\documentclass[a4paper,11pt]{book}

%-----------------PREAMBLE---------------%
%Graphics
\usepackage{graphicx}
\usepackage{color}
\DeclareGraphicsRule{.pdftex}{pdf}{.pdftex}{}
\usepackage[]{subfigure}
\usepackage{algorithm}
\usepackage{algorithmic}
\usepackage{enumerate}
\usepackage{bbold}

\usepackage{tikz}
\usetikzlibrary{decorations.markings,matrix,arrows}
\usetikzlibrary{trees}
\usetikzlibrary{decorations.pathmorphing}
\usetikzlibrary{decorations.markings}


%Text
\usepackage{appendix}
\usepackage{empheq}
\usepackage{listings}
\lstset{ %
language=C,                % the language of the code
basicstyle=\footnotesize,       % the size of the fonts that are used for 
numbersep=5pt,                  % how far the line-numbers are from the code
backgroundcolor=\color{white},  % choose the background color. You must add \usepackage{color}
showspaces=false,               % show spaces adding particular underscores
showstringspaces=false,         % underline spaces within strings
showtabs=false,                 % show tabs within strings adding particular underscores
frame=single,                   % adds a frame around the code
tabsize=2,                      % sets default tabsize to 2 spaces
captionpos=b,                   % sets the caption-position to bottom
breaklines=true,                % sets automatic line breaking
breakatwhitespace=false,        % sets if automatic breaks should only happen at whitespace
title=\lstname,                 % show the filename of files included with \lstinputlisting;
                                % also try caption instead of title
escapeinside={\%*}{*)},         % if you want to add a comment within your code
morekeywords={*,...}            % if you want to add more keywords to the set
}

%For tables
\usepackage{multirow} 

%---------------------------------------------------
%customizing page body layout using the geometry package:
%---------------------------------------------------
\usepackage[vscale=0.8,left=3cm,right=2cm,footskip=2.5cm, headsep=1cm]{geometry}
%-----------------------------------------------------------
%customizing headers and footers using the fancyhdr package:
%-----------------------------------------------------------
\usepackage{fancyhdr}
%redefining plain style,(necessary to remove default page-style setings):
\fancypagestyle{plain}{
\fancyhf{}
\renewcommand{\headrulewidth}{0pt}
\renewcommand{\footrulewidth}{0pt}}
%seting new style
\pagestyle{fancy}
\fancyhf{} %clears both head and foot
\renewcommand{\chaptermark}[1]{\markboth{\chaptername\ \thechapter.\ #1}{}}
\renewcommand{\sectionmark}[1]{\markright{\thesection\ #1}{}}
\fancyhead[LE]{\footnotesize \bfseries \leftmark}
\fancyhead[RO]{\footnotesize \bfseries \rightmark}
\headheight 15pt
\fancyfoot[LE,RO]{\bfseries \thepage}

%----------------------------------------
%redefining the \cleardoublepage command
%----------------------------------------
%automatically clears the (even) blank side following each chapter
\makeatletter
\def\cleardoublepage{\clearpage\if@twoside \ifodd\c@page\else
  \hbox{}
  %\vspace*{\fill}
  %If you want the blank page following the chapters
  %to have some content, this is added here. 
  %\vspace{\fill}
  \thispagestyle{empty}
  \newpage
  \if@twocolumn\hbox{}\newpage\fi\fi\fi}
\makeatother

%-------------------------------
%section-setings,no indentation:
%-------------------------------
%\usepackage{parskip}
%\setlength{\parindent}{0pt}
%\setlength{\parskip}{\baselineskip}

%-------------
%math setings:
%-------------
\usepackage{amsmath}
\usepackage{amsfonts}
\usepackage{amssymb}
\usepackage{simplewick} 
\usepackage{amsthm}  %Gives us abillity to begin proof

%changes refrences into links:
\usepackage{hyperref}

%Tabulars
\usepackage{booktabs,tabularx}

%---------------ENDING PREAMBLE---------------%

%----Note----%
%\footnote{}
%\emph{}
%\label{} - \ref{} - \pageref{} 

%--------------COSTUMIZE NEWCOMMANDS-----------------%
\newcommand{\bra}[1]{\langle{#1}|}
\newcommand{\ket}[1]{|{#1}\rangle}
\newcommand{\braket}[2]{\langle{#1}|{#2}\rangle}
\newcommand{\ql}{\left}
\newcommand{\qr}{\right}
\newcommand{\hh}[1]{\widehat{#1}}
\newcommand{\p}[1]{\hat{#1}}
\newcommand{\da}[1]{{#1}^{\dagger}}
%\newcommand{\bb}[1]{{\bf{#1}}}
\newcommand{\dd}[1]{\delta_{#1}}
\newcommand{\vv}[2]{\bra{#1}v\ket{#2}}
\newcommand{\vva}[2]{\bra{#1}|\hat{v}|\ket{#2}}


\renewcommand{\vec}[1]{\mathbf{#1}}
\newcommand{\bb}[1]{\, \ensuremath{\mathbf{#1}}}
%--------------COSTUMIZE amsthm package-----------------%
\newtheorem{theorem}{Theorem}[section]
\newtheorem{lemma}[theorem]{Lemma}
\newtheorem{proposition}[theorem]{Proposition}
\newtheorem{corollary}[theorem]{Corollary}


\providecommand{\e}[1]{\ensuremath{\times 10^{#1}}}

\begin{document}

\begin{titlepage}
\begin{center}
{\huge \bf Coupled-Cluster Studies of Double Quantum Dots}\\
\ \\ 
{\Large by}\\
\ \\
{\Large \bf Yang Min Wang}\\
\ \\
\ \\
\ \\
{\large \bf THESIS}\\
{\large for the degree of}\\
{\large \bf MASTER OF SCIENCE}\\
\ \\
(Master in Computational Physics)\\
\ \\
\ \\
%\includegraphics[width=60mm,height=60mm]{uiologo}
%\includegraphics[width=0.5\textwidth]{uio-logo.jpg}
\begin{figure}[h!]
\begin{center}
\includegraphics[width=0.3\textwidth]{uiologo}
%\includegraphics[width=0.8\textwidth]{UiO_Segl_black.eps}
\end{center}
\end{figure}
\ \\
\Large{\rm Faculty of Mathematics and Natural Sciences}\\
{\rm Department of Physics}\\
{\rm University of Oslo}\\
\ \\
\ \\
{\rm June 2011}

%\author{Marte Hoel J\o rgensen}
%\title{Thesis for the degree of Master of Science}
%\affiliation{Faculty of Mathematical and Natural Sciences}
%\maketitle
%\thispagestyle{empty}
\end{center}
\end{titlepage}

\chapter*{Acknowlegdement}

I am very thankful to my supervisor, professor Morten Hjort-Jensen, who has taught me ever since 2008. He has been a mentor and helped me far beyond what I would expected. His enthusiasm and humor have been a motivation. I am very lucky to have met Morten, and I hope I have in return fullfiled his expectations. I am deeply grateful to the many hours Simen Kvaal have spent helping me to implement the transformations and writing replies to my questions. 

It has been interesting to work with such a likable person as Marte Hoel J\o rgensen, whom I worked with on the optimization of the CCSD code. I admire her professionalism and discipline, and she has helped me to see things from different perspectives.

I would also thank Gustav Jansen and \O yvind Jensen for their insights and suggestions that have helped me understand the nature of my problems. Magnus Pedersen Lohne has been very helpful and responded every time on email. Thanks to my fellow students David Sk\aa lid Amundsen, Andreas Nakkerud, J\o rgen Tr\o mborg, Christoffer Hirth, J\o rgen H\o gberget, who have helped me along the way and made my time at the University of Oslo very enjoyable.
 
I have enjoyed the physics as well as non-physics discussions with my friends Henrik and Michael, they have broadened my perspectives on what a fruit is and is not. The discussions have been very fruitful to say the least. It has been a challenging year, and skiing has been very relaxing, thanks to Patrick, Vladimir, Morten, Christoffer, Thomas, Henrik and Michael for joining me for some memorable trips in Lillomarka. Shout-out to Henrik Holte for proof-reading my thesis.

Finally, I want to give credit to my mother, father, and my sister. Without their encouragement and support there would be no thesis. 

\begin{quotation}
\hfill Yang Min Wang
\end{quotation}


\clearpage
\begin{figure}
\begin{quotation}
\it ``The highest forms of understanding we can achieve is laughter and human compassion.''
\end{quotation}
\hfill R. P. Feynman.
\end{figure}

\tableofcontents

%\maketitle
%\thispagestyle{empty}

%\chapter*{Abstract}

\section*{Abbreviations}
\begin{table}[ht]
\begin{tabular}{l l}
EPV & Exclusion-Principle-Violationg\\
Eq & Equation\\
Fig & Figure\\
BCH & Baker-Campbell-Hausdorff expansion\\
SD & Slater Determinant\\
PEP & Pauli Exclusion Principle\\ 
HF & Hartree-Fock Method\\
RHF & Restriced Hartree-Fock Method\\
DFT & Density functional Theory\\
MBPT & Many-Body Perturbation Theory\\
CC & Coupled-Cluster\\
CCM & Coupled-Cluster Method\\
CCS & Coupled-Cluster Method with Single Excitations\\
CCD & Coupled-Cluster Method with Double Excitations\\
CCSD & Coupled-Cluster with Single and Double Excitations\\
CCSDT & Coupled-Cluster with Single, Double and Triple Excitations\\
CCSDTQ & Coupled-Cluster with Single, Double,Triple and Quadrouple Excitations\\
DMC & Diffusion Monte Carlo \\
\end{tabular}
\label{table:abbreviations}
\end{table}

%%%%%%%%%%%%%%%%%%%%%%%%%%% INTRODUCTION %%%%%%%%%%%%%%%%%%%%%%%%%%%%%%%%%%%%%%%%%%%%%%%%%%%
\chapter*{Introduction}
The goal of this thesis is to show how we can study  double quantum dots (and more complicated as well) systems using coupled custer theory. During this process we have used Magnus Pedersen Lohnes Master's Thesis \cite{lohne} as a starting point, and a lot of the work has been devoted in optimizing his code. I have been working with Marte Hoel Joergensen \cite{marte} with the optimalization, and we have developed a method of using a Hartree-Fock calculation as an input to improve the convergence of our coupled-cluster calculations. 

This thesis could be regarded as a academic curiosity. However we will show that the versatility of the CCSD method and the general features  of our confining potentials, can be used in linking theory with experiment. 
This is the main reason why we want to study quantum dots. People have applied Density Functional theories (DFT) to describe such systems \cite{introdft}, but we want to describe it from first principle methods. This is why we use the ab initio coupled-cluster method. Quantum dots form a very active and lively research field in condensedmatter theory. They have been used to develop devices such as single electron transistors, quantum dot lasers \cite{introqdotlaser} and ``artificial atoms'' \cite{ashoori}. They share similar properties such as shell structure and magic numbers as seen in atoms and nuclei. Lately one of the important applications is the use of colloidal quantum dot for detecting cancer cells \cite{gao}. 
To address such system, we need a proper theoretical description of the electronic system, exchange coupling, correlation energies, and ground state energies.

The quantum dot is essentially a device that can trap electrons. The typical size of these devices are between a micrometer to few nanometers. In these quantum dots we can see 
quantum mechanical effects such as discrete energy levels. Our quantum mechanical model for the quantum dot is a parabolic quantum well which traps the electrons in two dimensions. Such a model is an idealization of quantum dots that serve as  a starting point to understand realistic quantum dots. They are in fact crystalline and have periodic potentials. That is the reason why we are interested in the double quantum dot. This could give us insights of the physics in more real life systems. 
 
\subsubsection{Structure of Thesis}
The thesis is divided into three parts. Part one is an introduction to quantum mechanics, I start by covering the one-body Schr\"{o}dinger equation and then the development of the many-body theory which we need in part two.

In part two we cover the different many-body theories, \emph{Hartree-Fock} and \emph{CCSD}. Finally, part three contains the results and implementation of our code. We finish off with discussions and then present out conclusion.

%%%%%%%%%%%%%%%%%%%%%%%%%%%%%%%%%%%%%%%%%%%%%%%%%%%%%%%%%%%%%%%%%%%%%%%%%%%%%%%%%%%%%%%%%%%%
%%%%%%%%%%%%%%%%%%%%%%%%%%%%%%%%%%%% PART FUNDAMENTALS %%%%%%%%%%%%%%%%%%%%%%%%%%%%%%%%%%%%%
%%%%%%%%%%%%%%%%%%%%%%%%%%%%%%%%%%%%%%%%%%%%%%%%%%%%%%%%%%%%%%%%%%%%%%%%%%%%%%%%%%%%%%%%%%%%
 
\part{FUNDAMENTALS}

\chapter{Quantum Mechanics for single-particle systems}
\begin{quotation}
\it
``Quantum mechanics is that mysterious, confusing discipline, which really no one of us understands but which we know how to use. It works perfectly, as far as we can tell, in describing physical reality, but it is a counter-intuitive discipline,as the social scientists would say. Quantum mechanics is not a theory, but rather a framework within which we believe any correct theory must fit.''
\end{quotation}
\begin{quote}
   Murray Gell-Mann, in Mulvey (1981)
\end{quote}



\section{Introduction}
In the late 19th century scientists had problems with their classical understanding of the ways things are. Some experiments showed up that did not coicide with the current theories. The beginning of quantum mechanics was when Max Planck published in a theory of black-body radiation (1900). He explained that atoms can absorb and emit discrete quanta of radiation with energy $\epsilon = h f$, where $f$ is the radiation frequency and $h$ is the fundamental constant called Planck's constant.

\begin{equation}
h \approx 6.626 \times 10^{-34} \mbox{Js}
\end{equation}
%
In classical physics we distinguish between the particle and waves, a classical particle cannot be wavelike and particle like at the same time. But as some experiments have shown, this does not reflect the reality (Young's famous double slit experiment). Quantum mechanics can sometimes be counterintuitive in that regard, but then again, Newton's theory of gravity must have been difficult to grasp at first. How do we know about these ``invisible'' forces which we cannot see? 

Another problem people had was that the atoms would be unstable if we had electrons that were particle-like and orbited around defined orbits. The orbiting electrons would radiate electromagnetic energy and eventually fall into the nucleus. We need to have a wavelike description of the atomic electrons in order to explain their stability. Wave-like electrons are confined inside the atom, and at the lowest state, the ground state, the electron cannot radiate away its energy and fall into the nucleus. This gave us a whole new way of looking at nuclear physics. It revolutionized the 20th century physics.

\subsection{Measurement and Observables}
Measurements are done by a subject, usually called an \emph{observer}, that has an instrument which takes measurements on an \emph{object}. During this process we have disturbances and never ideal conditions. In classical physics the disturbance is directly associated with measurement itself and can be made arbitrarily small, depending of how good the engineers are. If we wanted to measure length with 5 decimal precision, we could just get a more precise ruler. But this is not the case in quantum mechanics. Each time we do measurements in quantum mechanics we have a probability distribution of different outcomes. This has nothing to do with the instrument itself, it is a part of the intrinsic nature of formulation of the quantum mechanics called the \emph{Heisenberg uncertainty principle}.

\begin{equation}
 \Delta x \Delta p \geq \frac{\hbar}{2}
 \label{heisenberg}
\end{equation}
%
This means that if we determine the momentum exactly $\Delta x = 0$, then the momentum is totally uncertain. The concept that particles exist with definite position and momentum is an idealistic classical concept.   
%
\subsection{The Schr\"{o}dinger equation} 
The Schr\"{o}dinger equation is the quantum mechanical eqvivalent to Newton's Second Law. It describes the motion of a quantum mechanical particle. This is a partial differential equation which describes how the wave function, representing the particle, \emph{flows}. The time-dependent Schr\"{o}dinger equation for a particle moving in a three-dimensional potential is:

\begin{equation}
i \hbar \frac{\partial \Psi}{\partial t} = \hat{H}\Psi
\label{schrodinger}
\end{equation}
%
Where $H$ is called the Hamilton operator and $\Psi$ is the wavefunction:

\begin{equation}
 \hat{H} = -\frac{\hbar^2}{2m} \nabla^2 + V(\vec{r})
\end{equation}
%
It is important to point out that this equation describes a non-relativistic motion of a quantum particle i.e. $E >> m_0 c^2$, where $m_0$ is the rest mass of the quantum particle. 

\subsection{Probabilities}
Max Born introduced an interpretation of the Schr\"{o}dinger wave function $\Psi(\vec{r},t)$. He pointed out that the probability of detecting a particle at a certain location and time is proportional to $|\Psi(\vec{r},t)|^2$. Thus $|\Psi(\vec{r},t)|^2$ is often viewed as a probability density at the position $r$ and time $t$. And the wavefunction is often referred to as a probability amplitude.

We can normalize our probability density by summing up all the possible position of the particle to one:

\begin{equation}
 \int |\Psi(\vec{r},t)|^2 d^3 \vec{r} = 1
\label{normalization}
\end{equation}
%
The wavefunction could also be a function of momentum and thereby describe the probability for finding a momentum in a certain range.

\section{Postulates of Quantum Mechanics}
We could argue that Newton's law of motion is the postulate in classical mechanics. What are Newton's laws derived from?. So likewise quantum mechanics is based on some fundamental ``laws of nature'' that must be underivable.  

A postulate is a statement made without any proof, an ``underived'' statement. In physics a postulate could be translated into a proposal which could either be verified or falsified based on experiment \cite{bowman}.

\begin{quotation}
{\bf Postulate 1}{\it We describe a system by its state vector $\ket{s}$, an observable q and it's a hermitian operator $\hat{Q}$ which operates on any $\ket{s}$}
\end{quotation}

\begin{quotation}
{\bf Postulate 2 }{\it The time evolution of the quantum state $\Psi$ is goverened by the time-dependent Schr\"{o}dinger equation. $\Psi$ is a function of position coordinates $q_n$ and the time $t$. $\hat{H}$ is the Hamiltonian}
\begin{equation}
 i\hbar\frac{\partial \Psi}{\partial t} = \hat{H} \Psi
\label{eq:tdseq}
\end{equation}
\end{quotation}

\begin{quotation}
{\bf Postulate 3 }{\it The possible results of measurement of an observable $Q$, are its eigenvalues $q_i$ of its operator $\hat{Q}$.}
\end{quotation}

\begin{quotation}
{\bf Postulate 4 }{\it If the results of the measurements is found to be $q_i$, then after the measurement the system will ''collapse`` into a corresponding eigenstate $\ket{q_i}$.}
\end{quotation}

\begin{quotation}
{\bf Postulate 5 }{\it The expectation value of an observable is given by}
\begin{equation}
\langle F \rangle = \int \psi^* \hh{F} \psi \, d \tau \qquad \mbox{(coordinate representation)}
\end{equation}
\end{quotation}


\section{Different Representations of Quantum Mechanics}
As we have seen, originally the state $\Psi$ was a solution to the Schr\"{o}dinger equation, which was a partial differential equation, second order in space and first order in time. So the ''state function`` $\Psi$ has to be a function of time and space coordinates. We will refer to this as the Schr\"{o}dinger representation. Quantum mechanics could be more abstract and more convenient as shown by Dirac \cite{dirac}. Another way to represent it is by matrix mechanics introduced by Heisenberg.
Those are the main ways that we can use to represent quantum mechanics, and they are mathematically equivalent.

\subsection{Dirac's Notation}
An abstract quantum mechanical state $\Psi$ is represented by a ''bra`` vector $\bra{\Psi}$ or a ''ket`` vector $\ket{\Psi}$. The distinction between the forms lies in the context in which they are used and will become clearer when we show this. The scalar product in Diracs notation is the ''bracket``, which is very convenient compared to writing the integral in the Schr\"{o}dinger representation.

\begin{equation}
\bra{\Phi}\ket{\Psi} \equiv \int \Phi^*(\vec{r})\Psi(\vec{r}) d\vec{r}
\label{eq:bracket}
\end{equation}
%
The expectation value of an operator is:

\begin{equation}
\bra{\Phi}\hat{Q}\ket{\Psi} \equiv \int \Phi^*(\vec{r})\hat{Q}\Psi(\vec{r})d\vec{r}
\label{eq:dirac_expect} 
\end{equation}
%
The time-independent Schr\"{o}dinger equation in the Dirac notation:

\begin{equation}
\hat{H} \ket{E_i} = E_i \ket{E_i}
\end{equation}
%
In this example we see that the label $E_i$ inside the ket tells us about its eigenvalue. The spectra of the eigenvalues to an operator could be discrete or continuous like the operator $\hat{x}$: 

\begin{equation}
 \hat{x}\ket{x} = x \ket{x}
\label{eq:xrep}
\end{equation}
%
If we wish to get back to the Schr\"{o}dinger representation $\Psi(x)$, we could do a ''projection`` of $\ket{\Psi}$ on the eigenfunction $\ket{x}$:

\begin{equation}
 \Psi(x) \equiv \braket{x}{\Psi}
\label{eq:coordinatrep}
\end{equation}
%
Similary the complex conjugate:
\begin{equation}
 \Psi^*(x) \equiv \braket{\Psi}{x}
\label{eq:ccoordinatrep}
\end{equation}

\subsection{Heisenberg's matrix formulation}
\label{Heisenberg}

The state vector $\ket{\Psi}$ is represented as a vector by its projection on a complete set of basis states $\ket{E_1},\ket{E_2}...$:

\begin{equation}
 \ket{\Psi} \equiv \left(\begin{array}{c} \braket{E_1}{\Psi} \\ \braket{E_2}{\Psi} \\ ... \\ ... \end{array}\right)
\label{eq:heisenbergstate}
\end{equation}
%
Similary we represent the operators by:

\begin{equation}
 \hat{H} \equiv 
\left(\begin{array}{cccc} 
\bra{E_1}\hat{H}\ket{E_1} & \bra{E_1}\hat{H}\ket{E_2} & \bra{E_1}\hat{H}\ket{E_3} & ... \\ 
\bra{E_2}\hat{H}\ket{E_1} & \bra{E_2}\hat{H}\ket{E_2} & ... & ... \\
... & ... & ... & ... \\
\end{array}\right)
\label{eq:heisenbergopeator}
\end{equation}
%
Since the basis states $\ket{E_n}$ are the eigenstates of $\hat{H}$, its matrix representation is diagonal, and the diagonal elements are eigenvalues of $\hat{H}$. To diagonalize an operator in matrix representation is therefore the same as solving the time-independent Schr\"{o}dinger equation. 
\section{Angular Momentum and Spin}

In classical physics angular momentum is defined as:
\begin{equation}
 \vec{L} = \vec{r} \times \vec{p}
 \label{angular momentum}
\end{equation}
$\vec{r}$ is the displacement vector from the origin and $\vec{p}$ is the linear momentum. 
We can write out the components of the angular momentum operator in quantum mechanics by using the substitution
\[
 p_x \rightarrow -i\hbar \frac{\partial}{\partial x}
\]
%
Gouldsmit and Uhlenbeck (in 1925) introduced the concept of interal, purely quantum mechanical, angular momentum called \emph{spin}. This was later experimentally confirmed by the Stern-Gerlach experiment (1922).
\newline
\newline
The spin eigenstates are:
\begin{equation}
\hat{S}_z \ket{l,m} = m \hbar \ket{l,m} 
 \label{zspin}
\end{equation}

\begin{equation}
\hat{S}^2 \ket{l,m} = l(l+1) \hbar^2 \ket{l,m}
 \label{totalspin}
\end{equation}
%
And their commutation relation reads

\begin{equation}
  \left[\hh{S}_x,\hh{S}_y \right] = i\hbar\hh{S}_z  \qquad \left[ \hh{S}_y,\hh{S}_z \right] = i\hbar\hh{S}_x \qquad \left[ \hh{S}_z, \hh{S}_x \right] = i\hbar \hh{S}_y
\end{equation}
%
This basically means that we cannot determine eigenvalues of two different components, for example $\hh{S}_z$ and $\hh{S}_x$ simultaneously. They are \emph{incompatible} observables. But we can determine one of the directions and $\hh{S}^2$. The spin quantum numbers are \cite{griffiths}

\begin{align}
  s = 0,\frac{1}{2},1,\frac{3}{2},... \\
  m_s = -s,-s+1,...,s-1,s
\end{align}
%
Where $s$ is defined as the \emph{spin} of the particle. Electrons have a spinvalue of $1/2$ spin while the photons have a spinvalue of $1$.

\subsection{The wavefunction}
The total wavefunction of a state vector mentioned in Postulate 1 are composed of a spatial part $\phi(x,y,z)$ and a spin part $\ket{\chi}$. They are from different Hilbert spaces and mathematically the state vector is a tensor product of these two. 

\begin{equation}
 \ket{\psi} \equiv \psi(r) \equiv \phi(x,y,z) \otimes \ket{\chi}.
\end{equation}


\section{Simple systems}

\subsection{Particle in a Infinite Potential Well}
One of the simplest single-particle systems we could solve exactly the infinite one-dimensional potential well, defined by:

\begin{equation}
V(x) = \left\{ \begin{array}{rl}
 -x &\mbox{ if $0\geq x \geq a$} \\
  \infty, &\mbox{ otherwise}
       \end{array} \right.
 \label{def:infinitepotential}
\end{equation}
%
and we want to solve the Schr\"{o}dinger equation with respect to this potential,

\begin{equation}
-\frac{\hbar^2}{2m}\frac{d^2 \psi}{dx^2} + V(x) = E \psi.
 \label{eq:schrodingerinfinitewell}
\end{equation}
% 
We can rewrite this equation as

\begin{equation}
\frac{d^2 \psi}{dx^2} = \frac{2m}{\hbar^2}\left[V(x)-E\right]\psi.
 \label{eq:rewriteschrodingerinfinitewell}
\end{equation}
%
For the case $x>a$ and $x<0$ where $E < V(x)$, we see that the second derivate always has the same sign and therefore cannot be normalized (ref postulate). Such wavefunctions cannot exist in this range. Therefore the wavefunctions only exist in the range $0\geq x \geq a$ and the continuity of $\psi$ requires the boundary conditions $\psi(0)=0$ and $\psi(L)=0$. 

Using standard methods of solving second order differential equations (see \cite{boas} for details), we get 

\begin{equation}
\psi(x) = A\sin kx
 \label{eq:infinitewellsolution}
\end{equation}

\begin{equation}
k \equiv \frac{\sqrt{2mE}}{\hbar} 
 \label{def:k}
\end{equation}
%
Normalizing in order to get the constant $A$:

\begin{equation}
\int_0^a |A|^2 \sin^2 kx \, dx = |A|^2 \frac{a}{2} = 1 \Rightarrow |A|^2 = \frac{2}{a},
 \label{deriv:normalizing}
\end{equation}
%
We then choose to use the positive root 

\begin{equation}
\psi(x) = \sqrt{\frac{2}{a}} \sin kx,
 \label{eq:infinitewellsolution2}
\end{equation}
%
Next we want to determine the konstant $k$. From the boundary condition we know that $\sin ka = 0$. This means that the possible values for $k$ are

\begin{equation}
 k = \frac{n \pi}{a}, \qquad \mbox{with $n = 1,2,3...$}
 \label{eq:kvalues}
\end{equation}
%
and hence the possible values of $E$ from Eq.~(\ref{def:k}) are discrete:

\begin{equation}
E_n = \frac{\hbar^2 k^2}{2m} = \frac{n^2\pi^2\hbar^2}{2ma^2}.
 \label{eq:energies}
\end{equation}


%%------Under construction--------

\subsection{Particle in a Harmonic Oscialltor Potential Well}
\label{onebodyharmonicoscillatorpotential}
The next simplest problem to solve is the \emph{harmonic oscillator} in one-dimensional case. The Hamiltonian we have is

\begin{equation}
  \hh{H} = \frac{\hat{p}^2}{2m} + \frac{1}{2}\omega^� \hat{x}^2,
\end{equation}
%
where $\hat{x} = x$ is the position operator and the $\hat{p}$ is the momentum operator, given by

\begin{equation}
  \hat{p} = -i\hbar \frac{d}{dx}.
\end{equation}
%
There are two main approaches to solve this. One of them is the analytical approach \cite{griffiths}, in which we get the solutions in Schr\"{o}dinger representation,

\begin{equation}
  \psi_n(x) = \sqrt{\frac{1}{2^n n!}} \left(\frac{m\omega}{\pi \hbar}\right)^{1/4} e^{-\beta^2 /2} H_n(\beta), \qquad \beta = \sqrt{\frac{m\omega}{\hbar}}x
\end{equation}
%
And the eigenvalues,

\begin{equation}
  E_n = \left(n + \frac{1}{2}\right) \hbar \omega, \qquad n = 0,1,2,3...
\end{equation}
%
The other way is algebraically \cite{griffiths}, with which we get solutions in Dirac formalism:

\begin{equation}
  \ket{\psi_n} = \frac{1}{\sqrt{n!}} (a_+)^n \ket{0},
\end{equation}
%
where $\ket{0}$ is our ground state and $a_+ = (a_-)^\dagger$ is the creation operator
%
\begin{equation}
  a_{\pm} = \frac{1}{\sqrt{2\hbar m \omega}} (\mp i\hat{p} + m\omega \hat{x}).
\end{equation}

%%-----------------  MATHEMATICAL DESCRIPTION OF THE QD -------------------
\chapter{The Quantum Mechanics behind Quantum Dots}
In this chapter we will give a quantum mechanical description for the two-dimensional quantum dot. We will start by solving the Schr\"{o}dinger equation for a single-electron parabolic quantum dot with an applied magnetic field. 

Recent developments in techniques of quantum dot growth have made it possible to create better solar cells, as well as other exotic phenomena such as quantum computing. Quantum dot is a semiconductor on the nanoscale. It can trap one or several electrons in a spatially confined potential. The size of quantum dot ranges from a few hundreds to many thousand atoms \cite{qdotshur}. And it can confine everything from one electron to hundreds. 

\section{Description of the Quantum Dot}
As we have seen in the chapter 1, quantum dots are artificially created. And there are various techniques and methods for creating quantum dots, which give them different properties. In this thesis we will concentrate on the quantum dots created inside the Gallium Arsenide (GaAs) semiconductors. The semiconductor is sandwiched between layers of Aluminum Gallium Arsenide (AlGaAs) semiconductor material which has a bigger bandgap. This acts like an insulator and results in a confinement in the vertical direction. Our choice for the confinement potential is a parabolic harmonic oscillator $\omega_x = \omega_y = \omega$. Both numerical \cite{numerical1, numerical2} and experimental \cite{experimental1, experimental2} studies have shown that this is a reasonable approximation. And in our case the electrons inside will only feel the Coulomb interaction. The Hamiltonian then becomes 
\begin{equation}
\hh{H} = \sum_{i=1}^N \left(-\frac{\hbar^2}{2 m^*} \nabla_i^2 + \frac{1}{2} m^* \omega^2 r_i^2 \right) + \frac{e^2}{4\pi \epsilon_0 \epsilon_r} \sum_{i<j}^N \frac{1}{r_{ij}}.
  \label{Hamiltonian for QD}
\end{equation}
The $e$ is the electron charge, $\epsilon_0$ is the vacuum permittivity, $\epsilon_r$ is the relative permittivity, and $r_{ij} = |r_i - r_j|$, $\omega$ is the oscillator frequency, and $r_i$ is the distance from electron $i$ to the potential minimum $(r=0)$.
It is important to notice that the $m^*$ here is the effective electron mass and not to be mistaken for the Newtonian reduced mass which is a classical phenomena. This is a simplification we have made to our problem. The effective electron mass differs from the free-electron mass $m$ and is isotropic and independent of both the position and the energy of the electron. The effective mass is the result of the motion of an electron in a periodic potential \cite{quantumheterostructures}. For example in GaAs the electrons appear to carry mass that is only $7\%$ of the free-electron mass \cite{ashoori}.  

\section{The One-electron Quantum Dot}
We start by deriving the wavefunctions for a general spherical symmetric potential $V(r) = V(-r)$ in two dimensions. Our Hamiltonian in cartesian coordinates reads
\begin{equation}
  -\frac{\hbar^2}{2m^*} \nabla^2 \psi(x,y) + V(x,y) \psi(x,y) = E \psi(x,y).
   \label{Hamiltonian cartesian coordinates}
\end{equation}
%
We want to rewrite this using polar coordinates
%
\begin{align}
x &= r \cos \theta, \\
y &= r \sin \theta, \\
r &= \sqrt{x^2 + y^2}.
\end{align}
The Laplacian then becomes
\begin{equation}
\nabla^2 = \frac{1}{r} \frac{\partial}{\partial r}\left(r\frac{\partial}{\partial r}\right) + \frac{1}{r^2} \frac{\partial}{\partial \theta^2}.
  \label{laplacian in 2D}
\end{equation}
%
Inserting this into Eq.~(\ref{Hamiltonian cartesian coordinates}) gives us the Schr\"{o}dinger equation in polar coordinates
\begin{equation}
  -\frac{\hbar^2}{2m^*} \left(\frac{\partial^2}{\partial r^2} + \frac{1}{r} \frac{\partial}{\partial r} + \frac{1}{r^2} \frac{\partial^2}{\partial \theta^2}\right) \psi(r,\theta) + V(r) \psi(r,\theta) = E \psi(r,\theta).
   \label{Hamiltonian polar coordinates}
\end{equation}
Introducing a solution on the form $\psi(r,\theta) = R(r) Y(\theta)$ and multiplying by $\frac{2m^*}{\hbar^2 R(r)Y(\theta)}r^2$ in Eq.~(\ref{Hamiltonian polar coordinates}), we get

\begin{equation}
\frac{r^2}{R(r)}\left[\frac{d^2 R}{d r^2} + \frac{1}{r} \frac{dR}{dr} + \frac{2m^*}{\hbar^2}(E-V(r))R(r) \right] = -\frac{1}{Y(\theta)}\frac{\partial^2 Y(\theta)}{\partial \theta^2},
  \label{separable Hamiltonian 2}
\end{equation}
the left side of this equation depends on $r$ while the right side depends on $\theta$. This can only be satisfied if each term is equal to a constant $k=m^2$
%
\begin{align}
\frac{r^2}{R(r)}\left[\frac{d^2 R}{d r^2} + \frac{1}{r} \frac{dR}{dr} + \frac{2m^*}{\hbar^2}(E-V(r))R(r) \right] &= -m^2 \label{separation constants1} \\
\frac{1}{Y(\theta)}\frac{\partial^2 Y(\theta)}{\partial \theta^2} &= m^2 .
  \label{separation constants2}
\end{align}
%
The solution to the angular part Eq.~(\ref{separation constants2}) is 
\begin{equation}
Y(\theta)  = Ce^{im\theta}.
\end{equation}
%
Normalization gives us the constant $C$
\begin{equation}
C^2 = \frac{1}{\int_0^{2\pi}Y(\theta)^2 \, d\theta} = \frac{1}{2\pi}.
  \label{constant C}
\end{equation}
%
The normalized solution for the angular part is
\begin{equation}
Y(\theta) = \frac{1}{\sqrt{2\pi}} e^{im\theta}.
  \label{normalized angular part}
\end{equation}
%
The total wavefunction must satisfy the physical condition that $\psi(r,\theta) = \psi(r,\theta+2\pi)$. This makes a restriction on the quantum number $m$ which can take integral values
\begin{equation}
m = 0, \pm 1, \pm 2, ...
  \label{m quantum number}
\end{equation}
For the Eq.~(\ref{separation constants1}) we can simplify by defining
\begin{equation}
  u(r) = \sqrt{r}R(r) \Rightarrow R(r) = \frac{u(r)}{\sqrt{r}},
  \label{radialequation function}
\end{equation}
which yields
\begin{equation}
  -\frac{\hbar}{2m^*}\frac{d^2 u}{dr^2} + \left[ V(r) + \frac{\hbar^2}{2m^*} \frac{m^2-\frac{1}{4}}{r^2} \right] u(r) = E u(r).
\end{equation}
This is the radial equation. It has the same form as the one-dimensional time-independent Schr\"{o}dinger equation with an effective potential
\begin{equation}
V_{eff} = V(r) + \frac{\hbar^2}{2m^*}\frac{m^2-\frac{1}{4}}{r^2}.
  \label{effective potential}
\end{equation}
The radial function Eq.~(\ref{radialequation function}) must satisfy the normalization conditions
\begin{equation}
\int_0^\infty |u(r)|^2 dr = 1.
  \label{normalization radial}
\end{equation}
For this to be normalizable we require the boundary conditions 
\begin{equation}
  \text{$u(0) = C$ and $u(\infty) = 0$}, \quad \text{where $C$ is a constant}.
\end{equation}
%
Finally the general solutions to the spherical symmetrical potential is
\begin{equation}
\psi(r,\theta) = R(r) \frac{1}{\sqrt{2\pi}} e^{im\theta}.
  \label{final solution to spherical}
\end{equation}

\subsection{Parabolic Quantum Dot with Influence of an Electromagnetic Field}
\label{parabolic quantum dot}

As an academic exercise we shall solve the one-electron Schr\"{o}dinger equation for a two dimensional quantum dot. Most of the derivation here is taken from \cite{jackson}. The classical Hamiltonian of a charged electron in an electromagnetic field reads
  
\begin{equation}
   H = \frac{1}{2m}(\bb{p} - e\bb{A})^2 + e\phi ,
   \label{electromagnetic Hamiltonian}
 \end{equation}
%
with $\bf{A}$ and $\phi$ as the electromagnetic vector and scalar potentials, $m$ is the electron mass, $e$ is the charge, and $\bb{p}$ is the momentum vector. The electromagnetic fields are related to the potentials 
\begin{align}
\bb{E} &= - \frac{1}{c} \frac{\partial \bb{A}}{\partial t} - \nabla \phi, \\ 
\bb{B} &= \nabla \times \bb{A},
  \label{elmag potentials}
\end{align}
where $\bb{E}$ is electric field, and $\bb{B}$ is the magnetic field, which satisfies Maxwells equations. The quantum mechanical Hamiltonian consists of and additional term which couples the spin to the electromagnetic field $-\mu\cdot \bb{B}$, where $\mu$ is the magnetic moment of the electron. 
\begin{equation}
  \hh{H} = \frac{1}{2m^*} (\hh{p} - e\bb{A})^2 + e\phi +\frac{1}{2}m^* \omega_0^2 r^2 - \hh{\mu} \bb{B}.
  \label{Total Hamiltonian}
\end{equation}
Here the $\hh{p}$ is a quantum mechanical momentum operator. The time-independent Schr\"{o}dinger we want to solve is 
\begin{align}
\frac{1}{2m^*} (\hh{p} - e\bb{A})^2 + e\phi +\frac{1}{2}m^* \omega_0^2 r^2 - \hh{\mu} \cdot \bb{B} \psi(\bb{r}) & = E \psi(\bb{r}),
  \label{Time-independent SE}
\end{align}
\begin{align}
H_r &= \frac{1}{2m^*} (\hh{p} - e\bb{A})^2  +\frac{1}{2}m^* \omega_0^2 r^2\label{spatial hamiltonian},\\
H_s &=  - \hh{\mu} \cdot \bb{B} \label{spin hamiltonian}.
\end{align}
%
where the wavefunction $\psi(\bb{r})$ also include the spin part which is decoupled with the spatial part, i.e. their expectation are uncorrelated, and we can therefore write the wavefunction as a product of each.
\begin{equation}
 \psi(\bb{r}) = R(r) \otimes \ket{m_s}.
  \label{wavef spin}
\end{equation}
%
Here the quantum number is $m_s=\pm \frac{1}{2}$, since electrons are fermions. Inserting this into Eqs. (\ref{spatial hamiltonian}) and (\ref{spin hamiltonian}) we obtain two equations, one spatial and one spin dependent
\begin{align}
  \left(\frac{1}{2m^*} (\hh{p} - e\bb{A})^2 + e\phi +\frac{1}{2}m^* \omega_0^2 r^2 \right)R(r) &= E_r R(r), \\
  -(\mu \cdot \bb{B}) \ket{m_s} &= E_s \ket{m_s}.
\end{align}
%
Which gives us the total energy 
%
\begin{equation}
E = E_r + E_s + e\phi.
  \label{total energy}
\end{equation}
%
We want to do a gauge transformation on the potential $\bb{A}$ with the Coulomb gauge condition. This will not change the potentials $\bb{E}$ and $\bb{B}$.
%
\begin{equation}
\bb{\nabla} \cdot \bb{A} = 0.
  \label{gauge}
\end{equation}
%
A choice that satisfies this condition is 
\begin{equation}
\vec{A} = \frac{1}{2} \bb{B} \times \left(x\,\bb{i} + y \, \bb{j} \right).
  \label{coulomb gauge}
\end{equation}
%
We want to expand the first term in the spatial Hamiltonian $\hh{H}_r$ Eq.~(\ref{spatial hamiltonian}) using the condition Eq.~(\ref{coulomb gauge})

\begin{align}
(\hh{p}^2 - e\bb{A})^2 &= \hh{p}^2 - e(\hh{p}\bb{A} + \bb{A}\hh{p}) + e^2 \bb{A},  \\
& = \hh{p}^2 - 2e \bb{A}\cdot\hh{p}+ e^2 \bb{A}^2, \label{second deriv} \\
& = \hh{p}^2 - e\mathbf{B}\cdot \hh{L} + \frac{e^2}{4}\left(\mathbf{B} \times (x \bb{i} + y \bb{j}) \right)^2, \label{thirdeq}  
\end{align}
%
In Eq.~(\ref{second deriv}) we have used that $\hh{p}$ and $\bb{A}$ commute because of the Coulomb gauge. 
\begin{equation}
\hh{p} \cdot \mathbf{A} \psi = -i\hbar\nabla \cdot \left(\mathbf{A} \psi \right) = -i\hbar\left( \underbrace{\nabla \cdot \mathbf{A}}_{=0} + \mathbf{A} \cdot \nabla \psi \right) = \mathbf{A} \cdot \left(-i\hbar\nabla \psi \right) = \bb{A}\cdot \hh{p} \psi
\end{equation}
And in Eq.~(\ref{thirdeq}) we have used the relation 
\begin{equation}
 \left( \mathbf{B} \times \vec{r} \right) \cdot \mathbf{p} = \mathbf{B} \cdot \left( \vec{r}  \times \mathbf{p} \right) =  \mathbf{B} \cdot \mathbf{L}.
  \label{math relation cross}
\end{equation}
%
The applied magnetic field is constant and homogeneous along the z-axis. $\bb{B} = B_0 \bb{k}$. Then our spatial Hamiltonian $\hh{H}_r$ simplifies to 
\begin{equation}
\hh{H}_r = \frac{1}{2m^*} \left[\hh{p}^2 - e\mathbf{B_0}\left(x\hh{p}_y - y \hh{p}_x \right) + \frac{e^2B_0^2}{4}\left( x^2 + y^2 \right)\right]
 + \frac{1}{2} m^* \omega_0^2 (x^2+y^2).
  \label{total hamiltonian 2}
\end{equation}
%
Introducing
\begin{equation}
  \omega_B \equiv \frac{eB_0}{2m^*},
\end{equation}
%
and 
\begin{equation}
  \omega \equiv \omega_0 + \omega_B^2,
\end{equation}
%
the Hamiltonian becomes 
\begin{equation}
\hh{H}_r = \frac{1}{2m^*} \left(\hh{p}^2 - eB_0\hh{L}_z \right) + \frac{1}{2} m^* \omega^2 \left(x^2+y^2\right) \label{final Hamil}.
\end{equation}
%
Where $\hh{L}_z$ is the angular momentum in the $z$-direction. $L_z = x\hh{p}_y - y \hh{p}_x$ \cite{hemmer}. The next step is to transform the Hamiltonian Eq.~(\ref{final Hamil}) from a cartesian coordinate representation to a polar coordinate representation. Then the angular momentum can be expressed
as \cite{hemmer} 
\begin{equation}
\hh{L}_z = -i\hbar \frac{\partial}{\partial \theta},
  \label{Lz in polar}
\end{equation}
%
which yields the time-independent Schr\"{o}dinger equation for $\hh{H}_r$ Eq.~(\ref{spatial hamiltonian})
\begin{equation}
\left[ - \frac{\hbar^2}{2m^*} \left(\frac{\partial^2}{\partial r^2} + \frac{1}{r} \frac{\partial}{\partial r} + \frac{1}{r^2}\frac{\partial^2}{\partial \theta^2} - \frac{ieB_0}{\hbar} \frac{\partial}{\partial \theta}\right) + \frac{1}{2}m^*\omega^2 r^2 \right] \psi(r,\theta) = E \psi(r,\theta).
  \label{Hamilton in polar}
\end{equation}
This is almost the same as Eq.~(\ref{Hamiltonian polar coordinates}) which we have solved, except we now have an additional term $\frac{ieB_0}{\hbar} \frac{\partial}{\partial \theta}$  caused by the magnetic field. But this can be separated from the radial equation, since it only depends on the angle $\theta$. 
%
The wavefunction can still be separated in an angular and a radial part, and we will use the same ansatz as Eq.~(\ref{final solution to spherical})
\begin{equation}
\psi(r,\theta) = R(r) \frac{1}{\sqrt{2\pi}} e^{im\theta}, \qquad m=0,\pm 1, \pm 2,... 
  \label{wave function ansatz}
\end{equation}
%
Inserting this into Eq.~(\ref{Hamiltonian polar coordinates}) 
\begin{align}
\left[ - \frac{\hbar^2}{2m^*} \left(\frac{\partial^2}{\partial r^2} + \frac{1}{r} \frac{\partial}{\partial r} + \frac{1}{r^2}\frac{\partial^2}{\partial \theta^2} + \frac{meB_0}{\hbar} \frac{\partial}{\partial \theta}\right) + \frac{1}{2}m^*\omega^2 r^2 \right] R(r) = E_r R(r),
  \label{Hamilton in polar 2}
\end{align}
%
the solution of this radial equation is
%
\begin{equation}
R_{nm}(r) = \sqrt{\frac{2n!}{(n+|m|)!}}\beta^{\frac{1}{2}(|m|+1)}r^{|m|}e^{-\frac{1}{2}\beta r^2} L_n^{|m|} (\beta r^2).
  \label{Laguerre solutions}
\end{equation}
%
Here the subscript $n$ denotes the principal quantum number, and $m$ is the angular momentum number
\begin{align}
  n &= 0, 1, 2, 3, ... \\
  m &= 0, \pm 1, \pm, 2, \pm 3, ...
\end{align}

$L_n^{|m|}$ is the associated Laguerre polynomials \cite{arfken}, and $\beta$ is defined as
\begin{equation}
\beta = \frac{m^*\omega}{\hbar}.
  \label{definition beta}
\end{equation}
%
The final eigenfunction to the spatial Hamiltonian $H_r$ is then
%
\begin{equation}
 \psi(r,\theta) = \sqrt{\frac{n!}{\pi(n+|m|)!}}\beta^{\frac{1}{2}(|m|+1)}r^{|m|}e^{-\frac{1}{2}\beta r^2} L_n^{|m|} (\beta r^2) e^{im\theta},
 \label{final for spatial Hamiltonian}
\end{equation}
%
with the corresponding eigenvalues
%
\begin{equation}
E_r = (1+|m|+2n)\hbar \omega + m\hbar \omega_B.
  \label{corresponding eigenvalues}
\end{equation}
%
See Appendix in \cite{lohne} and \cite{drosdal} for details. 

Now we will consider the spin Hamiltonian $H_s$ Eq.~(\ref{spin hamiltonian}). The quantum mechanical magnetic moment $\hh{mu}$ is given by \cite{griffiths}

\begin{equation}
\hh{\mu} = \frac{eg}{2m^*} \hh{S}.
  \label{magnetic moment}
\end{equation}
%
where $g$ is the $g$-factor and approximate $2$ for the electron. $\hh{S}$ is the spin operator. Since the magnetic field is $\mathbf{B} = B_0 \bb{k}$, the Schr\"{o}dinger equation reads
%
\begin{equation}
-\frac{egB_0}{2m^*}S_z\ket{m_s} = E_s \ket{m_s}.
  \label{SL for the spin part}
\end{equation}
%
$\hh{S}_z$ is the $z$-component of the total spin $\hh{S}$. The eigenvectors and eigenvalues of this operator are given by \cite{hemmer}
%
\begin{equation}
E_s = -\frac{eg\hbar B_0}{2m^*}m_s = gm_s \hbar \omega_B.
  \label{eigenvalue for spin}
\end{equation}
%
The total eigenvalue for the system becomes
\begin{equation}
E = (1+|m|+2n)\hbar \omega + m \hbar \omega_B + g m_s \hbar \omega_B + e\phi,
  \label{total energyeigenvalue}
\end{equation}
and the corresponding eigenstate
\begin{equation}
\psi_{nmm_s}(\bb{r},\theta) = \sqrt{\frac{n!}{\pi(n+|m|)!}}\beta^{\frac{1}{2}(|m|+1)}r^{|m|}e^{-\frac{1}{2}\beta r^2} L_n^{|m|} (\beta r^2) e^{im\phi} \otimes \ket{m_s}.
  \label{total eigenstate}
\end{equation}
%
Without any external magnetic field $\bb{B} = 0$, the energy becomes spin-independent because $\omega_B = 0$
\begin{equation}
E_{nm}^0 = (1+|m|+2n) \hbar \omega_0 
  \label{spin independent energy}
\end{equation}
%
In the spirit of perturbation theory we denote this with a superscript $0$. We will have a degeneracy in spin since the Hamiltonian is spin-independent. For each pair of the quantum numbers $\{n,m\}$ we have two different quantum states, one with $m_s=-\frac{1}{2}$ and the other with $m_s=\frac{1}{2}$.
\begin{equation}
R \equiv (1+|m|+2n).
  \label{shellnumber}
\end{equation}
%
$R$ is defined as the \emph{shellnumber}, it corresponds to the energy level and the degeneracy for each level $R$ is
\begin{equation}
  g(R) = 2R.
\end{equation}
%
This system has a \emph{shell structure}, Figure \ref{fig:shellnumber}, i.e. the energy levels are equidistant from each other and we have a defined degeneracy. This is similar to the shell model in nuclear physics, for which Goeppert-Mayer, Wiger and Jensen was awarded the Nobel Prize in Physics (1963).
%
\begin{figure}
\centering
\scalebox{0.6}{\input{shell.pdftex_t}}
\caption{Shell structure of a single-electron parabolic quantum dot for $\bb{B=0}$, where $R$ is the shell number defined in Eq.~ (\ref{shellnumber}) and $m$ is the angular quantum number. The arrows $\uparrow \downarrow$ denote the spin quantum number $m_s=\pm \frac{1}{2}$}
\label{fig:shellnumber}
\end{figure}
%
The total number of spin-orbitals for a given shellnumber $R$ is 
%
\begin{equation}
  N = \sum_{R=0}^{R'} g(R) = 2R' + 2(R'-1) + 2(R'-3) + ... + 2.
\end{equation}
%
We have tabulated some of the values in Table~\ref{tab:shellfilling} 

\begin{table}[H]
\centering 
\begin{tabular}{rrr}
\toprule
$R$ & $g(R)$ & $N$\\
\midrule
1	& 2 	 &  2 \\
2	& 4 	 &  6 \\
4	& 8 	 &  20  \\
8	& 16 	 &  72 \\
10	& 20 	 &  110 \\
15	& 30 	 &  240 \\
20	& 40 	 &  420 \\
\bottomrule
\end{tabular}
\caption{This table shows some values for different shellnumbers $R$, where $g(R)$ is the degeneracy. $N$ is the total number of single-electron spin-orbitals occupied in $R$-number of shelles. This is often referred to as \emph{magic numbers} which indicate the number of spin-orbitals needed to complete the shells.}
\label{tab:shellfilling}
\end{table}
%
If we take a look at Eq.~(\ref{total energyeigenvalue}), we see that in presense of a magnetic field, the degenerate energy levels would split because of the sign of $m$. If we simplify by neglecting the spin quantum number $m_s = 0$, and setting the constant $e\phi = 0$, we can express the energy in Eq.~(\ref{total energyeigenvalue}) by 

\begin{equation}
\frac{E_{nm}}{\hbar \omega_0} = (1+|m|+2n) \sqrt{1+\frac{\omega_B^2}{\omega_0^2}} + m\frac{\omega_B}{\omega_0}.
  \label{fock darwin}
\end{equation}
%
If we plot $\frac{E_{nm}}{\hbar \omega_0}$ as a function of $\frac{\omega_B}{\omega_0}$ we get the Fock-Darwin energy spectrum, which was first solved by V.Fock \cite{fock} and later by C.G. Darwin \cite{darwin}. 

\begin{figure}[H]
\centering
\scalebox{0.7}{\includegraphics{fockdarwin.eps}}
\caption{This is the Fock-Darwin energy spectrum for a single-electron quantum dot. When $\omega_B = 0$, all the quantum states with equal $R$ are degenerate, after $\omega_B$ increases the levels begin to split due to the angular momentum contribution from $m\hbar \omega_B$. The degeneracy can reappear for certain levels and certain $\omega_B$'s. The energy levels will shift back and forth between $(n,m)$ pairs, but they all appear to reach an asymptote in the high field limit, forming the famous Landau levels \cite{landaulevels}, indicated by the red line. The landau levels are plotted for $N_L =0,1,2$ in Eq.~(\ref{landaulevelsplot})}
\end{figure}
%
The energy of the different states will then decrease or increase with stronger magnetic field $\omega_B$, depending on $m$. States that belong to different shells for $\bb{B}=0$ will become degenerate, and when the magnetic field increases even more, we would reach an asymptote for the energies, as we clearly see in Figure \ref{fock darwin}. These energy levels are popularly called Landau levels, which Landau discovered in 1930 at an age of 22 \cite{landauonline}. 

We are interested in the energy levels when $\omega_B \rightarrow \infty$ for the lowest energy levels $m<0$

\begin{equation}
\lim_{\omega_B \rightarrow \infty} E_{nm} = (1+2n)\hbar\omega_B.
\end{equation}
%
The Landau levels appear when $N_L \equiv n = 0,1,2,3...$. And the energy only depends on $n$ in the high limit of $\bb{B}$

\begin{equation}
E_L \approx (1+2N_L) \hbar \omega_B.
\label{landaulevelsplot}
\end{equation}

\subsection{Scaling the Hamiltonian}
\label{section:scaling the hamiltonian}
When we do computations it is a good thing to have dimensionless parameters. The fewer things that can go wrong in the calculation the better. Therefore we want to rescale our many-body Hamiltonian Eq.~(\ref{Hamiltonian for QD}) so that its dimensionless. The following derivation is based on examples \cite{scaling,lohne}

\begin{align}
  \omega  & = \omega_c \bar{\omega}, \\
  \bb{r}  & = r_c \bb{\bar{r}}, \\
  \bb{\nabla}  & = \frac{1}{r_c} \bb{\bar{\nabla}}, \\
  r_i^2 & = r_c^2 \bar{r}_i^2, \\
  r_{ij} & = r_c \bar{r}_{ij}.
\end{align}
%
Here the variables with the subscript $c$ is just a constant with the same dimensions as the variable we want to rescale. The variables with a bar is the dimensionless variables that we want. Inserting these into our Hamiltonian Eq.~(\ref{Hamiltonian for QD}) gives

\begin{equation}
\hh{H} = -\frac{\hbar^2}{2m^*r_c^2} \sum_{i=1}^N \bar{\nabla}_i^2 + \frac{1}{2}m^* \omega_c^2 \bar{\omega}^2 r_c^2 \sum_{i=1}^N \bar{r}_i^2 + \frac{\hbar}{\epsilon r_c} \sum_{i<j}^N \frac{1}{\bar{r}_{ij}},
  \label{rescaled Qdot Hamiltonian}
\end{equation}
%
where 
\begin{equation}
  \epsilon = \frac{4\pi \epsilon_0 \epsilon_r}{e^2}.
\end{equation}
%
Furthermore we define the oscillator length to be
\begin{equation}
r_c = \sqrt{\frac{\hbar}{m^* \omega}}
  \label{oscillator length}
\end{equation}
Inserting this into our Hamiltonian Eq.~(\ref{rescaled Qdot Hamiltonian})
\begin{equation}
\hh{H} = -\frac{\omega_c \bar{\omega} \hbar}{2} \sum_{i=1}^N \bar{\nabla}^2_i + \frac{\hbar}{2} \omega_c \bar{\omega} \sum_{i=1}^N \bar{r}^2_i + \frac{\hbar}{\epsilon} \sqrt{\frac{m^*\omega_c \bar{\omega}}{\hbar}} \sum_{i<j}^N \frac{1}{\bar{r}_{ij}}
  \label{rerescaled Hamiltonian2}.
\end{equation}
%
We want to scale the Hamiltonian as well, so that it has units of the Hartree energy $E_h$ \cite{hartrees}
\begin{equation}
  \bar{H} =  \hh{H} / E_h.
\end{equation}
%
where the Hartree energy is defined as
%
\begin{equation}
E_h = m^* \left(\frac{e^2}{4\pi \epsilon_0 \epsilon_r \hbar}\right)^2 = \frac{m^*}{\epsilon^2}.
  \label{E hartree}
\end{equation}
%
The scaled Hamiltonian becomes
\begin{equation}
\bar{H} = -\frac{\omega_c\bar{\omega} \hbar \epsilon^2}{2m^*} \sum_{i=1}^N \bar{\nabla}^2_i + \frac{\hbar \epsilon^2}{2m^*} \omega_c \bar{\omega} \sum_{i=1}^N \bar{r}^2_i + \frac{\hbar\epsilon}{m^*} \sqrt{\frac{m^*\omega_c \bar{\omega}}{\hbar}} \sum_{i<j}^N \frac{1}{\bar{r}_{ij}}.
\end{equation}
%
To make it dimensionless we have to define
%
\begin{equation}
\omega_c = \frac{m^*}{\hbar \epsilon},
  \label{dimensional condition}
\end{equation}
%
which is fine since $\omega_c$ has the dimension of $[1/s]$ which is the same as $\omega$. Our final dimensionless N-electron scaled Hamiltonian reads

\begin{equation}
\bar{H} = -\frac{\omega_c}{2} \sum_{i=1}^N \bar{\nabla}^2 + \frac{1}{2} \omega_c \sum_{i=1}^N \bar{r}_i^2 + \sqrt{\omega_c} \sum_{i<j}^N \frac{1}{\bar{r}_{ij}}.
  \label{N-body Hamiltonian scaled}
\end{equation}


\section{The Double Quantum Dot}
\label{sec:The double quantum dot}
\label{section:the double dot potential}
The rapid development of nanotechnology has made the quantum computer a more realistic achievement. One of the possible candidates to the quantum bit is a double quantum dot. Experiments of this have been done both with a trapped nucleus \cite{doublenuclei} and electrons \cite{doubleelectron}. We will investigate a model based on \cite{doubledot}. Our one-electron Hamiltonian is the same as before but now with a change in the potential

\begin{equation}
  \hh{H} = -\frac{\hbar^2}{2m^*} \nabla_i^2 + V_c(x,y),
  \label{one electron hamilton double}
\end{equation}
where the confinement potential is
\begin{equation}
V_c(x,y) = \frac{1}{2} m^* \omega_0^2 \cdot \left[x^2+y^2 - 2L_x |x| + L_x^2 \right].
  \label{double potential}
\end{equation}
%
From \cite{doubledot} we use GaAs material parameters $m^* = 0.067m_e$, and the confinement strength $\hbar \omega_0 = 3.0$ meV. Which correspond to a harmonic oscillator length of $\sqrt{\hbar/\omega_0 m^*} \approx 5.3$ nm. And with the minima separated by a distance of $2L_x$ from each other. The values of $V_c(x,0)$ are plotted in Fig.~\ref{doubledotfigure}.

\begin{figure}[H]
\centering
\scalebox{0.7}{\includegraphics{doubledot.eps}}
\caption{Confinement potential for $V_c(x,0)$ and $L_x=50$ nm.}
\label{doubledotfigure}
\end{figure}
%
We will then solve the two-dimensional Schr\"{o}dinger equation by using a finite difference method with a three-point Laplacian \cite{computerjensen}. This will give us a Hermitian matrix for which the eigenvalues and eigenvectors can be solved using the QR-algorithm \cite{faires}. This algorithm was independently introduced in 1960s by Kublanovskaya \cite{QR2} and Francis \cite{francis}. And have been recognized as one the most important algorithms of the 21st century \cite{mostimportant}. 


\subsection{The QR-Algorithm}
\label{section:The QR-Algorithm}
The idea is surprisingly simple. First step is to factor the matrix $A$ into a product of an orthogonal matrix $Q_1$ and positive upper triangular matrix $R_1$. This also referred to as an QR-decomposition and solved by using the Gram-Schmidt algorithm \cite{linalg}.
\begin{equation}
A_1 = Q_1 R_1,
  \label{first iteration}
\end{equation}
next step is to multiply $Q$ and $R$ in a reversed order
\begin{equation}
A_2 = R_1 Q_1,
  \label{second iteration}
\end{equation}
%
we repeat this process by finding the $Q$ and $R$ values to $A_2$
\begin{equation}
A_2 = Q_2 Q_2.
\end{equation}
%
The complete algorithm can be written as
\begin{equation}
  A = Q_1R_1,\qquad R_kQ_k = A_{k+1} = Q_{k+1}R_{k+1}, \qquad k=1,2,3,...
\end{equation}
%
Where $ R_k,Q_k$ is from the previous steps, and $Q_{k+1}$ is still an orthogonal matrix ($Q^TQ=1$) and $R_{k+10}$ is a positive upper triangular.
%
\begin{algorithm}
\caption{\emph{QR-Algorithm}}
\begin{algorithmic}
\FOR{$i=1 \to n$}
\STATE $Q_iR_i = A_i$ 
\STATE $A_{i+1} = R_iQ_i$ 
\ENDFOR
\end{algorithmic}
\end{algorithm}
%
The iteration will finally create a matrix $\tilde{A}$ whose diagonal entries are eigenvalues of $A$. The reason why this works is because all the $A_k$ are similar to each other and therefore they have a common set of eigenvalues with different eigenvectors, i.e. 
\begin{proof}
\bf{If} $Ax = \lambda x$ and $\tilde{A} = S^TA S $ ($S S^T=1$) $\Rightarrow S^T A S x = \lambda x \Rightarrow  A(Sx) = \lambda (Sx)$
\label{similartransformation}
\end{proof}
%
And 
%
\begin{equation}
A_{k+1} = R_k Q_k = Q_k^T (Q_k R_k) Q_k = Q_k^T A_k Q_k .
  \label{shows the similarity transform}
\end{equation}
%
We will not go in to rigorous details for why this works, but readers are recommended to read \cite{qrproof}. Instead we will give an numerical example. Given a Hermitian matrix 
\begin{equation}
A = 		\left[\begin{array}{ccc}
			5 & 4 & 1\\
			4 & 3 & 2\\
			1 & 2 & 1\\
		\end{array}	\right],
		\label{Hermitian A}
\end{equation}
the \emph{exact} eigenvalues found by $\texttt{eig(A)}$ in \textsc{matlab} are 
%
\begin{equation}
\begin{matrix}
 \lambda_1 & = &  8.6625 \\
 \lambda_2 & = &  1.1444 \\
 \lambda_3 & = & -0.8070
\end{matrix}.
\label{eigenvalues}
\end{equation}
%
Then the initial QR-factorization $A_1 = Q_1R_1$ produces 
\begin{equation}
Q_1 = 		\left[\begin{array}{ccc}
    0.7715  & -0.0392 &  -0.6350 \\
    0.6172 & -0.1960  &  0.7620 \\
    0.1543  &  0.9798  &  0.1270
		\end{array}\right] 
				\qquad 
		R_1 = 		\left[\begin{array}{ccc}
    6.4807  &  5.2463  &  2.1602 \\
         0  &  1.2150  &  0.5487 \\
         0  &       0  &   1.0160
		\end{array}\right],
\end{equation}
which in turn gives us the new $A_2$
\begin{equation}
A_2 = R_1Q_1 = \left[\begin{array}{ccc}
    8.5714 &   0.8346  &  0.1568 \\
    0.8346 &   0.2995  &  0.9955 \\
    0.1568  &  0.9955  &  0.1290		
		\end{array}	\right].
\end{equation}
%
We continue this procedure
%
\begin{equation}
A_3 = R_2Q_2 = \left[\begin{array}{ccc}
    8.6611  &  0.1033 &   0.0168 \\
    0.1033  &  0.5392 &   0.9035 \\
    0.0168  &  0.9035 &  -0.2003		
		\end{array}	\right],
\end{equation}
%
\begin{equation}
A_4 = R_3Q_3 = \left[\begin{array}{ccc}
    8.6625  &  0.0131  &  0.0017 \\ 
    0.0131 &   0.7871  &  0.7548 \\
    0.0017  &  0.7548 &  -0.4496 
		\end{array}	\right].
\end{equation}
%
After $17$ iterations the off-diagonal elements are practically zero, and the eigenvalues on the diagonal correspond to the eigenvalues found in Eq.~(\ref{eigenvalues}). 
\begin{equation}
A_{17} = R_{16}Q_{16} = \left[\begin{array}{ccc}
    8.6625  &   0.0000  & -0.0000 \\
    0.0000  &  1.1444   & 0.0098 \\
    0.0000  &  0.0098  & -0.8069
		\end{array}	\right].
\end{equation}
%
But if we want to increase the precision, this method becomes rather slow. As we can see, the convergence of this method is not that impressive, assuming \textsc{matlab} is using the Householder's QR factorization method \cite{bible}, we are going to have $O(n^3)$ flops per iteration, in addition we have a matrix-matrix multiplication.

One of the optimalization we could do is to use Householder's method for tridiagonalization \cite{computerjensen}, i.e. we want to find a tridiagonal matrix $T$ that is a similar transformed of the matrix $A$.

\begin{equation}
  T = S^T A S, \qquad S = S_1S_2...S_{n-2} 
  \label{similar}
\end{equation}
%
This would speed up the iteration since we have fewer off-diagonal elements to worry about. A second optimazation is the QR-algorithm. We could improve the convergence by introducing a \emph{shift} on the diagonal. This is popularly called \emph{The accelerated QR-algorithm} or \emph{The shifted QR-algorithm} \cite{kincaid}

For a general tridiagonal matrix

\begin{equation}
T_m = \left[\begin{array}{cccccc}
      \alpha_1^{m}  &  \beta_1^{m}   & 0 	    & 0      & \hdots & 0  \\ 
      \beta_1^{m}   &  \alpha_2^{m}  & \beta_2^{m}  & 0 \\ 
                 &  \beta_2^{m}   & \alpha_3^{m} & \beta_3 & & \vdots  \\ 
       \vdots       &                &  \ddots      & \ddots &  \ddots  \\
           &      	     &              &  \beta_{n-2}^m & \alpha_{n-1}^m  & \beta_{n-1}^m \\ 
        0	    &      &      \hdots          &       & \beta_{n-1}^n & \alpha_n^m  \\ 
		\end{array}	\right].
\end{equation}

\begin{algorithm}
\caption{\emph{The accelerated QR-Algorithm}}
\begin{algorithmic}
\FOR{$i=1 \to n$}
\STATE $T_m - \alpha_n^m I = Q_m R_m $ 
\STATE $T_{m+1} = R_m Q_m + \alpha_{n}^m I$ 
\ENDFOR
\end{algorithmic}
\end{algorithm}
%
It was interesting to see how well this algorithm works compared to a, and we have therefore mad a plot for comparison in Fig. \ref{fig:QRsummary}

\begin{figure}[H]
\centering
\scalebox{0.7}{\includegraphics{QR.eps}}
\caption{This is a plot of different eigenvalues as a function of iteration. The blue line with diamond points represents our normal QR-algorithm on the matrix $A$ Eq.~(\ref{Hermitian A}), and the red line with points represents the accelerated QR-algorithm on a tridiagonalized matrix $A$. The improved QR-algorithm converges faster for the lower eigenvalues.}
\label{fig:QRsummary}
\end{figure}
The final issue is the numerical precision, or more correctly, the numerical imprecision. It is important to know the error when we do this type of calculations. The details of the mathematics can be read in \cite{lyche}. We will take use of the Hoffman-Wielandt Theorem which states that

\begin{theorem}
  Let $A$ and $E$ be a real symmetric $n\times n$. And let $T = A + E$ with eigenvalues $\{ \gamma_i \}$. And   $\{ \lambda_i \}$ the eigenvalues of $A$, arranged in increasing order. Then
  \begin{equation}
      \left[\sum_{j=1}^n \left(\lambda_i - \gamma_i \right) \right]^{\frac{1}{2}} \leq F(E),
  \end{equation}
 where $F(E) = \left(\sum_{ij}^n |a_{ij}|^2 \right)^{1/2}$ is the Frobenius norm of $E$
\end{theorem}
%
Let $T$ be a tridiagonal matrix and let $\tilde{T}$ be the new matrix obtained by deleting $\beta_{n-1}$ from the off-diagonal positions $(n-1,n)$ and $(n,n-1)$ of T. And let $\{\lambda_j\}$ and $\{\tilde{\lambda}_i\}$ denote the eigenvalues respectively. Then from the Wielandt-Hoffman theorem

\begin{equation}
  \left[\sum_{j=1}^n \left(\lambda_i - \tilde{\lambda}_i \right) \right]^{\frac{1}{2}} \leq F(T-\tilde{T}) = \sqrt{2}|\beta_{n-1}|,
\end{equation}
%
since $\beta_{n-1}$ is the only term that was left after the subtraction. The conclusion of this is that we are closest to the exact eigenvalues when the off-diagonals are smallest.

\subsection{Discretizing the Schr\"{o}dinger Equation}
Our one-electron Hamiltonian from Eq.~(\ref{one electron hamilton double}) reads

\begin{equation}
  \hh{H} = -\frac{\hbar^2}{2m^*} \nabla_i^2 +  \frac{1}{2} m^* \omega_0^2 \cdot \left[x^2+y^2 - 2L_x |x| + L_x^2 \right].
  \label{rerescaled Hamiltonian}
\end{equation}
Using the rescaled parameters in section \ref{section:scaling the hamiltonian} we get
\begin{equation}
  \bar{H} = -\frac{\omega_c}{2} \nabla^2 +  \frac{1}{2} \omega_c^2 \cdot V_c(\bar{x},\bar {y}) .
\end{equation}
%
The time-independent Schr\"{o}dinger equation of this Hamiltonian is then
\begin{equation}
-\frac{\omega_c}{2} \left[\frac{\partial}{\partial \bar{x}^2} +\frac{\partial}{\partial \bar{y}^2} \right]u(\bar{x},\bar{y}) +  \frac{1}{2} \omega_c^2 \cdot V_c(\bar{x},\bar{y}) u(\bar{x},\bar{y})  = E u(\bar{x},\bar{y}) 
  \label{tisl scaled hamilton}
\end{equation}
%
Where $u(\bar{x},\bar{y})$ is the single-electrons eigenfunctions. From here on the dimensionless coordinates $\bar{x}$ and $\bar{y}$ would be referred to as $x$ and $y$. And we would set the constant $\omega_c = 1$. The differential equation can be solved as a matrix diagonalization problem. By subtracting two Taylor series we get the numerical second derivate \cite{computerjensen}

\begin{align}
  u(x+h) + u(x-h) &= 2u(x) + h^2 f''(x) + O(h^4), \\
  u(y+h) + u(y-h) &= 2u(y) + h^2 f''(y) + O(h^4), 
\end{align}

\begin{equation}
 f''(x) = \frac{u(x+h)-2u(x)+u(x-h)}{h^2} + O(h^2), \\
 f''(y) = \frac{u(y+h)-2u(y)+u(y-h)}{h^2} + O(h^2).
 \label{discretize}
\end{equation}
%
For the two-dimensional case the Laplacian becomes
\begin{equation}
  \nabla^2  \approx \frac{1}{h_x^2} \left(u_{i-1,j} - 2u_{i,j} + u_{i+1,j} \right) +  \frac{1}{h_y^2} \left(u_{i,j-1} - 2u_{i,j} + u_{i,j+1} \right) + O(h^2).
  \label{poisson insert}
\end{equation}
%
Where we have used the more compact way of writing $u(x\pm h,y) = u_{i \pm 1, j}$, $u(x,y\pm h) = u_{i, j\pm1}$. An example could be a grid with integration points $n_x = 3$, and step $h=h_x=h_y=\frac{x_{\max}-x_{\min}}{n_x} = \frac{1}{3}$, $x=x_{\min}+ih$, $y=y_{\min}+jh$, where $(i,j)\in\{1,2,3\}$. See Figure \ref{unitgrid}
%
\begin{figure}[H]
\centering
\input{unitgrid}
\caption{The discretized unit grid in our example. The indices are arranged in lexical order}
\label{unitgrid}
\end{figure}
%
Inserting Eq.~(\ref{poisson insert}) into the Schr\"{o}dinger equation Eq.~(\ref{tisl scaled hamilton}), gives
\begin{equation}
  d_{ij} u_{i,j} - e u_{i,j-1} - e u_{i,j+1} - e u_{i-1,j} - e u_{i+1,j} = \lambda u_{i,j},
\end{equation}
%
where we have defined
\begin{equation}
  d_{ij} = \frac{4}{h^2} + V_{ij} \qquad e = -\frac{1}{h^2}.
\end{equation}
%
The left side is the discretized Poisson equation in two dimensions. This gives us a set of linear equations that can be written as a matrix eigenvalue problem.

\begin{equation}
 \left[\begin{array}{ccccccccc}
    -4   &  1  &    0  &   1  &      &      &      &      &     \\
     1   & -4  &    1  &   0  &   1  &      &      &      &     \\
     0   &  1  &   -4  &   0  &   0  &   1  &      &      &     \\
     1   &  0  &    0  &  -4  &   1  &   0  &   1  &      &     \\
         &  1  &    0  &   1  &  -4  &   1  &   0  &   1  &     \\
         &     &    1  &   0  &   1  &  -4  &   0  &   0  &   1 \\
         &     &       &   1  &   0  &   0  &  -4  &   1  &   0 \\
         &     &       &      &   1  &   0  &   1  &  -4  &   1 \\
         &     &       &      &      &   1  &   0  &   1  &  -4
		\end{array}	\right]
    \left[\begin{array}{c}
     u_{1,1}    \\
     u_{2,1}    \\
     u_{3,1}    \\
     u_{1,2}    \\
     u_{2,2}    \\
     u_{3,2}    \\
     u_{1,3}    \\
     u_{2,3}    \\
     u_{3,3}   
	\end{array}\right]		
	=
	\lambda
	  \left[\begin{array}{c}
     u_{1,1}    \\
     u_{2,1}    \\
     u_{3,1}    \\
     u_{1,2}    \\
     u_{2,2}    \\
     u_{3,2}    \\
     u_{1,3}    \\
     u_{2,3}    \\
     u_{3,3}   
	\end{array}\right].
\end{equation}
%
This is a sparse symmetric matrix which is almost tridiagonal. In reality we would choose larger integration points and the number of entries in the matrix grows with $n^4$. A lot of those entries are zero.

Since the $x$-coordinates are separable, i.e.

\begin{align}
  \hh{H}(x,y) &= \hh{H}_X(x) + \hh{H}_Y(y) \\ \Rightarrow E &= E_X + E_Y,
\end{align}
%
we do not need to discretize in two dimensions, since we already know the eigenvalues for $\hh{H}_Y$, they are the same as for the one-dimensional case, see section \ref{onebodyharmonicoscillatorpotential}. 
The Laplacian in one dimension is Eq.~(\ref{discretize})

\begin{equation}
  \frac{d^2}{dx^2} = \frac{u_{i-1} - 2u_i + u_{i+1} }{h^2} + O(h^2),
  \label{onedimdiscretized}
\end{equation}
%
where $h$ is our step defined as 
\begin{equation}
h = \frac{x_{max}-x_{max}}{N},
  \label{step}
\end{equation}
%
where $N$ is the number of steps or gridpoints. Inserting Eq.~(\ref{onedimdiscretized}) into the one-dimensional Schr\"{o}dinger equation
%
\begin{equation}
-\frac{u_{i+1} - 2u_{i} + u_{i-1}}{h^2} + V_i u_i = E_X u_i.
\end{equation}
%
This can be written as a tridiagonal matrix eigenvalue equation.
%
\begin{equation}
  d_iu_i + e_{i-1}u_{i-1} + e_{i+1}u_{i+1} = \lambda u_i,
  \label{onedim discrete}
\end{equation}
%
where
%
\begin{align}
  e_i  = -\frac{1}{h^2}, \qquad  d_i  = \frac{1}{h^2} + V_i,
\end{align}

\begin{equation}
 \left[\begin{array}{ccccccc}
    d_1   &  e_1   &  0  &   0  &  \cdots  &   0  &   0\\
    e_1   & d_2  &  e_2  &   0  &  \cdots  &   0   &  0 \\
     0    & e_2  &  d_3  &  e_3  &  \cdots  &   0   &  0 \\
 \vdots   & \vdots  &  \ddots  &   \ddots  &  \ddots  &   0   &  0 \\
     0    & 0  &  0  &   0  &  \cdots  &   d_{N-2}   &  e_{N-1}  \\		
     0    & 0  &  0  &   0  &  \cdots  &   e_{N-1}   &  d_N	
		\end{array}	\right]
    \left[\begin{array}{c}
     u_{1}    \\
     u_{2}    \\
     u_{3}    \\
    \vdots    \\	
     u_{N-1}   
	\end{array}\right]		
	=
	E_X
	  \left[\begin{array}{c}
    u_{1}    \\
     u_{2}    \\
     u_{3}    \\
    \vdots    \\	
     u_{N-1}   
	\end{array}\right].
\label{tridiagonal}
\end{equation}

%%%%%%%%%%%%%%%%%%%%%%%%%%%%%%%%%%%%%%%%%%%%%%%%%%%%%%%%%%%%%%%%%%%%%%%%%%%%%%%%%%%%%%%%
%%%%%%%%%%%%%%%%%%%%%%%%%%%%%%%  QM FOR MANY-BODY SYSTEMS       %%%%%%%%%%%%%%%%%%%%%%%%
%%%%%%%%%%%%%%%%%%%%%%%%%%%%%%%%%%%%%%%%%%%%%%%%%%%%%%%%%%%%%%%%%%%%%%%%%%%%%%%%%%%%%%%%
% \chapter{Quantum Mechanics for Many-Body Systems}
In this chapter, we introduce the notation that is common in many-body physics. We will also introduce the formalism of second-quantization and Wick's theorem.

\section{Introduction}
The underlying fundamentals of quantum mechanics as we know today have not changed much since its birth. We still need to know the Hamiltonian and solve the Schr\"{o}dinger equation. But how we use this theory is a different story. In the real world we have to deal with systems of more than one quantum particle, we therefore need to expand our Schr\"{o}dinger equation to include those particles, and this we call the many-body Schr\"{o}dinger equation. The degrees of freedom increase as our system size gets bigger. We are not able to solve our manybody Schr\"{o}dinger equation with conventional techniques, not analytically nor numerically. Another limiting factor is our knowlegde of the interactions, as the exact Hamiltonian is not known.

Therefore we have to make some assumptions and use approximations, and here is where the \emph{many-body methods} comes in. The first approximation we use is to the Hartree-Fock method on closed shells, assuming that all our electrons interact in the same way. The next method we are going to use is the coupled cluster (CC) method. They are all different in the sense that they have their regions of effectiveness. HF is very fast and gives reasonable result compared to  when we have small systems for very closed shells and the ground state is stable. The CC method, on the other hand, has applications for systems up to 40 electrons, but has problems with convergence and non-variational energies.


%26.01.2011
%Many of manybody methods we are called \emph{ab initio} methods. That means the calculation start from some fundamental interaction. These methods are not based on any empirical assumptions or parametrizations %[ask Morten] What about when we fit trial wave functions to STO's and parametrize the trial wave function. The whole HF is based on educated guess? empirical guess? some knowelgde about empiri gives us the guess? and the Electron Correlation energy Diagram to show how good HF is compared to other methods see page 32 in Patrick.
\section{The Many-Body Problem}
\label{sec:many-body problem}
Let us assume that we have a non-relativistic isolated system of $N$ particles. And assume we can describe the system with a time-independent Hamilton operator $\hat{H}$, then we could reduce the problem to solve the time-independent Schr\"{o}dinger equation:

\begin{equation}
 \hat{H}(r_1, r_2,...,r_N)\psi_\lambda(r_1,r_2,...,r_N) = E_\lambda \psi_\lambda(r_1, r_2,...,r_N).
\label{def:manybodyschrodinger}
\end{equation}
%
where the $r_i$ represents particle $i$ with spin $\ket{m_s}$. $\lambda$ denotes the set of quantum numbers for particles 1,...N. 

The many-body wavefunction $\Psi_\lambda$ is a N-body \emph{vector} in the composite Hilbert space:

\begin{equation}
\psi_\lambda \in \mathcal{H}_N := \mathcal{H}_1 \oplus \mathcal{H}_1 \oplus \ldots \oplus \mathcal{H}_1,
  \label{maybodyinhilbert}
\end{equation}
%
or
\begin{equation}
\ket{\Psi_\lambda} = \ket{\psi_1}  \oplus \ket{\psi_2} \oplus ... \oplus \ket{\psi_N} \equiv \ket{\psi_1\psi_2 ... \psi_N},
 \label{def:productstates}
\end{equation}
%
where $\ket{\psi_i}$ is a state in a single-particle Hilbert space $\mathcal{H}_1$, which is the space of square integrable function over  
$\mathbb{R}^d \oplus (\sigma)$, or formally:

\begin{equation}
\mathcal{H}_1:= L^2(\mathbb{R}^d \oplus (\sigma)).
\label{hilbertspacedim}
\end{equation}
%
\subsection{The Electronic Hamiltonian}
We want to describe our physical system with an \emph{ab initio} method, which basically means that we want our Hamiltonian to include the basic forces with no parametrization, in atomic units ($\hbar = c = m_e = 1$) the Hamiltonian is:

\begin{equation}
\hat{H}  = \hat{T} + \hat{V}.
 \label{eq:tplussv}
\end{equation}
%
Where $\hat{T}$ is the total kinetic energy operator, and $\hat{V}$ is the total potential energy operator.
 
\begin{equation}
\hat{T} = \sum_{k} \hat{t}_k.
 \label{eq:kinetic}
\end{equation}
%
$\hat{t}_k$ is the kinetic energy operator for particle $k$. 

In general we have:

\begin{equation}
 \hat{V} = \hat{V}_1 + \hat{V}_2 + ...  
 \label{eq:potential}
\end{equation}
 %
where 

\begin{equation}
 \hat{V}_n = \frac{1}{n!} \sum_{abc..z} \hat{v}^{(n)}_{abc...z}
\end{equation}
$n!$ is because we have indistinguishable particles. For systems of electrons like quantum dots, we will truncate our total potential operator to include up to the two-body potential operator $\hat{V}_2$. But some papers in nuclear physics (ref(papers of three body force)) have proven that the three-body force is a important contributor to the binding energy. Then the electronic Hamiltonian the reads (in atomic units)

\begin{equation}
\hat{H} = \sum_{k} \hat{t}_k + \frac{1}{2} \sum_{ij} \hat{v}_{ij},
 \label{def:electronic hamiltonian}
\end{equation}
%
where (in atomic units)

\begin{equation}
\hat{h}_k = -\frac{1}{2}\nabla^2_k + \sum_{A} \hat{c}_{kA}, \qquad \hat{v}_{ij} = \frac{1}{r_{ij}}.
 \label{def:electronic hamiltonian parts}
\end{equation}
%
Where the last sum is the Coulomb contribution of interaction of the single-particles with the \emph{core} particles $A$, e.g. electrons around a proton core. 


\subsection{Identical Particles}
In quantum mechanics particles are indistinguishable, and thus we can not tell which of the electrons are in which state. This means that the expectation value would have to be the same when we interchange the coordinates of particle $i$ and $j$. 

\begin{equation}
|\Psi_{\lambda}(r_1,r_2,..,r_i,..,r_j,..,r_N)|^2 = |\Psi_{\lambda}(r_1,r_2,..,r_j,..,r_i,..,r_N)|^2,
 \label{eq:interchange}
\end{equation}
%
gives us possible antisymmetric ($-$) and symmetric ($+$) solutions

\begin{equation}
\Psi_{\lambda}(r_1,r_2,..,r_i,..,r_j,..,r_N) = \pm \Psi_{\lambda}(r_1,r_2,..,r_j,..,r_i,..,r_N).
\end{equation}
%
We will later refer to the symmetric solution as bosons, and the other as fermions.
Introducing the Permutation operator 
%
\begin{equation}
\hat{P}_{ij}\Psi_\lambda(r_1,r_2,..,r_i,..,r_j,..,r_N) = \Psi_\lambda(r_1,r_2,..,r_j,..,r_i,..,r_N) 
 \label{eq:poperator}
\end{equation}
%
The eigenvalue equation for $\hat{P}$ gives 

\begin{equation}
 \hat{P}_{ij}\Psi_\lambda(r_1,r_2,..,r_i,..,r_j,..,r_N) = \beta \Psi_\lambda(r_1,r_2,..,r_i,..,r_j,..,r_N),
 \label{eq:poperatoreq}
\end{equation}

\begin{equation}
\beta = \pm 1 \qquad \text{Since} \qquad \hat{P}^2_{ij} = 1.
 \label{eq:beta} 
\end{equation}
%
wavefunctions with $(\beta = +1)$ are the bosons, and $(\beta = -1)$ are the fermions. The Hamiltonian is invariant under the interchange of particles and therefore commutes with the permutation operator.

\begin{proof}
\begin{equation}
\hat{P}_ {jk} \hat{H} \Psi = \hat{P}_{jk}(\hat{H}_{jk}\Psi_{jk}) =  \hat{H}_{kj}\Psi_{kj} 
\label{eq:overst}
\end{equation}
%
\begin{equation}
\hat{H} \hat{P}_{jk} \Psi = \hat{H}_{jk}\hat{P}_{jk}\Psi_{jk} = \hat{H}_{jk} \Psi_{kj}
\label{eq:nederst}
\end{equation}
%
We then subtract equation (\ref{eq:overst}) with (\ref{eq:nederst}) 
%
\begin{equation}
 \hat{P}_{jk}(\hat{H}\Psi) - \hat{H}(\hat{P}_{jk}\Psi) = \hat{H}_{kj}\Psi_{kj} - \hat{H}_{jk} \Psi_{kj}
\end{equation}
%
or equivalently
%
\begin{equation}
 \left[\hat{P}_{jk},\hat{H} \right] = \left(\hat{H}_{kj} - \hat{H}_{jk}\right)\Psi_{kj}
\end{equation}
\end{proof}
%
\noindent The permutation operator commutes with the Hamiltonian if and only if $\hat{H}_{kj} = \hat{H}_{jk}$.

According to (\ref{eq:kinetic}) and (\ref{eq:potential}) Hamiltonian is a sum of of onebody and two-body operators. The sums converge and are therefore interchangeable with respect to particles without changing our Hamiltonian.

\begin{equation}
H_{jk} = H_1... + H_j + ... + H_k + ... = H_1... + H_k + ... + H_j + ... = H_{kj}.
 \label{eq:HkjequalsHjk}
\end{equation}
%
Then it follows that $\hat{H}$ and $\hat{P}$ are compatible observable (\cite{griffiths}). i.e. there exists eigenfunctions for $\hat{H}$ that are also eigenfunctions of $\hat{P}$. We know we need to construct symmetric wavefunctions for the bosons and antisymmetric for the fermions. One way of doing this is using a symmetrization operator for bosons

\begin{equation}
\hat{S} = \frac{1}{N!} \sum_p \hat{P}
 \label{def:symmetrizer}
\end{equation}
%
where $p$ is the called the permutation number and is the \emph{set} of all possible permutations including the empty set. $p = \{\emptyset,[1,2],[1,3],[2,3]\}$  for three particles.
The normalized symmetric state $\Phi_S$ is then given by

\begin{equation}
\Phi_S(r_1,r_2,...,r_N) = \sqrt{\frac{N!}{n_\alpha!n_\beta!...n_\gamma!}} \hat{S} \phi_\alpha(r_1)\phi_\beta(r_2)...\phi_\gamma(r_N)
 \label{def:symmetricwave}
\end{equation}
%
Similarly we have the anti-symmetrization operator for fermions

\begin{equation}
\hat{A} = \frac{1}{N!} \sum_p (-1)^p \hat{P},
 \label{def:antisymmetrizer}
\end{equation}
%
and the normalized antisymmetric states

\begin{equation}
\Phi_{AS}(r_1,r_2,...,r_N) = \sqrt{N!} \hat{A} \psi_\alpha(r_1)\psi_\beta(r_2)...\psi_\gamma(r_N).
 \label{def:antisymmetricwave}
\end{equation}
%
or equivalently

\begin{equation}
  \Phi_{\alpha\beta...\gamma}(r_1,r_2,...,r_N) = \frac{1}{\sqrt N!} 
\left|\begin{array}{cccc} 
\phi_\alpha(r_1) & \phi_\beta(r_1) & ... & \phi_\gamma(r_1) \\
\phi_\alpha(r_2) & \phi_\beta(r_2) & ... & \phi_\gamma(r_2) \\
\vdots & \vdots & \vdots & \vdots \\
\phi_\alpha(r_N) & \phi_\beta(r_N) & ... & \phi_\gamma(r_N) \\
\end{array} \right|.
\label{def:slaterdeterminant}
\end{equation}
%
This was first introduced by J.C. Slater \cite{slaterdeterminant} in 1929 and is popularly called a Slater determinant. It obeys the Pauli Exclusion Principle (PEP): The determinant would be zero if two of the single-particle wavefunctions have the same quantum numbers $\alpha, \beta, ... \gamma$. 


The most general way of writing our wavefunction of the N-fermion system is to have a linear combination of the Slater determinants (\ref{def:slaterdeterminant}).

\begin{equation}
\Psi_\lambda(r_1,r_2,....,r_N) = \sum_{\alpha\beta...\gamma} C_{\alpha\beta...\gamma}^{\lambda}  \Phi_{\alpha\beta...\gamma}(r_1,r_2,...,r_N)
 \label{def:generalslaterdeterminant}
\end{equation}

\section{Second Quantization}
\label{sec:second quantization}
The second-quantization formalism was first introduced by Dirac (1927) and extended to fermion systems by Jordan and Klein (1927) and by Jordan and Wigner (1928) \cite{bartlett}. The formalism of \emph{second quantization} is just a simplification in the description of a many-body system, a reformulation of the original Schr\"{o}dinger equation. The quantum mechanical states are represented by annihilation and creation operators working on the physical vacuum state. 

We will look at fermionic systems, therefore we will restrict the many-particle functions to be antisymmetric and choose the \emph{Slater determinant} (\ref{def:slaterdeterminant}) as our candidate. And introduce the occupancy notation for Slater determinants
%
\begin{equation}
\Phi_{\alpha_1 \alpha_2...\alpha_N} \equiv \ket{\alpha_1\alpha_2...\alpha_N}.
 \label{def:state}
\end{equation}
%
Note: this is not the same as the product states in (\ref{def:productstates}). This is antisymmetrized 

\begin{equation}
\ket{\alpha_1..\alpha_i\alpha_j..\alpha_N} = -\ket{\alpha_1..\alpha_j\alpha_i..\alpha_N}.
 \label{def:antisymmetricstate}
\end{equation}
%
And the state ''lies`` in what we called the Fock space,  which is a tensor product space of antisymmetric Hilbert spaces:

\begin{equation}
\mathcal{F}_N = \bigoplus_{n=0}^N \mathcal{H}_n^{AS}.
 \label{def:fockspace}
\end{equation}
%
In this case we have a state in an $N$-dimentional Fock space. 

\subsection{Creation and Annihilation Operators}
The creation and annihilation operators are mappings between different $N$ and $N\pm1$ dimentional Hilbert spaces,

\begin{equation}
a_{\alpha}^{\dagger}: \mathcal{H}_{N}^{AS} \rightarrow \mathcal{H}_{N+1}^{AS},
 \label{def:creation}
\end{equation}

\begin{equation}
a_{\alpha}: \mathcal{H}_N^{AS} \rightarrow \mathcal{H}_{N-1}^{AS},
 \label{def:annihilation}
\end{equation}
%
where 

\begin{equation}
 \alpha \in \mathcal{H}_1.
\end{equation}
%
A creation operator $a_{\alpha}^{\dagger}$ will create a fermion with quantum number(s) $\alpha$ from the antisymmetric state (\ref{def:state})
 

\begin{equation}
a_{\alpha}^{\dagger} \ket{0} = \ket{\alpha}.
 \label{ex:creation}
\end{equation}
%
$\ket{0}$ is the vacuum state. If $\alpha$ is already occupied, the result is zero due to PEP. 

\begin{equation}
a_{\alpha}^{\dagger} \ket{\alpha} = 0.
 \label{ex:creationpep}
\end{equation}
%
An annihilation operator $a_{\alpha}$ will remove a fermion with quantum number(s) $\alpha$ from the antisymmetric state (\ref{def:state})

\begin{equation}
a_{\alpha}\ket{\alpha} = \ket{0},
 \label{ex:annihilation}
\end{equation}

\begin{equation}
a_{\alpha}\ket{0} = 0.
 \label{ex:annihilationvacuum}
\end{equation}
%
If $\alpha$ does not exist, the result is zero due to annihilation of a vacuum state. 

\begin{equation}
a_{\alpha} \underbrace{\ket{\alpha_1\alpha_2...\alpha_N}}_{\alpha \notin} = 0.
 \label{ex:annihilationpep}
\end{equation}
%
Our Slater determinant (\ref{def:state}) can now be written as a product of creation operators

\begin{equation}
\ket{\alpha_1\alpha_2...\alpha_N} = \prod_{i=1}^N a_{\alpha_i}^\dagger \ket{0}.
 \label{eq:slatercreationproduct}
\end{equation}
%
Using the antisymmetry of the states (\ref{def:antisymmetricstate}) we can show that

\begin{equation}
a_{\alpha_i}^\dagger a_{\alpha_k}^\dagger = - a_{\alpha_k}^\dagger a_{\alpha_i}^\dagger,
 \label{eq:creationcommute}
\end{equation}
%
leading to the anticommutation rule for creation operators

\begin{equation}
\{a_{\alpha}^\dagger ,a_{\beta}^\dagger \} = a_{\alpha}^\dagger a_{\beta}^\dagger  + a_{\beta}^\dagger a_{\alpha}^\dagger = 0.
 \label{eq:creationanticommute}
\end{equation}
%
Note: If $\alpha = \beta$ we would also get zero because of PEP. The hermitian conjugate (adjoint) of $a_{\alpha}^\dagger$ is the annihilation operator,

\begin{equation}
\left(a_{\alpha}^\dagger \right)^\dagger = a_\alpha
 \label{def:creationadjoint]}
\end{equation}
%
We have the following anticommutation relation for the annihilation operators (see \cite{bartlett} for details)

\begin{equation}
\{a_\alpha, a_\beta \} = a_\alpha a_\beta + a_\beta a_\alpha = 0
 \label{eq:annihilationanticommute}
\end{equation}
%
and 

\begin{equation}
\{a_{\alpha}^\dagger, a_\beta \} = \{a_{\alpha}, a_\beta^\dagger \}  = \delta_{\alpha \beta}
 \label{eq:creationannihilationanticommute}
\end{equation}
%
where $\delta_{\alpha \beta}$ is $0$ if $\alpha \neq \beta$ and $1$ if $\alpha = \beta$.

\subsection{Representation of Operators}
Now that we have a formalism for our states, we want to calculate matrix elements and expectation values of our \emph{many-body} operators. Starting with the number-operator. It is a way to test that our many-body formalism conserves the particle number. 

\begin{equation}
\hat{N} = \sum_\alpha a_{\alpha}^\dagger a_{\alpha},
 \label{def:numberoperator}
\end{equation}
%
and operating this on a state gives us the eigenvalue of $n$, which is the number of fermions in that state.

\begin{equation}
\hat{N} \ket{\alpha_1 \alpha_2... \alpha_N} = \sum_\alpha a_{\alpha}^\dagger a_{\alpha} \ket{\alpha_1 \alpha_2... \alpha_N} = n  \ket{\alpha_1 \alpha_2... \alpha_N} 
 \label{def:numberoperatorstate}
\end{equation}
%
Because from  (\ref{ex:creation}),(\ref{ex:annihilationpep}) and (\ref{ex:annihilationpep}) we get 

\begin{equation}
 a_{\alpha}^\dagger a_{\alpha} \ket{\alpha_1 \alpha_2... \alpha_N}= \left\{ \begin{array}{ll}
 0 & \mbox{$\alpha \notin \{\alpha_i\}$} \\
  \ket{\alpha_1 \alpha_2... \alpha_N} & \mbox{$\alpha \in \{\alpha_i\}$}
       \end{array} \right..
 \label{def:becausenumberoperators}
\end{equation}
% 
The number operator is a one-body operator since it acts on one single-particle state at a time. Another type is the one-body operator \cite{bartlett}.

\begin{equation}
\hat{F} = \sum_{\alpha \beta} \bra{\alpha} \hat{f} \ket{\beta} \ket{\alpha}\bra{\beta}
 \label{def:onebodyoperator}
\end{equation}
%
where $\ket{\alpha},\ket{\beta}$ is the chosen single-particle basis. It can be rewritten and expressed with creation and annihilation operators. The second quantisized form of $\hat{F}$

\begin{equation}
\hat{F} = \sum_{\alpha \beta} \bra{\alpha} \hat{f} \ket{\beta} a_{\alpha}^\dagger a_{\beta}
 \label{def:rewrittenonebodyoperator}
\end{equation}


%One spesial case is the one-body operator for the kinetic energy plus an external one-body %potential.

%\begin{equation}
%\hat{H}_0 = \sum_{\alpha \beta} \bra{\alpha} \hat{h}_0 \ket{\beta} a_{\alpha}^\dagger a_{\beta}
% \label{def:spesialonebodyoperator}
%\end{equation}

%where $\hat{h}_0 = \hat{t} + \hat{u}$. 

The operator $\hat{F}$ removes a fermion from the state $\beta$ and creates a new one in state $\alpha$. This transition is given by the probability amplitude $\bra{\alpha}\hat{f}\ket{\beta}$.

Generally we can do this for an N-body operator. But we will only consider a two-body operator:

\begin{equation}
\hat{V} = \sum_{\alpha \beta \gamma \delta} \bra{\alpha \beta} v \ket{\gamma \delta}   \ket{\gamma \delta} \bra{\alpha \beta}.
 \label{def:twobodyoperator}
\end{equation}
%
For an $N$-particle system we have

\begin{equation}
V_N = \sum_{i<j=1}^N \hat{v}_{ij} = \frac{1}{2} \sum_{i\neq j}^N \hat{v}_{ij}.
 \label{def:N-twobodyoperator}
\end{equation} 
%
This can be used to rewrite $\hat{V}$ to

\begin{align}
\hat{V} &= \frac{1}{2} \sum_{\alpha\beta\gamma\delta}\bra{\alpha \beta} v \ket{\gamma \delta} a_{\alpha}^\dagger a_{\beta}^\dagger a_{\delta}, a_{\gamma} \\
        &= \frac{1}{4} \sum_{\alpha\beta\gamma\delta} \bra{\alpha \beta}|v|\ket{\gamma \delta} \, a_{\alpha}^\dagger a_{\beta}^\dagger a_{\delta}, a_{\gamma} 
 \label{eq:secondquantized}
\end{align}
%
where we have defined the antisymmetric matrix element to be
\begin{align}
\bra{\alpha \beta} |v| \ket{\gamma \delta} = \bra{\alpha \beta} v \ket{\gamma \delta} - \bra{\alpha \beta} v \ket{\delta\gamma}.
 \label{def:antisymmetrizedmatrix}
\end{align}
%Note that if alpha = beta, gamma, delta we would automatically get 0, no need to have the condition alpha != beta like in the def:N-twobodyoperator. That is restriction of the summation index are not needed because of PEP
see \cite{manybodyjensen} and \cite{dick} for details of this derivation. The interpretation of the operator $\hat{V}$ is that it removes two fermions in the states $\gamma$ and $\delta$, and creates two others in states $\alpha$,$\beta$. This is done with probability amplitude  $\frac{1}{4}\bra{\alpha \beta} v \ket{\gamma \delta}_{\text{AS}}$. 
%
But what is interesting here is to calculate expectation values of the operator (\ref{eq:secondquantized}). Let us find the expectation value of $\hat{V}$ with respect to the two-particle product states $\ket{\alpha_1 \alpha_2}$ and $\ket{\beta_1}{\beta_2}$

\begin{align}
\bra{\alpha_1\alpha_2}\hat{V}\ket{\beta_1\beta_2} &= \frac{1}{4} \sum_{\alpha\beta\gamma\delta} 
 \bra{\alpha \beta} |v| \ket{\gamma \delta}\bra{\alpha_1 \alpha_2} a_{\alpha}^\dagger a_{\beta}^\dagger a_{\delta} a_{\gamma} \ket{\beta_1\beta_2}, \\ 
  & = \frac{1}{4} \sum_{\alpha\beta\gamma\delta} \bra{\alpha \beta} |v| \ket{\gamma \delta}\bra{0} a_{\alpha_1} a_{\alpha_2} a_{\alpha}^\dagger a_{\beta}^\dagger a_{\delta} a_{\gamma} a_{\beta_1}^\dagger a_{\beta_2}^\dagger \ket{0}.
\end{align}
%Derivate the expectation value
Using the anticommutation relations (\ref{eq:creationanticommute}),(\ref{eq:annihilationanticommute}) and (\ref{eq:creationannihilationanticommute}), we get

\begin{align}
 &\bra{0} a_{\alpha_1} a_{\alpha_2} a_{\alpha}^\dagger a_{\beta}^\dagger a_{\delta} a_{\gamma} a_{\beta_1}^\dagger a_{\beta_2}^\dagger \ket{0} \nonumber \\
&= \bra{0} a_{\alpha_1} a_{\alpha_2} a_{\alpha}^\dagger a_{\beta}^\dagger \left( a_{\delta} \delta_{\gamma\beta_1} a_{\beta_2}^{\dagger} -  a_\delta a_{\beta_1}^\dagger a_\gamma a_{\beta_2}^\dagger \right) \ket{0} \\ 
&= \bra{0} a_{\alpha_1} a_{\alpha_2} a_{\alpha}^\dagger a_{\beta}^\dagger \left( \delta_{\gamma\beta_1}\delta_{\delta \beta_2} - \delta_{\gamma \beta_1} a_{\beta_2}^{\dagger} a_{\delta} - a_{\delta} a_{\beta_1}^{\dagger} \delta_{\gamma \beta_2} + a_\delta a_{\beta_1}^{\dagger}a_{\beta_2}^{\dagger} a_\gamma \right) \ket{0}\\
&=\bra{0} a_{\alpha_1} a_{\alpha_2} a_{\alpha}^\dagger a_{\beta}^\dagger \left(\delta_{\gamma\beta_1}\delta_{\delta \beta_2} - \delta_{\gamma \beta_1} a_{\beta_2}^{\dagger} a_{\delta} - \delta_{\delta\beta_1}\delta_{\gamma\beta_2} \delta_{\gamma \beta_2} + \delta_{\gamma\beta_2}a_{\beta_1}^{\dagger}a_\delta + a_\delta a_{\beta_1}^{\dagger}a_{\beta_2}^{\dagger} a_\gamma \right) \ket{0}
\end{align}
% 
The only terms that survive are the terms with only kronecker deltas, because all the other terms have an annihilation operator to the left, which yields zero with vacuum state, (\ref{ex:annihilationvacuum}).

\begin{equation}
\bra{0} a_{\alpha_1} a_{\alpha_2} a_{\alpha}^\dagger a_{\beta}^\dagger a_{\delta} a_{\gamma} a_{\beta_1}^\dagger a_{\beta_2}^\dagger \ket{0} = \left(\delta_{\gamma\beta_1}\delta_{\delta\beta_2} - \delta_{\delta\beta_1}\delta_{\gamma\beta_2}\right)\bra{0} a_{\alpha_2} a_{\alpha_1} a^{\dagger}_{\alpha} a^{\dagger}_{\beta} \ket{0}
\label{eq:vacuumexpectation}
\end{equation}
%
Similariy we can rewrite 

\begin{equation}
\bra{0} a_{\alpha_2}a_{\alpha_1}a_{\alpha}^{\dagger}a_{\beta}^{\dagger} \ket{0} = \delta_{\alpha \alpha_1}\delta_{\beta \alpha_2} - \delta_{\beta \alpha_1}\delta_{\alpha \alpha_2}
 \label{eq:vacuumexpectation2}
\end{equation}
%
This gives us the following expectation value

\begin{align}
\bra{\alpha_1\alpha_2}\hat{V}\ket{\beta_1\beta_2} &= \frac{1}{2} \left[ \bra{\alpha_1\alpha_2}v\ket{\beta_1\beta_2} - \bra{\alpha_1\alpha_2}v\ket{\beta_2\beta_1} - \bra{\alpha_2\alpha_1}v\ket{\beta_1\beta_2} + \bra{\alpha_2\alpha_1}v\ket{\beta_2\beta_1} \right]\\
 &= \bra{\alpha_1\alpha_2}v\ket{\beta_1\beta_2} - \bra{\alpha_1\alpha_2}v\ket{\beta_2\beta_1}\\
 &= \bra{\alpha_1\alpha_2}v\ket{\beta_1\beta_2}_{\text{AS}} 
 \label{res:twobodyexpectation}
\end{align}
%
As we see, this can be very tedious and inefficient as we have to write out contributions that gives us zero. But we can use Wick's theorem to more easily find those terms that give us contribution. This will be our next topic. The second-quantized form of electronic Hamiltonian from (\ref{def:electronic hamiltonian},\ref{def:electronic hamiltonian parts}) is then %badness underfull 

\begin{equation}
\hat{H} = \sum_{ij} \bra{i}\hat{h}\ket{j}a_i^{\dagger} a_j+ \frac{1}{4} \sum_{ijkl} \bra{ij}|\hat{v}|\ket{kl} a_i^{\dagger} a_j^{\dagger} a_l a_k
 \label{def:secondquantized hamiltonian}
\end{equation}
%
And its vacuum expectation value:
\begin{equation}
\bra{0}\hat{H}\ket{0} = \sum_{i} \bra{i}h\ket{i} + \frac{1}{4} \sum_{ij} \bra{ij}|v|\ket{ij}
 \label{eq:second-qunatized vacuum expectation value}
\end{equation}


\subsection{Wick's Theorem}
Originally Gian-Carlo Wick established this method (1950) in order to evalute the $S$-matrix in quantum field theory (see for \cite{wick} and \cite{chang} for details). He introduced two concepts: \emph{normal ordering} and \emph{contractions}. Normal ordering is just a way to write products of annihilation and creation operators in a systematic manner. 

The operators $\hat{A},\hat{B},\hat{C},...$ represents both creation and annihilation operators. Then the \emph{normal ordering} of the operators $\{\hat{A}\hat{B}\hat{C}...\}$ are defined as the rearrangement such that all of the annihilation operators are to the left of the creation operators, multiplied with a phase factor which is $-1$ for each permutation of the nearest neighboor operators.

\begin{equation}
\left \{ \hat{A}\hat{B}...\hat{U}\hat{V} \right \} \equiv (-1)^{p} u^{\dagger}v^{\dagger}w^{\dagger}...cba 
 \label{def:normalorder}
\end{equation}
%
The superscript $p$ denotes the number of permutations needed to bring the original operator product into the normal ordered form. 
\\   %How to get Example \newline without badness???
Example:
\begin{align}
\{a^{\dagger}b\} &= a^{\dagger}b, \qquad \{ab^{\dagger}\} = -b^{\dagger}a, \nonumber \\
\{ab\} &= ab = -ba, \\
\{a^{\dagger}bc^{\dagger}d\} & = a^{\dagger} c^{\dagger} db = c^{\dagger}a^{\dagger} bd = -a^{\dagger} c^{\dagger}bd =  -c^{\dagger} a^{\dagger}db \nonumber
  \label{ex:normalorder}
\end{align}
Note that the normal ordered form of operators is not unique since creation and annihilation operators can permute among themselves. Also note that one of the important properties of normal ordered operators is that its vacuum expectation value is zero. 
\begin{equation}
\bra{0}\ql\{\hh{A}\hh{B}...\qr\}\ket{0} = 0
 \label{show:vacuumexpectnormalorder}
\end{equation}
%
Because of 
\begin{equation}
a_\alpha\ket{0} = 0
\label{eq:becauseannihilation} 
\end{equation}
\begin{equation}
\bra{0}a^\dagger_\alpha = 0
\label{eq:becausecreation} 
\end{equation}
%
A contraction between two operators is defined as
\begin{equation}
\contraction{}{\hh{A} }{}{\hh{B}}
\hh{A}\hh{B} \equiv \hh{A} \hh{B} - \ql \{ \hh{A} \hh{B} \qr \}
 \label{def:contraction}
\end{equation}
%
And we have only four possible contractions
\begin{align}
\contraction{}{\hat{a}}{{}_{\alpha}^{\dagger}}{a}
a_{\alpha}^{\dagger}a_{\beta}^{\dagger} &= a_{\alpha}^{\dagger}a_{\beta}^{\dagger} - a_{\alpha}^{\dagger}a_{\beta}^{\dagger} = 0 \\
\contraction{}{a}{{}_{\alpha}}{a}
a_{\alpha}a_{\beta} &= a_{\alpha}a_{\beta} = 0 \\
\contraction{}{\hat{a}}{{}_{\alpha}^{\dagger}}{a}
a_{\alpha}^{\dagger}a_{\beta} &= a_{\alpha}^{\dagger}a_{\beta} - a_{\alpha}^{\dagger}a_{\beta} = 0 \\
\contraction{}{\hat{a}}{{}_{\alpha}^{\dagger}}{a}
a_{\alpha}a_{\beta}^{\dagger} &= a_{\alpha}a_{\beta}^{\dagger} - ( - a_{\beta}^{\dagger}a_{\beta}^{\dagger} ) = \delta_{\alpha\beta} \qquad \text{from (\ref{eq:creationannihilationanticommute})}
 \label{eq:contractionpossibilities}
\end{align}
%
We can have contractions between operators inside a normal ordered product,
\begin{equation}
\contraction{\ql \{ \qr. \hh{A}\hh{B}\hh{C},.}{\hh{P}}{}{\hh{Q}} 
\contraction{\ql \{ \qr. \hh{A}\hh{B}\hh{C}...\hh{P}\hh{Q},.}{\hh{X}}{}{\hh{Y}} 
\contraction{\ql \{ \hh{A}\hh{B}\hh{C}...\hh{P}\hh{Q}...\hh{X}\hh{Y}...\qr \} = (-1)^{p} }{\hh{P}}{}{\hh{Q}} 
\contraction{\ql \{ \hh{A}\hh{B}\hh{C}...\hh{P}\hh{Q}...\hh{X}\hh{Y}...\qr \} = (-1)^{p} \hh{P}\hh{Q}}{\hh{X}}{}{\hh{Y}} 
\ql \{ \hh{A}\hh{B}\hh{C}...\hh{P}\hh{Q}...\hh{X}\hh{Y}...\qr \} = (-1)^{p} \hh{P}\hh{Q} \hh{X}\hh{Y} \ql \{ \hh{A}\hh{B}\hh{C}...\qr \}
 \label{def:contrationinsidenormalorder}
\end{equation}
%
Wick's theorem states that we can express any product of creation and annihilation operators as sum of normal ordered products with all possible ways of contractions, i.e.

\begin{align}
\hh{A}\hh{B}\hh{C}\hh{D}...\hh{V}\hh{X}\hh{Y}\hh{Z} &= \ql\{ \hh{A}\hh{B}\hh{C}\hh{D}...\hh{V}\hh{X}\hh{Y}\hh{Z} \qr \}  \nonumber\\
\contraction{ + \sum_{(1)} \ql\{ \qr.}{\hh{A}}{}{..}
& + \sum_{(1)} \ql\{ \hh{A}\hh{B}\hh{C}\hh{D}...\hh{V}\hh{X}\hh{Y}\hh{Z} \qr \} \nonumber\\
\contraction{ + \sum_{(1)} \ql\{ \qr.}{\hh{A}}{,.}{\hh{C}}  %A little cheating whith ,.
\contraction[2ex]{ + \sum_{(1)} \ql\{ \qr.\hh{A}}{\hh{B}}{,.}{\hh{D}}  %A little cheating whith ,.
& + \sum_{(2)} \ql\{ \hh{A}\hh{B}\hh{C}\hh{D}...\hh{V}\hh{X}\hh{Y}\hh{Z} \qr \} \nonumber\\
&+...\nonumber\\
\contraction{ + \sum_{(N/2)} \ql\{ \qr.}{\hh{A}}{,.}{\hh{C}}  %A little cheating whith ,.
\contraction[2ex]{ + \sum_{(N/2)} \ql\{ \qr.\hh{A}}{\hh{B}}{,.}{\hh{D}}  %A little cheating whith ,.
\contraction{ + \sum_{(N/2)} \ql\{ \qr.\hh{A}\hh{B}\hh{C}\hh{D}...}{\hh{V}}{.,}{\hh{Y}}
\contraction[2ex]{ + \sum_{(N/2)} \ql\{ \qr.\hh{A}\hh{B}\hh{C}\hh{D}...\hh{V}}{\hh{X}}{,.}{\hh{Z}}
& + \sum_{(N/2)} \ql \{ \hh{A}\hh{B}\hh{C}\hh{D}...\hh{V}\hh{X}\hh{Y}\hh{Z} \qr \} 
\label{def:wicksteorem}
\end{align}

$\sum_{(m)}$ means sum over all terms with $m$ representing the number of contractions. $N$ is the total number of creation and annihilation operators. If there are different numbers of creation and annihilation operators, the vacuum expectation value would be zero, because of Eq. (\ref{eq:becauseannihilation}) and Eq. (\ref{eq:becausecreation}). If $N$ is odd, one of the operators would not be contracted and we would get zero as well. In order to get contribution one must contract all of the operators. For details of the proof see \cite{peskin} or \cite{bartlett}. 

The generalized Wick's theorem follows directly from Wick's theorem and states that the normal ordered product of operators strings $\{...\}$ are the same as the sum of the normal ordered product of the total group with all possible ways of contractions, i.e. 

\begin{align}
\ql\{\hh{A}\hh{B}\hh{C}\hh{D}..\qr \} \ql \{\hh{V}\hh{X}\hh{Y}\hh{Z}..\qr \} & = \ql \{ \hh{A}\hh{B}\hh{C}\hh{D}..\hh{V}\hh{X}\hh{Y}\hh{Z}\qr\} \nonumber \\
\contraction{+\sum_{(1)} \ql \{ \qr .}{\hh{A}}{\hh{B}\hh{C}\hh{D},}{\hh{V}} 
& +\sum_{(1)} \ql \{ \hh{A}\hh{B}\hh{C}\hh{D}..\hh{V}\hh{X}\hh{Y}\hh{Z}\qr\} \nonumber \\
\contraction{+\sum_{(2)} \ql \{ \qr .}{\hh{A}}{\hh{B}\hh{C}\hh{D},}{\hh{V}} 
\contraction[2ex]{+\sum_{(2)} \ql \{ \qr . \hh{A}}{\hh{B}}{\hh{C}\hh{D},\hh{V}}{\hh{X}} 
& +\sum_{(2)} \ql \{ \hh{A}\hh{B}\hh{C}\hh{D}..\hh{V}\hh{X}\hh{Y}\hh{Z}\qr\} \nonumber \\
& + ... \nonumber \\
\contraction{+\sum_{(N/2)} \ql \{ \qr .}{\hh{A}}{\hh{B}\hh{C}\hh{D}.}{V}
\contraction[2ex]{+\sum_{(N/2)} \ql \{ \qr .\hh{A}}{\hh{B}}{\hh{C}\hh{D}..\hh{V}}{\hh{X}}
\contraction[3ex]{+\sum_{(N/2)} \ql \{ \qr .\hh{A}\hh{B}}{\hh{C}}{\hh{D}..\hh{V}\hh{X}}{Y}
\contraction[4ex]{+\sum_{(N/2)} \ql \{ \qr .\hh{A}\hh{B}\hh{C} }{\hh{D}}{..\hh{V}\hh{X}\hh{Y}}{Z}
& +\sum_{(N/2)} \ql \{ \hh{A}\hh{B}\hh{C}\hh{D}..\hh{V}\hh{X}\hh{Y}\hh{Z}\qr\} \nonumber \\
 \label{def:generalizednormalorder}
\end{align}
%
Note that the only contribution to the vacuum expectation value comes from the full contractions, only the last sum will give contribution. There are no \emph{internal}  contractions, i.e. contraction between pairs of operators inside each operator string $\{..\}$. 

As an example, we can now use Wick's theorem to find the vacuum expectation value of the following products:

\begin{align}
\contraction{\bra{0} a_{i}a_{j}^{\dagger} \ket{0} .;;  \ql \{ \qr.}{a}{{}_{i}}{a}
\bra{0} a_{i}a_{j}^{\dagger} \ket{0} &=   \ql \{ a_{i}a_{j}^{\dagger} \qr \} =\delta_{ij}  \label{eq:vacuumexpectation2redo one-body} \\
%
\contraction{\bra{0} a_{\alpha_2}a_{\alpha_1}a_{\alpha}^{\dagger}a_{\beta}^{\dagger} \ket{0} = \ql \{ \qr.}{a}{{}_{\alpha_2}a_{\alpha_1} }{a}
%
\contraction[2ex]{\bra{0} a_{\alpha_2}a_{\alpha_1}a_{\alpha}^{\dagger}a_{\beta}^{\dagger} \ket{0} = \ql \{ a_{\alpha_2} \qr.}{a}{{}_{\alpha_1}a_{\alpha}^{\dagger}}{a}
%
\contraction{\bra{0} a_{\alpha_2}a_{\alpha_1}a_{\alpha}^{\dagger}a_{\beta}^{\dagger} \ket{0} = \ql \{a_{\alpha_2}a_{\alpha_1}a_{\alpha}^{\dagger}a_{\beta}^{\dagger} \qr \} + \ql \{ \qr. a_{\alpha_2}}{a}{{}_{\alpha.}}{a}
%
\contraction[2ex]{\bra{0} a_{\alpha_2}a_{\alpha_1}a_{\alpha}^{\dagger}a_{\beta}^{\dagger} \ket{0} = \ql 
\{a_{\alpha_2}a_{\alpha_1}a_{\alpha}^{\dagger}a_{\beta}^{\dagger} \qr \} + \ql \{ \qr. }{a}{{}_{\alpha_2}a_{\alpha_1}a^{\dagger}_{\alpha}}{a}
%
\contraction{\bra{0} a_{\alpha_2}a_{\alpha_1}a_{\alpha}^{\dagger}a_{\beta}^{\dagger} \ket{0} = \ql \{a_{\alpha_2}a_{\alpha_1}a_{\alpha}^{\dagger}a_{\beta}^{\dagger} \qr \} + \ql \{a_{\alpha_2}a_{\alpha_1}a_{\alpha}^{\dagger}a_{\beta}^{\dagger} \qr \} 
+ ,,.}{a}{{}_{\alpha_1}}{a}
%
\contraction{\bra{0} a_{\alpha_2}a_{\alpha_1}a_{\alpha}^{\dagger}a_{\beta}^{\dagger} \ket{0} = \ql \{ a_{\alpha_2}a_{\alpha_1}a_{\alpha}^{\dagger}a_{\beta}^{\dagger} \qr \} + \ql \{a_{\alpha_2}a_{\alpha_1}a_{\alpha}^{\dagger}a_{\beta}^{\dagger} \qr \} 
+ \ql \{ \qr. a_{\alpha_2}a_{\alpha.}.}{a}{{}^{\dagger}_{\alpha}}{a}
%
\bra{0} a_{\alpha_2}a_{\alpha_1}a_{\alpha}^{\dagger}a_{\beta}^{\dagger} \ket{0}&= \ql \{a_{\alpha_2}a_{\alpha_1}a_{\alpha}^{\dagger}a_{\beta}^{\dagger} \qr \} + \ql \{a_{\alpha_2}a_{\alpha_1}a_{\alpha}^{\dagger}a_{\beta}^{\dagger} \qr \} 
+ \ql \{\underbrace{a_{\alpha_2}a_{\alpha_1}a_{\alpha}^{\dagger}a_{\beta}^{\dagger}}_{=0} \qr \}  \nonumber \\
%
& = \delta_{\alpha \alpha_1}\delta_{\beta \alpha_2} - \delta_{\beta \alpha_1}\delta_{\alpha \alpha_2}  \label{eq:vacuumexpectation2redo}
\end{align}


\subsection{Particle-Hole Formalism}
\label{sec:particle-hole formalism}
One of the advantages with second quantization is that we can easily introduce a reference SD: $\ket{c}$ instead of using the physical vacuum state $\ket{0}$. This reduces the dimensionality of the problem and the new reference state would be defined by a boldfaced zero

\begin{equation}
\ket{\bf 0} \equiv \ket{\Phi_0} = \ket{ijk...n} 
 \label{def:referencestate}
\end{equation}
%
The new reference state $\ket{\bf 0}$ is also referred to as the \emph{Fermi vacuum}, this is now our Fermi level which is the level of our last occupied quantum state, usually the highest occupied orbital.
\\
In this representation we have hole states in addition to particle states. We will define what the hole-states are after the following example. Assume that we have three states that are successively filled with $n-1$,$n$ and $n+1$ single-particle states $\alpha_i$

\begin{align}
\ket{\Phi_0} &\equiv \ket{\alpha_1\alpha_2...\alpha_{n}} &\qquad \text{(reference state)} \label{def:reference state} \\
\ket{\Phi_{\alpha_1}} &\equiv \ket{\alpha_2\alpha_3...\alpha_{n}} = a_{\alpha_1} \ket{\Phi_0} &\qquad \text{(creation of a hole)} \label{def:creation of a hole} \\
\ket{\Phi^{\alpha}} &\equiv \ket{\alpha\alpha_1\alpha_2...\alpha_{n}} = a_{\alpha}^{\dagger} \ket{\Phi_0} &\qquad \text{(creation of a particle)}
 \label{def:creation of a particle}
\end{align}
%
And assume that the energies of the single-particle orbitals is such that (see Figure (\ref{fig:particleholerepresentation}))
\begin{equation} \epsilon_{\alpha_{n+1}}>\epsilon_{\alpha_n}>\epsilon_{\alpha_{n-1}}>...\epsilon_{\alpha_{2}}>\epsilon_{\alpha_{1}}
 \label{def:singleparticleenergies} 
\end{equation}
%
Let us then define our Fermi level to be $\alpha_n$, a \emph{hole} is then a state that is below or equal to the Fermi level $\alpha_i \leq \alpha_n$. And a particle is a state above $\alpha_i > \alpha_n$. 

When we change our reference state from the physical vacuum state $\ket{0}$ to a particle-hole vaccum $\ket{\bf 0}$, we have to introduce new operators as well. 

\begin{equation}
a_{\alpha} \ket{\bf 0} \neq 0 
 \label{show:neednewoperators}
\end{equation}
%
since $\alpha \in \ket{\bf 0}$, while for the physical vacuum we have $a_\alpha \ket{0} = 0$ for all $\alpha$. The new operators need to have the relation $b_\alpha\ket{\bf 0} = 0$.

The new operators are called \emph{quasi-} annihilation and creation operators, viz

\begin{equation}
b_{\alpha}^{\dagger} = \left\{ \begin{array}{rl}
 a_\alpha^{\dagger}, & \alpha > \alpha_F\\
 a_\alpha, & \alpha \leq \alpha_F
       \end{array} \right.
 \label{def:quasi-creation}
\end{equation} 

\begin{equation}
b_{\alpha} = \left\{ \begin{array}{rl}
 a_\alpha, & \alpha > \alpha_F \\
 a_\alpha^{\dagger}, & \alpha \leq \alpha_F
       \end{array} \right.
 \label{def:quasi-annihilation}
\end{equation}
%
Where $\alpha_F$ is the Fermi level representing the last occupied single-particle orbit of the reference state $\ket{c}$. 
%
And we have the following anticommutation relations
\begin{align}
\ql \{ b_{\alpha}, b_{\beta} \qr \} &= 0  \label{def:quasioperator1}\\
\ql \{ b_{\alpha}^{\dagger}, b_{\beta}^{\dagger} \qr \} &= 0  \label{def:quasioperator2}\\
\ql \{ b_{\alpha}^{\dagger}, b_{\beta}  \qr \} &= \delta_{\alpha \beta}  \label{def:quasioperator3}
\end{align}
%
The reference state is normalized

\begin{equation}
 \braket{c}{c} = 1
\end{equation}
%
A quasi particle state is defined by a state which has one or more particles/holes  added to the reference state $\ket{c}$.
%
\begin{equation}
\ket{abcd...ijkl...pqrs...} \equiv b_a^{\dagger}b_b^{\dagger}b_c^{\dagger}b_d^{\dagger}...b_i^{\dagger}b_j^{\dagger}b_k^{\dagger}b_l^{\dagger}...b_p^{\dagger}b_q^{\dagger}b_r^{\dagger}b_s^{\dagger}\ket{c}
 \label{def:manybody-quasiparticle}
\end{equation}
%
The convention is that indices $i,j,k,l...$ indicate states which are occupied by holes. Indices $a,b,c,d...$ indicate the states which are occupied by particles. And $p,q,r,s...$ indicate any state. We are going to simplify the notation further

\begin{align}
 b_i^{\dagger} &= a_i           = i         \,\,\,\,\,\qquad \text{(creation of a hole = removing a state below $\alpha_F$)} \nonumber \\
 b_i           &= a_i^{\dagger} = i^\dagger \,\,\qquad \text{(annihilation of a hole = creating a state below $\alpha_F$)} \nonumber \\
 b_a^{\dagger} &= a_a^{\dagger} = a^\dagger \qquad \text{(creation of a particle = adding a state above $\alpha_F$) } \nonumber \\
 b_a           &= a_a           = a         \,\,\qquad \text{(annihilation of a particle = removing a state above $\alpha_F$)}
\label{def:simplyfied notation of quasi-operators}
\end{align}
%
with the following contractions
\begin{align}
\contraction{}{p}{{}^{\dagger}}{q}
\contraction{p^{\dagger}q^{\dagger} = }{p}{}{q}
p^{\dagger}q^{\dagger} &= pq = 0  \label{eq:contraction 1} \\
\contraction{}{i}{{}^{\dagger}}{j}
i^{\dagger}j &= \delta_{ij} \label{eq:contraction 2} \\
\contraction{}{i}{}{j}
ij^{\dagger} &= 0 \label{eq:contraction 3} \\
\contraction{}{\hat{a}}{}{b}
ab^{\dagger} &= \delta_{ab} \label{eq:contraction 4}\\
\contraction{}{\hat{a}}{{}^{\dagger}}{b}
a^{\dagger}b &= 0 \label{eq:contraction 5}
\end{align}

%\begin{equation}
% \text{\scalebox{0.5}{\input{occupation_orbitals2.pdftex_t}}}
%\end{equation}

\begin{figure}
\centering
     \scalebox{0.7}{\input{occupation_orbitals2.pdftex_t}}
  \caption{The dashed line represents our Fermi level}
  \label{fig:particleholerepresentation}
\end{figure}

\section{The Normal-Ordered Hamiltonian}
We want to rewrite the second-quantized form of the electronic Hamiltonian (\ref{def:secondquantized hamiltonian}) 

\begin{equation}
\hat{H} = \sum_{pq} \bra{p}h\ket{q}p^{\dagger}q + \frac{1}{4}\sum_{pqrs} \bra{pq}|v|\ket{rs}p^{\dagger}q^{\dagger}sr 
 \label{def:normal-ordered hamiltonian}
\end{equation}
%
using Wick's theorem, 

\begin{align}
\contraction{p^{\dagger}q.;; \ql \{ p^{\dagger}q \qr \} + \ql \{ \qr.}{p}{{}^{\dagger}}{q}
 p^{\dagger}q &= \ql \{ p^{\dagger}q \qr \} +  \underbrace{\ql \{ p^{\dagger}q \qr \}}_{ p=i,q=j}  = \ql \{ p^{\dagger}q \qr \}  + \delta_{ij}  \label{eq:normalordere one-particle} \\
%
\contraction{p^{\dagger}q^{\dagger}sr = \ql \{ p^{\dagger}q^{\dagger}sr \qr \} ;; \ql \{ \qr.}{p}{{}^\dagger q^\dagger}{s}
%
\contraction{p^{\dagger}q^{\dagger}sr = \ql \{ p^{\dagger}q^{\dagger}sr \qr \} + \ql \{ p^{\dagger}q^{\dagger}sr \qr \} ;; \ql \{ \qr .}{p}{{}^\dagger q^\dagger s}{r}
%
\contraction{p^{\dagger}q^{\dagger}sr = \ql \{ p^{\dagger}q^{\dagger}sr \qr \} + \ql \{ p^{\dagger}q^{\dagger}sr \qr \} + \ql \{ p^{\dagger}q^{\dagger}sr \qr \} .. \ql \{ p^{\dagger} \qr .}{q}{{}^\dagger}{s}
%
\contraction{p^{\dagger}q^{\dagger}sr = \ql \{ p^{\dagger}q^{\dagger}sr \qr \} + \ql \{ p^{\dagger}q^{\dagger}sr \qr \} + \ql \{ p^{\dagger}q^{\dagger}sr \qr \} .. \ql \{ p^{\dagger} \qr .}{q}{{}^\dagger}{s}
%
\contraction{p^{\dagger}q^{\dagger}sr = \ql \{ p^{\dagger}q^{\dagger}sr \qr \} + \ql \{ p^{\dagger}q^{\dagger}sr \qr \} + \ql \{ p^{\dagger}q^{\dagger}sr \qr \} + \ql \{ p^{\dagger}q^{\dagger}sr \qr \}.. \ql \{ p^{\dagger} \qr .}{q}{{}^\dagger s}{r}
%%%%
p^{\dagger}q^{\dagger}sr &= \ql \{ p^{\dagger}q^{\dagger}sr \qr \} + \ql \{ p^{\dagger}q^{\dagger}sr \qr \}  + \ql \{ p^{\dagger}q^{\dagger}sr \qr \}  + \ql \{ p^{\dagger}q^{\dagger}sr \qr \} +  \ql \{ p^{\dagger}q^{\dagger}sr \qr \}  \nonumber \\
%%%%
\contraction{... \ql \{ \qr.}{p}{{}^{\dagger} \da{q}}{s}
\contraction[2ex]{... \ql \{ p^\dagger \qr. }{q}{{}^\dagger s}{r}
%
\contraction[2ex]{... \ql \{ p^{\dagger}q^{\dagger}sr \qr \}  + \ql \{ \qr.}{p}{{}_{\dagger} \da{q}	s}{r}
\contraction{... \ql \{ p^{\dagger}q^{\dagger} sr \qr \} +  \ql \{ p^\dagger \qr. }{q}{.}{S}
%
&+ \ql \{ p^{\dagger}q^{\dagger}sr \qr \}  + \ql \{ p^{\dagger}q^{\dagger}sr \qr \} \nonumber  \\
%%%%
\contraction{ = \ql \{ p^{\dagger}q^{\dagger}sr \qr \} - ,,}{p}{\da{}}{s}
\contraction{ = \ql \{ p^{\dagger}q^{\dagger}sr \qr \} - ,,,,,p^\dagger s \ql \{ q^\dagger r \qr \} + }{p}{\da{}}{r}
\contraction{ = \ql \{ p^{\dagger}q^{\dagger}sr \qr \} - ,,,,,,,...p^\dagger s \ql \{ q^\dagger r \qr \} + p^\dagger r \ql \{ q^\dagger s \qr \} + }{q}{\da{}}{s}
\contraction{ = \ql \{ p^{\dagger}q^{\dagger}sr \qr \} - ,,,,,,,,,,,,p^\dagger s \ql \{ q^\dagger r \qr \} + p^\dagger r \ql \{ q^\dagger s \qr \} + q^\dagger s \ql \{ p^\dagger r \qr \} - }{p}{\da{}}{s}
%%%
& = \ql \{ p^{\dagger}q^{\dagger}sr \qr \} - \underbrace{p^\dagger s}_{p=i,s=j} \ql \{ q^\dagger r \qr \} + \underbrace{p^\dagger r}_{p=i, r=j} \ql \{ q^\dagger s \qr \} + \underbrace{q^\dagger s}_{q=i,s=j} \ql \{ p^\dagger r \qr \} - \underbrace{q^\dagger r}_{q=i,r=j} \ql \{ p^\dagger s \qr \} \nonumber \\
%%%
\contraction{- }{p}{\da{}}{s}
\contraction{.... \da{p}s}{q}{\da{}}{r}
%
\contraction{- \da{p}s\da{q}r + }{p}{\da{}}{r}
\contraction{- \da{p}s\da{q}r + \da{p}r}{q}{\da{}}{s}
& - \da{p}s\da{q}r + \da{p}r\da{q}s \qquad \text{only contribution when $p=i,s=j,q=k,r=l$} \nonumber  \\
%%%
& = \ql \{ p^{\dagger}q^{\dagger}sr \qr \} - \delta_{ij}\ql \{ q^\dagger r \qr \} + \delta_{ij}\ql \{ q^\dagger s \qr \} + \delta_{ij} \ql \{ p^\dagger r \qr \} - \delta_{ij}\ql \{ p^\dagger s \qr \} \nonumber \\
& = \delta_{ij}\delta_{kl} + \delta_{ij}\delta_{kl}
\end{align}
%
Here we have done contractions relative to a reference state $\ket{\bf 0}$ and followed the relations Eq. (\ref{eq:contraction 1}-\ref{eq:contraction 5}). Then the normal-ordered one-body Hamiltonian is

\begin{equation}
\hat{H}_1 = \sum_{pq} \bra{p}h\ket{q} \ql \{ p^{\dagger}q \qr \}  + \sum_{i}\bra{i}h\ket{i}
 \label{def:normal-ordered hamiltonian part 1}
\end{equation}
%
And the two-body 
\begin{align}
\hat{H}_2 &= \frac{1}{4} \sum_{pqrs} \bra{pq}|v|\ket{rs}  \ql \{ p^{\dagger}q^{\dagger}sr \qr \} - \frac{1}{4} \sum_{qri} \bra{iq}|v|\ket{ri}  \ql \{ p^{\dagger} r \qr \}  + \frac{1}{4}  \sum_{qsi} \bra{iq}|v|\ket{is} \ql \{ q^{\dagger}s \qr \} \nonumber \\
%
& + \frac{1}{4} \sum_{pri} \bra{pi}|v|\ket{ri} \ql \{ p^{\dagger} r \qr \} - \frac{1}{4} \sum_{psi} \bra{pi}|v|\ket{is}\ql \{ p^{\dagger} s \qr \} + \frac{1}{4}\sum_{ij}\ql [ \bra{ij}|v|\ket{ij} - \bra{ij}|v|\ket{ji}\qr ] 
 \label{def:normal-ordered hamiltonain part 2}
\end{align}
%
The second term is equal to the fourth term. And the third terms is equal to the fifth term.  

\begin{align}
 \hat{H} &= \sum_{pq} \bra{p}h\ket{q} \ql \{ p^{\dagger}q \qr \} + \sum_{i}\bra{i}h\ket{i} + \frac{1}{4} \sum_{pqrs} \bra{pq}|v|\ket{rs} \ql \{ p^{\dagger}q^{\dagger}sr \qr \}   \nonumber \\ 
&+ \frac{1}{2} \sum_{pri} \bra{ip}|v|\ket{ri}\ql \{ p^{\dagger} r \qr  \} - \frac{1}{2} \sum_{psi} \bra{pi}|v|\ket{is}\ql \{ p^{\dagger} s \qr \} + \frac{1}{2} \sum_{ij} \bra{ij}|v|\ket{ij}
\label{def:second-ordered Hamiltonian1}
\end{align}
%
We have changed summation variables $s,r\rightarrow q$

\begin{align}
 \hat{H} &= \sum_{pq} \bra{p}h\ket{q} \ql \{ p^{\dagger}q \qr \} + \sum_{i}\bra{i}h\ket{i} + \frac{1}{4} \sum_{pqrs} \bra{pq}|v|\ket{rs} \ql \{ p^{\dagger}q^{\dagger}sr \qr \}   \nonumber \\ 
& + \sum_{pq} \ql (  \bra{p}h\ket{q} + \sum_{i}\bra{ip}|v|\ket{qi}  \qr ) \ql \{ p^{\dagger} q \qr \} + \frac{1}{2} \sum_{ij} \bra{ij}|v|\ket{ij}
 \label{def:second-ordered Hamiltonian2}
\end{align}
%
And define the following
\begin{align}
f_q^p         &\equiv  \bra{p}h\ket{q} + \sum_{i}\bra{ip}|v|\ket{qi}  \label{def:f_p^q} \\
\hh{F}_N      &\equiv  \sum_{pq} f_q^p  \ql \{ p^{\dagger} q \qr \}    \label{def:F_N}     \\
\hh{V}_N      &\equiv  \frac{1}{4} \sum_{pqrs} \bra{pq}|v|\ket{rs} \ql \{ p^{\dagger}q^{\dagger}sr \qr \}\label{def:V_N}  \\
\hh{H}_N      &\equiv  \hh{F}_N + \hh{V}_N \label{def:H_N} 
\end{align}
%
From Eq. (\ref{def:second-ordered Hamiltonian2}) we see that only the last term survives and expectation value would be
\begin{equation}
\bra{\bf 0}\hh{H}\ket{\bf 0} = \sum_i \bra{i}h\ket{i} + \frac{1}{2} \sum_{ij} \bra{ij}|v|\ket{ij} 
 \label{eq:second-qunatized reference expectation value}
\end{equation}
%
Finally the Hamiltonian reads
\begin{equation}
\hh{H} = \hh{H}_N + \bra{\bf 0}\hh{H}\ket{\bf 0}
 \label{eq:final normal ordered hamiltonian}
\end{equation}
%
And the normal-ordered electronic Hamiltonian 
\begin{equation}
\hh{H}_N = \hh{H} - \bra{\bf 0}\hh{H}\ket{\bf 0}
 \label{eq:normal-ordered hamiltonian}
\end{equation}
%
This is nothing but a shift by a constant for the expectation value. The usefulness of this will be clearer when we introduce the Coupled Cluster Theory.




 
%%%%%%%%%%%%%%%%%%%%%%%%%%%%%%%%%%%%%%%%%%%%%%%%%%%%%%%%%%%%%%%%%%%%%%%%%%%%%%%%%%%%%%%%
%%%%%%%%%%%%%%%%%%%%%%%%%%%%%%%   PART MANY-BODY METHODS        %%%%%%%%%%%%%%%%%%%%%%%%
%%%%%%%%%%%%%%%%%%%%%%%%%%%%%%%%%%%%%%%%%%%%%%%%%%%%%%%%%%%%%%%%%%%%%%%%%%%%%%%%%%%%%%%%
% \part{MANY-BODY METHODS}
% \chapter{Hatree-Fock Method}
%In this chapter we shall...................
%Hartree-Fock theory and DFT are simplifications of the manybody problem with moving particles in a potential field. Physical systems such as atoms, molecules and  nuclei. In such systems particles are not only affected by the main potential set up by a \emph{core}, but aswell as the field generated by the other particles. 

%The first approximation is to assume that the kinetic energy of the \emph{core} is negligible compared to the particles surrounding it. Then we just have to solve the Schr\"{o}dinger equation for particles instead of the whole system, particles$+$core. This is the philosophy of \emph{Born-Oppenheimer approximation}.

The exact solution to the Schr\"{o}inger equation cannot be obtained in most of the systems of chemical interest (see \cite{helgaker}). The number of particles involved in those problems are just to big for computers to solve. Therefore approximations are needed. The approximations would of course depend on the physical system. We could just try some potential and interactions that works. But the many-body theories have hierachcal structure when it comes to approximations. The advantages of this is the fact that we build up experience about the given models. And we can do benchmarks and comparison tests. This allows too differentiate between different type of \emph{effects} in the theories. In this chapter we are going to look into one of the first approximations, the Hartree-Fock method.

\section{Introduction}
\label{sec:Hatree-Fock Method:Introduction}
The independent particle picture is an assumption that the electron-electron interaction is rather weak. And each electron could be viewed as independent particles which sees an effective field set up by the other electrons. This is the philosophies of Hartree-Fock theories and Density functional theories. 

The Hartree-Fock method is an optimization problem using Lagrange multipliers as the mathematical tool. The integral we want to minimize are in general

\begin{equation}
E[\Phi] = \int_a^b f(\Phi(r),\frac{\partial \Phi}{\partial r}, r)d^3 r
 \label{def:integral we want to minimize}
\end{equation}
%
We want to find a function $\Phi$ that minimizes the \emph{functional} $E[\Phi]$. The set of functions $V$ have the following conditions: $V = \ql \{ \Phi:[-\infty,\infty] \rightarrow \mathbb{R}:\text{$\Phi$ is continous and differentiable}: \braket{\Phi}{\Phi} = 1 \qr \}$. The unkowns are the functions in $V,\frac{\partial \Phi}{\partial r}$ and $r$. Although the integral limits $a,b$ are defined, the integration path is not. We want to find a path with a given set of unknowns such that $\delta E = 0$. This will give us minima, maxima or saddle points. So we have check if we have a minima after finding the solution.

The quantum mechanical functional we want to minimize is

\begin{equation}
E[\Phi] = \frac{\bra{\Phi}\hat{H}\ket{\Phi}}{\braket{\Phi}{\Phi}} =  \frac{\int \Phi^* \hat{H} \Phi d\tau}{\int \Phi^* \Phi d \tau}
 \label{def:energyfunctional}
\end{equation}
If $\Phi$ is an eigenstate of the $\hat{H}$, then the functional will be stationary
\begin{equation}
\hat{H} \ket{\Phi} = E_0 \ket{\Phi} \Rightarrow E[\Phi] = E_0 \Rightarrow \delta E[\Phi] = 0 
 \label{def:eigenstate then stationary}
\end{equation}
%
Conversely we can show a stationary point is an eigenstates of the Hamiltonian (\ref{def:eigenstate then stationary}), and therefore a solutions of the Schr\"{o}dinger equation. 
%
\begin{proof}
From (\ref{def:energyfunctional}) and the fact that setting $\delta E[\Phi]$ = 0 gives us
\begin{align}
0 &= \delta [ \bra{\Phi}\hat{H}\ket{\Phi} - E[\Phi] \ql (\braket{\Phi}{\Phi} - 1 \qr ) ] \nonumber \\
0 &= \delta [ \bra{\Phi} \ql( \hat{H}-E[\Phi] \qr ) \ket{\Phi}+ E[\Phi] ] \nonumber \\
0 &=  \bra{\delta \Phi} \ql( \hat{H}-E[\Phi] \qr ) \ket{\Phi}+  \bra{\Phi} \ql( \hat{H}-E[\Phi] \qr ) \ket{\delta \Phi}+ \delta E[\Phi] \nonumber \\
0 &=  \bra{\delta \Phi} \ql( \hat{H}-E[\Phi] \qr ) \ket{\Phi}+  \bra{\Phi} \ql( \hat{H}-E[\Phi] \qr ) \ket{\delta \Phi} 
\label{proof:stationary points are eigenvalues}
\end{align}
We can add a phase $\delta \Phi \rightarrow i\delta \Phi$. Adding a phase to a wave function will not change the expectation value. Then we get two se of equations 
\begin{align}
I. \qquad & \bra{\delta \Phi} \ql( \hat{H}-E \qr ) \ket{\Phi}+  \bra{\Phi} \ql( \hat{H}-E \qr ) \ket{\delta \Phi}  = 0\\
II. \qquad& \bra{\delta \Phi} \ql( \hat{H}-E \qr ) \ket{\Phi}-  \bra{\Phi} \ql( \hat{H}-E \qr ) \ket{\delta \Phi} = 0
 \label{proof:arbritrary phase}
\end{align}
We have multiplied with $i$ in (II) and remember that $i \ket{\delta \Phi} = -i\bra{\delta \Phi}$. And adding them together
\begin{equation}
\bra{\delta \Phi} \ql( \hat{H}-E \qr ) \ket{\Phi} = 0
 \label{proof:adding I with II}
\end{equation}
The choice of $\bra{\delta \Phi}$ is arbritrary, so the equation must hold for all possible $\bra{\delta \Phi}$'s. Therefore this can only be satisfied if $\Phi$ is an eigenfunction to $\hat{H}$.
\end{proof}
%
\noindent This is equivalent to Rayleigh-Ritz principle \cite{arfken} that tells us that the functional $\Omega_{RR} = \bra{\Psi}\hat{H}-E\ket{\Psi}$  is stationary at any eigenfunctions $H\ket{\Phi_m} = E_m \ket{\Phi_m}$. 
%
One important feature of the functional (\ref{def:energyfunctional}) is that errors are in the second order in $\delta \Phi$. 
 
\begin{proof}
Our trial wavefunction is $\ket{\phi} = \ket{\psi_0} + \lambda\ket{\psi}$. $\ket{\psi_0}$ is a stationary state, $\ket{\psi}$ is our guess, and $\lambda$ is complex number. Using the relations
\begin{align}
(\hat{H}-E_0) \ket{\psi_0} &= 0 \nonumber \\
(\hat{H}-E_0) [\ket{\phi} - \lambda \ket{\psi} ] &= 0 \nonumber \\
(\hat{H}-E_0) \ket{\phi} &= \lambda (\hat{H}-E_0 ) \ket{\psi} 
 \label{proof:need relations}
\end{align}
%
We get
\begin{align}
E[\Phi] - E_0 & = \frac{\bra{\phi}(\hat{H}-E_0)\ket{\phi}}{\braket{\phi}{\phi}} \nonumber \\ 
&= \frac{\bra{\lambda\psi}(\hat{H}-E_0)\ket{\lambda\psi}}{\braket{\phi}{\phi}} \nonumber \\
&= |\lambda|^2 \frac{\bra{\psi}(\hat{H}-E_0)\ket{\psi}}{\braket{\phi}{\phi}}
\end{align}
\end{proof}

\section{Derivation of the Hartree-Fock equations}
The HF ansatz reads

\begin{equation}
\Phi_{\text{HF}} = \ket{pqrs...}
 \label{def:Hartree Fock ansatz}
\end{equation}
%
Where $pqrs...$ are the HF spin-orbitals. $\Phi_{\text{HF}}$ is a SD with $N$ states written in second quantization formalism. We could vary these directly, giving us the Hartree-Fock-Roothaan method \cite{roothaan}. Another possibility is to expand our spin-orbitals $\ket{a}$ as a linear combination of a finite number of basis states $\ket{\alpha}$. (A unitary transformation of the states $\ket{a}$?) The sum goes to infinite in general but we do a truncation. (Why does this work? with the truncation).
\begin{equation}
 \ket{a} = \sum_{\alpha}^n C_{a\alpha} \ket{\alpha}
 \label{def:spin-orbital expantion}
\end{equation}
%
$C_{a\alpha}$ is the expansion coefficient of an unitary matrix $C^{\dagger}C = 1$. Writing the Hamiltonian in this basis up to two-body interactions gives

\begin{equation}
\hat{H} = \sum_{ab} \bra{a}h\ket{b}a_a^{\dagger}a_b + \frac{1}{2} \sum_{abcd} \bra{ab}v\ket{cd} a_a^{\dagger}a_b^{\dagger} a_d a_c
 \label{def:Hamiltonian in second quantization}
\end{equation}
Note that the Hamiltionian is in second quantization form (see section \ref{sec:second quantization} and \cite{bartlett} for details). $h$ is just the one-body operator, and $v$ is the two-body operator. 
%
We will restrict our system to a closed shell system (RHF-equations, see cite...), i.e. all possible spin-orbital levels are filled. Our trial wavefunction $\ket{\Phi_{\text{HF}}}$ will be the \emph{Fermi vacuum} $\ket{\bf0}$ (see section \ref{sec:particle-hole formalism}). We will use Wick's theorem to evaluate $\bra{\bf 0}\hat{H}\ket{\bf 0}$. First the Hamiltonian (\ref{def:Hamiltonian in second quantization}) have to be rewritten in the quasi-operator representation
\begin{equation}
\hat{H} = \sum_{ij} \bra{i}h\ket{j}b_ib_j^{\dagger} + \frac{1}{2} \sum_{ijkl} \bra{ij}v\ket{kl} b_ib_j b_l^{\dagger} b_k^{\dagger}
 \label{def:Hamiltonian in quasi-operator representation}
\end{equation}
Wick's theorem on quasi-operators
\begin{align}
\bra{\bf 0} b_{i}b_{j}^{\dagger} \ket{\bf 0} &= \delta_{ij} \qquad &\text{from (\ref{eq:vacuumexpectation2redo one-body})} \\
\bra{\bf 0} b_ib_j b_l^{\dagger} b_k^{\dagger} \ket{\bf 0} &= \delta_{kj}\delta_{li} - \delta_{lj}\delta_{ki} \qquad &\text{from (\ref{eq:vacuumexpectation2redo})} 
 \label{eq:evaluate quasi-operators}
\end{align}
Lead us to the following energy functional
\begin{align}
E[\Phi_{\text{HF}}] &= \bra{\bf 0}\hat{H}\ket{\bf 0} = \sum_{i} \bra{i}h\ket{i} + \frac{1}{2}   \sum_{ijkl} [\bra{ij}v\ket{ij} - \bra{ij}v\ket{ji}] \nonumber \\
        &= \sum_{i} \bra{i}h\ket{i} + \frac{1}{2}   \sum_{ij} \bra{ij}v\ket{ij}_{\text{AS}}
 \label{eq:energyfunctional}
\end{align}
%
Inserting the new basis states (\ref{def:spin-orbital expantion}) in this expression yields 
%
\begin{equation}
E[\Phi_{\text{HF}}] = \sum_{a}^N \sum_{\alpha \beta}^n C_{a\alpha}^*C_{b\beta} \bra{\alpha}h\ket{\beta} + \frac{1}{2} \sum_{ab}^N\sum_{\alpha\beta\gamma\delta}^n C_{a\alpha}^*C_{b\beta}^* C_{a\gamma}C_{b\delta} \bra{\alpha\beta}v\ket{\gamma\delta}_{\text{AS}}
 \label{eq:energyfunctional2}
\end{equation}
%
The second-quantized form of an operator is unchanged under a unitary transformation of the basis (see \cite{bartlett}). 

As mention the method of choice for minimizing $E[\Phi_{HF}]$ is the Lagrange multipliers method (see \cite{boas}).

In general we want to find a stationary point $p=(x_1,x_2,...,x_n)$ of a function $f(p)$ with multiple constraints
\begin{align}
g_1(p) &= 0 \nonumber \\
g_2(p) &= 0 \nonumber \\
\vdots  \\
g_N(p) &= 0  \nonumber 
 \label{def:constraints}
\end{align}
%
Then a point $p$ is a stationary if and only if 
%
\begin{equation}
\nabla_{\nu} f(p) - \sum_{i=1}^N \lambda_i \nabla_\nu g_i(p) = 0, \qquad \forall \nu \in p
 \label{eq:langrange multiplier}
\end{equation}
%
$\lambda_i$ is the Langrangian multipliers. We get a system of $N\cdot n$ equations to be solved for the $(N+n)$-set of variables  $\lambda_1,\lambda_2...\lambda_N,x_1,x_2,...x_n$

In our case the constraint is the HF spin-orbitals are orthonormal (Because we have done a unitary transformation?)(which they have to be? No Since a og b can have same spin, depends on the physical system)

\begin{equation}
\braket{a}{b} = \sum_{\alpha}^n C_{a\alpha}^*C_{b\alpha} = \delta_{ab}
 \label{def: Hartree Fock constraint}
\end{equation}
%
We want to find the stationary point of the energy functional (\ref{eq:energyfunctional2}) with respect to a specific coefficient $C_{k\mu}^*$ in (\ref{def:spin-orbital expantion}). We dont have to do the same for $C_{k\mu}$ since they only differ by a phase. This means that for our HF spin-orbitals $\ket{a} \rightarrow i\ket{a}$ will not change the expectation value. 

We define
\begin{align}
F &\equiv E[\Phi_{\text{HF}}] - \sum_a^N \lambda_a g_a \nonumber \\
  &= E[\Phi_{\text{HF}}] - \sum_a^N \lambda_a \sum_\alpha^n C_{a\alpha}^*C_{a\alpha}
 \label{def:Big F}
\end{align}
%
Then we take the partial derivate of $F$ with respect to $C_{k\mu}^*$ using Eq. (\ref{eq:langrange multiplier}), yielding

\begin{align}
\frac{\partial}{\partial C_{k\mu}^*}  &\ql(  \sum_{a}^N\qr.  \sum_{\alpha \beta}^n C_{a\alpha}^*C_{b\beta} \bra{\alpha}h\ket{\beta} + \frac{1}{2} \sum_{ab}^N\sum_{\alpha\beta\gamma\delta}^n C_{a\alpha}^*C_{b\beta}^*  C_{a\gamma}C_{b\delta} \bra{\alpha\beta}v\ket{\gamma\delta}_{\text{AS}} \nonumber - \ql. \sum_a^N \lambda_a \sum_\alpha C_{a\alpha}^*C_{a\alpha}  \qr ) \\
 &= \sum_{\alpha \beta}^n C_{k\beta} \bra{\alpha}h\ket{\beta} + \frac{1}{2} \sum_{a}^N \sum_{\alpha\beta\gamma\delta}^n C_{a\beta}^*C_{k\gamma}  C_{a\delta} \bra{\alpha\beta}v\ket{\gamma\delta}_{\text{AS}} - \lambda_k \sum_\alpha^n C_{k\alpha} \\
&= \sum_{\alpha}^n\ql (\sum_\beta^n C_{k\beta} \bra{\alpha}h\ket{\beta} + \frac{1}{2} \sum_{a}^N \sum_{\beta\gamma\delta}^n C_{a\beta}^*C_{k\gamma}  C_{a\delta} \bra{\alpha\beta}v\ket{\gamma\delta}_{\text{AS}} - \lambda_k C_{k\alpha} \qr ) = 0  \label{eq:partial of F} 
\end{align}
%
The factor $1/2$ dissapears because of the product rule when we derivate $C_{a\alpha}^*C_{b\beta}^*  C_{a\gamma}C_{b\delta}$, we get two cases, when $\{a = k, \alpha = \mu \}$ and $\{b = k, \beta = \mu\}$. We rewrite Eq. (\ref{eq:partial of F})

\begin{equation}
\sum_\gamma^n C_{k\gamma} \bra{\alpha}h\ket{\gamma} + \frac{1}{2} \sum_{a}^N \sum_{\beta\gamma\delta}^n C_{a\beta}^*C_{k\gamma}  C_{a\delta} \bra{\alpha\beta}v\ket{\gamma\delta}_{\text{AS}}  = \lambda_k C_{k\alpha}
  \label{eq: rewritten equation} 
\end{equation}
Where we have changed the summation index $\beta \rightarrow \gamma$ in the first sum. And this than be simplified further

\begin{equation}
\sum_\gamma^n \ql ( \bra{\alpha}h\ket{\gamma} + \frac{1}{2} \sum_{a}^N \sum_{\beta\delta}^n C_{a\beta}^*  C_{a\delta} \bra{\alpha\beta}v\ket{\gamma\delta}_{\text{AS}} \qr ) C_{k\gamma} = \lambda_k C_{k\alpha}
 \label{def: simplified HF equation}
\end{equation}
%
By defining the Hartree-Fock Hamiltonian as 
\begin{equation}
h_{\alpha \gamma}^{\text{HF}} \equiv \bra{\alpha}h\ket{\gamma} + \frac{1}{2} \sum_{a}^N \sum_{\beta\delta}^n C_{a\beta}^*  C_{a\delta} \bra{\alpha\beta}v\ket{\gamma\delta}_{\text{AS}}
 \label{def: HF Hamiltonian}
\end{equation}
We finally obtain the Hartree-Fock equation
\begin{equation}
\sum_\gamma^n h_{\alpha\gamma}^{\text{HF}} C_{k\gamma} = \lambda_k C_{k\alpha}
 \label{def: The final HF equation}
\end{equation}
%
This is a non-linear set of equations, meaning that the set of equations changes with solutions we find. One way to solve this is by the \emph{self-consisten field method}. 



%
%\emph{the Hylleraas-Undheim theorem} \cite{hylleraas} and \cite{helgaker}
%
%More Hartree-Fock Approximation (see Heinnonen 59 kapt 7) And about Restricted Hartree-Fock approximation and the symmetry dilemma.
%
%The quasi-particle concept Heinnonen 313 kapt 26
%
%Time-independent perturbation theory Heinnonen 197 kapt 17
%
%Particle and hole operators and Wick's theorem kapt 19,20,21,22...
\section{Self-Consistent Field Method}
This is an iterative method that involves diagonalizing the $h_{\alpha\beta}^{\text{HF}}$-matrix for each iteration and summing the eigenvalues. It reaches selfconsitensy  

In a quantum mechanical system of many particles, particles are interacting with each other and there is \emph{correlation} between them. Therefore we have to find a wavefunction which describe the system as an whole. Our guess is that this wavefunction could have the form of an Slater-Determinant which is product of the single-particle wavefunctions, which depends on both spin and positional coordinates. 

The single-particle wavefunctions must be \emph{self-consistent}, i.e. they are solutions of the  Schr\"{o}dinger equations for the average field produced by the other particles, and they determine the average potential of the field.

One of the earliest methods relating to a self-consistent field was the Thomas-Fermi method (L. Thomas and E. Fermi) (ref). It did not use wave functions but the density distribution of the particles. The self-consistent Thomas-Fermi field explains the sequence in which electron shells are filled in atoms. 

The Hartree method does not take into account the PEP, an improved self-consistent field was given by the Hartree-Fock method, which was introduced by V.A. Fock in 1930 (ref). He used the SD instead of just the product of single-particle wavefunctions. We use this method to find the single-particle states that gives us minimum in energy.

Hartree-Fock method gives contribution for pairs of particles with opposite spins, but not the case of when the pairs of the spins are parallel? (Attraction?). A generalization of the Hartree-Fock method that take this into account aswell was developed by N. N. Bogoliubov in 1958 (ref?). He used this in the theory of superconductivity and the theory of heavy nuclei.

Historically, the concept of a self-consistent field was introduced in 1907 by the French physiscist P. Weiss \cite{weiss}. He found that the magnetic moment of each atom in a ferromagnetic material is proportional to magnetic moment and is thus \emph{self-consistent}. %(Vi har en partikkel som ser et gjennomsnittlig felt dannet av andre partikler, men som er med p� � skape dette feltet.)

The interaction, averaged in a certain way, between a given particle and the other particles of a quantum-mechanical system consisting of many particles. Because the problem of many interacting particles is very complex and has no exact solution, calculations are done by approximate methods. One of the most often used approximate methods of quantum mechanics is based on the introduction of a self-consistent field, which permits the many-particle problem to be reduced to the problem of a single particle moving in the average self-consistent field produced by the other particles. 

A concept used to find approximate solutions to the many-body problem in quantum mechanics. The procedure starts with an approximate solution for a particle moving in a single-particle potential, which derives from its average interaction with all the other particles. This average interaction is determined by the wave functions of all the other particles. The equation describing this average interaction is solved and the improved solution obtained is used in the calculation of the interaction term. This procedure is repeated for wave functions until the wave functions and associated energies are not significantly changed in the cycle, self-consistency having been attained. In atomic theory the Hartree-Fock procedure makes use of self-consistent fields. 

We would refer to this as the restricted Hatree-Fock (RHF) method. Which is used in calculations of a \emph{closed-shell} system, i.e. all the levels are filled with electrons. For \emph{open-shell} systems unrestricted Hartree-Fock (UHF) is the method of choice (see \cite{uhf}).

We would like to know how well the HF method preform. Therefore we introduce the term \emph{correlation energy}. 

\begin{equation}
E_{corr} = E_{exact} - E_{HF}
 \label{def:correlation energy}
\end{equation}

It is the difference between the exact nonrelativistic energy (i.e. the FCI energy) and the Hartree-Focl energy. This gives us some indication of how much the electron-electron interaction contributes to our system.

\begin{figure}
\centering
  \scalebox{0.75}{
    \input{electron-correlation.pdftex_t}
}
\caption{A schematic overview over different energies for different approximations}
 \label{fig:Electron correlation energy}
\end{figure}


\section{Outline of the Algorithm}

\begin{figure}
\centering
\includegraphics{HF-flowchart.pdf}
\caption{Flowchart of Hartree-Fock Algorithm. Self-consistency is the condition that $E[\Phi_{new}] - E[\Phi_{old}] = 0$ }
 \label{fig:HF flow chart}
\end{figure}


%%%%%%%%%%%%%%%%%%%%%%%%%%%%%%%%%%%%%%%%%%% Comments %%%%%%%%%%%%%%%%%%%%%%%%%%%%%%%%%%%%%%

% Thjiissen:
% Hartree-Fock method a variational method? We have a SD.
% Individual one-electron wave functions are called \emph{orbitals}.
% The Schr\"{o}dinger equation have a potential determined by the electrons. A coupling between the orbitals via the potential makes the equation (SL?) nonlinear in the orbitals?. 
% 
% The solution to this \emph{nonlinear in the orbitals} is found iteratively in a self-consistency procedure??? 
% 
% The Hartree-Fock method is close the mean-field approach in Statmech???
% 
% Hartree-Fock method is very popular among chemists, and it has also been applied to solids. [See referanser...]
% 
% Hartree-Fock method used in fermionic systems, 
% 
% The Coulomb repulsion between the electrons is represented in an averaged way. 
% The Pauli Principle is included in the HF-theory.
% 
% Even with the position of the nuclei kept fixed (Born-Oppenheimer approximation). But still we have problems with too many degrees of freedom. Especially the the interaction term makes the problem difficult. (In what way?)
% 
% One common method is to ''uncouple`` the problem in which the interaction of one electron with the remaining ones is incorporated in an averaged way into a potential felt by the electron 
% 
% Def: Spin-orbitals are functions depending on the spatial and spin coordinates of one electron (page 48)
% 
% However the Hamiltonian does not act on the spin-dependent part of the spin-orbital (a quantum state |alpha>). Then spatial and the spin part can be seperated into two products. 
% 
% Hartree potential/Hartree term is something unphysical: it contains a coupling between orbital \emph{k} and itself.
% 
% Hartree derived his equation 4.11 in 1927, it neglects exchange as well as other correlations (which?)
% 
% Fock took the antisymmetry into account, exchange term. 
% 
% It is clear that the nonlinear Hartree-Fock equation, must be solved by a selfconsistency iterative procedure. (Describe) This means that we start with a finite number of states N. A guess. These are used to generate a new Fock operator which is diagonalized again and we get new states. This procedure is repeated until the wanted convergence is achieved.
% 
% CI (Configuration Interaction) is a systematic way of improving upon Hartree-Fock theory.
% 
% The exchange term vanished for orthogonal states, pair of opposite spin do not feel this term.
% 
% The derivation of the Hartree-Fock equation is based on a varitaional calculation for the Schr\"{o}dinger equation. 
% 
% We can use Koopman's theorem to approximate excitation energies from a ground state calculation.
% 
% In closed shells systems, all the levels are occupied by two electrons with opposite spin, but in open-shell system there some levels that are partially filled. 
% 
% Calculation with spin-orbitals paired are called restricted Hartree-Fock (RHF). Those in which all spin-orbitals are allowed to have a different spatial dependence are called unrestricted Hartree-Fock (UHF).
% 
% Magnus: \cite{lohne}  The foundation of HF is the variational principle of Rayleigh and Ritz (RR) (See E.K.U Gross)
% 
% Derived the Hartree Fock Equations
% 
% Patrick: \cite{merlot}
% Quantum Dots can be seen as an many-particle problem. (Show how?)
% 
% Hartree-Fock method is minimazation problem based on a mathematical technique known as the Lagrange multipliers, we want to minimize the energy.

% 
% Assume that we have a hermitian Hamiltonian operator that can be separated into two parts
% 
% \begin{equation}
% \hat{H} = \hat{H}_0 + \lambda\hat{V} 
%  \label{def: two part Hamiltonian}
% \end{equation}
% %
% Where $H_0$ is the zero-order and $\lambda\hat{V}$ the perturbation. $\lambda$ is called an \emph{order parameter} which is used to classify the order of the contributions. We assume that $\Psi_n$ is the eigenstate to (\ref{def: two part Hamiltonian}), and $\Phi_n$ the eigenstate for the zero-order $H_0$. And that they are ortonormal.
% 
% \begin{theorem}
% \emph{(Bartlett)}
% \label{thm:Lagrange}
% If $\Phi_n$ is nondegenerate, then its possible to find the solutions 
% \begin{equation}
%  \lim_{\lambda \rightarrow 0} \Psi_n = \Phi_n, \qquad  \lim_{\lambda \rightarrow 0} E_n = E_n^{(0)}
% \end{equation}
% \end{theorem}
% 
% \subsection{Projection Operators}
% We define the projection operators
% 
% \begin{align}
% \hat{P} &= \ket{\Phi_0}\bra{\Phi_0} \\
% \hat{Q} &= \hat{1} - \hat{P} = \sum_{i=1} \ket{\Phi_i}\bra{\Phi_i}
%  \label{def:projection operators}
% \end{align}
% %
% Rayleigh-Schr\"{o}dinger perturbation theory (RSPT) is extensive order by order. This is related to the fact that RSPT is a true order-by-order expansion, unlike BWPT, which has the infinite-order energy in all denominators \cite{bartlett} ????.
% 
% Time-independent perturbation theory Heinnonen 197 kapt 17
% 
% \subsection{Rayleigh-Schr\"{o}dinger Perturbation Theory}
% \subsection{Brioullin-Wigner Perturbation Theory}
% \subsection{Introduction to Diagrammatic notation}
% % 
% 
% %%%%%%%%%%%%%%%%%%%%%%%%%%%%%%%%%%%%%%%%%%%%%%%%%%%%%%%%%%%%%%%%%%%%%%%%%
% %%%%%%%%%%%%%%%%%%%%%%%%%%%%%% CCSD %%%%%%%%%%%%%%%%%%%%%%%%%%%%%%%%%%%%%
% %%%%%%%%%%%%%%%%%%%%%%%%%%%%%%%%%%%%%%%%%%%%%%%%%%%%%%%%%%%%%%%%%%%%%%%%%
% \chapter{Coupled Cluster Theory}
In this chapter we will introduce the Coupled Cluster Theory \cite{crawford} and derive the Coupled Cluster Equations for Single and Double Exitations (CCSD). 

\section{Introduction}
The theoretical framework of Coupled Cluster Theory was developed in the late 1960s by Coester and K\"{u}mmel. He applied it to problems in nuclear and subnuclear physics. Later it was introduced into quantum chemistry in the 1960s by \v{C}\'i\v{z}ek \cite{cizek1,cizek2}. K\"{u}mmel said that he found it remarkable that a quantum chemist would open an issue of a nuclear physics journal \cite{kummel-ccsd}. Then Monkhorts \cite{monkhorst} developed the CC response theory for calculating the molecular properties. And with the revolution of microprocessors the computational power increased rapidly. People like Pople et al. \cite{pople-ccsd} and Bartlett et al. \cite{bartlett-ccsd} began to look into more realistic systems by developing the spin-orbital CCD programs. Since then people have tried to develop efficient CCSD energy codes, and inclusion of higher excitations in the wavefunction. And there is still onging work with develping methods for open shell calculations (Equation of Motion CCSD) \cite{krylov}.


Later Bishop applied and developed the CCM theory for use in condensed matter physics, electron gas at both high and low densities. (See link) 

The CCM is a numerical method for solving the many-body problem (see section \ref{sec:many-body problem}). It is a popular choice with respect to computability ($\sim 100$ nucleons) (cite). Other \emph{ab intion} many-body methods such as Green's Function Monte Carlo have great precision compared to experiment but have expoential growth in computational need \cite{jensen-ccsd}.

\section{The Coupled Cluster Formalism}
In this section we will discuss the notation that are going to be used and derive the theory behind the CC equations.

\subsection{Reference state and operators}
We will use the particle-hole formalism (see section \ref{sec:particle-hole formalism}). The reference state are defined as

\begin{equation}
\ket{\bf 0} \equiv \ket{\Phi_0}
 \label{def:ccsd reference state}
\end{equation}

Where $\ket{\Phi_0}$ is a SD with single-particle orbitals up to the Fermi-level,

\begin{equation}
 \ket{\Phi_0} = a^{\dagger}_{\alpha_1}a^{\dagger}_{\alpha_2}...a^{\dagger}_{\alpha_F}\ket{0}
\label{def:old reference state}
\end{equation}

Where $\alpha_F$ denotes the quantum number that is on the Fermi-level. annihiliation and creation operators with subscripts $ijk...$ will denote hole states, and the subscripts $abc...$ will denote the particle states. Subscripts that are $pqr...$ are in either of these groups. For the creation and annihilation operators we will use the standard ones $a_\alpha$ and $a_\alpha^{\dagger}$ with respect to the physical vacuum $\ket{0}$. But we will keep the \emph{quasi}-annihilation $b_\alpha$ and creation operator $b_\alpha^{\dagger}$.   (\ref{def:quasi-creation},\ref{def:quasi-annihilation}) in mind when we use Wick's theorem.

Products of annihiliation and creation operators can be used to construct different states $\ket{\Phi^a}$ with respect to our reference state $\ket{\Phi_0}$ as shown in (\ref{def:creation of a hole},\ref{def:creation of a particle}). The subscripts will denote the hole states, and the superscripts will denote the particle states.  

%\begin{figure}
% \label{def:}
%\end{figure}

A general state can be represented by 

\begin{equation}
\ket{\Phi_{ijk...}^{abc...}}
 \label{def:general reference state}
\end{equation}
%
Where total number of particles in a system are given by 
%
\begin{equation}
N = N' + n_p - n_h
 \label{eq:total number of particles}
\end{equation}
%
$N'$ is the number of of particles in the reference state. And $n_p$ is the number of particles, $n_h$ is the number of holes.
We define the single-orbital exitation operator (cluster operator) as 

\begin{equation}
t_i \equiv \sum_a t_i^a a_a^{\dagger} a_i 
 \label{def:single-orbital cluster operator}
\end{equation}
%
And the two-orbital exitation operator  
\begin{equation}
t_{ij} \equiv \frac{1}{2} \sum_a t_{ij}^{ab} a_a^{\dagger}a_b^{\dagger} a_j a_i 
 \label{def:two particle orbital cluster operator}
\end{equation}
%
In genereal the $n$-orbital exitation operator is defined as 
\begin{equation}
t_{ijk...}^{abc...} \equiv \frac{1}{n!} \sum_{\underbrace{\mbox{\scriptsize $abc...$}}_n} t_{ijk...}^{abc...} a_a^{\dagger}a_b^{\dagger}a_c^{\dagger}...a_j a_i a_k...
 \label{def:n-ortbital exitation operator}
\end{equation}
%
$t_{ij..}^{ab..}$ is the $n$-orbital excitation amplitude. The factor $\frac{1}{n!}$ is due to the fact that we have an unrestricted summation over the particle states, and there are $n!$ ways of permuting $n$-particles which give rise to the same final state. 
%
The total excitation amplitdes is the sum of all possible excitations
\begin{align}
\hh{T}_1 &\equiv \sum_i \hat{t_i} \label{def:T_1} \\
\hh{T}_2 &\equiv \frac{1}{2} \sum_{ij} t_{ij}^{ab} \label{def:T_2} \\
& \vdots \qquad \vdots \nonumber \\
\hh{T}_n &\equiv \sum_{\underbrace{\mbox{\scriptsize $ijk...$}}_n} t_{\underbrace{ \mbox{\scriptsize $ijk...$}}_n}^{\overbrace{\mbox{\scriptsize $abc...$}}^{n}}
 \label{def:T_n}
\end{align}

\subsection{The Exponential Ansatz}
Our ansatz is that the exact wavefunction $\ket{\Psi}$ can be written as 

\begin{equation}
\ket{\Psi} = e^{\hat{T}} \ket{\Phi_0}
 \label{def:exponential ansatz}
\end{equation}
%
Where $\hat{T}$ is the total excitation operator
\begin{equation}
\hh{T} \equiv \sum_n^{\infty} \hh{T}_n 
 \label{def:total excitation operator}
\end{equation}
%
and $\ket{\Phi_0}$ is the refernce state. (The exponential ansatz manifest the extensitivity property??? see Bartlett). The motivation for this ansatz aswell???? 

This would be the exact solution to our many-body Schr\"{o}dinger equation and Why??? Proof??? (ref). But we have to truncate this sum and these truncations give rise to errors in the CC calculation. The fewer term we truncate the closer we are to the exact solution. Different CC schemes are determined by the level of truncation,

\begin{align}
\hh{T} &= \hh{T}_1 \qquad \mbox{(CCS)} \\
\hh{T} &= \hh{T}_1 + \hh{T}_2 \qquad \mbox{(CCSD)} \\
\hh{T} &= \hh{T}_1 + \hh{T}_2 + \hh{T}_3 \qquad \mbox{(CCSDT)} \\ 
\hh{T} &= \hh{T}_1 + \hh{T}_2 + \hh{T}_3 + \hh{T}_3 \qquad \mbox{(CCSDTQ)}
 \label{def:CC schemes}
\end{align}
%
We will focus on the CCSD scheme in this thesis.


\section{Energy equation for the CCSD}  

%Natural Truncation of the Exponential Ansatz......... page 22 Crawford

The problem we want to solve is the manybody Sch\"{o}dinger equation defined in Eq. (\ref{def:manybodyschrodinger})
\begin{equation}
\hh{H} \ket{\Psi} = E \ket{\Psi}
 \label{def:CC schrodinger}
\end{equation}
%
With our Exponential ansatz

\begin{equation}
\hh{H} e^{\hh{T}} \ket{\Phi_0} = E e^{\hh{T}}\ket{\Phi_0}
 \label{def:CC schrodinger ansatz}
\end{equation}
%
We can left-multiply to get the energy eigenvalue

\begin{equation}
\bra{\Phi_0}\hh{H}e^{\hh{T}}\ket{\Phi_0} = E
 \label{def:intermediate normalization}
\end{equation}
%
Where $\braket{\Phi_0}{\Phi_{CC}}=1$, because of $\braket{\Phi_0}{\Phi_0}=1$ and $\braket{\Phi_{ijk...}^{abc...}}{\Phi_0}=0$ . This can be taken further by multiplying with a general excited sates 
 
\begin{equation}
\bra{\Phi_{ijk...}^{abc..}}\hh{H}\ket{\Phi_0} = \bra{\Phi_{ijk...}^{abc..}}E e^{\hh{T}}\ket{\Phi_0}
\label{def:intermediate normalization excited} 
\end{equation}
%
This gives us a set of equations involving the excitation amplitude $t_{ijk..}^{abc..}$. And they are non-linear because of $e^{\hh{T}}$. But formally they are solved exactly when $\hat{T}$ is not truncated (ref?????). We want to decouple the energy and the excitation amplitudes. This can be done by multiplying with $\bra{\Phi_0} e^{\hat{T}}$.

\begin{align}
\bra{\Phi_0}e^{-\hh{T}} \hh{H} e^{-\hh{T}} \ket{\Phi_0}&=E
 \label{eq:mathematical foresight 1} \\
\bra{\Phi_{i}^{a}} e^{-\hh{T}} \hh{H} e^{\hh{T}} \ket{\Phi_0} &= 0
 \label{eq:mathematical foresight 2}  \\
\bra{\Phi_{ij}^{ab}} e^{-\hh{T}} \hh{H} e^{\hh{T}} \ket{\Phi_0} &= 0  \label{eq:mathematical foresight 22}\\
\vdots \nonumber\\
  \bra{\Phi_{ijk...}^{abc..}}e^{-\hh{T}} \hh{H} e^{\hh{T}} \ket{\Phi_0} &=0 \label{eq:mathematical foresight 3}
\end{align}
%
Which are the CC equations. The first one Eq. (\ref{eq:mathematical foresight 1}) is the CC energy equation. The next ones are the \emph{amplitude} equations, which we have to solve to get the excitation amplitudes $t_{ijk...}^{abc...}$.

We can write the similarity transformed Hamiltonian as a nested sum of commutators using the Baker-Campbell-Hausdorff relation (see Appendix \ref{dix:derivation of baker-campbell-hausdorff}).

\begin{equation}
e^{-\hh{T}} \hat{H} e^{\hh{T}} =  \hh{H} + \ql[\hh{H}, \hh{T}  \qr] + \frac{1}{2!} \ql[\ql[\hh{H}, \hh{T}  \qr],\hh{T} \qr] + \frac{1}{3!} \ql[ \ql[\ql[\hh{H}, \hh{T}  \qr],\hh{T} \qr],\hh{T} \qr] + ...
 \label{def:BCH formula}
\end{equation}
%
Instead of the Hamiltonian will are going to use the normal-ordered $\hh{H}_N$ to derive the CCSD equations, and truncate such that $\hh{T} = \hh{T}_1 + \hh{T}_2$. We define the similarity transformed normal-ordered  Hamiltonian, $\bb{H} \equiv e^{-\hh{T}} \hh{H}_N e^{\hh{T}}$. %Inserting for the normal-ordered Hamiltonian Eq. (\ref{eq:final normal ordered hamiltonian}) gives
% \begin{align}
% \bb{H} &=  e^{-\hh{T}}\ql[\hh{H}_N + \bra{\Phi_0} \hh{H} \ket{\Phi_0}\qr] e^{\hh{T}} \nonumber \\
% &= e^{-\hh{T}}\hh{H}_Ne^{\hh{T}} +  \bra{\Phi_0} \hh{H} \ket{\Phi_0}
%  \label{def:inserting normal ordered}
% \end{align}
%
Inserting this into Eq. (\ref{eq:mathematical foresight 1})-(\ref{eq:mathematical foresight 22}) 

\begin{align}
\bra{\Phi_0} \bb{H} \ket{\Phi_0} + \bra{\Phi_0}\hh{H}\ket{\Phi_0} &= E
 \label{eq:inserting this into 1} \\
\bra{\Phi_i^a}\bb{H} \ket{\Phi_0} & = 0  \label{eq:inserting this into 2} \\
\bra{\Phi_{ij}^{ab}}\bb{H} \ket{\Phi_0} & = 0  \label{eq:inserting this into 3} 
\end{align}
%
We define the CCSD energy as 

\begin{equation}
E_{\text{CCSD}} = \bra{\Phi_0} \bb{H} \ket{\Phi_0} + E_0
 \label{def:CCSD energy} 
\end{equation}
%
Where $E_0 \equiv \bra{\Phi_0}\hh{H}\ket{\Phi_0}$ is the SCF energy (???). There are two ways to derive the set of energy equations (\ref{eq:inserting this into 1})-(\ref{eq:inserting this into 2}), the diagramatic and the algebraic. We are going to do the latter first. 

\subsection{The Algebraic Approach}
Now we want to do a BCH-expansion on $\bb{H}$
%
\begin{equation}
\bb{H} = \hh{H}_N + \ql[\hh{H}_N, \hh{T}\qr] + \frac{1}{2!}\ql[\ql[\hh{H}_N, \hh{T}\qr], \hh{T}\qr] 
 \label{eq: bb H BCH-expansion}
\end{equation}
%
The expansion goes to infinity but terminates since $\hh{H}_N$ (Eq. \ref{def:second-ordered Hamiltonian2}) is a two-particle operator and therefore can at most de-excite a state that have been twofold excited. This is important because the ortonormality of the reference state $\ket{\Phi_0}$. Therefore we will never get contributions from $\hh{T}^3$ and higher, because $\bra{\Phi_0}\hh{H}_N T_1^3 \ket{\Phi_0} = 0$. This is due to the type of Hamilton operator and not electron numbers. Inserting for $\hh{T} = \hh{T}_1 + \hh{T}_2$
\begin{align}
\bb{H} &= \hh{H}_N + \ql[\hh{H}_N, \hh{T}_1\qr] + \ql[\hh{H}_N, \hh{T}_2\qr] + \frac{1}{2} \ql[\ql[\hh{H}_N, \hh{T}_1\qr] , \hh{T}_1\qr] \nonumber \\ 
&+  \frac{1}{2} \ql[\ql[\hh{H}_N, \hh{T}_2\qr] , \hh{T}_2\qr] + \frac{1}{2}\ql[\ql[\hh{H}_N, \hh{T}_1\qr] , \hh{T}_2\qr] + \frac{1}{2}\ql[\ql[\hh{H}_N, \hh{T}_2\qr] , \hh{T}_1\qr] 
 \label{eq:bar H similarity series}
\end{align}
%
We know that the $\hh{T}_n$ operators commute with each other beacause of commutation relations Eq.  (\ref{def:quasioperator1})-(\ref{def:quasioperator3}). This gives us the following relation

\begin{equation}
\ql[ \ql [\hh{H}_N,\hh{T}_2 \qr],\hh{T}_1 \qr] = \ql[ \ql [\hh{H}_N,\hh{T}_1 \qr],\hh{T}_2 \qr]
 \label{def:commutation relation}
\end{equation}
%
And the following Hamiltonian
\begin{align}
\bb{H} &= \overbrace{\hh{H}_N}^{\textbf{1th term}} + \overbrace{\ql[\hh{H}_N, \hh{T}_1\qr]}^{\textbf{2nd term}}+ \overbrace{\ql[\hh{H}_N, \hh{T}_2\qr]}^{\textbf{3rd term}}+  \overbrace{\frac{1}{2} \ql[\ql[\hh{H}_N, \hh{T}_1\qr] , \hh{T}_1\qr]}^{\textbf{4th term}} \nonumber \\ 
&+  \underbrace{\frac{1}{2}  \ql[\ql[\hh{H}_N, \hh{T}_2\qr] , \hh{T}_2\qr]}_{\textbf{5th term}} + \underbrace{\ql[\ql[\hh{H}_N, \hh{T}_1\qr] , \hh{T}_2\qr]}_{\textbf{6th term}} 
 \label{eq:following hamiltonian}
\end{align}
%
We need to express these commutation relations in a second-quantized form which leads us to the algebraic approach.

\subsubsection*{First term}
The first term gives us the contribution 

\begin{empheq}[box=\fbox]{equation}
\bra{\Phi_0} \hh{H}_N \ket{\Phi_0} = 0 \rightarrow E_{\text{CCSD}}^{(1)}  
 \label{eq:first term}
\end{empheq}


\subsubsection*{Second term}
\begin{equation}
\ql[\hh{H}_N, \hh{T}_1\qr] =  \ql[\hh{F}_N, \hh{T}_1\qr] +  \ql[\hh{V}_N, \hh{T}_1\qr]   
 \label{def:second commutator} 
\end{equation}
%
\begin{align}
\hh{V}_N\hh{T}_1&= \frac{1}{4}\sum_{pqrs}\sum_{ia} \bra{pq}|v|\ket{rs} t_i^a \{a_p^{\dagger}a_q^{\dagger}a_s a_r\}\{a_a^\dagger a_i\}
 \label{def:V_N T_1}
\end{align}
%
There are four operators in the first operator string and two in the second. Using the general Wick's theorem we see that we cannot obtain fully contracted terms and thus does not contribute to the CCSD energy. Contractions within the operator strings gives zero since they are already on a normalordered form. 

\begin{equation}
\bra{\Phi_0}[ \hh{V}_N,\hh{T}_1 ]\ket{\Phi_0} = 0 \rightarrow E_{\text{CCSD}}^{(2)} 
 \label{eq:second term T_1 V_1 energy}
\end{equation}
%
Let us find a secondquantized form of the first commutator 

\begin{align}
\hh{F}_N\hh{T}_1 &= \sum_{pq}\sum_{ia}f_q^p t_i^a \{a_p^{\dagger}a_q\}\{a_a^\dagger a_i\}  \label{def:F_N T_1}\\
\hh{T}_1\hh{F}_N &= \sum_{pq}\sum_{ia}f_q^p t_i^a \{a_a^\dagger a_i\}\{a_p^{\dagger}a_q\}
 \label{def:T_1 F_1}
\end{align}
%
From section \ref{sec:particle-hole formalism} we have the following contraction relations

\begin{align}
\contraction{}{a}{{}^\dagger}{a} 
a_{i}^\dagger a_j &= \delta_{ij} \label{rel:1}\\
\contraction{}{a}{{}_a}{a}            
a_{a} a_b^\dagger &= \delta_{ab} \label{rel:2}\\
\contraction{}{a}{{}^\dagger}{a}
a_{a}^\dagger a_b &= 0           \label{rel:3}\\
\contraction{}{a}{{}^\dagger}{a}
a_{i} a_j^\dagger &= 0           \label{rel:4}
\end{align}
%
Using the generalized form of Wick's theorem from Eq. (\ref{def:generalizednormalorder}) on the operator strings
%
\begin{align}
\contraction{\{a_p^{\dagger}a_q\}\{a_a^{\dagger} a_i\} = \{a_p^{\dagger}a_q a_a^\dagger a_i\} + \{}{a}{{}^{\dagger}a_q a_a^{\dagger} }{a}
\contraction{\{a_p^{\dagger}a_q\}\{a_a^{\dagger} a_i\} = \{a_p^{\dagger}a_q a_a^\dagger a_i\} + \{a_p^{\dagger} a_q a_a^{\dagger} a_i\} + \{a_p^{\dagger} }{a}{{}_{}^\dagger}{a}
\contraction[2ex]{\{a_p^{\dagger}a_q\}\{a_a^{\dagger} a_i\} = \{a_p^{\dagger}a_q a_a^\dagger a_i\} + \{a_p^{\dagger} a_q a_a^{\dagger} a_i\} + \{a_p^{\dagger} a_q a_a^{\dagger} a_i\}  + \{}{a}{{}_{}^\dagger a_q a_a}{a}
\contraction{\{a_p^{\dagger}a_q\}\{a_a^{\dagger} a_i\} = \{a_p^{\dagger}a_q a_a^\dagger a_i\} + \{a_p^{\dagger} a_q a_a^{\dagger} a_i\} + \{a_p^{\dagger} a_q a_a^{\dagger} a_i\}  + \{a_p^{\dagger}}{a}{{}_q}{a}
%
\{a_p^{\dagger}a_q\}\{a_a^{\dagger} a_i\} &= \{a_p^{\dagger}a_q a_a^\dagger a_i\} + \{a_p^{\dagger} a_q a_a^{\dagger} a_i\} + \{a_p^{\dagger} a_q a_a^{\dagger} a_i\}  + \{a_p^{\dagger}a_q a_a^\dagger a_i\} \nonumber \\
%
&= \{a_p^{\dagger}a_q a_a^\dagger a_i\} + \delta_{pi}\{a_q a_a^\dagger\} + \delta_{qa} \{a_p^{\dagger} a_i \} + \delta_{pi}\delta_{qa} \label{eq:generalized wick on F_N T_1} \\
%
\{a_a^{\dagger} a_i\}\{a_p^{\dagger} a_q\} &= \{a_a^{\dagger} a_i a_p^{\dagger} a_q\} = \{a_p^{\dagger} a_q a_a^{\dagger} a_i\} \label{eq:T_1 F_1 contraction}
\end{align}
%
The non-contracted terms cancels and the commutation relation of reads

\begin{equation}
\ql[\hh{F}_N,\hh{T}_1\qr] = \sum_{qia} f_q^i t_i^a \{a_q a_a^{\dagger} \} + \sum_{pia} f_a^p t_i^a \{a_p^{\dagger} a_i\} +  \sum_{ia} f_i^a t_a^i 
 \label{eq:second term F_N T_1 commuation}
\end{equation}
%
And the contribution from the second term to the energy 

\begin{empheq}[box=\fbox]{equation}
\bra{\Phi_0} [\hh{H}_N,\hh{T}_1] \ket{\Phi_0} = \sum_{ia}f_a^i t_i^a \rightarrow E_{\text{CCSD}}^{(2)}  
 \label{eq:second term energy}
\end{empheq}


\subsubsection*{Third term}
\begin{equation}
\ql[\hh{H}_N, \hh{T}_2\qr] =  \ql[\hh{F}_N, \hh{T}_2\qr] +  \ql[\hh{V}_N, \hh{T}_2\qr]   
 \label{def:third commutator} 
\end{equation}
%
Using the same argumentation we had for $[\hh{V}_N,\hh{T}_1]$, we see that $[\hh{F}_N,\hh{T}_2]$ does not contribute to the energy

\begin{align}
\hh{F}_N\hh{T}_2 &= \frac{1}{4} \sum_{abji}\sum_{pq} f_p^q t_{ij}^{ab} \{a_p^\dagger a_q\}\{a_a^{\dagger}a_b^{\dagger}a_j a_i\}
 \label{eq:F_N T_2}
\end{align}

\begin{equation}
\bra{\Phi_0}[\hh{F}_N,\hh{T}_2] \ket{\Phi_0} = 0 \rightarrow E_{\text{CCSD}}
 \label{eq:third term T_1 V_1}
\end{equation}
% 
The second part of the commutation relation of the third term gives

\begin{align}
\hh{V}_N\hh{T}_2 &= \frac{1}{16}\sum_{pqrs}\sum_{ijab} t_{ij}^{ab} \bra{pq}|v|\ket{rs} \{a_p^{\dagger}a_q^\dagger a_s a_r \}\{a_a^\dagger a_b^\dagger a_j a_i\} \\
\hh{V}_2\hh{V}_N &= \frac{1}{16}\sum_{pqrs}\sum_{ijab} t_{ij}^{ab} \bra{pq}|v|\ket{rs} \{a_a^\dagger a_b^\dagger a_j a_i\} \{a_p^{\dagger}a_q^\dagger a_s a_r \}
 \label{eq: third term commuation V_N T_2}
\end{align}
%
Using the generalized Wick's theorem 
%
\begin{align}
\contraction[4ex]{a_p^{\dagger}a_q^\dagger a_s a_r \}\{a_a^\dagger a_b^\dagger a_j a_i\} =   \{a_p^{\dagger}a_q^\dagger a_s a_r a_a^\dagger a_b^\dagger a_j a_i\} + \{a }{a}{{}_p^{\dagger}a_q^\dagger a_s a_r a_a^\dagger a_b^\dagger a_j }{a}
\contraction[3ex]{a_p^{\dagger}a_q^\dagger a_s a_r \}\{a_a^\dagger a_b^\dagger a_j a_i\} =   \{a_p^{\dagger}a_q^\dagger a_s a_r a_a^\dagger a_b^\dagger a_j a_i\} + \{a a_p^{\dagger} }{a}{{}_q^\dagger a_s a_r a_a^\dagger a_b^\dagger}{a}
\contraction[2ex]{a_p^{\dagger}a_q^\dagger a_s a_r \}\{a_a^\dagger a_b^\dagger a_j a_i\} =   \{a_p^{\dagger}a_q^\dagger a_s a_r a_a^\dagger a_b^\dagger a_j a_i\} + \{a a_p^{\dagger} {a}_q^\dagger}{a}{{}_s a_r a_a^\dagger}{a}
\contraction{a_p^{\dagger}a_q^\dagger a_s a_r \}\{a_a^\dagger a_b^\dagger a_j a_i\} =   \{a_p^{\dagger}a_q^\dagger a_s a_r a_a^\dagger a_b^\dagger a_j a_i\} + \{a a_p^{\dagger} {a}_q^\dagger a_s}{a}{{}_r }{a}
%
%
\contraction[4ex]{a_p^{\dagger}a_q^\dagger a_s a_r \}\{a_a^\dagger a_b^\dagger a_j a_i\} =   \{a_p^{\dagger}a_q^\dagger a_s a_r a_a^\dagger a_b^\dagger a_j a_i\} + \{a_p^{\dagger}a_q^\dagger a_s a_r a_a^\dagger a_b^\dagger a_j a_i\} + \{a }{a}{{}_p^{\dagger}a_q^\dagger a_s a_r a_a^\dagger a_b^\dagger a_j }{a}
\contraction[3ex]{a_p^{\dagger}a_q^\dagger a_s a_r \}\{a_a^\dagger a_b^\dagger a_j a_i\} =   \{a_p^{\dagger}a_q^\dagger a_s a_r a_a^\dagger a_b^\dagger a_j a_i\} + \{a_p^{\dagger}a_q^\dagger a_s a_r a_a^\dagger a_b^\dagger a_j a_i\} + \{a a_p^{\dagger} }{a}{{}_q^\dagger a_s a_r a_a^\dagger a_b^\dagger}{a}
\contraction[2ex]{a_p^{\dagger}a_q^\dagger a_s a_r \}\{a_a^\dagger a_b^\dagger a_j a_i\} =   \{a_p^{\dagger}a_q^\dagger a_s a_r a_a^\dagger a_b^\dagger a_j a_i\} + \{a_p^{\dagger}a_q^\dagger a_s a_r a_a^\dagger a_b^\dagger a_j a_i\} + \{a a_p^{\dagger} {a}_q^\dagger}{a}{{}_s a_r}{a}
\contraction{a_p^{\dagger}a_q^\dagger a_s a_r \}\{a_a^\dagger a_b^\dagger a_j a_i\} =   \{a_p^{\dagger}a_q^\dagger a_s a_r a_a^\dagger a_b^\dagger a_j a_i\} + \{a_p^{\dagger}a_q^\dagger a_s a_r a_a^\dagger a_b^\dagger a_j a_i\} + \{a a_p^{\dagger} {a}_q^\dagger a_s}{a}{{}_r a_a^\dagger}{a}
%
%
\{a_p^{\dagger}a_q^\dagger a_s a_r \}\{a_a^\dagger a_b^\dagger a_j a_i\} &=   \{a_p^{\dagger}a_q^\dagger a_s a_r a_a^\dagger a_b^\dagger a_j a_i\} +  \{a_p^{\dagger}a_q^\dagger a_s a_r a_a^\dagger a_b^\dagger a_j a_i\} + \{a_p^{\dagger}a_q^\dagger a_s a_r a_a^\dagger a_b^\dagger a_j a_i\}  \nonumber\\
%
%
\contraction[4ex]{+\{ }{a}{{}_p^{\dagger}a_q^\dagger a_s a_r a_a^\dagger a_b^\dagger }{a}
\contraction[3ex]{+\{ a_p^\dagger }{a}{{}_q^\dagger a_s a_r a_a^\dagger a_b^\dagger a_j}{a}
\contraction[2ex]{+\{ a_p^\dagger a_q^\dagger}{a}{{}_s a_r a_a^\dagger}{a}
\contraction{+\{ a_p^\dagger a_q^\dagger a_s}{a}{{}_r }{a}
%
%
\contraction[4ex]{+\{a_p^{\dagger}a_q^\dagger a_s a_r a_a^\dagger a_b^\dagger a_j a_i\} + \{ }{a}{{}_p^{\dagger}a_q^\dagger a_s a_r a_a^\dagger a_b^\dagger }{a}
\contraction[3ex]{+\{a_p^{\dagger}a_q^\dagger a_s a_r a_a^\dagger a_b^\dagger a_j a_i\}+\{ a_p^\dagger }{a}{{}_q^\dagger a_s a_r a_a^\dagger a_b^\dagger a_j}{a}
\contraction[2ex]{+\{a_p^{\dagger}a_q^\dagger a_s a_r a_a^\dagger a_b^\dagger a_j a_i\}+\{ a_p^\dagger a_q^\dagger}{a}{{}_s a_r}{a}
\contraction{+\{a_p^{\dagger}a_q^\dagger a_s a_r a_a^\dagger a_b^\dagger a_j a_i\}+\{ a_p^\dagger a_q^\dagger a_s}{a}{{}_r a_a^\dagger}{a} 
%
%
&+ \{a_p^{\dagger}a_q^\dagger a_s a_r a_a^\dagger a_b^\dagger a_j a_i\}  + \{a_p^{\dagger}a_q^\dagger a_s a_r a_a^\dagger a_b^\dagger a_j a_i\} + ...
 \label{eq: third term generalized wick} \\
%
% 
&= \{a_p^{\dagger}a_q^\dagger a_s a_r a_a^\dagger a_b^\dagger a_j a_i\} + \dd{pi}\dd{qj}\dd{sb}\dd{ra} - \dd{pi}\dd{qj}\dd{sa}\dd{rb}  \nonumber \\ 
&- \dd{pj}\dd{qi}\dd{sb}\dd{ra} + \dd{pj}\dd{qi}\dd{sa}\dd{rb}  + ... \nonumber \\
%
%
\{a_a^\dagger a_b^\dagger a_j a_i\} \{a_p^{\dagger}a_q^\dagger a_s a_r \}  &= \{a_a^\dagger a_b^\dagger a_j a_i a_p^{\dagger}a_q^\dagger a_s a_r \} = \{a_p^{\dagger}a_q^\dagger a_s a_r a_a^\dagger a_b^\dagger a_j a_i\}
%
%
\end{align}
%
The non-contracted term cancels, but we are only interested in the fully contracted terms
\begin{align}
[\hh{V}_N,\hh{T}_2] &= \frac{1}{16} \sum_{ijab} \ql[\vva{ij}{ab} - \vva{ij}{ba} - \vva{ji}{ab} + \vva{ji}{ba} \qr] t_{ij}^{ab} \nonumber \\
&= \frac{1}{4} \sum_{ijab} \vva{ij}{ab} t_{ij}^{ab}
 \label{eq:commutation V_N T_2 final}
\end{align}
An the contribution to the energy

\begin{empheq}[box=\fbox]{equation}
\bra{\Phi_0} [\hh{H}_N,\hh{T}_2] \ket{\Phi_0} =\frac{1}{4} \sum_{ijab} \vva{ij}{ab}t_{ij}^{ab} \rightarrow E_{\text{CCSD}}^{(3)} 
 \label{eq:third term energy}
\end{empheq}

\subsubsection*{Fourth term}
\begin{equation}
\ql[\ql[\hh{H}_N, \hh{T}_1\qr], \hh{T}_1\qr] =  \ql[\ql[\hh{F}_N, \hh{T}_1\qr], \hh{T}_1\qr] +  \ql[\ql[\hh{V}_N, \hh{T}_1\qr], \hh{T}_1\qr]  
 \label{def:fourth commutator} 
\end{equation}
%
From Eqs. (\ref{def:T_2}) and (\ref{eq:second term F_N T_1 commuation}) we obtain these expressions

\begin{equation}
\ql[\hh{F}_N,\hh{T}_1\qr]\hh{T}_1 = \sum_{qia}\sum_{jb}  f_q^i t_i^a t_j^b \{ a_q a_a^{\dagger}\}\{a_b^{\dagger} a_j\} + \sum_{pia}\sum_{jb}   f_a^p t_i^a t_j^b\{a_p^{\dagger} a_i\} \{a_b^{\dagger} a_j\} +  \sum_{ia}\sum_{jb}  f_i^a t_a^i t_j^b \{a_b^{\dagger} a_j \} 
 \label{eq:[F_N,T_1]T_1}
\end{equation}
\begin{equation}
\hh{T}_1 \ql[\hh{F}_N,\hh{T}_1\qr]= \sum_{qia}\sum_{jb}  f_q^i t_i^a t_j^b \{a_b^{\dagger} a_j\} \{a_q a_a^{\dagger}\}+ \sum_{pia}\sum_{jb}   f_a^p t_i^a t_j^b\{a_b^{\dagger} a_j\} \{a_p^{\dagger} a_i\} +  \sum_{ia}\sum_{jb}  f_i^a t_a^i t_j^b \{a_b^{\dagger} a_j \} 
 \label{eq:T_1[F_N,T_1]}
\end{equation}
%
The last two terms cancels each other in a commutation, and all possible contractions is zero. In Eq.  (\ref{eq:T_1[F_N,T_1]}) the particle creation operator $a_b^\dagger$ gives zero in contraction with a particle or hole operator from the left, because of Eq. (\ref{rel:3}). And a in Eq. (\ref{eq:[F_N,T_1]T_1}) the contractions between a particle and hole operator pairs always gives zero. Therefore there is no contribution to the energy from this part of the fourth term.

\begin{align}
\ql[\hh{V}_N, \hh{T}_1\qr]\hh{T}_1  &=  \frac{1}{4}\sum_{pqrs}\sum_{ia}\sum_{jb} \bra{pq}|v|\ket{rs} t_i^at_j^b \{a_p^{\dagger}a_q^{\dagger}a_s a_r\}\{a_a^\dagger a_i\}\{a_b^\dagger a_j\} \\
&-\frac{1}{4} \sum_{pqrs}\sum_{ia}\sum_{jb} \vva{pq}{rs} t_i^at_j^b \{a_a^\dagger a_i\} \{a_p^{\dagger}a_q^{\dagger}a_s a_r\} \{a_b^\dagger a_j\} \label{eq:[V_N,T_1]T_1 1}\\
%
%
%
\hh{T}_1\ql[\hh{V}_N, \hh{T}_1\qr]  &= \frac{1}{4}\sum_{pqrs}\sum_{ia}\sum_{jb} \bra{pq}|v|\ket{rs} t_i^at_j^b \{a_b^\dagger a_j\} \{a_p^{\dagger}a_q^{\dagger}a_s a_r\}\{a_a^\dagger a_i\}\\ 
&-  \frac{1}{4}\sum_{pqrs}\sum_{ia}\sum_{jb} \bra{pq}|v|\ket{rs} t_i^at_j^b \{a_a^\dagger a_i\}\{a_b^\dagger a_j\} \{a_p^{\dagger}a_q^{\dagger}a_s a_r\}
 \label{eq:[V_N,T_1]T_1 2}
\end{align}
%
We see that the only term that contribute to the expectation value is the first term in Eq. (\ref{eq:[V_N,T_1]T_1 1}). All the other terms have a particle creation operator in the leftmost operator string, and the contrations with this operator therefore leads to zero from Eq. (\ref{rel:3}). Using the generalized Wick's theorem on the first term in Eq. (\ref{eq:[V_N,T_1]T_1 1}) 
%
\begin{align}
\contraction[4ex]{\{a_p^{\dagger}a_q^\dagger a_s a_r \}\{a_a^\dagger a_i\}\{a_b^\dagger a_j\} =  \{ }{a}{{}_p^{\dagger}a_q^\dagger a_s a_r a_a^\dagger a_i a_b^\dagger }{a}
\contraction[3ex]{\{a_p^{\dagger}a_q^\dagger a_s a_r \}\{a_a^\dagger a_i\}\{a_b^\dagger a_j\} = \{a_p^{\dagger}}{a}{{}_q^\dagger a_s a_r a_a^\dagger}{a}
\contraction[2ex]{\{a_p^{\dagger}a_q^\dagger a_s a_r \}\{a_a^\dagger a_i\}\{a_b^\dagger a_j\} = \{ a_p^{\dagger} a_q^\dagger}{a}{a_r a_a^\dagger a_i .}{a}
\contraction{\{a_p^{\dagger}a_q^\dagger a_s a_r \}\{a_a^\dagger a_i\}\{a_b^\dagger a_j\} = \{ a_p^{\dagger} a_q^\dagger a_s}{a}{{}_r}{a}
%
%
\contraction[4ex]{\{a_p^{\dagger}a_q^\dagger a_s a_r \}\{a_a^\dagger a_i\}\{a_b^\dagger a_j\} = \{a_p^{\dagger}a_q^\dagger a_s a_r a_a^\dagger a_i a_b^\dagger a_j\} + \{ }{a}{{}_p^{\dagger}a_q^\dagger a_s a_r a_a^\dagger a_i a_b^\dagger }{a}
\contraction[3ex]{\{a_p^{\dagger}a_q^\dagger a_s a_r \}\{a_a^\dagger a_i\}\{a_b^\dagger a_j\} = \{a_p^{\dagger}a_q^\dagger a_s a_r a_a^\dagger a_i a_b^\dagger a_j\} + \{a_p^{\dagger}}{a}{{}_q^\dagger a_s a_r a_a^\dagger}{a}
\contraction[2ex]{\{a_p^{\dagger}a_q^\dagger a_s a_r \}\{a_a^\dagger a_i\}\{a_b^\dagger a_j\} = \{a_p^{\dagger}a_q^\dagger a_s a_r a_a^\dagger a_i a_b^\dagger a_j\} + \{ a_p^{\dagger} a_q^\dagger}{a}{a_r .}{a}
\contraction{\{a_p^{\dagger}a_q^\dagger a_s a_r \}\{a_a^\dagger a_i\}\{a_b^\dagger a_j\} = \{a_p^{\dagger}a_q^\dagger a_s a_r a_a^\dagger a_i a_b^\dagger a_j\} + \{ a_p^{\dagger} a_q^\dagger a_s}{a}{a_a^\dagger a_i .}{a}
%
%
\{a_p^{\dagger}a_q^\dagger a_s a_r \}\{a_a^\dagger a_i\}\{a_b^\dagger a_j\} &=   \{a_p^{\dagger}a_q^\dagger a_s a_r a_a^\dagger a_i a_b^\dagger a_j\} + \{a_p^{\dagger}a_q^\dagger a_s a_r a_a^\dagger a_i a_b^\dagger a_j\} \\
%
%
\contraction[4ex]{+ \{}{a}{{}_p^{\dagger}a_q^\dagger a_s a_r a_a^\dagger}{a}
\contraction[3ex]{+ \{ a_p^\dagger}{a}{{}_q^\dagger a_s a_r  a_a^\dagger a_i a_b^\dagger}{a}
\contraction[2ex]{+ \{a_p^\dagger a_q^\dagger}{a}{a_r  a_a^\dagger a_i .}{a}
\contraction{+ \{ a_p^\dagger a_q^\dagger a_s}{a}{{}_r}{a}
%
%
\contraction[4ex]{+ \{a_p^{\dagger}a_q^\dagger a_s a_r a_a^\dagger a_i a_b^\dagger a_j\} + \{}{a}{{}_p^{\dagger}a_q^\dagger a_s a_r a_a^\dagger}{a}
\contraction[3ex]{+ \{a_p^{\dagger}a_q^\dagger a_s a_r a_a^\dagger a_i a_b^\dagger a_j\} + \{ a_p^\dagger}{a}{{}_q^\dagger a_s a_r  a_a^\dagger a_i a_b^\dagger}{a}
\contraction[2ex]{+ \{a_p^{\dagger}a_q^\dagger a_s a_r a_a^\dagger a_i a_b^\dagger a_j\} +  \{a_p^\dagger a_q^\dagger}{a}{a_r . }{a}
\contraction{+ \{a_p^{\dagger}a_q^\dagger a_s a_r a_a^\dagger a_i a_b^\dagger a_j\} + \{ a_p^\dagger a_q^\dagger a_s}{a}{{a}_r a_i^\dagger . }{a}
%
%
&+ \{a_p^{\dagger}a_q^\dagger a_s a_r a_a^\dagger a_i a_b^\dagger a_j\} + \{a_p^{\dagger}a_q^\dagger a_s a_r a_a^\dagger a_i a_b^\dagger a_j\} + ...\\
%
%
&= -\dd{pj}\dd{qi}\dd{ra}\dd{sb} + \dd{pj}\dd{qi}\dd{rb}\dd{sa} + \dd{pi}\dd{qj}\dd{ra}\dd{sb} - \dd{pi}\dd{qj}\dd{rb}\dd{sa} + ...
\end{align}
%
The expectation value of the fourth term is then

\begin{align}
\bra{\Phi_0} [[\hh{V}_N,\hh{T}_1 ],\hh{T}_1 ]\ket{\Phi_0} &= \frac{1}{4} \sum_{ia}\sum_{jb} t_i^a t_j^b \ql[\vva{ij}{ba} - \vva{ij}{ab} - \vva{ji}{ba} + \vva{ji}{ab} \qr] \\
&= \sum_{ijab} t_i^a t_j^b\vva{ij}{ab} 
 \label{eq:commutator expectation 4th}
\end{align}
%
The contribution to the energy is then

\begin{empheq}[box=\fbox]{equation}
\bra{\Phi_0} \frac{1}{2} [[\hh{H}_N,\hh{T}_1 ],\hh{T}_1 ] \ket{\Phi_0} =\frac{1}{2} \sum_{ijab} t_i^a t_j^b\vva{ij}{ab}  \rightarrow E_{\text{CCSD}}^{(4)}
 \label{eq:fourth term energy}
\end{empheq}

\subsubsection*{Fifth term}

\begin{equation}
[[\hh{H}_N,\hh{T}_2 ],\hh{T}_2] = [[\hh{F}_N,\hh{T}_2 ],\hh{T}_2] +  [[\hh{V}_N,\hh{T}_2 ],\hh{T}_2] 
 \label{eq:fith term}
\end{equation}
%
The expectation value of $[\hh{F}_N,\hh{T}_2 ]$ is zero Eq. (\ref{eq:F_N T_2}), and we are left with finding $[\hh{V}_N,\hh{T}_2 ]\hh{T}_2$,

\begin{align}
[\hh{V}_N,\hh{T}_2 ]\hh{T}_2 &= \frac{1}{16} \sum_{pqrs}\sum_{abij}\sum_{cdkl}\vva{pq}{rs}t_{ij}^{ab}t_{kl}^{cd} \{a_p^\dagger a_q^\dagger a_s a_r\}\{a_a^\dagger a_b^\dagger a_j a_i\}\{a_c^\dagger a_d^\dagger a_l a_k\} \\
&- \frac{1}{16} \sum_{pqrs}\sum_{abij}\sum_{cdkl}\vva{pq}{rs}t_{ij}^{ab}t_{kl}^{cd} \{a_a^\dagger a_b^\dagger a_j a_i\}\{a_p^\dagger a_q^\dagger a_s a_r\}\{a_c^\dagger a_d^\dagger a_l a_k\} 
\end{align}
%
There are in total four hole-annihilation operators $a_s,a_r,a_k,a_l$ and only two-hole creation operators $a_p^\dagger,a_q^\dagger$, this will not be fully contracted. For the same reason $\hh{T}_2[\hh{V}_N,\hh{T}_2 ]$ will not contribute to the energy.

\begin{empheq}[box=\fbox]{equation}
\bra{\Phi_0} \frac{1}{2} [[\hh{H}_N,\hh{T}_2 ],\hh{T}_2] \ket{\Phi_0} = 0  \rightarrow E_{\text{CCSD}}^{(5)}
 \label{eq:fifth term energy}
\end{empheq}

\subsubsection*{Sixth term}
The mixed term $\hh{H}_N\hh{T}_1\hh{T}_2$ have three annihilation-creation pairs (one from $\hh{T}_1$ and two from $\hh{T}_2$), while $\hh{H}_N$ only have two. This term can therefore never be fully contracted and gives no contribution to the energy

\begin{empheq}[box=\fbox]{equation}
\bra{\Phi_0} \frac{1}{2} [[\hh{H}_N,\hh{T}_1 ],\hh{T}_2 ] \ket{\Phi_0} = 0 \rightarrow E_{\text{CCSD}}^{(6)}
 \label{eq:sixth  term energy}
\end{empheq}
%

\subsubsection{Final term}
The final expression becomes 

\begin{empheq}[box=\fbox]{equation}
E_{\text{CCSD}} - E_0 =  \sum_{ia}f_a^i t_i^a + \frac{1}{4} \sum_{ijab} \vva{ij}{ab}t_{ij}^{ab} + \frac{1}{2} \sum_{ijab} \vva{ij}{ab} t_i^a t_j^b 
 \label{eq:final energy}
\end{empheq}
%
This equation is valid for CCSDT and other schemes. Since the higher terms $\hh{T}_3$ and $\hh{T}_4$ etc. cannot produce fully contracted terms with the two-body Hamiltonian as we have seen. However the higher order excitation operators contributes through the amplitude equations Eqs. (\ref{eq:mathematical foresight 2})-(\ref{eq:mathematical foresight 3}).

\section{Introduction to Coupled Cluster Diagrams}
Even though Wick's theorem have simplyfied the method of evaluating expectation values. However the algebraic approach to the amplitude equations would involve even longer operator strings and be very time consuming. Another method is the digramatic approach popularized by Kucharski and Bartlett \cite{bartlett-kucharski}. We will use their approach to derive the energy equation Eq. (\ref{eq:final energy}) and the amplitude equations.
   The particle-hole formalism is still in use, lines represent a particle or a hole with respect to the reference state $\ket{\Phi_0}$. Lines with the downard arrows represent the hole states. And lines with the upward arrows represent the particle states (see Fig.  \ref{fig:particle-hole representation}).  

\begin{figure}[h!]
\centering
\input{Figurer/Tikz/particle-hole-representation}  %latex'en klarer ikke mellomrom!!!!
\caption{Representation of particle and holes}
 \label{fig:particle-hole representation}
\end{figure}
%
Notice that the covention is to write the rightmost particle/hole first, this have no physical significance since we have a sum over all particles and holes. Its just a phase factor of (-1) difference, $\ket{\Phi_{ij}^{ab}} = -\ket{\Phi_{ji}^{ab}} = \ket{\Phi_{ji}^{ab}}$. Since we already have defined the directions of particles and holes, we want to be consistent with the algebraic ordering of the operators, and have a convecntion that the rightmost operators starts always at the bottom this page which is the reference state $\ket{\Phi_0}$. 

The matrix elements of an interaction with a one-body operator $\bra{b}h\ket{a}$ is represented with a vertex (the black dot), $\bullet---\times$. Where the dashed line indicate our interaction potential.

\begin{figure}[h!]
\centering
\subfigure[$\equiv \bra{b}u\ket{a}\{a_b^\dagger a_a \}$]{
\input{Figurer/Tikz/oneparticle-pp}   %HUSK ALLTID \input() P� .TEX FILER!!!!!
\label{fig:oneparticle-pp}
}
\subfigure[$\equiv \bra{j}u\ket{i}\{a_i^\dagger a_j \}$]{
\input{Figurer/Tikz/oneparticle-hh}
\label{fig:oneparticle-hh}
}
\subfigure[$\equiv \bra{a}u\ket{i}\{a_a^\dagger a_i \}$]{
\input{Figurer/Tikz/oneparticle-ph}
\label{fig:oneparticle-ph}
}
\subfigure[$\equiv \bra{i}u\ket{a}\{a_i^\dagger a_a \}$]{
\input{Figurer/Tikz/oneparticle-hp}
\label{fig:oneparticle-hp}
}

\caption[Optional caption for list of figures]{Diagramatic representation of one-body operators. As we can see from Fig. \ref{fig:one particle operators} the quasi-creation operators lie above the interaction line while the quasi-annihilation operators is below.}
\label{fig:one particle operators}
\end{figure}

\subsection*{Diagram rules part 1}

\begin{quote}
{\bf Rule 1:} \em When expressing operators: Any unlabeled particle/hole-line is summed over
\end{quote}

\begin{quote}
{\bf Rule 2:} \em Incoming lines in a vertex $\bullet---\times$ represent an annihilation operator and correspond to the $\ket{}$ $ket$-part of the matrix element. Outgoing lines represent an creation operator and correspond to the $\bra{}$ $bra$-part of the matrix element. %And the rightmost incoming/outgoing line is always written first, 

\begin{equation}
\bra{\text{leftmost-out ... rightmost-out}}u\ket{\text{leftmost-in ...  rightmost-in}}
\end{equation}

\end{quote}

\begin{quote}
{\bf Rule 3:} \em Lines that follow the same direction can be contracted. And a phase factor of $-1$ is multiplied when contracting two holes
\end{quote} 

\begin{figure}[h!]
\centering
\subfigure[{$\equiv \contraction{}{a}{{}_a}{a}a_a a_b^\dagger$}]{
\input{Figurer/Tikz/contraction-pp}  %HUSK ALLTID \input() P� .TEX FILER!!!!!
\label{fig:contraction-pp}
}
\hspace{4cm}
\subfigure[{$\equiv \contraction{}{a}{{}_i^\dagger}{a} a_i^\dagger a_j$}]{
\input{Figurer/Tikz/contraction-hh}
\label{fig:contraction-hh}
}
\label{fig:one particle contractions}
\caption{Non-zero contractions}
\end{figure}
%
\noindent Our one-body operator $\hat{F}_N$ Eq. (\ref{def:F_N}) is then a sum of all the different diagrams in Fig. \ref{fig:one particle operators} and over all the indices. 

\begin{align}
\hh{F}_N &=& \sum_{ab} f_b^a \{a_a^\dagger a_b \} \quad &+& \sum_{ij} f_j^i \{a_i^\dagger a_j \} \quad &+& \sum_{ia} f_a^i \{a_i^\dagger a_a \}\quad &+& \sum_{ai} f_i^a \{a_a^\dagger a_i \}   \nonumber \\
%%
&\equiv& \begin{matrix}\scalebox{0.6}{\input{Figurer/Tikz/oneparticle-pp-noindex}}\\ \epsilon_1 = 0\end{matrix} \quad &+& \begin{matrix}\scalebox{0.6}{\input{Figurer/Tikz/oneparticle-hh-noindex}} \\ \epsilon_2 = 0 \end{matrix} \quad  &+&
\begin{matrix}\scalebox{0.6}{\input{Figurer/Tikz/oneparticle-hp-noindex}}\\ \\ \\\epsilon_3 = -1\end{matrix} \quad  &+&
\begin{matrix}\scalebox{0.6}{\input{Figurer/Tikz/oneparticle-ph-noindex}}\\ \\ \\ \epsilon_4 = +1\end{matrix} \label{diagram:F_N} \\
\end{align}
%
$\epsilon_n$ is the excitation level of diagram nr $n$ and defined by the difference between the number of quasi-creation operators $\mathcal{N}_c$ and the quasi-annihilation operators $\mathcal{N}_a$ divided by 2.

\begin{equation}
\epsilon = \frac{\mathcal{N}_c - \mathcal{N}_a}{2}
 \label{def:excitation level}
\end{equation}
%
A two-body operator will be represented by two vertices, $\bullet --- \bullet$ and the matrix elements are defined as before by diagram the rule 2. The diagramatic representation of $\hh{V}_N$ then becomes.

\begin{align}
\hh{V}_N &=& \frac{1}{4} \sum_{abcd} \vva{ab}{cd} \{a_a^\dagger a_b^\dagger a_d a_c \} \,\, &+& \frac{1}{4} \sum_{ijkl} \vva{ij}{kl} \{a_i^\dagger a_j^\dagger a_l a_k \} \, &+&  \sum_{iabj} \vva{ia}{bj} \{a_i^\dagger a_a^\dagger a_j a_b \} \nonumber \\
%%
&+& \frac{1}{2} \sum_{aibc} \vva{aj}{bc} \{ a_a^\dagger a_i^\dagger a_c a_b \} \,\,  &+& \frac{1}{2} \sum_{ijka} \vva{ij}{ka} \{ a_i^\dagger a_j^\dagger a_a a_k \} \,\,   &+& \frac{1}{2} \sum_{abci} \vva{ab}{ci} \{ a_a^\dagger a_b^\dagger a_i a_c \} \nonumber \\
%%
&+& \frac{1}{2} \sum_{iakl} \vva{ia}{kj} \{ a_i^\dagger a_a^\dagger a_k a_j \} \,\,   &+& \frac{1}{4} \sum_{abij} \vva{ab}{ij} \{ a_a^\dagger a_b^\dagger a_j a_i \} \,\,  &+& \frac{1}{4} \sum_{ijab} \vva{ij}{ab} \{ a_i^\dagger a_j^\dagger a_b a_a \}  \nonumber \\
%%
&\equiv& \begin{matrix}\scalebox{0.6}{\input{Figurer/Tikz/V1}}\\ \epsilon_1 = 0 \end{matrix} \quad  &+& 
\begin{matrix}\scalebox{0.6}{\input{Figurer/Tikz/V2}}\\ \epsilon_2 = 0 \end{matrix} \quad  &+&
\begin{matrix}\scalebox{0.6}{\input{Figurer/Tikz/V3}}\\ \epsilon_3 = 0 \end{matrix} \quad  \nonumber \\
%%
&+& \begin{matrix}\scalebox{0.6}{\input{Figurer/Tikz/V4}}\\ \epsilon_4 = -1 \end{matrix} \quad  &+& 
\begin{matrix}\scalebox{0.6}{\input{Figurer/Tikz/V5}}\\ \epsilon_5 = -1 \end{matrix} \quad  &+&
\begin{matrix}\scalebox{0.6}{\input{Figurer/Tikz/V6}}\\ \epsilon_6 = +1 \end{matrix} \quad  \nonumber \\
%%
&+& \begin{matrix}\scalebox{0.6}{\input{Figurer/Tikz/V7}}\\ \epsilon_7 = +1 \end{matrix} \quad  &+& 
\begin{matrix}\scalebox{0.6}{\input{Figurer/Tikz/V8}}\\\\\\ \epsilon_8 = +2 \end{matrix} \quad  &+&
\begin{matrix}\scalebox{0.6}{\input{Figurer/Tikz/V9}}\\\\\\ \epsilon_9 = -2 \end{matrix} \quad  \nonumber \\
\label{diagram:V_N}
\end{align}
%
Notice that the $\hh{V}_N$ is a sum over indices $p,q,r,s$ which could either be a particle or hole, we would have $2^4=16$ sums, but some of the sums are equivalent. For instance we have used that 

\begin{align}
\frac{1}{4} \sum_{iabj} \vva{ia}{bj} \{a_i^\dagger a_a^\dagger a_j a_b \} \quad &+& \frac{1}{4} \sum_{iabj} \vva{ai}{bj} \{a_a^\dagger a_i^\dagger a_j a_b \} \quad  &+& \frac{1}{4} \sum_{iabj} \vva{ia}{jb} \{a_i^\dagger a_a^\dagger a_b a_j \} \nonumber \\
%%
&+& \frac{1}{4} \sum_{iabj} \vva{ai}{bj} \{a_a^\dagger a_i^\dagger a_b a_j \} \quad  &=&  \sum_{iabj} \vva{ai}{jb} \{a_i^\dagger a_a^\dagger a_j a_b \}
 \label{proof:V_N term 3} 
\end{align}
%
Which is term 3 in $\hh{V}_N$. Each vertex is unique $A\bullet --- \bullet B$, i.e. we can therefore permute the lines leaving or entering the vertex $A$ or $B$, it will give different diagrams but the same matrix element with a phase factor. For example, the third diagram in $\hh{V}_N$ which is the sum over $\vva{ia}{jb}\{a_i^\dagger a_a^\dagger a_j a_b \}$ can be written in four different ways which is the diagramatic equivalent of Eq. (\ref{proof:V_N term 3}).

\begin{align}
\begin{matrix}\scalebox{0.6}{\input{Figurer/Tikz/V3proof1}}\end{matrix} =
\begin{matrix}\scalebox{0.6}{\input{Figurer/Tikz/V3proof4}}\end{matrix} = -
\begin{matrix}\scalebox{0.6}{\input{Figurer/Tikz/V3proof2}}\end{matrix} = - 
\begin{matrix}\scalebox{0.6}{\input{Figurer/Tikz/V3proof3}}\end{matrix} 
\label{diagram:V_N proof}
\end{align}
%
The diagrams are antisymmetric with respect to permutation of hole and particle pairs. Permutation means placing them in different vertices. 

In addition we have the cluster operator $\hh{T} = \hh{T}_1 + \hh{T}_2$ 

\begin{align}
\hh{T} &=& \sum_{ia} t_i^a \{ a_a^\dagger a_i \} \qquad &+& \frac{1}{4} \sum_{ijab} t_{ij}^{ab} \{ a_a^\dagger a_b^\dagger a_j a_i \} \nonumber \\
&\equiv& \begin{matrix}\scalebox{0.6}{\input{Figurer/Tikz/T1}}\\ \epsilon_{\hh{T}_1} = +1 \end{matrix}  \qquad &+& 
\begin{matrix}\scalebox{0.6}{\input{Figurer/Tikz/T2}}\\ \epsilon_{\hh{T}_2} = +2 \end{matrix}
 \label{diagram:cluster operator}
\end{align}
%
Where the interaction is represented by a solid bar ''\, --- \,``. 

Now that we have established diagrams for the operators, we can now show how the matrix elements of operatos between SDs are represented.

\begin{align}
\bra{\Phi_i^a}\hh{T}_1\ket{\Phi_0} = \quad \begin{matrix}\scalebox{0.8}{\input{Figurer/Tikz/T1matrix1}}\end{matrix} \rightarrow \begin{matrix}\scalebox{0.8}{\input{Figurer/Tikz/T1matrix2}}\end{matrix}
 \label{diagram:T1matrix}
\end{align}
%
As mentioned earlier the diagrams are read from bottom to top. We have the reference state on bottom, which is \emph{white space} followed by an N-body operator with an interaction in the middle and the finally the $bra$-state on top.

\begin{align}
\bra{\Phi_{ij}^{ab}}\hh{T}_1\ket{\Phi_0} = \quad \begin{matrix}\scalebox{0.8}{\input{Figurer/Tikz/T2matrix1}}\end{matrix} \rightarrow \begin{matrix}\scalebox{0.8}{\input{Figurer/Tikz/T2matrix2}}\end{matrix}
 \label{diagram:T2matrix}
\end{align}
%
All the lines in the N-body operator have to be contracted in order to give nonzero matrix elements, additional particle/hole lines which are not involved in the interaction will be written on the right side.

\begin{align}
\bra{\Phi_{ikl}^{acd}}\hh{V}_N\ket{\Phi_{jkl}^{bcd}} = \quad \begin{matrix}\scalebox{0.8}{\input{Figurer/Tikz/V3matrix}}\end{matrix}
 \label{diagram:V3matrix}
\end{align}
%
Diagram 3 in $\hh{V}_N$ is therefore the only nonzero contribution for the given the SD. Diagrams that represents the energy equations can be created in the same way. We must remember our reference state is white space, so we cannot have lines above or below our interaction, i.e. the \emph{excitation level}  have to be zero. This will be the condition that gives us the natural truncation to include only $\hh{T}^2$ diagrams.0 Since $\hh{T}^3$ have at least an excitation level of 3, but $\hh{H}_N$ have at an excitation level between $-2$ and $2$. 

\subsection{Energy Equation}

We can write out the commutators in Eq. (\ref{eq:bar H similarity series}) and only include those terms in which the Hamiltonian is to the left of the cluster operators. The reason for this is because of Wicks theorem, we cannot get a full contraction when we have a  particle creation operator $a_a^\dagger$ to the left, as we have seen an example of in Eqs. (\ref{def:T_1 F_1}) and (\ref{eq:T_1 F_1 contraction}), i.e. $\bra{\Phi_0}\hh{T}_1\hh{H}_N \ket{\Phi_0} = 0$.

\begin{align}
\bb{H}_c = \hh{H}_N + \hh{H}_N\hh{T}_1 + \hh{H}_N\hh{T}_2 + \frac{1}{2} \hh{H}_N \hh{T}_1^2 + \frac{1}{2} \hh{H}_N \hh{T}_2^2 + \hh{H}_N \hh{T}_1 \hh{T}_2
 \label{eq:rewrite bar H similartiy series}
\end{align}
%
The subscript $c$ indicates that we have a \emph{connceted cluster} form of the similarity-transformed Hamiltonian. 

Since both $\hh{T}_2$ and $\hh{T}_1\hh{T}_2$ have excitation levels bigger than $+2$, we would not get contributions from the last two terms.  We are going to consider the first nonzero contribution to the coupled cluster energy

\begin{equation}
E_{CCSD} - E_0 = \bra{\Phi_0} \hh{H}_N \hh{T}_1 \ket{\Phi_0} + \bra{\Phi_0} \hh{H}_N \hh{T}_2   \ket{\Phi_0} + \frac{1}{2} \bra{\Phi_0} \hh{H}_N \hh{T}_1^2 \ket{\Phi_0}   
 \label{def:CCSD energy c}
\end{equation}
%
A $\hh{T}_1$-diagram have an excitation level of $+1$, we have to combine this with a diagram in $\hh{H}_N$ that have have an excitation level of $-1$ to get an nonzero contribution. There are several diagrams in $\hh{F}_N$ and $\hh{V}_N$ that fullfill this criteria, but only diagram 3 in $\hh{F}_N$ that have the reference state on top.

\begin{align}
\bra{\Phi_0} \hh{F}_N \hh{T}_1 \ket{\Phi_0} = \quad \begin{matrix}\epsilon = -1 \\ \scalebox{0.6}{\input{Figurer/Tikz/F_NT_1part1}}\\ \epsilon = +1\end{matrix} \quad \rightarrow \qquad \begin{matrix}\scalebox{0.6}{\input{Figurer/Tikz/F_NT_1part2}} \\ \epsilon = 0\end{matrix}
\label{diagram:F_NT_1}
\end{align}
%
Next we consider $\bra{\Phi_0} \hh{H}_N \hh{T}_2 \ket{\Phi_0}$, here the $\hh{T}_2$ have an excitation level of $+2$, which has to be combined with diagram 9 in $\hh{V}_N$ with an excitation level of $-2$ and reference state on top.

\begin{align}
\bra{\Phi_0} \hh{V}_N \hh{T}_2 \ket{\Phi_0} = \quad \begin{matrix}\epsilon = -2 \\ \scalebox{0.6}{\input{Figurer/Tikz/V_NT_2part1}}\\ \epsilon = +2\end{matrix} \qquad \rightarrow \qquad \begin{matrix}\scalebox{0.6}{\input{Figurer/Tikz/V_NT_2part2}} \\ \epsilon = 0\end{matrix}
\label{diagram:V_NT_2}
\end{align}
%
The only contribution left is $\bra{\Phi_0} \hh{H}_N \hh{T}_1^2 \ket{\Phi_0}$, the excitation level of $\hh{T}_1$ is the same as the sum of each which is $+2$. This can only be coupled to diagram 9 in $\hh{V}_N$. 

\begin{align}
\frac{1}{2} \bra{\Phi_0} \hh{V}_N \hh{T}_1^2 \ket{\Phi_0} = \quad \begin{matrix}\epsilon = -2 \\ \scalebox{0.6}{\input{Figurer/Tikz/V_NT_1T_1part1}}\\ \epsilon = +2\end{matrix} \qquad \rightarrow \qquad \begin{matrix}\scalebox{0.6}{\input{Figurer/Tikz/V_NT_1T_1part2}} \\ \epsilon = 0\end{matrix}
\label{diagram:V_NT_1^2}
\end{align}
%
The factor $1/2$ seem to be very arbritrary, but the Diagrammatic rules are consistent with using Wick's theorem and this factor is taken care of when we imply the diagrammatic rules. The energy equation in the diagramatic form is then

\begin{align}
E_{CCSD} - E_0 = \quad \begin{matrix}\scalebox{0.6}{\input{Figurer/Tikz/CCSDenergypart1}}\end{matrix} \quad + \qquad \begin{matrix}\scalebox{0.6}{\input{Figurer/Tikz/CCSDenergypart2}}\end{matrix} \quad + \quad \begin{matrix}\scalebox{0.6}{\input{Figurer/Tikz/CCSDenergypart3}}\end{matrix}
\label{diagram:E_CCSD}
\end{align}
%
Haven written out the energy equation in a diagramatic form we now want to translate this into the algebraic equations we got from Eq. (\ref{eq:final energy}). First we have to introduce some rules to how diagrams can be interpreted.

\subsection*{Diagram rules part 2}

\begin{quote}
{\bf Rule 4:} \em Label the hole lines with indices \boldmath$ijk$.., and particle lines with indices  \boldmath$abc$..
\end{quote} 

\begin{quote}
{\bf Rule 5:} \em Each interaction line contributes with an matrix element $f_{in}^{out}$ or a amplitude $t_{ijk..}^{abc..}$, it is consistent with \text{\bf Rule 2} 
\end{quote} 

\begin{quote}
{\bf Rule 6:} \em Sum over all incdices that is associated with lines that begin and end at interaction lines (internal lines). 
\end{quote}

\begin{quote}
{\bf Rule 7:} \em The phase of the diagram is \boldmath$(-1)^{(l+h)}$, where \boldmath$h$ is number of hole lines. \boldmath$l$ is the number of \boldmath$loops$ that we have in our diagram. A \boldmath$loop$ is a route along a series of lines that returns to its beginning or begins at one external line and ends at another \cite{crawford}.
\end{quote}

\begin{figure}[h!]
\centering
\subfigure[$l=2$,\,\,$h=2$]{
  \scalebox{0.8}{\input{Figurer/Tikz/loopex1}}  
\label{fig:loop1}
}
\hspace{1cm}
\subfigure[$l=1$,\,\,$h=2$]{
 \scalebox{0.8}{\input{Figurer/Tikz/loopex2}}
\label{fig:loop2}
}
\caption[Optional caption for list of figures]{Examples of some loops}
\label{fig:loop examples}
\end{figure}

\begin{quote}
{\bf Rule 8:} \em Each pair of equivalent vertices and equivalent gives a factor of $\frac{1}{2}$ to be multiplied onto the algebraic expression. Particle/hole-lines that begins and end at the same interactions are  equivalent lines. Equivalent vertices are connected by two $\hh{T}_N$ operators to  $\hh{V}_N$ in the exact same way, i.e. same incoming and outgoing arrows for  
\end{quote}

\begin{quote}
{\bf Rule 9:} \em Each pair of \boldmath$unique$ external hole or particle lines give rise to a permutation operator $P(p,q)$. $Unique$ means that the external hole and particles enters/leaves different interactions lines. This rule makes the total diagram antisymmetric and includes the Pauli Principle which are already inclued in the interaction. 

\begin{equation} 
P(p,q)f(p,q) = f(p,q) - f(q,p)
 \label{def:permutation operator}
\end{equation}


\begin{figure}[h!]
\centering
  \scalebox{0.8}{ \input{Figurer/Tikz/unique-ex1}}  
\caption{\boldmath Here $i$ and $j$ enters two different interactions and they are therefore a $unique$ hole pair. We have to multiply with $P(i,j)$}
\label{fig:unique examples}
\end{figure}
%
\end{quote}
%
Now that we have established the diagram rules we can express them in an algebraic form

\begin{equation}
\begin{matrix}\scalebox{0.6}{\input{Figurer/Tikz/CCSDenergy1}}\end{matrix} = \sum_{ia} f_a^i t_i^a 
 \label{diagram:F_N algebraic}
\end{equation}
%
In this diagram we have two loops and two hole lines, one amplitude, one matrix elements and two internal indicies that have to be summed over which gives us a factor $+1$ in total.
\begin{equation}
\begin{matrix}\scalebox{0.6}{\input{Figurer/Tikz/CCSDenergy2}}\end{matrix} = \frac{1}{4} \sum_{ijab} \vva{ij}{ab} t_i^a 
 \label{diagram:V_NT_2 algebraic}
\end{equation}
%
In this diagram we have two loops and two hole lines, two amplitudes, one matrix element and four internal indicies that have to be summed over. And two pair of equivalent lines that gives us a factor $+\frac{1}{4}$ in total.

\begin{equation}
\begin{matrix}\scalebox{0.6}{\input{Figurer/Tikz/CCSDenergy3}}\end{matrix} = \frac{1}{2} \sum_{ijab} \vva{ij}{ab} t_i^a t_j^b
 \label{diagram:V_NT_1^2 algebraic}
\end{equation}
%
In this diagram we have two loops and two hole lines, two amplitudes, two matrix elements and four internal indicies that have to be summed over. And one pair of equivalent vertices that gives us a factor of $+\frac{1}{2}$
This is the same expressions as Eq. (\ref{eq:final energy}) when we used Wick's theorem.

\section{The Amplitude Equations}
We will use the same procedure for amplitude equations. The \emph{connected cluster} form of the similartiy-transformed Hamiltonian for both the singles and doubles equations

\begin{equation}
\bra{\Phi_i^a} \hh{H}_N \left(1 + \hh{T}_1 + \hh{T}_2 + \frac{1}{2} \hh{T}_1^2 + \hh{T}_1\hh{T}_2 + \frac{1}{3!} \hh{T}_1^3 \right) \ket{\Phi_0} = 0
\label{def:H_c for T_1}
\end{equation}

\begin{equation}
\bra{\Phi_{ij}^{ab}} \hh{H}_N \left(1 + \hh{T}_1 + \hh{T}_2 + \frac{1}{2} \hh{T}_1^2 + \hh{T}_1\hh{T}_2 + \frac{1}{2} \hh{T}_2^2 + \frac{1}{2} \hh{T}_1^2 \hh{T}_2 + \frac{1}{3!} \hh{T}_1^3 + \frac{1}{4!} \hh{T}_1^4 \right) \ket{\Phi_0} = 0
\label{def:H_c for T_2}
\end{equation}

\setlength{\fboxsep}{15pt}
\begin{figure}[h!]
  \centering
  \fbox{$\begin{matrix}\scalebox{0.4}{\input{Figurer/Tikz/S1}} 
& \scalebox{0.4}{\input{Figurer/Tikz/S2a}}
&\scalebox{0.4}{\input{Figurer/Tikz/S2b}}
&\scalebox{0.4}{\input{Figurer/Tikz/S2c}}
&\scalebox{0.4}{\input{Figurer/Tikz/S3a}}
\\S_1 & S_{2a} & S_{2b} & S_{2c} & S_{3a}\\
&\scalebox{0.4}{\input{Figurer/Tikz/S3b}}
&\scalebox{0.4}{\input{Figurer/Tikz/S3c}}
& \scalebox{0.4}{\input{Figurer/Tikz/S4a}}
& \scalebox{0.4}{\input{Figurer/Tikz/S4b}}
\\& S_{3b} & S_{3c} & S_{4a} & S_{4b}\\ 
&
& \scalebox{0.4}{\input{Figurer/Tikz/S4c}}
& \scalebox{0.4}{\input{Figurer/Tikz/S5a}}
& \scalebox{0.4}{\input{Figurer/Tikz/S5b}}
\\& & S_{4c} & S_{5a} & S_{5b}
\\
&
&
&\scalebox{0.4}{\input{Figurer/Tikz/S5c}}
&\scalebox{0.4}{\input{Figurer/Tikz/S6}}
\\& & & S_{5c} & S_{6}
\end{matrix}$}
  \caption{Diagrams of the $\hh{T}_1$ amplitude equation Eq. (\ref{def:H_c for T_1})}
  \label{fig:Diagrams of T1}
\end{figure}


\begin{table}[h!]
\caption{Algebraic expression of $\hh{T}_1$-amplitude diagrams from Fig. \ref{fig:Diagrams of T1}}
\centering
\begin{tabular}{c c c c l}
\hline\hline
Interaction & Contraction & $\epsilon$ & Diagram & Expression \\   
\hline
$\bra{\Phi_i^a}\hh{H}_N\ket{\Phi_0}$         &$\hh{F}_{N,4}$ & $+1$ & $S_1$     & $\quad \qquad f_i^a$\\
\hline
$\bra{\Phi_i^a}\hh{H}_N\hh{T}_1\ket{\Phi_0}$ &$\hh{F}_{N,1}$ & $0$  &$S_{2a}$   & $\begin{aligned}\quad \sum_c f_c^a t_i^c\end{aligned}$\\
$\bra{\Phi_i^a}\hh{H}_N\hh{T}_1\ket{\Phi_0}$ &$\hh{F}_{N,2}$ & $0$  &$S_{2b}$   & $\begin{aligned}\quad \sum_k f_i^k t_k^a\end{aligned}$\\
$\bra{\Phi_i^a}\hh{H}_N\hh{T}_1\ket{\Phi_0}$ &$\hh{V}_{N,3}$ & $0$  &$S_{2c}$   & $\begin{aligned}\quad \sum_{kc} \vva{ka}{ci} t_k^c\end{aligned}$\\
\hline
$\bra{\Phi_i^a}\hh{H}_N\hh{T}_2\ket{\Phi_0}$ &$\hh{F}_{N,3}$ & $-1$  &$S_{3a}$  & $\begin{aligned}\quad \sum_{kc}  f_c^k t_{ik}^{ac}\end{aligned}$\\
$\bra{\Phi_i^a}\hh{H}_N\hh{T}_2\ket{\Phi_0}$ &$\hh{V}_{N,4}$ & $-1$  &$S_{3b}$  & $\begin{aligned}\quad \frac{1}{2}\sum_{kcd}\vva{ka}{cd} t_{ki}^{cd}\end{aligned}$\\
$\bra{\Phi_i^a}\hh{H}_N\hh{T}_2\ket{\Phi_0}$ &$\hh{V}_{N,5}$ & $-1$  &$S_{3c}$  & $\begin{aligned}-\frac{1}{2}\sum_{klc}\vva{kl}{ci} t_{kl}^{ca}\end{aligned}$\\
\hline
$\bra{\Phi_i^a}\frac{1}{2}\hh{H}_N\hh{T}_1^2\ket{\Phi_0}$ &$\hh{F}_{N,3}$ & $-1$ & $S_{4a}$ & $\begin{aligned}-\sum_{kc} f_c^k t_i^c t_k^a \end{aligned}$\\
$\bra{\Phi_i^a}\frac{1}{2}\hh{H}_N\hh{T}_1^2\ket{\Phi_0}$ &$\hh{V}_{N,4}$ & $-1$ & $S_{4b}$ & $\begin{aligned}-\sum_{klc}\vva{kl}{ci} t_{k}^{c} t_{l}^{a}\end{aligned}$\\
$\bra{\Phi_i^a}\frac{1}{2}\hh{H}_N\hh{T}_1^2\ket{\Phi_0}$ &$\hh{V}_{N,5}$ & $-1$ & $S_{4c}$ & $\begin{aligned}\quad \sum_{kcd} \vva{ka}{cd} t_{k}^{c} t_{i}^{d}\end{aligned}$\\
\hline
$\bra{\Phi_i^a}\hh{H}_N\hh{T}_1\hh{T}_2\ket{\Phi_0}$ &$\hh{V}_{N,9}$ & $-2$ & $S_{5a}$ & $\begin{aligned} \quad \sum_{klcd} \vva{kl}{cd} t_{k}^{c} t_{li}^{da} \end{aligned}$\\
$\bra{\Phi_i^a}\hh{H}_N\hh{T}_1\hh{T}_2\ket{\Phi_0}$ &$\hh{V}_{N,9}$ & $-2$ & $S_{5b}$ & $\begin{aligned} -\frac{1}{2}\sum_{klcd} \vva{kl}{cd} t_{kl}^{ca} t_{d}^{i} \end{aligned}$\\
$\bra{\Phi_i^a}\hh{H}_N\hh{T}_1\hh{T}_2\ket{\Phi_0}$ &$\hh{V}_{N,9}$ & $-2$ & $S_{5c}$ & $\begin{aligned} -\frac{1}{2}\sum_{klcd} \vva{kl}{cd} t_{ki}^{cd} t_{l}^{a} \end{aligned}$\\
\hline
$\bra{\Phi_i^a}\frac{1}{3!}\hh{H}_N\hh{T}_1^3\ket{\Phi_0}$       &$\hh{V}_{N,9}$ & $-2$ &$S_{6}$   & $\begin{aligned} -\sum_{klcd} \vva{kl}{cd} t_{k}^{c} t_{i}^{d} t_{l}^{a} \end{aligned}$\\
\hline\hline
\end{tabular}
\label{table:T1 expressions}
\end{table}

\setlength{\fboxsep}{15pt}
\begin{figure}[h!]
  \centering
  \fbox{$\begin{matrix}\scalebox{0.35}{\input{Figurer/Tikz/D1}} 
& \scalebox{0.35}{\input{Figurer/Tikz/D2a}}
&\scalebox{0.35}{\input{Figurer/Tikz/D2b}}
&\scalebox{0.35}{\input{Figurer/Tikz/D2c}}
\\D_1 & D_{2a} & D_{2b} & D_{2c}\\
\scalebox{0.35}{\input{Figurer/Tikz/D2d}} 
&\scalebox{0.35}{\input{Figurer/Tikz/D2e}}
&\scalebox{0.35}{\input{Figurer/Tikz/D3a}}
&\scalebox{0.35}{\input{Figurer/Tikz/D3b}}
\\D_{2d} & D_{2e} & D_{3a} & D_{3b}\\
&\scalebox{0.35}{\input{Figurer/Tikz/D3c}} 
&\scalebox{0.35}{\input{Figurer/Tikz/D3d}}
&
\\& D_{3c} & D_{3d} & \\
\end{matrix}$}
  \caption{Diagrams of the $\hh{T}_2$ amplitude equation Eq. (\ref{def:H_c for T_2}), with $\hh{T}_2$ contractions only (CCD) }
  \label{fig:Diagrams of T2}
\end{figure}


\begin{table}[h!]
\caption{Algebraic expression of $\hh{T}_2$-amplitude diagrams from Fig. \ref{fig:Diagrams of T2}}
\centering
\begin{tabular}{c c c c l}
\hline\hline
Interaction & Contraction & $\epsilon$ & Diagram & \qquad Expression \\   
\hline
$\bra{\Phi_{ij}^{ab}}\hh{H}_N\ket{\Phi_0}$         &$\hh{V}_{N,8}$ & $+2$ & $D_1$     & $\quad \qquad \vva{ab}{ij}$\\
\hline
$\bra{\Phi_{ij}^{ab}}\hh{H}_N\hh{T}_2\ket{\Phi_0}$ &$\hh{F}_{N,1}$ & $0$  &$D_{2a}$   & $\begin{aligned}\quad P(ab) \sum_c f_c^b t_{ij}^{ac}\end{aligned}$\\
$\bra{\Phi_{ij}^{ab}}\hh{H}_N\hh{T}_2\ket{\Phi_0}$ &$\hh{F}_{N,2}$ & $0$  &$D_{2b}$   & $\begin{aligned} - P(ij) \sum_k f_j^k t_{ik}^{ab}\end{aligned}$\\
$\bra{\Phi_{ij}^{ab}}\hh{H}_N\hh{T}_2\ket{\Phi_0}$ &$\hh{V}_{N,1}$ & $0$  &$D_{2c}$   & $\begin{aligned}\quad \frac{1}{2} \sum_{cd} \vva{ab}{cd} t_{ij}^{cd}\end{aligned}$\\
$\bra{\Phi_{ij}^{ab}}\hh{H}_N\hh{T}_2\ket{\Phi_0}$ &$\hh{V}_{N,2}$ & $0$  &$D_{2d}$   & $\begin{aligned}\quad \frac{1}{2} \sum_{kl} \vva{kl}{ij} t_{kl}^{ab}\end{aligned}$\\
$\bra{\Phi_{ij}^{ab}}\hh{H}_N\hh{T}_2\ket{\Phi_0}$ &$\hh{V}_{N,3}$ & $0$  &$D_{2e}$   & $\begin{aligned} P(ij)P(ab) \sum_{kc} \vva{kb}{cj} t_{ik}^{ac}\end{aligned}$\\
\hline
$\bra{\Phi_{ij}^{ab}}\frac{1}{2}\hh{H}_N\hh{T}_2^2\ket{\Phi_0}$ &$\hh{V}_{N,9}$ & $-2$  &$D_{3a}$  & $\begin{aligned}\quad \frac{1}{4} \sum_{klcd} \vva{kl}{cd} t_{ij}^{ac} t_{kl}^{ab}\end{aligned}$\\
$\bra{\Phi_{ij}^{ab}}\frac{1}{2}\hh{H}_N\hh{T}_2^2\ket{\Phi_0}$ &$\hh{V}_{N,9}$ & $-2$  &$D_{3b}$  & $\begin{aligned}P(ij) \sum_{klcd}\vva{kl}{cd} t_{ik}^{ac}t_{jl}^{bd}\end{aligned}$\\
$\bra{\Phi_{ij}^{ab}}\frac{1}{2}\hh{H}_N\hh{T}_2^2\ket{\Phi_0}$ &$\hh{V}_{N,9}$ & $-2$  &$D_{3c}$  & $\begin{aligned}-\frac{1}{2}P(ij)\sum_{klcd}\vva{kl}{cd} t_{ik}^{dc}t_{lj}^{ab}\end{aligned}$\\
$\bra{\Phi_{ij}^{ab}}\frac{1}{2}\hh{H}_N\hh{T}_2^2\ket{\Phi_0}$ &$\hh{V}_{N,9}$ & $-2$  &$D_{3d}$  & $\begin{aligned}-\frac{1}{2}P(ab)\sum_{klcd}\vva{kl}{cd} t_{lk}^{ac}t_{ij}^{db}\end{aligned}$\\
\hline\hline
\end{tabular}
\label{table:T2 expressions}
\end{table}


\setlength{\fboxsep}{15pt}
\begin{figure}[h!]
  \centering
  \fbox{$\begin{matrix}\scalebox{0.35}{\input{Figurer/Tikz/D4a}} 
& \scalebox{0.35}{\input{Figurer/Tikz/D4b}}
&\scalebox{0.35}{\input{Figurer/Tikz/D5a}}
&\scalebox{0.35}{\input{Figurer/Tikz/D5b}}\\
%%
D_{4a} & D_{4b} & D_{5a} & D_{5b}\\
\scalebox{0.35}{\input{Figurer/Tikz/D5c}} 
& \scalebox{0.35}{\input{Figurer/Tikz/D6a}}
&\scalebox{0.35}{\input{Figurer/Tikz/D6b}}
&\scalebox{0.35}{\input{Figurer/Tikz/D6c}}\\
%%
D_{5c} & D_{6a} & D_{6b} & D_{6c}\\
\scalebox{0.35}{\input{Figurer/Tikz/D6d}} 
& \scalebox{0.35}{\input{Figurer/Tikz/D6e}}
&\scalebox{0.35}{\input{Figurer/Tikz/D6f}}
&\scalebox{0.35}{\input{Figurer/Tikz/D6g}}\\
%%
D_{6d} & D_{6e} & D_{6f} & D_{6g}\\
\scalebox{0.35}{\input{Figurer/Tikz/D6h}} 
& \scalebox{0.35}{\input{Figurer/Tikz/D7a}}
&\scalebox{0.35}{\input{Figurer/Tikz/D7b}}
&\scalebox{0.35}{\input{Figurer/Tikz/D7c}}\\
%%
D_{6h} & D_{7a} & D_{7b} & D_{7c}\\
\scalebox{0.35}{\input{Figurer/Tikz/D7d}} 
& \scalebox{0.35}{\input{Figurer/Tikz/D7e}}
&\scalebox{0.35}{\input{Figurer/Tikz/D8a}}
&\scalebox{0.35}{\input{Figurer/Tikz/D8b}}\\
%%
D_{7d} & D_{7e} & D_{8a} & D_{8b}\\
\scalebox{0.35}{\input{Figurer/Tikz/D9}}\\
D_9
\end{matrix}$}
  \caption{Diagrams of the $\hh{T}_2$ amplitude equation Eq. (\ref{def:H_c for T_2}), with $\hh{T}_1+\hh{T}_2$ contractions (CCSD) }
  \label{fig:Diagrams of T1 and T2}
\end{figure}


\begin{table}[h!]
\caption{Algebraic expression of $\hh{T}_1+\hh{T}_2$-amplitude diagrams from Fig. \ref{fig:Diagrams of T1 and T2}}
\centering
\begin{tabular}{c c c c l}
\hline\hline
Interaction & Contraction & $\epsilon$ & Diagram & \qquad Expression \\   
\hline
$\bra{\Phi_{ij}^{ab}}\hh{H}_N\hh{T}_1\ket{\Phi_0}$         &$\hh{V}_{N,6}$ & $+1$ & $D_{4a}$     & $\begin{aligned}P(ij)\sum_{c}  \vva{ab}{cj} t_i^c \end{aligned}$\\
$\bra{\Phi_{ij}^{ab}}\hh{H}_N\hh{T}_1\ket{\Phi_0}$         &$\hh{V}_{N,7}$ & $+1$ & $D_{4b}$     & $\begin{aligned}-P(ab)\sum_{k} \vva{kb}{ij  t_k^a}\end{aligned}$\\
\hline
$\bra{\Phi_{ij}^{ab}}\frac{1}{2}\hh{H}_N\hh{T}_1^2\ket{\Phi_0}$         &$\hh{V}_{N,1}$ & $0$ & $D_{5a}$     & $\begin{aligned}\frac{1}{2} P(ij)\sum_{cd} \vva{ab}{cd} t_i^c t_j^d\end{aligned}$\\
$\bra{\Phi_{ij}^{ab}}\frac{1}{2}\hh{H}_N\hh{T}_1^2\ket{\Phi_0}$         &$\hh{V}_{N,2}$ & $0$ & $D_{5b}$     & $\begin{aligned}\frac{1}{2} P(ab)\sum_{kl} \vva{kl}{ij} t_k^a t_l^b\end{aligned}$\\
$\bra{\Phi_{ij}^{ab}}\frac{1}{2}\hh{H}_N\hh{T}_1^2\ket{\Phi_0}$         &$\hh{V}_{N,3}$ & $0$ & $D_{5c}$     & $\begin{aligned}-P(ij|ab)\sum_{kc}  \vva{kb}{cj} t_i^c t_k^a\end{aligned}$\\
\hline
$\bra{\Phi_{ij}^{ab}}\hh{H}_N\hh{T}_1\hh{T}_2\ket{\Phi_0}$ &$\hh{F}_{N,1}$ & $-1$  &$D_{6a}$  & $\begin{aligned}-P(ij)    \sum_{kc} f_c^k t_i^c t_{kj}^{ab}\end{aligned}$\\
$\bra{\Phi_{ij}^{ab}}\hh{H}_N\hh{T}_1\hh{T}_2\ket{\Phi_0}$ &$\hh{F}_{N,1}$ & $-1$  &$D_{6b}$  & $\begin{aligned}-P(ab)    \sum_{kc} f_c^k t_a^k t_{ij}^{cb}\end{aligned}$\\
$\bra{\Phi_{ij}^{ab}}\hh{H}_N\hh{T}_1\hh{T}_2\ket{\Phi_0}$ &$\hh{V}_{N,4}$ & $-1$  &$D_{6c}$  & $\begin{aligned} P(ij|ab) \sum_{kcd}\vva{kd}{cd} t_k^a t_{ij}^{cd}\end{aligned}$\\
$\bra{\Phi_{ij}^{ab}}\hh{H}_N\hh{T}_1\hh{T}_2\ket{\Phi_0}$ &$\hh{V}_{N,5}$ & $-1$  &$D_{6d}$  & $\begin{aligned}-\frac{1}{2} P(ab) \sum_{kcd} \vva{kb}{cd} t_k^a t_{ij}^{cd}\end{aligned}$\\
$\bra{\Phi_{ij}^{ab}}\hh{H}_N\hh{T}_1\hh{T}_2\ket{\Phi_0}$ &$\hh{V}_{N,4}$ & $-1$  &$D_{6e}$  & $\begin{aligned} P(ab) \sum_{kcd} \vva{ka}{cd} t_k^c t_{ij}^{db}\end{aligned}$\\
$\bra{\Phi_{ij}^{ab}}\hh{H}_N\hh{T}_1\hh{T}_2\ket{\Phi_0}$ &$\hh{V}_{N,5}$ & $-1$  &$D_{6f}$  & $\begin{aligned} -P(ij|ab) \sum_{klc} \vva{kl}{ic} t_i^c t_{kl}^{ab}\end{aligned}$\\
$\bra{\Phi_{ij}^{ab}}\hh{H}_N\hh{T}_1\hh{T}_2\ket{\Phi_0}$ &$\hh{V}_{N,4}$ & $-1$  &$D_{6g}$  & $\begin{aligned}\frac{1}{2} P(ij)\sum_{klc} \vva{kl}{cj} t_i^c t_{kl}^{ab}\end{aligned}$\\
$\bra{\Phi_{ij}^{ab}}\hh{H}_N\hh{T}_1\hh{T}_2\ket{\Phi_0}$ &$\hh{V}_{N,5}$ & $-1$  &$D_{6h}$  & $\begin{aligned}- P(ij)\sum_{klc} \vva{kl}{ci} t_k^c t_{lj}^{ab}\end{aligned}$\\
\hline
$\bra{\Phi_{ij}^{ab}}\frac{1}{2}\hh{H}_N\hh{T}_1^2\hh{T}_2\ket{\Phi_0}$ &$\hh{V}_{N,9}$ & $-2$  &$D_{7a}$  & $\begin{aligned} \frac{1}{4}P(ij)\sum_{klcd} \vva{kl}{cd}t_i^c t_{kl}^{ab} t_j^d t_{lj}^{ab}\end{aligned}$\\
$\bra{\Phi_{ij}^{ab}}\frac{1}{2}\hh{H}_N\hh{T}_1^2\hh{T}_2\ket{\Phi_0}$ &$\hh{V}_{N,9}$ & $-2$  &$D_{7b}$  & $\begin{aligned} \frac{1}{4}P(ab)\sum_{klcd} \vva{kl}{cd}t_k^a t_{ij}^{cd} t_l^b t_{lj}^{ab}\end{aligned}$\\
$\bra{\Phi_{ij}^{ab}}\frac{1}{2}\hh{H}_N\hh{T}_1^2\hh{T}_2\ket{\Phi_0}$ &$\hh{V}_{N,9}$ & $-2$  &$D_{7c}$  & $\begin{aligned} -P(ij|ab)\sum_{klcd} \vva{kl}{cd} t_i^c t_a^k t_{lj}^{db} t_{lj}^{ab}\end{aligned}$\\
$\bra{\Phi_{ij}^{ab}}\frac{1}{2}\hh{H}_N\hh{T}_1^2\hh{T}_2\ket{\Phi_0}$ &$\hh{V}_{N,9}$ & $-2$  &$D_{7d}$  & $\begin{aligned} -P(ij)\sum_{klcd} \vva{kl}{cd} t_k^c t_i^d t_{lj}^{ab} t_{lj}^{ab}\end{aligned}$\\
$\bra{\Phi_{ij}^{ab}}\frac{1}{2}\hh{H}_N\hh{T}_1^2\hh{T}_2\ket{\Phi_0}$ &$\hh{V}_{N,9}$ & $-2$  &$D_{7e}$  & $\begin{aligned} -P(ab)\sum_{klcd} \vva{kl}{cd}t_k^c t_l^a t_{ij}^{db} t_{lj}^{ab}\end{aligned}$\\
\hline
$\bra{\Phi_{ij}^{ab}}\frac{1}{2}\hh{H}_N\hh{T}_1^2\hh{T}_2\ket{\Phi_0}$ &$\hh{V}_{N,9}$ & $-2$  &$D_{8a}$  & $\begin{aligned}\frac{1}{2}P(ij|ab)\sum_{klc} \vva{kl}{cd}t_i^c t_{kl}^{ab} t_j^d t_{lj}^{ab}\end{aligned}$\\
$\bra{\Phi_{ij}^{ab}}\frac{1}{2}\hh{H}_N\hh{T}_1^2\hh{T}_2\ket{\Phi_0}$ &$\hh{V}_{N,9}$ & $-2$  &$D_{8b}$  & $\begin{aligned}\frac{1}{2}P(ij|ab)\sum_{klc} \vva{kl}{cd}t_i^c t_{kl}^{ab} t_j^d t_{lj}^{ab}\end{aligned}$\\
\hline
$\bra{\Phi_{ij}^{ab}}\frac{1}{2}\hh{H}_N\hh{T}_1^2\hh{T}_2\ket{\Phi_0}$ &$\hh{V}_{N,9}$ & $-2$  &$D_{9}$  & $\begin{aligned}\frac{1}{4}P(ij|ab)\sum_{klcd} \vva{kl}{cd}t_i^c t_j^d t_k^a t_l^b t_{lj}^{ab}\end{aligned}$\\
\hline\hline
\end{tabular}
\label{table:T1 and T2 expressions}
\end{table}

\newpage

\section{Concluding Remarks}

\subsection{Issues with the Coupled Cluster Method}
Convergence etc. see \cite{kummel-ccsd}. The CC equations for the energy does not have any variational condition. Resulting in that we could in theory get lower energy than the \emph{exact} when we truncate $\hh{T}$ \cite{crawford}. 

%
One of the fundamental postulates in quantum mechanics states that observables are eigenvalues to an Hermitian operator, but the similartiy transformed Hamiltonian $e^{-\hh{T}}\hh{H}e^{\hh{T}}$ is not Hermitian for any $\hh{T}$. 
 
\begin{equation}
\ql (e^{-\hh{T}} \hh{H} e^{\hh{T}} \qr)^{\dagger} = e^{\hh{T}^{\dagger}} \hh{H} e^{-\hh{T}^{\dagger}}  \neq  e^{-\hh{T}} \hh{H} e^{\hh{T}}
\label{eq:not hermitian}
\end{equation}
%

The similarity transformed Hamiltonian $\hh{H}'$ and $\hh{H}$ same eigenvalues:

\small{
\begin{proof} 
\begin{equation*}
 \ket{\Psi} = e^{\hh{T}} \ket{\Phi_0}  
\end{equation*}
\begin{align*}
 \hh{H}\ket{\Psi} = \hh{H} (e^{\hh{T}} \ket{\Phi_0}) =
 e^{\hh{T}}e^{-\hh{T}}\hh{H}e^{\hh{T}}\ket{\Phi_0} &=e^{\hh{T}}\hh{H}'\ket{\Phi_0} = \lambda (e^{\hh{T}}\ket{\Phi_0}) = \lambda \ket{\Psi}
\end{align*}
 \label{proof:eigenvalues of a similarity transformed Hamiltonian}
\end{proof}
%
$\ket{\Phi_0}$ is the eigenfunctions of $\hh{H}'$ and $\ket{\Psi}$ is the eigenfunction of $\hh{H}$. Both with the same eigenvalues
%
The $\hh{T} = \hh{T}_1$ gives us the Hartree Fock approxiamtion,

\begin{equation}
E_{CCSD} - E_0 = \sum_{ia}f_a^i t_i^a + \frac{1}{2}\sum_{ijab} \vva{ij}{ab}t_i^a t_j^b
 \label{eq:ccs and Hartree Fock}
\end{equation}
%
and the expression for the $\hh{T}_1$-amplitude  
\begin{equation}
f_i^a + \sum_c f_c^a t_i^c + \sum_{kc}\vva{ka}{ci} t_k^c = 0
 \label{def:ccs and Hatree Fock amplitude equations}
\end{equation}


\subsection{Wavefunction Seperability and Size Consistency of the Energy}
In CC theory it is possible to $\Psi = e^{T_X + T_Y} \Phi_0$ and get $E = E_X + E_Y$ no can do in CI with linear terms in the ansatz



% 

% %%%%%%%%%%%%%%%%%%%%%%%%%%%%%%%%%%%%%%%%%%%%%%%%%%%%%%%%%%%%%%%%%%%%%%%%%
% %%%%%%%%%%%%%%%%%%%%%%%%%%%%%% Implementation%%%%%%%%%%%%%%%%%%%%%%%%%%%%
% %%%%%%%%%%%%%%%%%%%%%%%%%%%%%%%%%%%%%%%%%%%%%%%%%%%%%%%%%%%%%%%%%%%%%%%%%
% \part{Implementation and Results}
% \chapter{Implementations}
\begin{quote}
  "There's beauty not just at the dimension of one centimeter; there's also a beauty at smaller dimensions."  (From ``What do you care what other people think?'')
\end{quote}

\begin{quote}
  R. P. Feynman
\end{quote}



\section{Implementation of the Configuration Class}
\label{config}
The code-development made in this thesis is done in collaboration with Marte Hoel J\o rgensen. Before we start getting into the algorithms, we first want to introduce the ''Configuration``-Class or just refer to it as ''Config''. This is the class that is linked both to the Hartree-Fock and The Coupled-Cluster program. It keeps track of our single particle orbital quantum numbers. We choose the harmonic oscillator functions as the basis functions. We have the following mapping of our basis. 
\begin{align*}
  \ket{0} & = \{n=0, m=0, m_s = -0.5\} \\ 
  \ket{1} & = \{n=0, m=0, m_s = 0.5\} \\
  \ket{2} & = \{n=0, m=-1, m_s = -0.5\} \\
  \ket{3} & = \{n=0, m=-1, m_s = 0.5\} \\
  \ket{4} & = \{n=0, m=+1, m_s = -0.5\} \\
  \ket{5} & = \{n=0, m=+1, m_s = 0.5\}   
\end{align*}
%
Our basis is numbered such that the lowest subshells are filled first, and there is a special sequence for doing this. 
%
\begin{figure}[ht]
\centering
\scalebox{0.7}{\input{shellfilling}}
\caption{We fill the states such that the we start with the negative $m$ valued subshells then the positive $m$ ones, and last we label the $m=0$ subshells. The spin down state $m_s$ are always the even numbers.}
\end{figure}
% 
\begin{algorithm}[ht]
\caption{\emph{Algorithm for the mapping}}
\begin{enumerate}
\item Loop over $shellnumber$
\item Loop over $m=-(shellnumber-1);m<=0; m = m+2$
\item If $m \neq 0$ Then fill the $m<0$ first and then $m>0$ subshell
\item If $m = 0$ Fill this subshell
\end{enumerate}
\end{algorithm}
%
From Eq. (\ref{shellnumber}) we know that given shellnumber $R$ and $m$ we can find the radial quantum number $n$ by
%
\begin{equation}
n = \frac{R - 1 - |m|}{2}
\label{radialquantum number}  
\end{equation}
%
Next we want to tabulate pairs of states that have equal total $M = m_1 + m_2$ and $M_s = m_{s_1} + m_{s_2}$. This will be used later for the calculation of the  interaction elements. For example for two shells we have 
 
\begin{table}[H]
\centering 
\begin{tabular}{rrrr}
\toprule
$M$ & $M_s$ & $\alpha$ & State\\
\midrule
 -2 &  0 &  1  &$\ket{3,2}$ \\
 -1 & -1  & 3  &$\ket{2,0}$ \\
 -1 &  0 &  4  &$\ket{3,0}$ \\
 -1 &  0  & 4  &$\ket{2,1}$ \\
 -1 &  1 &  5  &$\ket{3,1}$ \\
  0 & -1 &  6  &$\ket{4,2}$ \\
  0 &  0 &  7  &$\ket{1,0}$ \\
  0 &  0 & 7   &$\ket{5,2}$ \\
  0 &  0 & 7   &$\ket{4,3}$\\
  0 &  1 &  8  &$\ket{5,3}$\\
  1 & -1 &  9  &$\ket{4,0}$ \\
  1 &  0 & 10  &$\ket{5,0}$\\
  1 &  0 & 10  &$\ket{4,1}$\\
  1 &  1 &  11 &$\ket{5,1}$ \\
  2 &  0 & 13  &$\ket{5,4}$ \\
\bottomrule
\end{tabular}
\caption{In this table we have paired orbitals with same $M$ and $M_s$. Where $\alpha \in \{0,...,N\}$ and $N$ is the number of pairs $\{M,M_s\}$. Notice that for $\alpha = 1$ we only have one pair, and for $\alpha = 4$, we have two pairs.}
\label{tab:configpair}
\end{table}
%
The algorithm for this is to loop over a pair of single-particle states and test if they have the same $\alpha$ value, i.e. the same $\{M,M_s\}$-pair. This information is stored in the array \texttt{c\_bc} which is a class member of the config class. 

In this class there are also methods that distinguish between a hole-hole pair \texttt{c\_hh}, particle-particle pair \texttt{c\_pp} and particle-hole pair \texttt{c\_ph} for a given $\alpha$. This becomes useful later when we store the two-particle interaction elements.
 
 Note that the states always are classified such that the biggest spin-orbital quantum numbers are labeled first, but this is not uniquely defined, so that the permutations must be considered as well. This is done in the calculations, where we make sure to permute where it needs to be permuted. 
 
 \section{Implementation of the Restriced Hartree-Fock Method}
In this chapter we shall present the alogrithms involved in the Restriced Hartree-Fock Method (RHF) for a two-dimentional parabolic quantum dot. The restricted Hartree-Fock means that we only consider the closed shell systems, ie all the orbitals in the last shell are filled. We have implemented the Hartree-Fock Scheme presented in Chapter \ref{chapter:hartree-fock}. The Hartree-Fock orbitals are expanded in a Harmonic Oscillator basis. The minima of the Hartree Fock Energy is found iteratively and we use the identiy-matrix as our first guess. A Hartree-Fock calculation is needed to accelerate the convergence of the coupled cluster energy. 
\\
\subsection{Input}
The input parameters of our system are
\begin{itemize}
  \item $nh$ is the number of holes in our systems
  \item $np$ is the number of particles in our systems
  \item $tol$ is the tolerance needed to stop the iteration
  \item $max\_iter$ is the maximum number of iteration
  \item $\bra{\alpha}h\ket{\beta}$ is the one-particle interaction element
  \item $\bra{\alpha\beta}v\ket{\gamma \delta}$ is the two-particle interaction element
\end{itemize}
%
The spin-orbital quantum numbers $\ket{\alpha}$ of our system are enumerated according to the config class. Next we need the single-particle interaction  elements $\bra{\alpha}h\ket{\beta}$, which we get from the config class. And we also need the two-particle interaction elements $\bra{\alpha\beta}v\ket{\gamma\delta}$. Note that the interaction elements $\bra{\alpha\beta}v\ket{\gamma\delta}$ are antisymmetrized.

\begin{equation}
\bra{\alpha \beta} v \ket{\gamma \delta}_{\text{AS}} = \bra{\alpha \beta}v\ket{\gamma \delta} - \bra{\alpha \beta}v\ket{\delta \gamma} 
  \label{antisymmetrized elements}
\end{equation}
%
We will from here on use the notation $\bra{\alpha \beta}v\ket{\gamma \delta}$ for the antisymmetric interaction $\bra{\alpha \beta}v\ket{\gamma \delta}_AS$, if nothing else is specified.


\subsection{Output}
The program prints the Hartree-Fock energy for each iteration. When the energy has converged it use the unitary coefficient-matrix $C_{\alpha \beta}$ to create new single-particle interaction elements $\bra{\alpha'}h\ket{\beta'}$
and two-particle interaction elements $\bra{\alpha'\beta'}v\ket{\gamma'\delta'}$.

\subsection{Validation of the Code}
We validate the the Hartree-Fock code by testing the code for a given situation where we know the answer analytically, i.e. the non-interacting system [ref]. For N = 2,6,12 and 20, the non-interacting energy is $2 \hbar \omega$, $10 \hbar \omega$, $28\hbar \omega$ and $60 \hbar \omega$, respectivly. Howerer, this is merely an indication that the code is correct. In order to be satisfyied we need to get the same results as \cite{lohne},\cite{merlot} and \cite{lohne-article}, for the interacting system. 

\subsection{Code Structure}
The Program is structured into a class, with class methods. This makes the structure of the program flexible and we can modify functions without changing the whole program. The code is tailored for a parabolic quantum dot, but can easily be changed to deal with other quantum systems. As mentioned the Hartree-Fock program is linked up to the config class which establishes the mapping:
\begin{equation}
  \ket{\alpha} \rightarrow \ket{n m m_s}
\end{equation}
And also pairs 
\begin{equation}
  \ket{\alpha \beta} \rightarrow \ket{M M_s}
\end{equation}
Where $M$ is the total angular momentum 
\begin{equation}
M = m_\alpha + m_\beta
  \label{totalmomentum}
\end{equation}
And the $M_s$ is the total spin
\begin{equation}
M_s = m_{s_1} + m_{s_2} 
  \label{spin}
\end{equation}
The single-particle interaction $\bra{\alpha}h\ket{\beta}$ is diagonal for the harmonic oscillator basis
\begin{equation}
\bra{\alpha}h\ket{\beta} = \delta_{\alpha \beta} \epsilon_\alpha
\end{equation}
%
These interaction elements are stored in an array \texttt{sp[n]}, where \texttt{n} is the index for the spin-orbital quantum number. \texttt{sp[n]} is an array from the config class. The two-particle interaction elements have to be feed in externally from a binary file generated by a program \texttt{tabulate.cc}, created by Simen Kvaal, see \cite{Kvaal2009}. The binary file is stored with the first four quantum numbers $\alpha,\beta,\gamma,\delta$ as \texttt{short int} (This is because we operate mostly with at most 3 digit numbers), the fifth number is stored as a double and represent our interaction value. The number of elements are therefore given by 

\begin{lstlisting}
    int number_of_elements = filestat.st_size / (sizeof (double) + 4 * sizeof (short int));
\end{lstlisting}
%
We store all the four quantum numbers in a c++ vector struct \texttt{<string,double>} called \texttt{mymap}. The quantum numbers are stored as a string type in \texttt{mymap}, while the interaction value are stored as double, this can be regared as a dictionary in Python language terms, where the quantum numbers are the \emph{key}, and the interaction is the \emph{value} from the key.  
\begin{lstlisting}
    for (int i = 0; i < number_of_elements; i++) {
        Ifile.read((char*) &q, sizeof (short int));
        Ifile.read((char*) &r, sizeof (short int));
        Ifile.read((char*) &s, sizeof (short int));
        Ifile.read((char*) &t, sizeof (short int));
        Ifile.read((char*) &value, sizeof (double));
        sprintf(key, "%d %d %d %d", q, r, s, t);
        mymap.insert(MapType::value_type(key, value));
        }
\end{lstlisting}
%
The reason we want to store the two-particle interaction in such a way is to be able to search for the elements faster. We know that the Hamiltonian does not change spin nor the angular momentum. Therefore the only contribution from the two-particle elements are when $M$ and $M_s$ are conserved. 

\begin{equation}
  \bra{M M_s}v\ket{M' M_s'} = \delta_{MM',M_s M_s'}
\end{equation}
%
\begin{lstlisting}
    for (int alpha = 0; alpha < Conf->dim_alpha; alpha++) {
        n1_map = 0;
        for (int n1 = 0; n1 < 2 * Conf->n_bc[alpha]; n1 += 2) {
            n2_map = 0;
            for (int n2 = 0; n2 < 2 * Conf->n_bc[alpha]; n2 += 2) {
            	p = Conf->c_bc[alpha][n1];
            	q = Conf->c_bc[alpha][n1 + 1];
            	r = Conf->c_bc[alpha][n2];
            	s = Conf->c_bc[alpha][n2 + 1];
                sprintf(configuration, "%d %d %d %d", p, q, r, s);
                search_value = mymap.find(configuration);
                v_matrix[alpha][n1_map][n2_map] = search_value->second;
                n2_map++;

            }
            n1_map++;
        }
    }
\end{lstlisting}
%
For each \texttt{alpha}, i.e. $\{M, M_s\}$ - pair, we loop over pairs $p,q$ and $r,s$ and store the value in the three-dimensional array \texttt{v\_matrix}, where \texttt{n1\_map} and \texttt{n2\_map} indicate which of the pairs that give contribution to our interaction elements. The number of pairs in each \texttt{alpha} would vary as we see in Table \ref{tab:configpair}. So it is necessary to know the number of pairs for each \texttt{alpha} in advance. This information is stored in the array \texttt{n\_bc} which is a class member of the config class. 

The main Hartree-Fock program is initiated from the constructor of the \texttt{HF-iter} class. 
\begin{lstlisting}
  HFinter::HFinter(int nh, int np, double tol, char *filename, Config *Conf) {
    this->Conf = Conf;
    this->iter = iter;
    this->filename = filename;
    this->tol = tol;
    shellnumber = int(-1 / 2 + sqrt(1 + 4 * (m_nh + m_np)) / 2);
    sp_states = shellnumber * (shellnumber + 1);
    max_mvalue = 2 * (shellnumber - 1);
    dim_alpha = 2 * (max_mvalue) * 3 + 3;

    //Set up the two-particle interaction
     setup_coulombmatrix();
    
    //Start Hartree Fock Iteration  
    setup_hf();
    
    //Create new single-particle and two-particle interaction elements
    setup_coulombmatrix_hf();

} //end Constructor
\end{lstlisting}
%
The philosophy of the Hartree-Fock Method is to diagonalize the Hartree-Fock Hamiltonian
\begin{equation}
H_{\alpha,\beta}^{HF} = \begin{bmatrix} H_{11} & H_{12} & \hdots & H_{1n} \\ H_{21} & H_{22} & \hdots & H_{2n} \\ \vdots & \vdots & \ddots & \vdots \\ H_{n1} & H_{n2} & \hdots & H_{nn} \end{bmatrix},
  \label{HF-Hamiltonian}
\end{equation}
%
where 
\begin{equation}
  H_{\alpha,\beta}^{HF} = \bra{\alpha}h\ket{\beta} + \sum_{a}^N \sum_{\gamma\delta}C_{a\gamma}^* C_{a\delta}^* \bra{\alpha \gamma} v \ket{\beta \delta}_{AS}.
\end{equation}
%
The eigenvalue equation gives 
\begin{equation}
H^{HF} \mathbf{C}_k = \lambda_k \mathbf{C}_k
  \label{HF-eigenvalue}
\end{equation}
%
Where $\mathbf{C}_k$ are our eigenvectors corresponding to the given eigenvalue $\lambda_k$. The eigenvector contains the expansion coefficients of the $k$-th HF-orbital Eq. \ref{def:spin-orbital expantion}. Since the HF-matrix are dependent on the coefficient vectors $\mathbf{C}_k$, the equation is non-linear and can be solved iteratively. 
\begin{equation}
\mathbf{C}_k = \begin{pmatrix}C_{k1} \\ C_{k2} \\ C_{k3} \\ \vdots \\ C_{kn} \end{pmatrix}
 \label{coeff}
\end{equation}
%

There are several LAPACK routines for diagonalization. And we could do this by brute force, but the HF-matrix is very sparse. The two-body operator $v$ and the onebody operator $h$ does not change the spin nor the angular momentum of the harmonic oscillator basis, therefore we can divide the HF-matrix into blocks of $\{m,m_s\}$. That is we can only couple spin-orbitals with the same $m$ and $m_s$-values. For three shells we have the following block matrices

\begin{table}[H]
\begin{minipage}[b]{0.5\linewidth}\centering 
\begin{tabular}{ccc}
\toprule
$m$ & $m_s$ & $Block$\\
\midrule
-2	& -0.5 	 &  $\begin{bmatrix}H_{6,6}\end{bmatrix}$ \\
-2	& 0.5 	 &  $\begin{bmatrix}H_{7,7}\end{bmatrix}$ \\
-1	& -0.5 	 &  $\begin{bmatrix}H_{2,2}\end{bmatrix}$ \\
-1	& 0.5 	 &  $\begin{bmatrix}H_{3,3}\end{bmatrix}$ \\
0	& -0.5 	 &  $\begin{bmatrix}H_{0,0} & H_{0,10} \\ H_{10,0} & H_{10,10} \end{bmatrix}$ \\
0	& 0.5 	 &  $\begin{bmatrix}H_{1,1} & H_{1,11} \\ H_{11,1} & H_{11,11} \end{bmatrix}$ \\
1	& -0.5 	 &  $\begin{bmatrix}H_{4,4}\end{bmatrix}$ \\
1	& 0.5 	 &  $\begin{bmatrix}H_{5,5}\end{bmatrix}$ \\
2	& -0.5 	 &  $\begin{bmatrix}H_{8,8}\end{bmatrix}$ \\
2	& 0.5 	 &  $\begin{bmatrix}H_{9,9}\end{bmatrix}$ \\
\bottomrule
\end{tabular}
\caption{Table of the block we have to diagonalize for three shells, note that the quantum numbers for our HF-basis are indicated as subscript. The corresponding eigenvectors are found and placed in }
\label{tab:blockdiag}
\end{minipage}
\hspace{0.5cm}
\begin{minipage}[b]{0.5\linewidth}
  \begin{tabular}{ccc}
\toprule
$m$ & $m_s$ & $Block$\\
\midrule
-2	& -0.5 	 &  $\begin{bmatrix}C_{6,6}\end{bmatrix}$ \\
-2	& 0.5 	 &  $\begin{bmatrix}C_{7,7}\end{bmatrix}$ \\
-1	& -0.5 	 &  $\begin{bmatrix}C_{2,2}\end{bmatrix}$ \\
-1	& 0.5 	 &  $\begin{bmatrix}C_{3,3}\end{bmatrix}$ \\
0	& -0.5 	 &  $\begin{bmatrix}C_{0,0} & C_{0,10} \\ C_{10,0} & C_{10,10} \end{bmatrix}$ \\
0	& 0.5 	 &  $\begin{bmatrix}C_{1,1} & C_{1,11} \\ C_{11,1} & C_{11,11} \end{bmatrix}$ \\
1	& -0.5 	 &  $\begin{bmatrix}C_{4,4}\end{bmatrix}$ \\
1	& 0.5 	 &  $\begin{bmatrix}C_{5,5}\end{bmatrix}$ \\
2	& -0.5 	 &  $\begin{bmatrix}C_{8,8}\end{bmatrix}$ \\
2	& 0.5 	 &  $\begin{bmatrix}C_{9,9}\end{bmatrix}$ \\
\bottomrule
\end{tabular}
\caption{Table of diagonalized blocks for 3 shells, note that eigenvectors $C_{ka}$ are sorted, so that the lowest eigenvalue is in the leftmost place}
\label{tab:blockdiag2}
\end{minipage}
\end{table}
%
The coefficient matrix is set to equal the identity matrix as our first guess $\mathbf{C}_{m,n} = \delta_{m,n}$. 


\begin{algorithm}
\caption{\emph{Hartree - Fock Algorithm}}
\begin{enumerate}[1.]
\item Calculate $\bra{\alpha}h\ket{\beta}$ and $\bra{\alpha\beta}v\ket{\gamma\delta}$
\item Initialize $C_{n,m} = \delta_{n,m}$
\item Check if $E_{new} - E_{old} < tolerance$
\begin{enumerate}[1.]
\item Calculate the HF matrix.
\item Block-diagonalize the HF matrix and find the eigenvalues and eigenvectors
\item Sort the eigenvalues and eigenvectors such that the lowest is $k=0$ for $\mathbf{C}_k$
\item Calculate the HF energy
\end{enumerate}
\end{enumerate}
Output: HF energy and new interaction elements $\bra{a}h\ket{b}$ and $\bra{ab}v\ket{cd}$
\end{algorithm}
%
The new one-body interaction elements are defined in the new basis as
%
\begin{equation}
 \bra{a}h\ket{b} = \sum_{\alpha}^N\sum_{\beta}^N C_{a\alpha}^*C_{b\beta} \bra{\alpha}h\ket{\beta}
\end{equation}
%
And the two-body interaction elements
%
\begin{equation}
 \bra{ab}v\ket{cd} = \sum_{\alpha}^N\sum_{\beta}^N\sum_{\gamma}^N\sum_{\delta}^N C_{a\alpha}^*C_{b\beta}^* C_{c\gamma}C_{d\delta}  \bra{\alpha}\bra{\beta}v\ket{\gamma\delta}
\end{equation}
%
The Hartree Fock energy is defined by
%
\begin{equation}
  E_{HF} = \sum_{k=1}^{nh} \lambda_k + \frac{1}{2} \sum_{ab}^{nh}\sum_{\alpha\beta\gamma\delta}^N  C_{a\alpha}^*C_{b\beta}^* C_{c\gamma}C_{d\delta}  \bra{\alpha}\bra{\beta}v\ket{\gamma\delta}
\end{equation}
% 
Where $nh$ is the number of hole states in our system.


\section{Implementation of the Coupled Cluster Singles and Doubles}
In the following section we will give a general outline of our CCSD code. It is based on Magnus Pedersen Lohnes originial code \cite{lohne}. We (Yang Min Wang and Marte Hoel Joergensen) have jointly optimized the code so that it in some cases has run 200 times faster. The optimization include rewriting the expressions into matrix multiplication so that we could use OpenMP to make it run in parallel CPUs, and storing the interaction elements in a better way, so that we loop over the configurations that give contribution. It is here the config class becomes useful. In comparison the code in \cite{lohne} used 36 hours to run 10 H.O. shells and 20 particles. While our optimized code are using 4 minutes and 20 seconds! Much of that is the overhead caused by Bliz++, and it's ineffective way of accessing elements in the memory.  

The code still only calculates the energy and the amplitudes $t_i^a$ and $t_{ij}^{ab}$ for a closed shell system. The coupled code have a option to run a Hartree-Fock calculation. If we run the Hartree-Fock calculation it will get new single-particle and two-particle elements for the CCSD input.

\subsection{Input}
The input parameters of our system are

\begin{itemize}
  \item $nh$ is the number of holes in our systems
  \item $np$ is the number of particles in our systems
  \item $tol$ is the tolerance needed to stop the iteration
  \item $max\_iter$ is the maximum number of iteration
  \item $\bra{\alpha}h\ket{\beta}$ is the one-particle interaction element
  \item $\bra{\alpha\beta}v\ket{\gamma \delta}$ is the two-particle interaction element
\end{itemize}
%
The matrix elements $\bra{\alpha}h\ket{\beta}$ and $\bra{\alpha\beta}v\ket{\gamma \delta}$ are tabulated in a binary file with the ordering

\begin{itemize}
  \item $\begin{matrix}\alpha & \beta & \bra{\alpha}h\ket{\beta} \end{matrix}$       
  \item $\begin{matrix}\alpha & \beta & \gamma & \delta & \bra{\alpha\beta}v\ket{\gamma \delta} \end{matrix}$      
\end{itemize}
%
$\alpha,\beta,\gamma,\delta$ are the single-particle quantum numbers mapped according to section \ref{config}

\subsection{Output}
The program prints the CCSD Energy for each iteraction. When the energy has converged it has the option to write the T-amplitudes $t_i^a$ and $t_{ij}^{ab}$ to a textfile.

\subsection{Validation of the Code}
One of the first checks we can do is to run the code for the non-interacting case, where analytical expressions can be obtained. The program reproduce these results. Next we need to validate the with interacting part of the program. 

The CCCSD code should be compared to an exact diagonalization of the Hamiltonian for the two-particle. The ground state energies should be equal since we only have single and double excitation. The exact diagonalization is popularly called Full Configuration Interaction method (FCI) \cite{Kvaal2008}. We will give a brief outline. 

In general the Schr\"{o}dinger equation can be solved as a matrix eigenvalue problem, section \ref{Heisenberg}. 

\begin{equation}
  H\ket{\Psi} = E \ket{\Psi}
  \label{Hket}
\end{equation}
%
Where $\ket{\Psi}$ is a linear combination of slater determinants 
%
\begin{equation}
\ket{\Psi} = \sum_{i=1}^{d} C_i \ket{\Phi_i}
  \label{d}
\end{equation}
%
Since our system is the parabolic dot, we will choose the Harmonic oscillator as our basis for the Slater determinant $\ket{\Phi}$
%
\begin{equation}
  \ket{\Phi} = \frac{1}{\sqrt{N!}} \sum_p (-1)^p \hh{P} \prod_{i=1}^d \ket{\alpha_i}
\end{equation}
%
The basis
\begin{equation}
  \mathcal{B} \equiv \{\ket{\alpha_i}\}_{i=1}^{d}
\end{equation}
%
The Hamiltonian in secondquantized form is given by Eq. 
%
\begin{equation}
\hat{H} = \sum_{ij} \bra{i}\hat{h}\ket{j}a_i^{\dagger} a_j+ \frac{1}{4} \sum_{ijkl} \bra{ij}\hat{v}\ket{kl} a_i^{\dagger} a_j^{\dagger} a_l a_k
\end{equation}
%
In principle $d \rightarrow \infty$ gives the exact energy eigenvalues and eigenvectors for Eq. \ref{Hket}, but we must truncate our model space, yielding an approximation to the eigenfunctions and eigenvalues. The diagonalization and the CCSD should give us the same energy given the same model space.

Using Wick's theorem \ref{def:wicksteorem}, the matrix elements $H_{ij} = \bra{\Phi_i}\hh{H}\ket{\Phi_j}$ can be evaluated in terms of $\bra{\alpha}h\ket{\beta}$ and $\bra{\alpha\beta}v\ket{\gamma \delta}$. The onebody part is defined by

\begin{equation}
\bra{\alpha}h\ket{\beta} = \delta_{\alpha \beta} \epsilon_\alpha,
  \label{onebodyinteraction}
\end{equation}
%
where
\begin{equation}
\epsilon_\alpha = 2n_\alpha + |m_\alpha| + 1
\end{equation}
%
The two-particle interaction elements $\bra{\alpha\beta}v\ket{\gamma \delta}$ are found analytically from \cite{rontani}. We are going to validate the CCSD code for the 2-electron case with 6 of the lowest basis functions $d=6$, i.e. 2 H.O. shells. As in section \ref{config} we have the following mapping

\begin{align*}
  \ket{0} & = \{n=0, m=0, m_s = -0.5\} \\ 
  \ket{1} & = \{n=0, m=0, m_s = 0.5\} \\
  \ket{2} & = \{n=0, m=-1, m_s = -0.5\} \\
  \ket{3} & = \{n=0, m=-1, m_s = 0.5\} \\
  \ket{4} & = \{n=0, m=+1, m_s = -0.5\} \\
  \ket{5} & = \{n=0, m=+1, m_s = 0.5\}   
\end{align*}
%
\begin{figure}[ht]
\centering
\scalebox{0.7}{\input{shellfilling2}}
\end{figure}
%
Since the Coulomb interaction does not depend on spin nor the angular momentum, the only nonzero elements are 
\begin{equation}
\bra{M, M_s} v \ket{M, M_s}
  \label{equation2}
\end{equation}
%
where
%
\begin{align}
  M & = m_\alpha + m_\beta = m_\gamma + m_\delta = 0 \\
  M_s & = m_{s_\alpha} + m_{s_\beta} = m_{s_\gamma} + m_{s_\delta} = 0
 \end{align}
%
Then we can diagonalize for the cases 

\begin{table}[H]
\centering
\begin{tabular}{cccc}
\toprule
$M$ & $M_s$ & $H$ & $E$\\
\midrule
-2  & 0    &  $\begin{bmatrix}\bra{2 3}H\ket{2 3}\end{bmatrix}$   &  $\begin{bmatrix}4.86165\end{bmatrix}$ \\
-1  & -1   &  $\begin{bmatrix}\bra{0 2}H\ket{0 2}\end{bmatrix}$   &  $\begin{bmatrix}3.62666\end{bmatrix}$ \\
-1  & 0    &  $\begin{bmatrix}\bra{0 3}H\ket{0 3} & \bra{0 3}H\ket{1 2}\\\bra{1 2}H\ket{0 3} & \bra{1 2}H\ket{1 2}\end{bmatrix}$   &  $\begin{bmatrix}3.62666 & 4.2533\end{bmatrix}$  \\
-1  & 1    &  $\begin{bmatrix}\bra{1 3}H\ket{1 3}\end{bmatrix}$   &  $\begin{bmatrix}3.6267\end{bmatrix}$ \\
 0  & -1   &  $\begin{bmatrix}\bra{2 4}H\ket{2 4}\end{bmatrix}$   &  $\begin{bmatrix}4.6267\end{bmatrix}$  \\
 0  & 0    &  $\begin{bmatrix}\bra{0 1}H\ket{0 1} & \bra{0 1}H\ket{2 5} & \bra{0 1}H\ket{3 4}\\ \bra{2 5}H\ket{0 1} & \bra{2 5}H\ket{2 5} & \bra{2 5}H\ket{3 4}\\ \bra{3 4}H\ket{0 1} & \bra{3 4}H\ket{2 5} & \bra{3 4}H\ket{3 4}\end{bmatrix}$ &  $\begin{bmatrix}3.1523 & 4.6267 & 5.1976 \end{bmatrix}$ \\
\bottomrule
\end{tabular}
\caption{Table of the blocks we have to diagonlize for two shells, where $E$ are the eigenvalues. The diagonalization for $M>0$ would be exactly the same as $M<0$}
\end{table}
%
In this case the lowest (groundstates) energy was
\begin{equation}
E = 3.1523
  \label{groundstates}
\end{equation}
%
Which our coupled cluster program reproduces, see \cite{lohne}. FCI could also be used to validate other CC schemes like CCSDT. For this we would need 3 particles in order to test the 3-particle-3-hole excitation. Since triples have at max 3-particle-3-hole excitations. 

\subsection{Code Structure}
The code is structured such that \texttt{main.cpp} is linked to three \emph{abstract} classes:
\begin{itemize}
  \item \texttt{CCalgo}: Class containing the CCSD algorithm
  \begin{itemize}
  \item ccsd1: A subclass of \texttt{CCalgo} 
  \end{itemize}
  \item \texttt{Fmatrix}: Class for handling the F-matrix
    \begin{itemize}
  \item f1: A subclass of \texttt{Fmatrix} 
  \end{itemize}
  \item \texttt{Interaction}: Class for handling the interaction elements $\bra{\alpha \beta}v\ket{\gamma \delta}$
      \begin{itemize}
  \item int1: A subclass of \texttt{Interaction} 
  \end{itemize}
\end{itemize}
In addition linked to 3 regular classes:
\begin{itemize}
  \item \texttt{Amplitudes}: Class for calculating the CCSD amplitudes $t_i^a$ and $t_{ij}^{ab}$. 
  \item \texttt{Config}: Class containing the one and two-particle basis map
  \item \texttt{HFiter}: Class for calculating the Hartree-Fock and new interaction elements 
  \item \texttt{DDot}: Class for calculating the interaction elements for our Double Dot
\end{itemize}
%
The reason we do this is to be able to expand the program to handle different physical problems in the future. The advantages of an abstract is that the subclasses can share the main class methods. The \texttt{CCAlgo} class runs the CCSD algorithm and it goes like this.
\begin{algorithm}
\caption{\emph{CCSD Algorithm}}
\begin{enumerate}[1.]
\item Setup the model space: $nh,np$
\item Setup the F-matrix $f_p^q$ and the two-particle interaction elements $\bra{\alpha\beta}v\ket{\gamma\delta}$
\item Setup the reference energy $E_0 = \bra{\Phi_0}\hh{H}\ket{\Phi_0}$
\item Initialize the amplitudes $t_i^a$ and $t_{ij}^{ab}$
\item Check if $E_{new} - E_{old} \leq tolerance$
\begin{enumerate}[1.]
\item Calculate the Intermediates
\item Calculate $t_i^a$-matrix elements
\item Calculate $t_{ij}^{ab}$-matrix elements
\item Calculate the energy $E_{new}$
\item Calculate Set $E_{old} = E_{new}$
\end{enumerate}
\end{enumerate}
Output: HF energy and new interaction elements $\bra{a}h\ket{b}$ and $\bra{ab}v\ket{cd}$
\end{algorithm}
%
The first thing we do in our CCSD code is setting up the interaction elements $\bra{pq}v\ket{rs}$, they are read from file and we want to structure those elements into six arrays: 

 \begin{align*}  
\texttt{ hhhh}  &= \bra{ij}v\ket{kl} \\
\texttt{ phhh} &= \bra{aj}v\ket{kl} \\
\texttt{ pphh} &= \bra{ab}v\ket{kl} \\
\texttt{ phph} &= \bra{aj}v\ket{bl} \\
\texttt{ ppph} &= \bra{ab}v\ket{cl} \\
\texttt{ pppp} &= \bra{ab}v\ket{cd}
\end{align*}
%
Where $ijkl...$ denote the hole states, and $abcd...$ denote the particle states. If we stored the matrix $\bra{pq}v\ket{rs}$ in a four dimentional array we would need $8np^4$ bytes of memory. In the case of 20 shells (420 one-particle states) we would need 232 GB! But many of the matrix-elements are zero and we do not need to store those. Since interaction-elements preserve total $M$ and $M_s$, we can use the Config class to help us find the two-particle pairs that give a nonzero element, viz.
 
\begin{equation}
  \texttt{hhhh[$\alpha$][1.pair][2.pair]} = \bra{\underbrace{\alpha \beta}_{1.pair}} v \ket{\underbrace{\gamma \delta}_{2.pair}}
  \label{takenintoaccount}
\end{equation}
%
where $\alpha$ points to the set $\{M,M_s\}$. This makes it possible to calculate for 20 H.O. shells, and we need about 8GB of memory space for this calculation, which is acceptable. Note that our interaction elements are antisymmetrized

\begin{equation*}
  \bra{pq}v\ket{rs} = -\bra{pq}v\ket{sr}
\end{equation*}
%
This needs to be taken into account, when we use those matrix elements \ref{takenintoaccount}.  
%
The next step is to setup the single particle interaction elements and the F-matrix:

\subsubsection{Fmatrix: CLASS implementation}

Fmatrix is an abstract base class with one subclass f1. This class tabulates the single-particle matrix elements $\bra{\alpha}h\ket{\beta}$ and they are used for calculating the F-matrix.

The sp-elements are provided in a text file, which is read and tabulated in the f1-class function \emph{read\_sp\_energy}. See listing \ref{list:impl:fmatrixread}. 
\begin{figure}[htb!]
\begin{lstlisting}[label={list:impl:fmatrixread},caption={Illustration of the implemented Fmatrix class function read\_sp\_energy. See the text for description.}]
void f1::read_sp_energy(char* filename){
  // open and reading from file
  ifstream file(filename,ios_base::in);
  while(!file.eof()){ 
    // read <bra| = <q|
    file >> q;
    // read |ket> = |r>
    file >> r; 
    // read single-particle energy <q|h_0|r>
    file >> value;
    
    if(q<(nh+np) && r<(nh+np)){
      if(q<nh && r<nh){
	s_hh[q][r] = value;
      }
      else if(q<nh && r>=nh){
	s_hp[q][r-nh] = value;
      }
      else if(q>=nh && r<nh){
	s_ph[q-nh][r] = value;
      }
      else if(q>=nh && r>=nh){
	s_pp[q-nh][r-nh] = value;
      }
    }
  }   
  file.close(); 
}// end read_sp_energy
\end{lstlisting}
\end{figure}
The class stores the sp-elements in different two-dimensional arrays according to the placement of holes and particles in the element. This gives four different arrays
\begin{align}  
{\bf s\_{hh}}&=\langle i|\hat{h}|j\rangle\nonumber\\
{\bf s\_{hp}}&=\langle i|\hat{h}|a\rangle\nonumber\\
{\bf s\_{ph}}&=\langle a|\hat{h}|i\rangle\nonumber\\
{\bf s\_{pp}}&=\langle a|\hat{h}|b\rangle
\end{align}

where the subscript $h$ denotes hole states, and the subscript $p$ denotes particle states. Since we have a orthogonal basis of harmonic oscillator functions, we only obtain contribution on the form
\begin{align}  
{\bf s\_{hh}}&=\langle i|\hat{h}|i\rangle\nonumber\\
{\bf s\_{pp}}&=\langle a|\hat{h}|a\rangle\nonumber
\end{align}
However all possibilities are implemented making the code more general. Note that the rescaling of particles numbers encountered in the implementation of the \emph{Interaction} class, is a necessary implementation also for this class, when creating the arrays.  

The f1-class function \emph{set\_up\_fmatrix} calculates the F-matrix defined as
\begin{equation}
 f_{q}^{p}= \langle p|h|q\rangle+\sum_i^d\langle pi|v|qi\rangle
\label{eq:impl:fockmatrix}
\end{equation} 
where $pq\ldots$ denotes both particle and hole states, $i$ denotes hole states, $d$ is the number of hole states, and the interaction elements are antisymmetrized. Listing \ref{list:impl:fmatrixcalc} illustrates the implementation of \emph{set\_up\_fmatrix}.
%
\begin{lstlisting}[label={list:impl:fmatrixcalc},caption={code-snippet illustrating the implementation of the Fmatrix calculation. See the text for description.}]
void f1::set_up_fmatrix(Interaction *interaction) {
    int i, j, k, l; // hole states
    int a, b; // particle states
    
    int n1_map, n2_map;
    
    // set up f_hh = <i|h_0|j> + SUM_k <i k||j k>
    for (i = 0; i < nh; i++) {
        for (j = 0; j < nh; j++) {
            f_hh[i][j] = s_hh[i][j];
        }
    }
    for (int alpha = C->alpha_hh; alpha < C->dim_alpha_hh; alpha++) {
        n1_map = 0;
        for (int n1 = 0; n1 < 2 * C->n_hh[alpha]; n1 += 2) {
            i = C->c_hh[alpha][n1];
            j = C->c_hh[alpha][n1 + 1];
            n2_map = 0;
            for (int n2 = 0; n2 < 2 * C->n_hh[alpha]; n2 += 2) {
                k = C->c_hh[alpha][n2];
                l = C->c_hh[alpha][n2 + 1];
                if (j == l)
                    f_hh[i][k] += interaction->hhhh[alpha][n1_map][n2_map];
                if (i == l)
                    f_hh[j][k] -= interaction->hhhh[alpha][n1_map][n2_map];
                if (j == k)
                    f_hh[i][l] -= interaction->hhhh[alpha][n1_map][n2_map];
                if (i == k)
                    f_hh[j][l] += interaction->hhhh[alpha][n1_map][n2_map];
                n2_map++;
            }
            n1_map++;
        }
    }
    for (int alpha = 0; alpha < C->dim_alpha; alpha++) {
        n1_map = 0;
        for (int n1 = 0; n1 < 2 * C->n_ph[alpha]; n1 += 2) {
            a = C->c_ph[alpha][n1];
            l = C->c_ph[alpha][n1 + 1];
            n2_map = 0;
            for (int n2 = 0; n2 < 2 * C->n_hh[alpha]; n2 += 2) {
                i = C->c_hh[alpha][n2];
                j = C->c_hh[alpha][n2 + 1];
                if (l == j) {
                    f_hp[i][a] += interaction->phhh[alpha][n1_map][n2_map];
                    f_ph[a][i] += interaction->phhh[alpha][n1_map][n2_map];
                } else if (l == i) {
                    f_hp[j][a] = f_hp[j][a] - interaction->phhh[alpha][n1_map][n2_map];
                    f_ph[a][j] = f_ph[a][j] - interaction->phhh[alpha][n1_map][n2_map];
                }
                n2_map++;
            }
            n1_map++;
        }
    }
    
    // set up f_pp = <a|h_0|b> + SUM_k <a k||b k>
    for (a = 0; a < np; a++) {
        f_pp[a][a] = s_pp[a][a];
    }

    for (int alpha = C->alpha_ph; alpha < C->dim_alpha_ph; alpha++) {
        n1_map = 0;
        for (int n1 = 0; n1 < 2 * C->n_ph[alpha]; n1 += 2) {
            a = C->c_ph[alpha][n1];
            i = C->c_ph[alpha][n1 + 1];
            n2_map = 0;
            for (int n2 = 0; n2 < 2 * C->n_ph[alpha]; n2 += 2) {
                b = C->c_ph[alpha][n2];
                j = C->c_ph[alpha][n2 + 1];
                if (i == j)
                    f_pp[a][b] += interaction->phph[alpha][n1_map][n2_map];
                n2_map++;
            }
            n1_map++;
        }
    }

} // end set_up_fmatrix
\end{lstlisting}
The F-matrix is also tabulated according to the placement of holes and particles, thus we obtain four arrays
 \begin{align}  
{\bf f\_{hh}}&=f_{i}^{j}\nonumber\\
{\bf f\_{hp}}&=f_{i}^{a}\nonumber\\
{\bf f\_{ph}}&=f_{a}^{i}\nonumber\\
{\bf f\_{pp}}&=f_{a}^{b}
\label{eq:impl:fmatrices}
\end{align}
In listing \ref{list:impl:fmatrixcalc} the implementation of ${\bf f\_{hh}}$ and ${\bf f\_{pp}}$ is shown explicit. The \texttt{C->} is a pointer to the Config class, and \texttt{interaction->} points to our Interaction class. Next step is to initialize the amplitudes $t_i^a$ and $t_{ij}^{ab}$.
%
\subsubsection{ Amplitudes: CLASS implementation}
\emph{Amplitudes} is an abstract base class with the derived class amp1. This class implements the amplitude equations for both $\hat{T}_1$ and $\hat{T}_2$. It calculates the amplitudes iteratively. In this section we present the implementation of the $\hat{T}_1$ and $\hat{T}_2$ amplitude equations.. 
%
\subsubsection{\bf $\hat{T}_1$-amplitude equation}
Consider the $\hat{T}_1$-amplitude equation, note that the summation-notation is omitted. We rearrange this equation as follows:
\begin{align}
0&=f_i^a+ \langle ja|v|bi \rangle t_{j}^{b}+\frac{1}{2}\langle aj|v|bc\rangle t_{ij}^{bc}+\left( f_{b}^{a}t_{i}^{b}+\langle aj|v|bc\rangle t_i^bt_j^c\right )\nonumber\\
&+\left(-f_{i}^{j}t_{j}^{a}-f_b^jt_j^at_i^b-\langle jk|v|ib\rangle t_{j}^{a}t_{k}^{b}-\langle jk|v|bc\rangle t_i^bt_j^at_k^c-\frac{1}{2}\langle jk|v|bc\rangle t_j^at_{ik}^{bc} \right)\nonumber\\
&+\left(-\frac{1}{2}\langle jk|v|ib\rangle t_{jk}^{ab}-\frac{1}{2}\langle jk|v|bc\rangle t_i^bt_{jk}^{ac} \right)+ \left(f_{b}^{j}t_{ij}^{ab}+\langle jk|v|bc\rangle t_j^bt_{ik}^{ac}\right)
\label{eq:impl:t1amplitude1}
\end{align}
which is equivalent to
\begin{align}
0&=f_i^a+ \langle ja|v|bi \rangle t_{j}^{b}+\frac{1}{2}\langle aj|v|bc\rangle t_{ij}^{bc}+\left( f_{b}^{a}+\langle aj|v|bc\rangle t_j^c\right )t_i^b\nonumber\\
&-\left(f_{i}^{j}+f_b^jt_i^b+\langle jk|v|ib\rangle t_{k}^{b}+\langle jk|v|bc\rangle t_i^bt_k^c+\frac{1}{2}\langle jk|v|bc\rangle t_{ik}^{bc} \right)t_{j}^{a}\nonumber\\
&+\frac{1}{2}\left(\langle jk|v|ic\rangle +\langle jk|v|bc\rangle t_i^b\right)t_{jk}^{ca} + \left(f_{c}^{k}+\langle jk|v|bc\rangle t_j^b\right)t_{ik}^{ac}
\label{eq:impl:t1amplitude2}
\end{align}
In the above equation we have relabeled some of the indices in the last line, in order to extract a common amplitude from the parenthesis. We simplify this expression further by defining the parenthesis in Eq. (\ref{eq:impl:t1amplitude2}) as intermediates. These intermediates are manipulated such that the matrix elements fit the ordering of the six matrices implemented in the \emph{Interaction} class.
\begin{align}
[I1]_b^a&=f_{b}^{a}+\langle aj|v|bc\rangle t_j^c\nonumber\\
&=f_{b}^{a}+\langle bc|v|aj\rangle t_j^c\label{eq:impl:t1intermediates1}\\
\nonumber\\
[I2]_c^k&=f_{c}^{k}+\langle jk|v|bc\rangle t_j^b\nonumber\\
&=f_{c}^{k}+\langle bc|v|jk\rangle t_j^b\label{eq:impl:t1intermediates2}\\
\nonumber\\
[I3]_i^j&=f_{i}^{j}+f_b^jt_i^b+\langle jk|v|ib\rangle t_{k}^{b}+\langle jk|v|bc\rangle t_i^bt_k^c+\frac{1}{2}\langle jk|v|bc\rangle t_{ik}^{bc}\nonumber\\
&=f_{i}^{j}+\langle jk|v|ib\rangle t_{k}^{b}+\frac{1}{2}\langle jk|v|bc\rangle t_{ik}^{bc}+\left(f_b^j+\langle cb|v|kj\rangle t_k^c\right)t_i^b\nonumber\\
&=f_{i}^{j}-\langle bi|v|jk\rangle t_{k}^{b}+\frac{1}{2}\langle bc|v|jk\rangle t_{ik}^{bc}+[I2]_b^j\phantom{.}t_i^b
\label{eq:impl:t1intermediates3}
\end{align}
%
\begin{align}
[I4]_{ic}^{jk}&=\langle jk|v|ic\rangle +\langle jk|v|bc\rangle t_i^b\nonumber\\
&=-\langle ci|v|jk\rangle+\frac{1}{2}\langle bc|v|jk\rangle t_i^b+\frac{1}{2}\langle bc|v|jk\rangle t_i^b\nonumber\\
&=[I5]_{ic}^{jk}+\frac{1}{2}\langle bc|v|jk\rangle t_i^b\label{eq:impl:t1intermediates4}\\
\nonumber\\
[I5]_{ic}^{jk}&=-\langle ci|v|jk\rangle+\frac{1}{2}\langle bc|v|jk\rangle t_i^b\label{eq:impl:t1intermediates5}
\end{align}
We insert these definitions into Eq. (\ref{eq:impl:t1amplitude2}), and obtain the $\hat{T}_1$ amplitude equation:
\begin{align}
0&=f_i^a - \langle aj|v|bi \rangle t_{j}^{b}+\frac{1}{2}\langle bc|v|aj\rangle t_{ij}^{bc}+[I1]_b^a\phantom{.}t_i^b\nonumber\\
&-[I3]_i^j\phantom{.}t_{j}^{a}+\frac{1}{2}[I4]_{ic}^{jk}\phantom{.}t_{jk}^{ca}+[I2]_c^k\phantom{.}t_{ik}^{ac}
\label{eq:impl:t1amplitude3}
\end{align}
Next we wish to obtain an equation for the $\hat{T}_1$ amplitude $t_i^a$. We accomplish this by performing the trick of adding and subtracting the $t_i^a$ amplitude inside eq. \ref{eq:impl:t1amplitude3}. Each term containing a $t_x^y$ amplitude is expressed in terms of the $t_i^a$ amplitude, and added to the equation. This added term is then subtracted by including the original expression multiplied by one or more delta-functions. Thus we obtain
\begin{align}
0&=f_i^a - \langle ai|v|ai \rangle t_{i}^{a}-(1-\delta_{ab}\delta_{ij})\langle aj|v|bi \rangle t_{j}^{b}+\frac{1}{2}\langle bc|v|aj\rangle t_{ij}^{bc}\nonumber\\
&+[I1]_a^a\phantom{.}t_i^a+(1-\delta_{ab})[I1]_b^a\phantom{.}t_i^b-[I3]_i^i\phantom{.}t_{i}^{a}-(1-\delta_{ij})[I3]_i^j\phantom{.}t_{j}^{a}\nonumber\\
&+\frac{1}{2}[I4]_{ic}^{jk}\phantom{.}t_{jk}^{ca}+[I2]_c^k\phantom{.}t_{ik}^{ac}
\label{eq:impl:t1amplitude4}
\end{align}
We are now able to extract an equation for the amplitude $t_i^a$, which reads
\begin{align}
t_i^a&=\frac{1}{D_i^a}\left(f_i^a+\frac{1}{2}\langle bc|v|aj\rangle t_{ij}^{bc}-(1-\delta_{ab}\delta_{ij})\langle aj|v|bi \rangle t_{j}^{b}\right.\nonumber\\
&\left. +(1-\delta_{ab})[I1]_b^a\phantom{.}t_i^b-(1-\delta_{ij})[I3]_i^j\phantom{.}t_{j}^{a}+\frac{1}{2}[I4]_{ic}^{jk}\phantom{.}t_{jk}^{ca}+[I2]_c^k\phantom{.}t_{ik}^{ac}\right)
\label{eq:impl:t1amplitude5}
\end{align} 
where $D_i^a$ is given as
\begin{equation}
D_i^a= \langle ai|v|ai \rangle-[I1]_a^a+[I3]_i^i
\label{eq:impl:t1denom}
\end{equation}
The Eq. (\ref{eq:impl:t1amplitude5}) is implemented in the class function \emph{t1\_uncoupled\_calc}, which is illustrated in Listing \ref{list:impl:calct1}. The implementation of the intermediates will be presented after the derivation of the $\hat{T}_2$ amplitude equation.
%
\begin{lstlisting}[label={list:impl:calct1},caption={Implementation of the amp1 class function t1\_uncoupled\_calc()}]
void Amplitudes::t1_uncoupled_calc(){
  // calculate t1 intermediates
  t1_uncoupled_intermediates();
  
  // f_i^a
  for(i=0; i<np; i++){
    for(j=0; j<nh; j++){
      t1[i][j] = F->f_ph[i][j];
    }
  }

  // -(1-delta_{ab}delta_{ij})<aj|v|bi>t_{j}^{b}
  t1_uncoupled_term2(t1);

  //(1-delta_{ab})[I1]_{b}^{a} t_{i}^{b}
  t1_uncoupled_term3(t1);
  
  // -(1-delta_{ij})[I3]_{i}^{j}t_{j}^{a}
  t1_uncoupled_term4(t1);

  //  0.5<bc|v|aj>t_ij^bc
  t1_uncoupled_term5(t1);

  // 0.5[I4]_{ic}^{jk}t_{jk}^{ca}
  t1_uncoupled_term6(t1);

  // [I2]_{c}^{k}t_{ik}^{ac}
  t1_uncoupled_term7(t1);

  // calculate t1 denominator
  t1_uncoupled_denom(t1);
  
} //end t1_uncoupled_calc
\end{lstlisting}
In Listing \ref{list:impl:calct1} \texttt{t1}, represents a two-dimensional array with dimension \texttt{t1[np][nh]}, containing all possible $t_i^a$ amplitudes. We proceed by illustrating the implementation of the term-functions in \emph{t1\_uncoupled\_calc}, which represent the terms in Eq. (\ref{eq:impl:t1amplitude5}). First the algebraic expression is given, followed by its implementation. Note how we in these implementations loop effectively over contributing interaction elements only. This technique saves much cpu-time compared with the brute force technique of looping over all interaction elements. Also note that the intermediates are calculated prior to the \texttt{t1} calculations. These intermediates are tabulated in matrices carrying names including \texttt{barh}. These names are not correspondent with the algebraic notation utilized above, however the connection is seen from the algebraic expressions above the implementation. Here we have \texttt{B->} which points to the Config class.
%
\begin{equation*}
t_i^aD_i^a\quad \leftarrow \quad-(1-\delta_{ab}\delta_{ij})\langle aj|v|bi \rangle t_{j}^{b}
\end{equation*}
%
\begin{lstlisting}[label={list:impl:t1term2},caption={implementation of the amp1 class function t1\_uncoupled\_term2()}]
void Amplitudes::t1_uncoupled_term2(double **t1) {
    for (alpha = 0; alpha < dim_alpha; alpha++) {
        map1 = 0;
        for (ph = 0; ph < phbcount[alpha]; ph += 2) {
            a = B->c_ph[alpha][ph];
            j = B->c_ph[alpha][ph + 1];
            map2 = 0;
            for (ph2 = 0; ph2 < phbcount[alpha]; ph2 += 2) {
                b = B->c_ph[alpha][ph2];
                i = B->c_ph[alpha][ph2 + 1];
                if (j != i || b != a) {
                    t1[a][i] = t1[a][i]-V->phph[alpha][map1][map2] * t1_old[b][j];
                }
                map2 += 1;
            }
            map1 += 1;
        }
    }
} // end t1_uncoupled_term2
\end{lstlisting}
%
\begin{equation*}
t_i^aD_i^a\quad \leftarrow \quad(1-\delta_{ab})[I1]_b^a\phantom{.}t_i^b
\end{equation*}
%
\begin{lstlisting}[label={list:impl:t1term3},caption={implementation of the amp1 class function t1\_uncoupled\_term3()}]
void Amplitudes::t1_uncoupled_term3(double **t1){
  
  for(a=0; a<np; a++){
    for(i=0; i<nh; i++){ 
      temp = 0.0;
      for(b=0; b<np; b++){
	if(b!=a){
	  temp = temp + barh_i02a[a][b]*t1_old[b][i];
	}
      }
      t1[a][i] = t1[a][i] + temp;
    }
  }

} // end t1_uncoupled_term3
\end{lstlisting}
%
\begin{equation*}
t_i^aD_i^a\quad \leftarrow \quad-(1-\delta_{ij})[I3]_i^j\phantom{.}t_{j}^{a}\nonumber
\end{equation*}
%
\begin{lstlisting}[label={list:impl:t1term4},caption={implementation of the amp1 class function t1\_uncoupled\_term4()}]
void Amplitudes::t1_uncoupled_term4(double **t1){
 
  for(a=0; a<np; a++){
    for(i=0; i<nh; i++){
      temp = 0.0;
      for(j=0; j<nh; j++){
	if(j!=i){
	  temp = temp + barh_03[j][i]*t1_old[a][j];
	}
      }
      t1[a][i] = t1[a][i] - temp;
    }
  }
  
} // end t1_uncoupled_term4
\end{lstlisting}
%
\begin{equation*}
t_i^aD_i^a\quad \leftarrow \quad\frac{1}{2}\langle bc|v|aj\rangle t_{ij}^{bc}\nonumber
\end{equation*}
%
\begin{lstlisting}[label={list:impl:t1term5},caption={implementation of the amp1 class function t1\_uncoupled\_term5()}]
void Amplitudes::t1_uncoupled_term5(double **t1){
  
  for(alpha=0; alpha<dim_alpha; alpha++){
    for(i=0; i<nh; i++){
      map1=0;
      for(p=0; p<ppbcount[alpha]; p+=2){
	b = B->c_pp[alpha][p];
	c = B->c_pp[alpha][p+1];
	map2=0;
	for(ph=0; ph<phbcount[alpha]; ph+=2){
	  a = B->c_ph[alpha][ph];
	  j = B->c_ph[alpha][ph+1];
	  t1[a][i] += 0.5*V->ppph[alpha][map1][map2]*t2_old[b][c][i][j];
	  t1[a][i] -= 0.5*V->ppph[alpha][map1][map2]*t2_old[c][b][i][j];
	  map2+=1;
	}
	map1+=1;
      }
    }
  }
} // end t1_uncoupled_term5
\end{lstlisting}


\begin{equation*}
t_i^aD_i^a\quad \leftarrow \quad\frac{1}{2}[I4]_{ic}^{jk}\phantom{.}t_{jk}^{ca}\nonumber
\end{equation*}

\begin{lstlisting}[label={list:impl:t1term6},caption={implementation of the amp1 class function t1\_uncoupled\_term6()}]
void Amplitudes::t1_uncoupled_term6(double **t1){
  
  for(a=0; a<np; a++){
    for(i=0; i<nh; i++){
      temp = 0.0;
      for(j=0; j<nh; j++){
	for(k=0; k<nh; k++){
	  for(c=0; c<np; c++){
	    temp = temp + barh_07[j][k][i][c]*t2_old[c][a][j][k];
	  }
	}
      }
      t1[a][i] = t1[a][i] + 0.5*temp;
    }
  }

} // end t1_uncoupled_term6
\end{lstlisting}


\begin{equation*}
t_i^aD_i^a\quad \leftarrow \quad [I2]_c^k\phantom{.}t_{ik}^{ac}
\end{equation*}
\begin{lstlisting}[label={list:impl:t1term7},caption={implementation of the amp1 class function t1\_uncoupled\_term7()}]
void Amplitudes::t1_uncoupled_term7(double **t1){
  
  for(a=0; a<np; a++){
    for(i=0; i<nh; i++){
      temp = 0.0;
      for(k=0; k<nh; k++){
	for(c=0; c<np; c++){
	  temp = temp + barh_01[k][c]*t2_old[a][c][i][k];
	}
      }
      t1[a][i] = t1[a][i] + temp;
    }
  }

} // end t1_uncoupled_term7
\end{lstlisting}

\begin{equation*}
t_i^a\quad \leftarrow \quad \frac{t_i^a}{D_i^a}=t_i^a/\left(\langle ai|v|ai \rangle-[I1]_a^a+[I3]_i^i\right)
\end{equation*}

\begin{lstlisting}[label={list:impl:t1denom},caption={implementation of the amp1 class function t1\_uncoupled\_denom()}]
void Amplitudes::t1_uncoupled_denom(double **t1){
  
  double **C;
  C = new double*[np];
  for(a=0; a<np; a++){
    C[a] = new double[nh];
    for(i=0; i<nh; i++){
      C[a][i] = 0.0;
    }
  }
  
  //C = [I3]-[I1]
  for(i=0; i<nh; i++){
    for(a=0; a<np; a++){
      C[a][i] = barh_03[i][i] - barh_i02a[a][a];
    }
  }
  //C += <ai|v|ai>
  for(int alpha=0; alpha<dim_alpha; alpha++){
    map1=0;
    for(ph=0; ph<phbcount[alpha]; ph+=2){
      a = B->c_ph[alpha][ph]];
      i = B->c_ph[alpha][ph+1];
      C[a][i] += V->phph[alpha][map1][map1];
      map1+=1;
    }
  }
  //t_i^a/D_i^a
  for(i=0; i<nh; i++){
    for(a=0; a<np; a++){
      t1[a][i] = t1[a][i]/C[a][i];
    }
  }
  
  //deallocating
  for(a=0; a<np; a++)
    delete[] C[a];
  delete[] C;
  
} // end t1_uncoupled_denom
\end{lstlisting}

\subsubsection{$\bf \hat{T}_2$-amplitude equation}

Consider the $\hat{T}_2$-amplitude equation. Again note that the summation-notation is omitted. We rearrange this equation as follows:
\begin{align} 
0&=\langle ij|v|ab\rangle+\frac{1}{2}\langle ab|v|cd\rangle t_{ij}^{cd}\\
&-\left( \hat{P}_{(ij)}f_j^kt_{ik}^{ab}+\hat{P}_{(ij)}f_c^kt_i^ct_{kj}^{ab}+\hat{P}_{(ij)}\langle kl|v|ci\rangle t_k^ct_{jl}^{ba}+\hat{P}_{(ij)}\langle kl|v|cd\rangle t_k^ct_i^dt_{jl}^{ba}\right.\\
&\phantom{++}\left.+\frac{1}{2}\hat{P}_{(ij)}\langle kl|v|cd\rangle t_{ki}^{cd}t_{lj}^{ab}\right)\\
&+\frac{1}{2}\left(\langle kl|v|ij\rangle t_{kl}^{ab}+\hat{P}_{(ij)}\langle kl|v|cj\rangle t_i^ct_{kl}^{ab}+\frac{1}{2}\langle kl|v|cd\rangle t_{kl}^{ab}t_{ij}^{cd}+\frac{1}{2}\hat{P}_{(ij)}\langle kl|v|cd\rangle t_i^ct_{kl}^{ab}t_j^d\right)\\
&+\left(\hat{P}_{(ab)}f_c^bt_{ij}^{ac}-\hat{P}_{(ab)}f_c^kt_k^at_{ij}^{cb}+\hat{P}_{(ab)}\langle ka|v|cd\rangle t_k^ct_{ij}^{db}-\hat{P}_{(ab)}\langle kl|v|cd\rangle t_k^ct_l^at_{ij}^{db}\right.\\
&\phantom{++}\left.-\frac{1}{2}\hat{P}_{(ab)}\langle kl|v|cd\rangle t_{kl}^{ca}t_{ij}^{db}\right)\\
&+\left(\hat{P}_{(ab)}\hat{P}_{(ij)}\langle kb|v|cj\rangle t_{ik}^{ac}+\hat{P}_{(ab)}\hat{P}_{(ij)}\langle ak|v|dc\rangle t_i^dt_{jk}^{bc}-\hat{P}_{(ab)}\hat{P}_{(ij)}\langle lk|v|ic\rangle t_l^at_{kj}^{cb}\right.\\
&\phantom{++}\left.+\hat{P}_{(ab)}\hat{P}_{(ij)}\langle kl|v|cd\rangle t_i^ct_k^at_{jl}^{bd}+\frac{1}{2}\hat{P}_{(ab)}\hat{P}_{(ij)}\langle kl|v|cd\rangle t_{ik}^{ac}t_{jl}^{bd}\right)\\
&+\left(-\hat{P}_{(ab)}\langle kb|v|ij\rangle t_k^a-\frac{1}{2}\hat{P}_{(ab)}\langle kb|v|cd\rangle t_k^at_{ij}^{cd}-\hat{P}_{(ab)}\hat{P}_{(ij)}\langle kb|v|cj\rangle t_k^at_i^c\right.\\
&\phantom{++}\left.-\frac{1}{2}\hat{P}_{(ab)}\hat{P}_{(ij)}\langle kb|v|cd\rangle t_k^at_i^ct_j^d+\frac{1}{2}\hat{P}_{(ab)}\langle kl|v|ij\rangle t_{k}^{a}t_l^b+\frac{1}{4}\hat{P}_{(ab)}\langle kl|v|cd\rangle t_k^at_{ij}^{cd}t_l^b\right.\\
&\phantom{++}\left.+\frac{1}{2}\hat{P}_{(ab)}\hat{P}_{(ij)}\langle kl|v|cj\rangle t_k^at_i^ct_l^b+\frac{1}{4}\hat{P}_{(ab)}\hat{P}_{(ij)}\langle kl|v|cd\rangle t_k^at_i^ct_l^bt_j^d\right)\\
&+\left(\hat{P}_{(ij)}\langle ab|v|cj\rangle t_i^c+\frac{1}{2}\hat{P}_{(ij)}\langle ab|v|cd\rangle t_i^ct_j^d\right)
\label{eq:impl:t2amplitude}
\end{align}
This is equivalent to:
\begin{align}
0&=\langle ij|v|ab\rangle+\frac{1}{2}\langle ab|v|cd\rangle t_{ij}^{cd}\nonumber\\
&-\hat{P}_{(ij)}\left(f_i^l + f_c^lt_i^c + \langle kl|v|ci\rangle t_k^c + \langle kl|v|cd\rangle t_k^ct_i^d+\frac{1}{2}\langle kl|v|cd\rangle t_{ki}^{cd}\right)t_{lj}^{ab}\nonumber\\
&+\frac{1}{2}\left(\langle kl|v|ij\rangle +\hat{P}_{(ij)}\langle lk|v|jc\rangle t_i^c+\frac{1}{2}\langle kl|v|cd\rangle t_{ij}^{cd}+\frac{1}{2}\hat{P}_{(ij)}\langle kl|v|cd\rangle t_i^ct_j^d\right)t_{kl}^{ab}\nonumber\\
&+\hat{P}_{(ab)}\left(f_d^a-f_d^kt_k^a+\langle ka|v|cd\rangle t_k^c-\langle lk|v|cd\rangle t_l^ct_k^a-\frac{1}{2}\langle kl|v|cd\rangle t_{kl}^{ca}\right)t_{ij}^{db}\nonumber\\
&+\hat{P}_{(ab)}\hat{P}_{(ij)}\left(\langle kb|v|cj\rangle + \langle bk|v|dc\rangle t_j^d-\langle lk|v|jc\rangle t_l^b-\langle kl|v|cd\rangle t_j^dt_l^b+\frac{1}{2}\langle kl|v|cd\rangle t_{jl}^{bd}\right)t_{ik}^{ac}\nonumber\\
&-\hat{P}_{(ab)}\left(\langle kb|v|ij\rangle +\frac{1}{2}\langle kb|v|cd\rangle t_{ij}^{cd}+\hat{P}_{(ij)}\langle kb|v|cj\rangle t_i^c+\frac{1}{2}\hat{P}_{(ij)}\langle kb|v|cd\rangle t_i^ct_j^d\right.\nonumber\\
&\phantom{++}\left.-\frac{1}{2}\langle kl|v|ij\rangle t_l^b-\frac{1}{4}\langle kl|v|cd\rangle t_{ij}^{cd}t_l^b-\frac{1}{2}\hat{P}_{(ij)}\langle kl|v|cj\rangle t_i^ct_l^b-\frac{1}{4}\hat{P}_{(ij)}\langle kl|v|cd\rangle t_i^ct_l^bt_j^d\right)t_k^a\nonumber\\
&+\hat{P}_{(ij)}\left(\langle ab|v|cj\rangle +\frac{1}{2}\langle ab|v|cd\rangle t_j^d\right)t_i^c
\label{eq:impl:t2amplitude2}
\end{align}
where we have relabeled some indices in the first, third and fourth parenthesis in order to extract common factors. We simplify this equation in a similar manner as for $\hat{T}_1$, by defining intermediates. The intermediates corresponding to $\hat{T}_2$ are given explicitly below. Note that we manipulate these expressions so that they correspond to the six matrices defined in the \emph{Interaction} class.

We recognize the first parenthesis of Eq. (\ref{eq:impl:t2amplitude2}) as the intermediate $[I3]$, already defined for the $\hat{T}_1$ amplitude Eq. (\ref{eq:impl:t1intermediates3}), viz...
\begin{align}
[I3]_{i}^{l} &= f_i^l + f_c^lt_i^c + \langle kl|v|ci\rangle t_k^c + \langle kl|v|cd\rangle t_k^ct_i^d+\frac{1}{2}\langle kl|v|cd\rangle t_{ki}^{cd}\nonumber\\
&=f_i^l+\langle ci|v|kl\rangle t_k^c+\frac{1}{2}\langle cd|v|kl\rangle t_{ki}^{cd}+(f_c^l+\langle dc|v|kl\rangle t_k^d)t_i^c\nonumber\\
&=f_i^l-\langle ci|v|lk\rangle t_k^c+\frac{1}{2}\langle cd|v|lk\rangle t_{ik}^{cd}+(f_c^l+\langle cd|v|lk\rangle t_k^d)t_i^c\nonumber\\
&=f_i^l-\langle ci|v|lk\rangle t_k^c+\frac{1}{2}\langle cd|v|lk\rangle t_{ik}^{cd}+[I2]_c^l\phantom{.}t_i^c
\label{eq:impl:t2intermediates3}
\end{align}
The second parenthesis reads
\begin{align}
[I6]_{ij}^{kl} &= \langle kl|v|ij\rangle +\hat{P}_{(ij)}\langle lk|v|jc\rangle t_i^c+\frac{1}{2}\langle kl|v|cd\rangle t_{ij}^{cd}+\frac{1}{2}\hat{P}_{(ij)}\langle kl|v|cd\rangle t_i^ct_j^d\nonumber\\
&=\langle kl|v|ij\rangle+\frac{1}{2}\langle cd|v|kl\rangle t_{ij}^{cd}+\hat{P}_{(ij)}\left(-\langle cj|v|lk\rangle+\frac{1}{2}\langle cd|v|kl\rangle t_j^d \right)t_i^c\nonumber\\
&=\langle lk|v|ji\rangle+\frac{1}{2}\langle dc|v|lk\rangle t_{ij}^{cd}+\hat{P}_{(ij)}\left(-\langle cj|v|lk\rangle+\frac{1}{2}\langle dc|v|lk\rangle t_j^d \right)t_i^c\nonumber\\
&=\langle lk|v|ji\rangle+\frac{1}{2}\langle dc|v|lk\rangle t_{ij}^{cd}+\hat{P}_{(ij)}[I5]_{jc}^{lk}\phantom{.}t_i^c\nonumber\\
&=\langle kl|v|ij\rangle+\frac{1}{2}\langle cd|v|kl\rangle t_{ij}^{cd}+\hat{P}_{(ij)}[I5]_{ic}^{kl}\phantom{.}t_j^c
\label{eq:impl:t2intermediates6}
\end{align}
where the intermediate $[I5]$ is defined in Eq. (\ref{eq:impl:t1intermediates5}). The third parenthesis reads
\begin{align}
[I7]_{d}^{a} &= f_d^a-f_d^kt_k^a+\langle ka|v|cd\rangle t_k^c-\langle lk|v|cd\rangle t_l^ct_k^a-\frac{1}{2}\langle kl|v|cd\rangle t_{kl}^{ca}\nonumber\\
&=\left(f_d^a +\langle dc|v|ak\rangle t_k^c\right)-\left(f_d^k+\langle cd|v|lk\rangle t_l^c\right)t_k^a-\frac{1}{2}\langle cd|v|kl\rangle t_{kl}^{ca}\nonumber\\
&=[I1]_d^a-[I2]_d^k\phantom{.}t_k^a-\frac{1}{2}\langle dc|v|kl\rangle t_{kl}^{ac}
\label{eq:impl:t2intermediates7}
\end{align}
where the intermediates $[I1]$ and $[I2]$ are defined in eqs. \ref{eq:impl:t1intermediates1} and \ref{eq:impl:t1intermediates2} respectively. The fourth parenthesis reads
\begin{align}
[I8]_{cj}^{kb} &=\langle kb|v|cj\rangle + \langle bk|v|dc\rangle t_j^d-\langle lk|v|jc\rangle t_l^b-\langle kl|v|cd\rangle t_j^dt_l^b+\frac{1}{2}\langle kl|v|cd\rangle t_{jl}^{bd}\nonumber\\
&=\left(-\langle bk|v|cj\rangle + \frac{1}{2}\langle dc|v|bk\rangle t_j^d\right)+\frac{1}{2}\langle dc|v|bk\rangle t_j^d\nonumber\\
&\phantom{=}-\left(-\langle cj|v|lk\rangle+\frac{1}{2}\langle cd|v|kl\rangle t_j^d+\frac{1}{2}\langle cd|v|kl\rangle t_j^d\right)t_l^b\nonumber\\
&\phantom{=}+\frac{1}{2}\langle cd|v|kl\rangle t_{jl}^{bd}\nonumber\\
&=[I9]_{cj}^{kb}+\frac{1}{2}\langle dc|v|bk\rangle t_j^d-[I4]_{jc}^{lk}\phantom{.}t_l^b+\frac{1}{2}\langle cd|v|kl\rangle t_{jl}^{bd}
\label{eq:impl:t2intermediates8}
\end{align}
where $[I4]$ corresponds to Eq. (\ref{eq:impl:t1intermediates4}) and $[I9]$ reads
\begin{align}
[I9]_{cj}^{kb} = -\langle bk|v|cj\rangle + \frac{1}{2}\langle dc|v|bk\rangle t_j^d
\label{eq:impl:t2intermediates9}
\end{align}
The fifth parenthesis yields
\begin{align}
[I10]_{ij}^{kb}&= \langle kb|v|ij\rangle +\frac{1}{2}\langle kb|v|cd\rangle t_{ij}^{cd}+\hat{P}_{(ij)}\langle kb|v|cj\rangle t_i^c+\frac{1}{2}\hat{P}_{(ij)}\langle kb|v|cd\rangle t_i^ct_j^d\nonumber\\
&\phantom{=}-\frac{1}{2}\langle kl|v|ij\rangle t_l^b-\frac{1}{4}\langle kl|v|cd\rangle t_{ij}^{cd}t_l^b-\frac{1}{2}\hat{P}_{(ij)}\langle kl|v|cj\rangle t_i^ct_l^b-\frac{1}{4}\hat{P}_{(ij)}\langle kl|v|cd\rangle t_i^ct_l^bt_j^d\nonumber\\
&=-\langle bk|v|ij\rangle-\frac{1}{2}\langle cd|v|bk\rangle t_{ij}^{cd}+\hat{P}_{(ij)}\left(-\langle bk|v|cj\rangle+ \frac{1}{2}\langle dc|v|bk\rangle t_j^d \right)t_i^c\nonumber\\
&\phantom{=}-\frac{1}{2}\left(\langle kl|v|ij\rangle+\frac{1}{2}\langle cd|v|kl\rangle t_{ij}^{cd}+\hat{P}_{(ij)}\left(\langle cj|v|kl\rangle+ \frac{1}{2}\langle cd|v|kl\rangle t_j^d\right)t_i^c\right)t_l^b\nonumber\\
&=-\langle bk|v|ij\rangle-\frac{1}{2}\langle cd|v|bk\rangle t_{ij}^{cd}+\hat{P}_{(ij)}[I9]_{cj}^{kb}\phantom{.}t_i^c\nonumber\\
&\phantom{=}-\frac{1}{2}\left(\langle lk|v|ji\rangle+\frac{1}{2}\langle dc|v|lk\rangle t_{ij}^{cd}+\hat{P}_{(ij)}\left(-\langle cj|v|lk\rangle+ \frac{1}{2}\langle dc|v|lk\rangle t_j^d\right)t_i^c\right)t_l^b\nonumber\\
&=-\langle bk|v|ij\rangle-\frac{1}{2}\langle cd|v|bk\rangle t_{ij}^{cd}+\hat{P}_{(ij)}[I9]_{cj}^{kb}\phantom{.}t_i^c-\frac{1}{2} [I6]_{ij}^{kl}\phantom{.}t_l^b
\label{eq:impl:t2intermediates10}
\end{align}
%
Finally the sixth parenthesis yields
\begin{equation}
[I11]_{cj}^{ab} = \langle ab|v|cj\rangle +\frac{1}{2}\langle ab|v|cd\rangle t_j^d
\label{eq:impl:t2intermediates11}
\end{equation}
We now reinsert these definitions into Eq. (\ref{eq:impl:t2amplitude2}), and obtain the $\hat{T}_2$ amplitude equation
\begin{align}
0&=\langle ab|v|ij\rangle+\frac{1}{2}\langle ab|v|cd\rangle t_{ij}^{cd}-\hat{P}_{(ij)}[I3]_{i}^{l}\phantom{.}t_{lj}^{ab}+\frac{1}{2}[I6]_{ij}^{kl}\phantom{.}t_{kl}^{ab}\nonumber\\
&+\hat{P}_{(ab)}[I7]_{d}^{a}\phantom{.}t_{ij}^{db}+\hat{P}_{(ab)}\hat{P}_{(ij)}[I8]_{cj}^{kb}\phantom{.}t_{ik}^{ac}-\hat{P}_{(ab)}[I10]_{ij}^{kb}\phantom{.}t_{k}^{a}+\hat{P}_{(ij)}[I11]_{cj}^{ab}\phantom{.}t_{i}^{c}
\label{eq:impl:t2amplitude3}
\end{align}
Similarly, as for the $\hat{T}_1$ amplitude, we wish to obtain an equation for the $\hat{T}_2$ amplitude $t_{ij}^{ab}$. This is obtained by using the same technique as in the case of $\hat{T}_1$, thus
\begin{align}
0&=\langle ab|v|ij\rangle+\frac{1}{2}\langle ab|v|ab\rangle t_{ij}^{ab}+\frac{1}{2}(1-\delta_{ca}\delta_{db})\langle ab|v|cd\rangle t_{ij}^{cd}\nonumber\\
&-\hat{P}_{(ij)}[I3]_{i}^{i}\phantom{.}t_{ij}^{ab}-(1-\delta_{il})\hat{P}_{(ij)}[I3]_{i}^{l}\phantom{.}t_{lj}^{ab}+\frac{1}{2}[I6]_{ij}^{ij}\phantom{.}t_{ij}^{ab}\nonumber\\
&+\frac{1}{2}(1-\delta_{ki}\delta_{lj})[I6]_{ij}^{kl}\phantom{.}t_{kl}^{ab}+\hat{P}_{(ab)}[I7]_{a}^{a}\phantom{.}t_{ij}^{ab}+\hat{P}_{(ab)}(1-\delta_{da})[I7]_{d}^{a}\phantom{.}t_{ij}^{db}\nonumber\\
&+\hat{P}_{(ab)}\hat{P}_{(ij)}[I8]_{cj}^{kb}\phantom{.}t_{ik}^{ac} -\hat{P}_{(ab)}[I10]_{ij}^{kb}\phantom{.}t_{k}^{a}+\hat{P}_{(ij)}[I11]_{cj}^{ab}\phantom{.}t_{i}^{c}
\label{eq:impl:t2amplitude4}
\end{align}
We extract the $t_{ij}^{ab}$ amplitudes and obtain the $t_{ij}^{ab}$-amplitude equation
\begin{align}
t_{ij}^{ab}&=\frac{1}{D_{ij}^{ab}}\left(\langle ab|v|ij\rangle+\frac{1}{2}(1-\delta_{ca}\delta_{db})\langle ab|v|cd\rangle t_{ij}^{cd}-\hat{P}_{(ij)}(1-\delta_{il})[I3]_{i}^{l}\phantom{.}t_{lj}^{ab}\right.\nonumber\\
&\left. +\frac{1}{2}(1-\delta_{ki}\delta_{lj})[I6]_{ij}^{kl}\phantom{.}t_{kl}^{ab}+\hat{P}_{(ab)}(1-\delta_{da})[I7]_{d}^{a}\phantom{.}t_{ij}^{db}+\hat{P}_{(ab)}\hat{P}_{(ij)}[I8]_{cj}^{kb}\phantom{.}t_{ik}^{ac}\right.\nonumber\\
&\left.-\hat{P}_{(ab)}[I10]_{ij}^{kb}\phantom{.}t_{k}^{a}+\hat{P}_{(ij)}[I11]_{cj}^{ab}\phantom{.}t_{i}^{c}\right)
\label{eq:impl:t2amplitude5}
\end{align}
where
\begin{align}
D_{ij}^{ab}=\hat{P}_{(ij)}[I3]_{i}^{i}-\frac{1}{2}\langle ab|v|ab\rangle-\frac{1}{2}[I6]_{ij}^{ij}-\hat{P}_{(ab)}[I7]_{a}^{a}
\label{eq:impl:t2denom}
\end{align}
Eq. \ref{eq:impl:t2amplitude5} is implemented in the \texttt{amp1} class function \texttt{t2\_uncoupled\_calc}, which is illustrated in Listing \ref{list:impl:calct2}
%
\begin{lstlisting}[label={list:impl:calct2},caption={Implementation of the amp1 class function t2\_uncoupled\_calc()}]
void Amplitudes::t2_uncoupled_calc() {
    //calculating t2 intermediates
    t2_uncoupled_intermediates();

    // <ab|v|ij>
    for (a = 0; a < np; a++)
        for (b = 0; b < np; b++)
            for (i = 0; i < nh; i++)
                for (j = 0; j < nh; j++)
                    t2[a][b][i][j] = 0.0;
    for (int alpha = 0; alpha < dim_alpha; alpha++) {
        map1 = 0;
        for (i = 0; i < ppbcount[alpha]; i += 2) {
            a = B->c_pp[alpha][i];
            b = B->c_pp[alpha][i + 1];
            map2 = 0;
            for (j = 0; j < hhbcount[alpha]; j += 2) {
                m = B->c_hh[alpha][j];
                n = B->c_hh[alpha][j + 1];
                t2[a][b][m][n] = V->pphh[alpha][map1][map2];
                t2[b][a][m][n] = -V->pphh[alpha][map1][map2];
                t2[a][b][n][m] = -V->pphh[alpha][map1][map2];
                t2[b][a][n][m] = V->pphh[alpha][map1][map2];
                map2 += 1;
            }
            map1 += 1;
        }
    }

    // 0.5(1-delta_{ca}delta_{db}) <ab|v|cd> t_{ij}^{cd}
    t2_uncoupled_term2(t2);

    // -P_(ij)(1-delta_{il})[I3]_{i}^{l}t_{lj}^{ab}
    t2_uncoupled_term3(t2);

    // 0.5(1-delta_{ki}delta_{lj})[I6]_{ij}^{kl}t_{kl}^{ab}
    t2_uncoupled_term4(t2);

    // P_(ab)(1-delta_{da})[I7]_{d}^{a}t_{ij}^{db}
    t2_uncoupled_term5(t2);

    // P_(ab)P_(ij)[I8]_{cj}^{kb}t_{ik}^{ac}
    t2_uncoupled_term6(t2);

    // -P_(ab)[I10]_{ij}^{kb}t_{k}^{a}
    t2_uncoupled_term7(t2);

    // P_(ij)[I11]_{cj}^{ab}t_{i}^{c}
    t2_uncoupled_term8(t2);

    // calculating t2 denominator
    t2_uncoupled_denom(t2);

} // end t2_uncoupled_calc
\end{lstlisting}
%
In Listing \ref{list:impl:calct2}, $t2$ represents the four-dimensional array with dimensions \texttt{[np][np][nh][nh]}, containing all possible $t_{ij}^{ab}$ amplitudes. Note how we are utilizing the \emph{Config} class (B) and \emph{Interaction} class (V) to get hold of the interaction elements in question. We proceed by illustrating the implementation of each term-function in \emph{t2\_uncoupled\_calc}, corresponding to the terms of Eq. (\ref{eq:impl:t2amplitude5}). In order to limit the implementation scope, standard procedures, like allocating arrays, are replaced by explanatory comments. The implementations are first presented in listings [\ref{list:impl:t2term2}-\ref{list:impl:t2denom}], followed by a more detailed description of \texttt{t2\_uncoupled\_term2} and \texttt{t2\_uncoupled\_term4}.
%
\begin{equation}
t_{ij}^{ab}D_{ij}^{ab}\quad \leftarrow \quad \frac{1}{2}(1-\delta_{ca}\delta_{db})\langle ab|v|cd\rangle t_{ij}^{cd}\nonumber
\end{equation}
%
\begin{lstlisting}[label={list:impl:t2term2},caption={implementation of the amp1 class function t2\_uncoupled\_term2()}]
void Amplitudes::t2_uncoupled_term2(double ****t2){

  /* allocate matrix C[dim_alpha][ppbcount[alpha]][nh*(nh-1)] */
  
  //filling of t2_pppp matrix = 2dimensional t2-old
  t2pppp_fill();
  
  //matrix multiplication
  for(alpha=0; alpha<dim_alpha; alpha++){
#pragma omp parallel default(shared) private(v,u,w) 
    {      
      nthreads = omp_get_num_threads();
#pragma omp for schedule (static)
      for(ab=0; ab<ppbcount[alpha]; ab++){
	for(ij=0; ij<temp; ij++){
	  for(cd=0; cd<ppbcount[alpha]; cd++){
	    C[alpha][ab][ij] += 0.5*v_pppp[alpha][ab][cd]
                                *t2_pppp[alpha][cd][ij];
	  }
	}
      } 
    }
  }
  //translation from 2dim to 4dim
  translste(t2,C);
  
  /* deallocate matrix C */
} // end t2_uncoupled_term2


void Amplitudes::t2pppp_fill()
{
  for(int a=0; a<dim_alpha ; a++){
    for(int cd=0; cd<ppbcount[a]; cd+=2){
      map2=0;
      for(int ij=0; ij<2*nh*(nh-1); ij+=2){
	t2_pppp[a][cd][map2] =t2_old[B->c_pp[a][cd]][B->c_pp[a][cd+1]]
                                      [ij_map[ij]][ij_map[ij+1]];
	t2_pppp[a][cd+1][map2]=t2_old[B->c_pp[a][cd+1]][B->c_pp[a][cd]]
                                     [ij_map[ij]][ij_map[ij+1]];
	map2+=1;
      }
    }
  }
}  //end fill


void Amplitudes::translste(double ****t2, double*** C)
{
   for(int alpha=0; alpha<dim_alpha; alpha++){
     for(int ab=0; ab<ppbcount[alpha]; ab+=2){
       map1=0;
       for(int ij=0; ij<2*nh*(nh-1); ij+=2){
	 t2[B->c_pp[alpha][ab]][B->c_pp[alpha][ab+1]]
           [ij_map[ij]][ij_map[ij+1]] += C[alpha][ab][map1];
	 t2[B->c_pp[alpha][ab+1]][B->c_pp[alpha][ab]]
           [ij_map[ij]][ij_map[ij+1]] += C[alpha][ab+1][map1];
	 map1+=1;
       }
     }
   }
}  //end translate
\end{lstlisting}
%
\begin{equation*}
t_{ij}^{ab}D_{ij}^{ab}\quad \leftarrow \quad -\hat{P}_{(ij)}(1-\delta_{il})[I3]_{i}^{l}\phantom{.}t_{lj}^{ab}
\end{equation*}

\begin{lstlisting}[label={list:impl:t2term3},caption={implementation of the amp1 class function t2\_uncoupled\_term3()}]
void Amplitudes::t2_uncoupled_term3(double ****t2){

  /*allocating new matrix barh03[nh][nh], where if(i!=l) 
  is implemented in form of zero elements along the diagonal*/
  for(i=0; i<nh-1; i++){
    for(l=i+1; l<nh; l++){
      barh03[i][l] = barh_03[i][l];
      barh03[l][i] = barh_03[l][i];
    }
  }
  for(i=0; i<nh; i++)
    barh03[i][i]=0.0;
  
  for(a=0; a<np-1; a++){
    for(b=a+1; b<np; b++){
      for(i=0; i<nh-1; i++){
	for(j=i+1; j<nh; j++){
	  temp=0.0;
          for(l=0; l<nh; l++){
	    temp -= barh03[l][i]*t2_old[a][b][l][j] 
                 -barh03[l][j]*t2_old[a][b][l][i];
	  }
	  t2[a][b][i][j] += temp;
	  t2[b][a][i][j] -= temp;
	  t2[a][b][j][i] -= temp;
	  t2[b][a][j][i] += temp;
	}
      }
    }
  }
  
  //deallocating matrix barh03
} // end t2_uncoupled_term3
\end{lstlisting}

\begin{equation*}
 t_{ij}^{ab}D_{ij}^{ab}\quad \leftarrow \quad \frac{1}{2}(1-\delta_{ki}\delta_{lj})[I6]_{ij}^{kl}\phantom{.}t_{kl}^{ab}
\end{equation*}

\begin{lstlisting}[label={list:impl:t2term4},caption={implementation of the amp1 class function t2\_uncoupled\_term4()}]
void Amplitudes::t2_uncoupled_term4(double ****t2){
  
  /* allocating matrix: A[np^2][nh^2], barh[nh^2][nh^2], C[np^2][nh^2]  */
  //filling matrices A, barh:
  for(a=0; a<np; a++)
    for(b=0; b<np; b++)
      for(k=0; k<nh; k++)
	for(l=0; l<nh; l++)
	  A[a*np+b][k*nh+l] = t2_old[a][b][k][l];
  for(k=0; k<nh; k++)
    for(l=0; l<nh; l++)
      for(i=0; i<nh; i++)
	for(j=0; j<nh; j++)
	  if(i!=k || j!=l)
	    barh[k*nh+l][i*nh+j] = barh_09[k][l][i][j];	    

  //matrix multiplication
#pragma omp parallel default(shared) private(v,u,w) 
  {      
    nthreads = omp_get_num_threads();
#pragma omp for schedule (static)
    for(ab=0; ab<np*np; ab++){
      for(ij=0; ij<nh*nh; ij++){
	for(kl=0; kl<nh*nh; kl++){
	  C[ab][ij] += 0.5*A[ab][kl]*barh[kl][ij];
	}
      }
    }
  }
  //translate
  for(a=0; a<np; a++){
    for(b=0; b<np; b++){
      for(i=0; i<nh; i++){
	for(j=0; j<nh; j++){
	  t2[a][b][i][j] = t2[a][b][i][j]+C[a*np+b][i*nh+j];
	}
      }
    }
  }
  //deallocating matrix A,barh,C
} // end t2_uncoupled_term4
\end{lstlisting}
%
\begin{equation*}
t_{ij}^{ab}D_{ij}^{ab}\quad \leftarrow \quad \hat{P}_{(ab)}(1-\delta_{da})[I7]_{d}^{a}\phantom{.}t_{ij}^{db}
\end{equation*}
%
\begin{lstlisting}[label={list:impl:t2term5},caption={implementation of the amp1 class function t2\_uncoupled\_term5()}]
void Amplitudes::t2_uncoupled_term5(double ****t2){

  /*allocating new matrix barh02[np][np], which implement 
  if(a!=d) by introducing zero elements along the diagonal*/
  for(a=0; a<np-1; a++){
    for(d=a+1; d<np; d++){
      barh02[a][d] = barh_02[a][d];
      barh02[d][a] = barh_02[d][a];
    }
  }
  for(a=0; a<np; a++)
    barh02[a][a]=0.0;
  
  for(a=0; a<np-1; a++){
    for(b=a+1; b<np; b++){
      for(i=0; i<nh-1; i++){
	for(j=i+1; j<nh; j++){	
	  temp = 0.0;
	  for(d=0; d<np; d++){
	    temp += barh02[a][d]*t2_old[b][d][i][j]
                   -barh02[b][d]*t2_old[a][d][i][j];
	  }
	  t2[a][b][i][j] += temp;
	  t2[b][a][i][j] -= temp;
	  t2[a][b][j][i] -= temp;
	  t2[b][a][j][i] += temp;
	}
      }
    }
  }
  
  //deallocating barh02
} // end t2_uncoupled_term5
\end{lstlisting}

\begin{equation*}
t_{ij}^{ab}D_{ij}^{ab}\quad \leftarrow \quad\hat{P}_{(ab)}\hat{P}_{(ij)}[I8]_{cj}^{kb}\phantom{.}t_{ik}^{ac}
\end{equation*}

\begin{lstlisting}[label={list:impl:t2term6},caption={implementation of the amp1 class function t2\_uncoupled\_term6()}]
void Amplitudes::t2_uncoupled_term6(double ****t2){
  
  /*allocating matrix: A[np*nh][np*nh], I[np*nh][np*nh], C[np*nh][np*nh] */

  //filling the matrices
  for(j=0; j<nh; j++)
    for(b=0; b<np; b++)
      for(c=0; c<np; c++)
	for(k=0; k<nh; k++)
	  I[k*np+c][j*np+b]= barh_i10c[k][b][c][j];

  for(i=0; i<nh; i++)
    for(a=0; a<np; a++)
      for(c=0; c<np; c++)
	for(k=0; k<nh; k++)
	  A[i*np+a][k*np+c] = t2_old[a][c][i][k];

  //multiplication
#pragma omp parallel default(shared) private(v,u,w) 
  {      
    nthreads = omp_get_num_threads();
#pragma omp for schedule (static)
    for(u=0; u<np*nh; u++){
      for(v=0; v<np*nh; v++){
	for(w=0; w<np*nh; w++){
	  C[u][v]+=A[u][w]*I[w][v];
	}
      }
    }
  }

  //translate
  for(a=0; a<np; a++){
    for(b=0; b<np; b++){
      for(i=0; i<nh; i++){
	for(j=0; j<nh; j++){
	  t2[a][b][i][j] += C[i*np+a][j*np+b];
	  t2[a][b][i][j] -= C[i*np+b][j*np+a];
	  t2[a][b][i][j] -= C[j*np+a][i*np+b];
	  t2[a][b][i][j] += C[j*np+b][i*np+a];
	}
      }
    }
  }
  //deallocate matrix: A, I, C 
} // end t2_uncoupled_term6
\end{lstlisting}

\begin{equation*}
t_{ij}^{ab}D_{ij}^{ab}\quad \leftarrow \quad-\hat{P}_{(ab)}[I10]_{ij}^{kb}\phantom{.}t_{k}^{a} 
\end{equation*}

\begin{lstlisting}[label={list:impl:t2term7},caption={implementation of the amp1 class function t2\_uncoupled\_term7()}]
void Amplitudes::t2_uncoupled_term7(double ****t2){
  
  /* t1 used instead if t1_old for a quicker convergence */
  for(a=0; a<np-1; a++){
    for(b=a+1; b<np; b++){
      for(i=0; i<nh-1; i++){
	for(j=i+1; j<nh; j++){	
	  temp = 0.0;
	  for(k=0; k<nh; k++){
	    temp = temp + (barh_i12a[k][b][i][j]*t1[a][k] 
                 - barh_i12a[k][a][i][j]*t1[b][k]); 
	  }
	  t2[a][b][i][j] -= temp;
	  t2[b][a][i][j] += temp;
	  t2[a][b][j][i] += temp;
	  t2[b][a][j][i] -= temp;
	}
      }
    }
  }
} // end t2_uncoupled_term7
\end{lstlisting}

\begin{equation*}
t_{ij}^{ab}D_{ij}^{ab}\quad \leftarrow \quad\hat{P}_{(ij)}[I11]_{cj}^{ab}\phantom{.}t_{i}^{c} 
\end{equation*}

\begin{lstlisting}[label={list:impl:t2term8},caption={implementation of the amp1 class function t2\_uncoupled\_term8()}]
void Amplitudes::t2_uncoupled_term8(double ****t2){
  
  /* t1 used instead if t1_old for a quicker convergence*/
  for(alpha=0; alpha<dim_alpha; alpha++){
    for(p=0; p<ppbcount[alpha]; p+=2){
      a=B->c_pp[alpha][p];
      b=B->c_pp[alpha][p+1];	  
      for(j=0; j<nh; j++){
	for(i=0; i<nh; i++){
	  temp1=0.0;
	  temp2=0.0;
	  for(c=0; c<np; c++){
	    temp1 += (barh_i11a[a][b][c][j]*t1[c][i] 
                  - barh_i11a[a][b][c][i]*t1[c][j]);
	  }
	  t2[a][b][i][j] +=temp1;
	  t2[b][a][i][j] -=temp1; 
	}
      }
    }
  }
} // end t2_uncoupled_term8
\end{lstlisting}

\begin{equation*}
t_{ij}^{ab}\quad \leftarrow \quad \frac{t_{ij}^{ab}}{D_{ij}^{ab}}=t_{ij}^{ab}/\left(\hat{P}_{(ij)}[I3]_{i}^{i}-\frac{1}{2}\langle ab|v|ab\rangle-\frac{1}{2}[I6]_{ij}^{ij}-\hat{P}_{(ab)}[I7]_{a}^{a}\right) 
\end{equation*}

\begin{lstlisting}[label={list:impl:t2denom},caption={implementation of the amp1 class function t2\_uncoupled\_denom()}]
void Amplitudes::t2_uncoupled_denom(double ****t2){

  /* allocating matrix C[np][np][nh][nh] */
    
  for(a=0; a<np; a++){
    for(b=0; b<np; b++){
      h2ab = barh_02[a][a] + barh_02[b][b];
      for(j=0; j<nh; j++){
	for(i=0; i<nh; i++){
	  h3ij = barh_03[i][i] + barh_03[j][j];
	  h9ijij = 0.5*barh_09[i][j][i][j];
	  C[a][b][i][j] = h3ij - h2ab - h9ijij;
	}
      }
    }
  }

  double temp;
  for(alpha=0; alpha<dim_alpha; alpha++){
    map1=0;
    for(p=0; p<ppbcount[alpha]; p+=2){
      a= B->c_pp[alpha][p];
      b= B->c_pp[alpha][p+1];
      temp=0.5*V->pppp[alpha][map1][map1];
      for(i=0; i<nh; i++){
	for(j=0; j<nh; j++){
	  C[a][b][i][j]-= temp;
	  C[b][a][i][j]-= temp;  
	}
      }
      map1+=1;
    }
  }
  
  for(b=0; b<np-1; b++){
    for(a=b+1; a<np; a++){
      for(j=0; j<nh-1; j++){
	for(i=j+1; i<nh; i++){
	  t2[b][a][j][i] =  t2[b][a][j][i]/C[b][a][j][i];
	  t2[a][b][j][i] = -t2[b][a][j][i];
	  t2[b][a][i][j] = -t2[b][a][j][i];
	  t2[a][b][i][j] =  t2[b][a][j][i];	  
	}
      }
    }
  }
  
  //deallocate matrix C
  
} // end t2_uncoupled_denom
\end{lstlisting}

Let us review the class function \texttt{t2\_uncoupled\_term2} and \texttt{t2\_uncoupled\_term4} in detail. We start with the former. The expression calculated in this function reads
\begin{equation*}
\frac{1}{2}(1-\delta_{ca}\delta_{db})\langle ab|v|cd\rangle t_{ij}^{cd}
\label{eq:impl:t2term2expression}
\end{equation*}
Originally this calculation was implemented with brute force by M. P. Lohne, see illustration in listing \ref{list:impl:oldt2term2}. 
\begin{lstlisting}[label={list:impl:oldt2term2},caption={M.P.Lohne's brute force implementation of amp1 class function t2\_uncoupled\_term2}]
void Amplitudes::t2_uncoupled_term2(Array<double,4> ans){

  for(j=0; j<nh; j++){
    for(i=0; i<nh; i++){
      for(b=0; b<np; b++){
	for(a=0; a<np; a++){
	  temp = 0.0;
	  for(d=0; d<np; d++){
	    for(c=0; c<np; c++){
	      if(a!=c || b!=d){
		temp = temp + V->pppp(a,b,c,d)*t2_old(c,d,i,j);
	      }
	    }
	  }
	  ans(a,b,i,j) = 0.5*temp;
	}
      }
    }
  }
} // end t2_uncoupled_term2  
\end{lstlisting}

The first simplification we make, is to utilize the two-particle basis embedded in the class \emph{Config} (\texttt{B->}), and the corresponding interaction matrices in the class \emph{Interaction} (\texttt{V->}). By coupling the a-loop and the b-loop into one loop over the \texttt{c\_pp} in \emph{Config}, and repeating this for the c-loop and the d-loop, we obtain the simplification illustrated in listing \ref{list:impl:t2term2firstsimplification}
\begin{lstlisting}[label={list:impl:t2term2firstsimplification},caption={illustrates the first simplification of M.P.Lohne's implementation of amp1 class function t2\_uncoupled\_term2}]
void Amplitudes::t2_uncoupled_term2(double ****t2){

  for(alpha=0; alpha<dim_alpha; alpha++){
    for(j=0; j<nh; j++){
      for(i=0; i<nh; i++){
	map1=0;
	for(p=0; p<ppbcount[alpha]; p+=2){
	  a = B->c_pp[alpha][p];
	  b = B->c_pp[alpha][p+1];
	  map2=0;
	  for(p2=0; p2<ppbcount[alpha]; p2+=2){
	    c = B->c_pp[alpha][p2];
	    d = B->c_pp[alpha][p2+1];
	    if(a!=e || b!=f){
	      t2[a][b][i][j] = ans[a][b][i][j] 
                + 0.5*V->pppp[alpha][map1][map2]*t2_old[c][d][i][j];
	      t2[b][a][i][j] = ans[b][a][i][j] 
                + 0.5*V->pppp[alpha][map1][map2]*t2_old[d][c][i][j];
	    }
	    if(b!=e || a!=f){
	      t2[b][a][i][j] = ans[b][a][i][j] 
                - 0.5*V->pppp[alpha][map1][map2]*t2_old[c][d][i][j];
	      t2[a][b][i][j] = ans[a][b][i][j] 
                - 0.5*V->pppp[alpha][map1][map2]*t2_old[d][c][i][j];
	    }
	    map2+=1;
	  }
	  map1+=1;
	}
      }
    }
  }
} // end t2_uncoupled_term2  
\end{lstlisting}
Next we notice in listing \ref{list:impl:oldt2term2} that the interaction matrix \texttt{pppp} have two common indices with the four-dimensional matrix \texttt{t2\_old}. Thus, this multiplication is in reality a matrix multiplication. Therefore, by mapping the four-dimensional \texttt{t2\_old} into a two-dimensional matrix \texttt{t2\_pppp}, we can replace these loops by a matrix multiplication. This matrix multiplication results in a significant speed up of our code. The new matrix \texttt{t2\_pppp} is created in correspondence with the three-dimensional interaction matrix \texttt{pppp}. For each possible $\{M,M_s\}$-value, we create a two-dimensional \texttt{t2\_pppp[cd][ij]}-matrix. The common dimension \texttt{[cd]} is determined by the \texttt{c\_pp} in \emph{Config} class, while the dimension of \texttt{[ij]} reads
\begin{equation*}
ij = nh \cdot (nh-1)
\label{eq:impl:ijdimension}
\end{equation*}
This dimension equals the number of all pairs of i and j, where i and j are different. If i and j are equal it violates the Pauli exclusion principle. The mapping of i and j into one variable is obtained as illustrated in listing \ref{list:impl:ijmap}. We then create an array \texttt{ij\_map} with twice the size of the \texttt{[ij]} dimension Eq. (\ref{eq:impl:ijdimension}), and tabulate the pairs of \texttt{i} and \texttt{j} values constituting the new \texttt{ij} values. This mapping is performed once in the \emph{Amplitudes} constructor.
\begin{lstlisting}[label={list:impl:ijmap},caption={Illustrates the map of the two quantum numbers i,j into one number ij}]
//mapping i,j -> ij  
int ij_map = new int[2*nh*(nh-1)];
temp=0;
for(int i=0; i<nh; i++){
  for(int j=0; j<nh; j++){
    if(i!=j){
      ij_map[tmp]=i;
      ij_map[tmp+1]=j;
      tmp+=2;
    }
  }
}
\end{lstlisting}
We utilize this mapping when filling \texttt{t2\_pppp}, and when we translate the new two-dimensional system back into the four-dimensional system of $t2$, see the functions \texttt{t2pppp\_fill} and \emph{translate} respectively, in listing \ref{list:impl:t2term2}. The matrix multiplication however, is not straightforward because of the if-statements in both listing \ref{list:impl:oldt2term2} and \ref{list:impl:t2term2firstsimplification}. We therefore create a new interaction matrix \texttt{v\_pppp}, where these if-statements are incorporated. We incorporate the if-statements by implementing the unwanted matrix elements as zero. The matrix multiplication between \texttt{t2\_pppp} and \texttt{v\_pppp} is then carried out, and $t2_{ij}^{ab}$ is obtained by translating back into a four-dimensional system, see listing \ref{list:impl:t2term2}.

Next we consider the class function \texttt{t2\_uncoupled\_term4}. This function performs the calculation of the following expression
\begin{equation*}
\frac{1}{2}(1-\delta_{ki}\delta_{lj})[I6]_{ij}^{kl}\phantom{.}t_{kl}^{ab}
\label{eq:impl:t2term4expression}
\end{equation*}
Originally this expression was implemented with brute force by M. P. Lohne, see illustration in listing \ref{list:impl:oldt2term4}
\begin{lstlisting}[label={list:impl:oldt2term4},caption={M.P.Lohne's brute force implementation of amp1 class function t2\_uncoupled\_term4} \cite{lohne}]
void Amplitudes::t2_uncoupled_term4(Array<double,4> ans){

  for(j=0; j<nh; j++){
    for(i=0; i<nh; i++){
      for(b=0; b<np; b++){
	for(a=0; a<np; a++){
	  temp = 0.0;
	  for(l=0; l<nh; l++){
	    for(k=0; k<nh; k++){
	      if(i!=k || j!=l){
		temp = temp + barh_09(k,l,i,j)*t2_old(a,b,k,l);
	      }
	    }
	  }
	  ans(a,b,i,j) = 0.5*temp;
	}
      }
    }
  }
  
} // end t2_uncoupled_term4
\end{lstlisting}
In this implementation there is no direct connections to an interaction element, thus simplification utilizing the two-particle basis embedded in the class \emph{Config} is not obtainable. However, we observe that the four-dimensional matrices \texttt{barh\_09} and \texttt{t2\_old} have two common indices \texttt{k} and \texttt{l}. This means that by mapping the matrices into a two-dimensional form, the performed calculation is equivalent to a matrix multiplication between \texttt{barh\_09[kl][ij]} and \texttt{t2\_old[ab][kl]}, which results in a matrix with dimension [ab][ij]. We create three new two-dimensional matrices, \texttt{barh} representing \texttt{barh\_09[kl][ij]}, \texttt{A} representing \texttt{t2\_old[ab][kl]}, and \texttt{C} for performing the multiplication. The mapping from four to two dimensions is performed with brute force when filling the matrices \texttt{A} and \texttt{barh} in listing \ref{list:impl:t2term4}. The dimensions thus read

\begin{align*}
ab &= a \cdot b=np^2 \\
ij &= i \cdot j=nh^2 \\
kl &= k \cdot l=nh^2
\label{eq:impl:t2term4dimensions}
\end{align*}

Note how the if-statement is incorporated in the new matrix \texttt{barh}. From here the matrix multiplication is straightforward, and the final result for \texttt{t2} is obtained by mapping back into the four-dimensional system, in the same manner which we mapped into the two-dimensional system.

In both listing \ref{list:impl:t2term2} and \ref{list:impl:t2term4} we are utilizing Open Multi-Processing (OMP), which is an application for parallellizing programs in a shared memory environment. OMP provides tools which create and manage threads automatically, and this makes parallellizing a much easier job. Parallel computing refers to computations where many calculations are carried out simultaneously, such that large and time consuming problems can be divided between threads into smaller ones, and then solved concurrently. Pragmas are special compiler commands, providing the compiler with additional information. OMP contains a set of pragmas that instructs the compiler to parallellize the code, but only if the compiler support OMP. The most basic pragma is the "$\#$pragma omp parallel", which denotes the region one wishes to parallellize. OMP is thereby a source to speed up our code. However, it is worth noticing that OMP can not run on different remote machines, like in a cluster of machines. This means that the speed up is limited by the number of processors on the single machine performing the calculation.  

\subsubsection{Intermediates}

In the following we present the implementation of the intermediates for both amplitudes $\hat{T}_1$ and $\hat{T}_2$. We calculate the intermediates since those terms appear more than one time in the CCSD equation. Then it would save us some computational time if we just calculate them once. We are not utilizing matrix multiplication when calculating the intermediates, even though this is possible in some of the calculations. We did not prioritize this implementation because the time consumption of calculating the intermediates did not imply that this would constitute a great speed up of our code.

\begin{equation*}
[I1]_{b}^{a}=f_{b}^{a}+\langle bc|v|aj\rangle t_j^c
\end{equation*}
\begin{lstlisting}[label={list:impl:intermediate1},caption={implementation of I1 in the amp1 class function ccsd\_uncoupled\_barh\_i02a\_store}]
void Amplitudes::ccsd_uncoupled_barh_i02a_store(){
 
  //initializing barh_i02a ([I1]) by including f_{b}^{a}
  for(int a=0; a<np; a++){
    for(int b=a; b<np; b++){
      barh_i02a[a][b] = F->f_pp[a][b];
      barh_i02a[b][a] = F->f_pp[b][a];
    }
  } 
  // calculating second term of[I1]
  for(alpha=0; alpha<dim_alpha; alpha++){
    map1=0;
    for(p=0; p<ppbcount[alpha]; p+=2){
      b = B->c_pp[alpha][p];
      c = B->c_pp[alpha][p+1];
      map2=0;
      for(ph=0; ph<phbcount[alpha]; ph+=2){   
	a = B->phbasis[alpha][ph];
	j = B->phbasis[alpha][ph+1];
	barh_i02a[a][b] += V->ppph[alpha][map1][map2]*t1[c][j];
	barh_i02a[a][c] -= V->ppph[alpha][map1][map2]*t1[b][j];
	map2+=1;
      }
      map1+=1;
    }
  }
}  //end [I1]
\end{lstlisting}

\begin{equation*}
[I2]_c^k =f_{c}^{k}+\langle bc|v|jk\rangle t_j^b
\end{equation*}
\begin{lstlisting}[label={list:impl:intermediate2},caption={implementation of I2 in the amp1 class function ccsd\_uncoupled\_barh\_01\_store}]
void Amplitudes::ccsd_uncoupled_barh_01_store(){

  //initializing barh_01 ([I2]) by including f_{c}^{k}
  for(int c=0; c<nh; c++){
    for(int k=0; k<np; k++){
      barh_01[c][k] = F->f_hp[c][k];
    }
  }
  //calculating second term of [I2]
  for(alpha=0; alpha<dim_alpha; alpha++){
    map1=0;
    for(p=0; p<ppbcount[alpha]; p+=2){
      c = B->c_pp[alpha][p];
      b = B->c_pp[alpha][p+1];
      map2=0;
      for(h=0; h<hhbcount[alpha]; h+=2){
	k = B->c_hh[alpha][h];
	j = B->c_hh[alpha][h+1];
	barh_01[k][c] +=  V->pphh[alpha][map1][map2]*t1[b][j];
	barh_01[k][b] -=  V->pphh[alpha][map1][map2]*t1[c][j];
	barh_01[j][c] -=  V->pphh[alpha][map1][map2]*t1[b][k];
	barh_01[j][b] +=  V->pphh[alpha][map1][map2]*t1[c][k];
	map2+=1;
      }
      map1+=1;
    }
  }

% }  //end [I2]
% \end{lstlisting}

\begin{equation*}
[I3]_i^j=f_{i}^{j}-\langle bi|v|jk\rangle t_{k}^{b}+\frac{1}{2}\langle bc|v|jk\rangle t_{ik}^{bc}+[I2]_b^j\phantom{.}t_i^b
\end{equation*}
\begin{lstlisting}[label={list:impl:intermediate3},caption={implementation of I3 in the amp1 class function ccsd\_uncoupled\_barh\_03\_store}]
void Amplitudes::ccsd_uncoupled_barh_03_store(){
 
  //initializing barh_03 ([I3]) by including f_{i}^{j}
  for(int i=0; i<nh; i++){
    for(int j=i; j<nh; j++){
      barh_03[i][j] = F->f_hh[i][j];
      barh_03[j][i] = F->f_hh[j][i];
    }
  }
  //calculating [I2]
  for(j=0; j<nh; j++){
    for(i=0; i<nh; i++){
      for(b=0; b<np; b++){ 
	barh_03[j][i] = barh_03[j][i] + barh_01[j][b]*t1[b][i];
      }
    }
  }
  //calculating second term of [I3]
  for(alpha=0; alpha<dim_alpha; alpha++){
    map1=0;
    for(ph=0; ph<phbcount[alpha]; ph+=2){
      b = B->phbasis[alpha][ph];
      i = B->phbasis[alpha][ph+1];
      map2=0;
      for(h=0; h<hhbcount[alpha]; h+=2){
	j = B->c_hh[alpha][h];
	k = B->c_hh[alpha][h+1];
	barh_03[j][i] -= V->phhh[alpha][map1][map2]*t1[b][k];
	barh_03[k][i] += V->phhh[alpha][map1][map2]*t1[b][j];
	map2+=1;
      }
      map1+=1;
    }
  }
  //calculating third term of [I3]
  for(alpha=0; alpha<dim_alpha; alpha++){
    map1=0;
    for(p=0; p<ppbcount[alpha]; p+=2){
      b = B->c_pp[alpha][p];
      c = B->c_pp[alpha][p+1];
      map2=0;
      for(h=0; h<hhbcount[alpha]; h+=2){
	j = B->c_hh[alpha][h];
	k = B->c_hh[alpha][h+1];
	for(i=0; i<nh; i++){
	  barh_03[j][i] += V->pphh[alpha][map1][map2]*t2[b][c][i][k];
	  barh_03[k][i] -= V->pphh[alpha][map1][map2]*t2[b][c][i][j];
	}
	map2+=1;
      }
      map1+=1;
    }
  }
} //end [I3]

\end{lstlisting}

\begin{equation*}
[I4]_{ic}^{jk}=[I5]_{ic}^{jk}+\frac{1}{2}\langle bc|v|jk\rangle t_i^b 
\end{equation*}
%
\begin{lstlisting}[label={list:impl:intermediate4},caption={implementation of I4 in the amp1 class function ccsd\_uncoupled\_barh\_07\_store}]
void Amplitudes::ccsd_uncoupled_barh_07_store(){

  //initializing barh_07 ([I4]) by including [I5]_{ic}^{jk}
  for(int j=0; j<nh; j++){
    for(int k=j; k<nh; k++){
      for(int i=0; i<nh; i++){
	for(int c=0; c<np; c++){
	  barh_07[j][k][i][c] = barh_i07a[j][k][i][c];
	  barh_07[k][j][i][c] = barh_i07a[k][j][i][c];
	}
      }
    }
  }
  //calculating second term of [I4]
  for(alpha=0; alpha<dim_alpha; alpha++){
    map1=0;
    for(p=0; p<ppbcount[alpha]; p+=2){
      b = B->c_pp[alpha][p];
      c = B->c_pp[alpha][p+1];
      map2=0;
      for(h=0; h<hhbcount[alpha]; h+=2){
	j = B->c_hh[alpha][h];
	k = B->c_hh[alpha][h+1];
	temp3 = 0.5 * V->pphh[alpha][map1][map2];
	for(i=0; i<nh; i++){
	  temp1 = temp3*t1[f][i];
	  temp2 = temp3*t1[e][i];
	  barh_07[j][k][i][c] += temp1;
	  barh_07[j][k][i][b] -= temp2;
	  barh_07[k][j][i][c] -= temp1;
	  barh_07[k][j][i][b] += temp2;
	}
	map2+=1;
      }
      map1+=1;
    }
  }
} //end[I4]
\end{lstlisting}
%
\begin{equation*}
[I5]_{ic}^{jk}=-\langle ci|v|jk\rangle+\frac{1}{2}\langle bc|v|jk\rangle t_i^b
\end{equation*}
\begin{lstlisting}[label={list:impl:intermediate5},caption={implementation of I5 in the amp1 class function ccsd\_uncoupled\_barh\_i07a\_store}]
void Amplitudes::ccsd_uncoupled_barh_i07a_store(){
  
  //calculating first term of [I5] (barh_i07a)
  for(alpha=0; alpha<dim_alpha; alpha++){
    map1=0;
    for(ph=0; ph<phbcount[alpha]; ph+=2){
      c = B->phbasis[alpha][ph];
      i = B->phbasis[alpha][ph+1];
      map2=0;
      for(h=0; h<hhbcount[alpha]; h+=2){
	j = B->c_hh[alpha][h];
	k = B->c_hh[alpha][h+1];
	barh_i07a[j][k][i][c] = -V->phhh[alpha][map1][map2];
	barh_i07a[k][j][i][c] = +V->phhh[alpha][map1][map2];
	map2+=1;
      }
      map1+=1;
    }
  }
  //calculating second term of [I5]
  for(alpha=0; alpha<dim_alpha; alpha++){
    map1=0;
    for(p=0; p<ppbcount[alpha]; p+=2){
      b = B->c_pp[alpha][p];
      c = B->c_pp[alpha][p+1];
      map2=0;
      for(h=0; h<hhbcount[alpha]; h+=2){
	j = B->c_hh[alpha][h];
	k = B->c_hh[alpha][h+1];
	temp3 = 0.5*V->pphh[alpha][map1][map2];
	for(i=0; i<nh; i++){
	  temp1 = temp3*t1[b][i];
	  temp2 = temp3*t1[c][i];
	  barh_i07a[j][k][i][c] += temp1;
	  barh_i07a[k][j][i][c] -= temp1;
	  barh_i07a[k][j][i][b] += temp2;
	  barh_i07a[j][k][i][b] -= temp2;
	}
	map2+=1;
      }
      map1+=1;
    }
  }
} //end [I5]
\end{lstlisting}

\begin{equation*}
 [I6]_{ij}^{kl}=\langle kl|v|ij\rangle+\frac{1}{2}\langle cd|v|kl\rangle t_{ij}^{cd}+\hat{P}_{(ij)}[I5]_{ic}^{kl}\phantom{.}t_j^c
\end{equation*}
\begin{lstlisting}[label={list:impl:intermediate6},caption={implementation of I6 in the amp1 class function ccsd\_uncoupled\_barh\_09\_store}]
void Amplitudes::ccsd_uncoupled_barh_09_store(){
 
  //calculating first term of [I6] (barh_09)
  for(alpha=0; alpha<dim_alpha; alpha++){
    map1=0;
    for(i=0; i<hhbcount[alpha]; i+=2){
      k = B->c_hh[alpha][i];
      l = B->c_hh[alpha][i+1];
      map2=0;
      for(j=0; j<hhbcount[alpha]; j+=2){
	i=B->c_hh[alpha][j];
	j=B->c_hh[alpha][j+1];
	barh_09[k][l][i][j] =  V->hhhh[alpha][map1][map2];
	barh_09[k][l][j][i] = -V->hhhh[alpha][map1][map2];
	barh_09[l][k][i][j] = -V->hhhh[alpha][map1][map2];
	barh_09[l][k][j][i] =  V->hhhh[alpha][map1][map2];
	map2+=1;
      }
      map1+=1;
    }
  }
  //calculating the [I5] term of [I6]
  for(k=0; k<nh; k++){
    for(l=0; l<nh; l++){
      for(i=0; i<nh; i++){ 
	for(j=0; j<nh; j++){	
	  temp = 0.0;
	  for(c=0; c<np; c++){
	    temp = temp + (barh_i07a[k][l][i][c]*t1[c][j] 
                 - barh_i07a[k][l][j][c]*t1[c][i]);
	  }
	  barh_09[k][l][i][j] += temp;
	}
      }
    }
  }
  //calculating second term of [I6]
  for(alpha=0; alpha<dim_alpha; alpha++){
    map1=0;
    for(p=0; p<ppbcount[alpha]; p+=2){
      c = B->c_pp[alpha][p];
      d = B->c_pp[alpha][p+1];
      map2=0;
      for(h=0; h<hhbcount[alpha]; h+=2){ 
	k = B->c_hh[alpha][h];
	l = B->c_hh[alpha][h+1];
	for(j=0; j<nh; j++){
	  for(i=0; i<nh; i++){ 
	    temp1 = V->pphh[alpha][map1][map2]*t2[c][d][i][j];
	    barh_09[k][l][i][j] += temp1;
	    barh_09[l][k][i][j] -= temp1;
	  }
	}
	map2+=1;
      }
      map1+=1;
    }
  }    
} //end [I6]
\end{lstlisting}

\begin{equation*}
[I7]_{d}^{a} =[I1]_d^a-[I2]_d^k\phantom{.}t_k^a-\frac{1}{2}\langle dc|v|kl\rangle t_{kl}^{ac}
\end{equation*}
%
\begin{lstlisting}[label={list:impl:intermediate7},caption={implementation of I7 in the amp1 class function ccsd\_uncoupled\_barh\_02\_store}]
void Amplitudes::ccsd_uncoupled_barh_02_store(){
 
  //initializing barh_02 ([I7]) by including [I1]
  for(int a=0; a<np; a++){
    for(int d=a; d<np; d++){
      barh_02[a][d] = barh_i02a[a][d];
      barh_02[d][a] = barh_i02a[d][a];
    }
  }
  //including the [I2] term
  for(a=0; a<np; a++){
    for(d=0; d<np; d++){ 
      for(k=0; k<nh; k++){
	barh_02[a][d] = barh_02[a][d] - barh_01[k][d] * t1[a][k];
      }
    }
  }
  //calculating the third term of [I7]
  for(alpha=0; alpha<dim_alpha; alpha++){
    for(a=0; a<np; a++){ 
      map1=0;
      for(p=0; p<ppbcount[alpha]; p+=2){
	d = B->c_pp[alpha][p];
	c = B->c_pp[alpha][p+1];
	map2=0;
	for(h=0; h<hhbcount[alpha]; h+=2){
	  k = B->c_hh[alpha][h];
	  l = B->c_hh[alpha][h+1];
	  temp1= V->pphh[alpha][map1][map2]*t2[a][c][k][l];
	  temp2= V->pphh[alpha][map1][map2]*t2[a][d][k][l];
	  barh_02[a][d] -=temp1;
	  barh_02[a][c] +=temp2;
	  map2+=1;
	}
	map1+=1;
      }
    }
  }
} //end[I7]
\end{lstlisting}

\begin{equation*}
[I8]_{cj}^{kb}=[I9]_{cj}^{kb}+\frac{1}{2}\langle dc|v|bk\rangle t_j^d-[I4]_{jc}^{lk}\phantom{.}t_l^b+\frac{1}{2}\langle cd|v|kl\rangle t_{jl}^{bd}  
\end{equation*}

\begin{lstlisting}[label={list:impl:intermediate8},caption={implementation of I8 in the amp1 class function ccsd\_uncoupled\_barh\_i10b\_store and ccsd\_uncoupled\_barh\_i10c\_store}]
void Amplitudes::ccsd_uncoupled_barh_i10b_store(){
  
  //initializing barh_i10b, equivalent to calculating [I9]
  for(int k=0; k<nh; k++){
    for(int b=0; b<np; b++){
      for(int c=b; c<np; c++){
	for(int j=k; j<nh; j++){
	  barh_i10b[k][b][c][j] = barh_i10a[k][b][c][j];
	  barh_i10b[j][b][c][k] = barh_i10a[j][b][c][k];
	  barh_i10b[k][c][b][j] = barh_i10a[k][c][b][j];
	  barh_i10b[j][c][b][k] = barh_i10a[j][c][b][k];
	}
      }
    }
  }
  //calculating second term of [I8]
  for(alpha=0; alpha<dim_alpha; alpha++){
    map1=0;
    for(p=0; p<ppbcount[alpha]; p+=2){
      d = B->c_pp[alpha][p];
      c = B->c_pp[alpha][p+1];
      map2=0;
      for(ph=0; ph<phbcount[alpha]; ph+=2){
	b = B->phbasis[alpha][ph];
	k = B->phbasis[alpha][ph+1];
	temp = 0.5 * V->ppph[alpha][map1][map2];
	for(j=0; j<nh; j++){
	  barh_i10b[k][b][c][j] += temp*t1[d][j];
	  barh_i10b[k][b][d][j] -= temp*t1[c][j];
	}
	map2+=1;
      }
      map1+=1;
    }
  }
  //calculating 0.5*[I4]
  for(k=0; k<nh; k++){ 
    for(b=0; b<np; b++){
      for(c=0; c<np; c++){
	for(j=0; j<nh; j++){
	  temp1=0.0;
	  for(l=0; l<nh; l++){ 
	    temp1 += 0.5 * barh_07[l][k][j][c]*t1[b][l];
	  }
	  barh_i10b[k][b][c][j]-=temp1;
	}
      }
    }
  }
} //end part of [I8]

void Amplitudes::ccsd_uncoupled_barh_i10c_store(){
 
 //initializing barh_i10c ([I8]) by including barh_i10b
  for(int k=0; k<nh; k++){
    for(int b=0; b<np; b++){
      for(int c=b; c<np; c++){
	for(int j=k; j<nh; j++){
	  barh_i10c[k][b][c][j] = barh_i10b[k][b][c][j];
	  barh_i10c[j][b][c][k] = barh_i10b[j][b][c][k];
	  barh_i10c[k][c][b][j] = barh_i10b[k][c][b][j];
	  barh_i10c[j][c][b][k] = barh_i10b[j][c][b][k];
	}
      }
    }
  }
  //calculating 0.5*[I4]
  for(k=0; k<nh; k++){ 
    for(b=0; b<np; b++){
      for(c=0; c<np; c++){
	for(j=0; j<nh; j++){
	  temp1=0.0;
	  for(l=0; l<nh; l++){ 
	    temp1 += 0.5 * barh_07[l][k][j][c]*t1[b][l];
	  }
	  barh_i10b[k][b][c][j]-=temp1;
	}
      }
    }
  }
  //calculating last term of [I8]
  for(b=0; b<np; b++){
    for(alpha=0; alpha<dim_alpha; alpha++){
      map1=0;
      for(p=0; p<ppbcount[alpha]; p+=2){
	d = B->c_pp[alpha][p];
	c = B->c_pp[alpha][p+1];
	map2=0;
	for(h=0; h<hhbcount[alpha]; h+=2){
	  l = B->c_hh[alpha][h];
	  k = B->c_hh[alpha][h+1];
	  temp1 = 0.5*V->pphh[alpha][map1][map2];
	  for(j=0; j<nh; j++){
	    barh_i10c[k][b][c][j] += temp1*t2[b][d][j][l];
	    barh_i10c[k][b][d][j] -= temp1*t2[b][c][j][l];
	    barh_i10c[l][b][c][j] -= temp1*t2[b][d][j][k];
	    barh_i10c[l][b][d][j] += temp1*t2[b][c][j][k];
	  }
	  map2+=1;
	}
	map1+=1;
      }
    }
  }
}//end [I8]
\end{lstlisting}

\begin{equation*}
[I9]_{cj}^{kb} = -\langle bk|v|cj\rangle + \frac{1}{2}\langle dc|v|bk\rangle t_j^d 
\end{equation*}
\begin{lstlisting}[label={list:impl:intermediate9},caption={implementation of I9 in the amp1 class function ccsd\_uncoupled\_barh\_i10a\_store}]
void Amplitudes::ccsd_uncoupled_barh_i10a_store(){
  
  //calculating first term of [I9] (barh_i10a)
  for(alpha=0; alpha<dim_alpha; alpha++){
    map1=0;
    for(ph=0; ph<phbcount[alpha]; ph+=2){ 
      b = B->phbasis[alpha][ph];
      k = B->phbasis[alpha][ph+1];
      map2=0;
      for(ph2=0; ph2<phbcount[alpha]; ph2+=2){
	c = B->phbasis[alpha][ph2];
	j =  B->phbasis[alpha][ph2+1];  
	barh_i10a[k][b][c][j] = -V->phph[alpha][map1][map2];
	map2+=1;
      }
      map1+=1;
    }
  }
  //calculating second term of [I9]
  for(alpha=0; alpha<dim_alpha; alpha++){
    map1=0;
    for(p=0; p<ppbcount[alpha]; p+=2){ 
      d = B->c_pp[alpha][p];
      c = B->c_pp[alpha][p+1];
      map2=0;
      for(ph=0; ph<phbcount[alpha]; ph+=2){
	b = B->phbasis[alpha][ph];
	k = B->phbasis[alpha][ph+1];
	temp = 0.5 * V->ppph[alpha][map1][map2];
	for(j=0; j<nh; j++){
	  barh_i10a[k][b][c][j] += temp*t1[d][j];
	  barh_i10a[k][b][d][j] -= temp*t1[c][j];
	}
	map2+=1;
      }
      map1+=1;
    }
  }
} //end [I9]
\end{lstlisting}

\begin{equation*}
[I10]_{ij}^{kb}=-\langle bk|v|ij\rangle-\frac{1}{2} \langle cd|v|bk\rangle t_{ij}^{cd}+\hat{P}_{(ij)}[I9]_{cj}^{kb}\phantom{.}t_i^c-\frac{1}{2} [I6]_{ij}^{kl}t_l^b 
\end{equation*}

\begin{lstlisting}[label={list:impl:intermediate10},caption={implementation of I10 in the amp1 class function ccsd\_uncoupled\_barh\_i12a\_store}]
void Amplitudes::ccsd_uncoupled_barh_i12a_store(){
  
  //calculating first term of [I10] (barh_i12a)
  for(alpha=0; alpha<dim_alpha; alpha++){
    map1=0;
    for(ph=0; ph<phbcount[alpha]; ph+=2){
      b = B->phbasis[alpha][ph];
      k = B->phbasis[alpha][ph+1];
      map2=0;
      for(h=0; h<hhbcount[alpha]; h+=2){
	i = B->c_hh[alpha][h];
	j = B->c_hh[alpha][h+1];
	barh_i12a[k][b][i][j] = -V->phhh[alpha][map1][map2];
	barh_i12a[k][b][j][i] = +V->phhh[alpha][map1][map2];
	map2+=1;
      }
      map1+=1;
    }
  }
  // calculating the [I9] term of [I10]
  for(k=0; k<nh; k++){
    for(b=0; b<np; b++){
      for(j=0; j<nh; j++){
	for(i=0; i<nh; i++){
    	  temp = 0.0;
	  for(c=0; c<np; c++){
	    temp = temp + (barh_i10a[k][b][c][j]*t1[c][i] 
                 - barh_i10a[k][b][c][i]*t1[c][j]);
	  }
	  barh_i12a[k][b][i][j] = barh_i12a[k][b][i][j] + temp;
	}
      }
    }
  }
  // calculating the [I6] term of [I10]
  for(k=0; k<nh; k++){
    for(b=0; b<np; b++){
      for(j=0; j<nh; j++){
	for(i=0; i<nh; i++){	
	  temp = 0.0;
	  for(l=0; l<nh; l++){
	    temp = temp + barh_09[k][l][i][j]*t1[b][l];
	  }
	  barh_i12a[k][b][i][j] = barh_i12a[k][b][i][j] - 0.5*temp;
	}
      }
    }
  }
  //calculating second term of [I10]
  for(alpha=0; alpha<dim_alpha; alpha++){
    map1=0;
    for(p=0; p<ppbcount[alpha]; p+=2){
      c = B->c_pp[alpha][p];
      d = B->c_pp[alpha][p+1];
      map2=0;
      for(ph=0; ph<phbcount[alpha]; ph+=2){
	b = B->phbasis[alpha][ph];
	k = B->phbasis[alpha][ph+1];
	for(i=0; i<nh; i++){
	  for(j=0; j<nh; j++){
	    barh_i12a[k][b][i][j] -= V->ppph[alpha][map1][map2]*t2[c][d][i][j];
	  }
	}
	map2+=1;
      }
      map1+=1;
    }
  }
} //end [I10]
\end{lstlisting}

\begin{equation*}
[I11]_{cj}^{ab} = \langle ab|v|cj\rangle +\frac{1}{2}\langle ab|v|cd\rangle t_j^d 
\end{equation*}
\begin{lstlisting}[label={list:impl:intermediate11},caption={implementation of I11 in the amp1 class function ccsd\_uncoupled\_barh\_i11a\_store}]
void Amplitudes::ccsd_uncoupled_barh_i11a_store(){
   
  //calculating first term of [11] (barh_i11a)
  for(alpha=0; alpha<dim_alpha; alpha++){
    map1=0;
    for(p=0; p<ppbcount[alpha]; p+=2){
      a = B->c_pp[alpha][p];
      b = B->c_pp[alpha][p+1];
      map2=0;
      for(ph=0; ph<phbcount[alpha]; ph+=2){
	c = B->phbasis[alpha][ph];
	j = B->phbasis[alpha][ph+1];
	barh_i11a[a][b][c][j] =  V->ppph[alpha][map1][map2];
	barh_i11a[b][a][c][j] = -V->ppph[alpha][map1][map2];
	map2+=1;
      }
      map1+=1;
    }
  }
  //calculating second term of [11]
  for(alpha=0; alpha<dim_alpha; alpha++){
    for(j=0; j<nh; j++){
      map1=0;
      for(p=0; p<ppbcount[alpha]; p+=2){
	a = B->c_pp[alpha][p];
	b = B->c_pp[alpha][p+1];
	map2=0;
	for(p2=0; p2<ppbcount[alpha]; p2+=2){ 
	  c = B->c_pp[alpha][p2];
	  d = B->c_pp[alpha][p2+1];
	  temp3 = 0.5*V->pppp[alpha][pmap][p2map];
	  temp1 = temp3*t1[f][i];
	  temp2 = temp3*t1[e][i]; 
	  barh_i11a[a][b][c][j] += temp1;
	  barh_i11a[b][a][c][j} -= temp1;
	  barh_i11a[a][b][d][j] -= temp2;
	  barh_i11a[b][a][d][j] += temp2;
	  map2+=1;
	}
	map1+=1;
      }
    }
  }
}//end [I11]
\end{lstlisting}


\section{Implementation of the double dot}
We will in this section outline the numerical methods and alogrithms related to finding the one-electron eigenvalues of double dot potential, see section \ref{section:the double dot potential}. The code can take a general potential in two dimensions, and find the eigenvalues and eigenvectors. But for our particular problem, we only diagonalize in $x$-direction, find the eigenvalues and eigenvectors. Then find the overlap coefficients with an harmonic oscillator basis
% 
\begin{equation}
  C_{n_x',n_X m_X} = \sum_{n_x}^N \braket{n_x'}{n_X m_X} \ket{n_X m_X} 
\end{equation}
%
Where $n_X$ and $m_X$ are the radial and angular quantum numbers from a polar basis solution, see Eq.~(\ref{total eigenstate}). Since there is no barrier in the $y$-direction, the hamiltonian \ref{rerescaled Hamiltonian} is separable and we get 
\begin{equation}
  C_{n_Y' m_Y', n_Y m_Y}  = \delta_{n_Y' m_Y', n_Y m_Y}  \ket{n_X m_X} 
\end{equation}
%
The relations between cartesian quantum numbers $(n_x,n_y)$ and $n,m$ are as follows 
\begin{equation}
 E = n_x + n_y + 1 = 2n + |m| + 1
\end{equation}
%
Then our new basis states can be expressed as linear combination of harmonic oscillator states 
%
\begin{equation}
 \ket{a} = \sum_{\alpha} C_\alpha \ket{\alpha} , \qquad C_\alpha =  C_{n_x',n_X m_X} C_{n_Ym_Y} 
 \label{new spinorbitals}
\end{equation}
%
Our goal is to find the two-particle interaction elements in the new basis
%
\begin{equation}
 \bra{ab}v\ket{cd} = \sum_{\alpha\beta\gamma\delta} C_\alpha^*C_\beta^*C_\gamma C_\delta \bra{\alpha\beta}v\ket{\gamma\delta}
 \label{new interactions}
\end{equation}
%
The interaction elements for the parabolic qunatum dot Eq.~\ref{Hamiltonian for QD} $\bra{\alpha\beta}v\ket{\gamma\delta}$ can be obtained from Simen Kvaal's Open FCI code for quantum dots \cite{kvaalOpenFCI}. 
%
This is a transformation of the interaction elements from a polar basis to a cartesian basis. And these basis transformations are costly to calculate, we have tried optimalize as much as we could. 

These new interaction elements Eq.~\ref{new interactions}, together with the new spinorbital quantum numbers Eq.~\ref{new spinorbitals} $\ket{a}$ and single-particle energies $\braket{a}{b}$ are used as an input in our CCSD and Hartree-Fock. 

\subsection{Scaling the Hamiltonian}
One of the potentials we wish to study \cite{doubledot}

\begin{equation}
V(x,y) = \frac{1}{2} m^* \omega_0^2 [x^2 + y^2 -2L_x|x| + L_x^2]
\end{equation}
%
We want our potential to dimensionless in our calculation. Introducing the scaled variables
\begin{align}
 L_x & = L_c \bar{L}_x \\
   x & = x_c \bar{x}   \\ 
   y & = y_c \bar{y}   \\
   \omega_0 &= \omega_c \bar{\omega}
\end{align}
Setting 
\begin{align}
 L_c & = x_c = y_c = l_0 \\
\end{align}
We get
\begin{equation}
  \bar{V} = \frac{1}{2} m^* \omega_0 l_0^2 \left[\bar{x}^2 + \bar{y}^2 - 2|\bar{x}| + \bar{L}_x \right]
\end{equation}
%
And setting
\begin{equation}
l_0 = \sqrt{\frac{\hbar}{m^* \omega_c \bar{\omega}}}
\end{equation}
%
Which gives
\begin{equation}
  \bar{V}_C  =  \frac{\hbar}{2} \omega_c \bar{\omega} \left[\bar{x}^2 + \bar{y}^2 - 2|\bar{x}| + \bar{L}_x \right]
    \label{Vabs}
\end{equation}
%
Defining $E_h$

\begin{equation}
 E_h = \frac{m^*}{\epsilon^2} \qquad \kappa = \frac{e^2}{4\pi \epsilon_0 \epsilon_r \hbar}
\end{equation}
%
Scaling the Hamiltonian $\hh{H} = E_h \bar{H}$, where $E_h$ is effective Hartrees.
 
\begin{equation}
  \bar{H} = -\frac{\omega_c \bar{\omega} \hbar \kappa^2}{2m^*} \bar{\nabla}^2 + \frac{\hbar\kappa^2}{2m^*}\omega_c\bar{\omega} \left[\bar{x}^2 + \bar{y}^2 - 2|\bar{x}| + \bar{L}_x \right]
\end{equation}
%
Setting 
\begin{equation}
 \frac{\hbar \kappa^2  \omega_c}{m^*} = 1
\end{equation}
%
We then get the one-body Hamiltonian
%
\begin{equation}
 \bar{H} = -\frac{\bar{\omega}}{2} \nabla + \frac{1}{2} \bar{\omega} \left[ \bar{x}^2 + \bar{y}^2 - 2|\bar{x}| + \bar{L}_x \right]
 \label{doublepotential abs}
\end{equation}
%
The two-body interaction is the same as in Eq.~\ref{N-body Hamiltonian scaled}. And the total Hamiltonian for our model becomes
%
\begin{equation}
 \bar{H} = -\frac{\bar{\omega}}{2} \nabla + \frac{1}{2} \bar{\omega} \left[ \bar{x}^2 + \bar{y}^2 - 2|\bar{x}| + \bar{L}_x \right] + \sqrt{\bar{\omega}}\sum_{i<j}^N \frac{1}{\bar{r}_{ij}}
  \label{total hamiltonian double dot}
\end{equation}
%
In the article \cite{doubledot}, they set $L_x = 50$ nm , then we can choose $L_c = 50$ nm and setting $\bar{L}_x = 1$, $\bar{\omega}=1$. 
%
\begin{figure}[H]
\centering
\scalebox{0.7}{\includegraphics{Doubledot3D.eps}}
\caption{Confinement potential for Eq.~\ref{Vabs} $V(\bar{x},\bar{y})$, $\bar{L}_x = 10$ and $\bar{\omega} = 1$}
\end{figure}
%
\begin{figure}[H]
\centering
\scalebox{0.7}{\includegraphics{Doubledot2D.eps}}
\caption{Confinement potential for Eq.~\ref{Vabs} $V(\bar{x},0)$, $\bar{L}_x = 10$ and $\bar{\omega} = 1$}
\end{figure}
%
The code is so general that we can change the potential, instead of the absolute value barrier $|\bar{x}|$, we can choose a gaussian curve barrier which is smoother. An example could be 
%
\begin{equation}
\bar{V}_G = \frac{1}{2}\bar{\omega}\left[\bar{x}^2 + \bar{y}^2 \right] + V_0 \exp\left(-\frac{\bar{x}^2}{2\sigma}\right)
  \label{Vgauss}
\end{equation}
%
\begin{figure}[H]
\centering
\scalebox{0.7}{\includegraphics{Doubledot3D_gaussian.eps}}
\caption{Confinement potential for Eq.~\ref{Vgauss} $V(\bar{x},\bar{y})$, $V_0 = 100$ and $\sigma=1, \omega = 1$}
\label{fig:Vgauss3D}
\end{figure}
%
\begin{figure}[H]
\centering
\scalebox{0.7}{\includegraphics{Doubledot2D_gaussian.eps}}
\caption{Confinement potential for Eq.~\ref{Vgauss} $V(x,0)$, $V_0 = 20,\sigma=1$ and $\bar{\omega} = 1$}
\label{fig:Vgauss2D}
\end{figure}
%

\subsection{Finding the eigenvectors and eigenvalues}
%
We diagonalize the one-body Hamiltonian in $x$-direction
\begin{equation}
\hh{H}_X = -\frac{\bar{\omega}}{2}\frac{d}{d\bar{x}^2} + V(\bar{x}) 
  \label{we want to diagonalize}
\end{equation}
%
In order to find the eigenvalues and the eigenvectors we can use the technique of discretization described in section \ref{section:the double dot potential}. We use the algorithm (Listing \ref{list:fillingtridiagonal}) to create the tridiagonal matrix Eq.~\ref{tridiagonal}

\begin{lstlisting}[label={list:fillingtridiagonal},caption={Filling the tridiagonal matrix}]
    double e = 0.5 * (-1.0) / (xstep * xstep);
    double f = 1.0 / (xstep * xstep);
    double V = 0;
    for (int i = 0; i < N - 1; i++) {
        x = xmin + i*xstep;
        V = potential(x);

        //Filling the diagonal
        T[i][i] = f + V;

        //Filling the off diagonal;
        T[i][i + 1] = e;
        T[i + 1][i] = e;
    }
    //Filling the last diagonal element
    x = xmin + (N - 1) * xstep;
    V = potential(x);
    T[N - 1][N - 1] = f + V;
\end{lstlisting}
 %
We use the \texttt{LAPACK} routine for solving a symmetric dense matrix, \texttt{dsyev()}, which basically use the accelerated QR-Algorithm, section \ref{section:The QR-Algorithm}. The eigenvalues $E_{n_x'}$ are then sorted in ascending order,

\begin{equation}
 n_{x}' = 0,1,2...N \qquad E_{0} \geq E_1 \geq E_2 \geq ... \geq E_{n_x' = N}
 \label{eigenvalues found by diag}
\end{equation}


\begin{equation}
  \psi_{n_x'} = \begin{bmatrix} \psi_{n_x'}(x_{\min}) \\ \psi_{n_x'}(x_1) \\ \vdots \\ \psi_{n_x}(x_N) \end{bmatrix}, \qquad x_j = x_{\min} + j\cdot h, \qquad j \in \{0,1,..,N\}                           
\end{equation}
%
\begin{figure}[h!]
\centering
$\begin{array}{cc}
\includegraphics[width=2.1in]{1D_Lx=5_0} & \includegraphics[width=2.1in]{1D_Lx=5_1} \\ \includegraphics[width=2.1in]{1D_Lx=5_2} & \includegraphics[width=2.1in]{1D_Lx=5_3}
\end{array}$
\caption{The probability density of the four lowest eigenstates for $L_x = 5,\, \omega = 1$.}
\label{wavefunctions1}
\end{figure}

\begin{figure}[h!]
\centering
$\begin{array}{cc}
\includegraphics[width=2.1in]{1D_V=20_0} & \includegraphics[width=2.1in]{1D_V=20_1} \\ \includegraphics[width=2.1in]{1D_V=20_2} & \includegraphics[width=2.1in]{1D_V=20_3}
\end{array}$
\caption{The probability density of the four lowest eigenstates for $V_0 = 20,\, \omega = 1$.}
\label{wavefunctions2}
\end{figure}


\subsection{Finding the coefficients}
We want to find the overlap coefficients
\begin{equation}
C_{n_x',n} = \braket{\psi_{n_x'}}{n} = \sum_{n}^N \psi_{n_x'}(x_i) \psi_n(x_i) \qquad n = 0,1,2,...N
  \label{overlap}
\end{equation}
%
\begin{equation}
x_i = x_{\min} + i \cdot h \qquad h = \frac{x_{\max} - x_{\min}}{N} 
\end{equation}

Where $\psi_n(x)$ are the harmonic oscialltor eigenfuncions in one dimension:
%
\begin{align}
  \psi_n(x) & = \left(\frac{\beta}{\pi}\right)^{1/4} \frac{1}{\sqrt{2^n n!}} H_n(x)e^{\frac{\beta x^2}{2}}
\end{align}
%
\begin{equation}
 \beta = \frac{m \omega}{\hbar}
\end{equation}
%
$H_n(x)$ are the Hermite polynomials \cite{boas}. Using the recursive relation for the Hermite polynomials 
%
\begin{align}
H_{n+1}(x) = 2xH_n(x) - 2nH_{n-1}(x)
\end{align}
%
We get
\begin{align}
  \psi_{n+1} & = \left(\frac{1}{\pi}\right)^{1/4} \frac{1}{\sqrt{2^{n+1} (n+1)!}} H_{n+1}e^{\frac{1 x^2}{2}} \\
             & = \left(\frac{1}{\pi}\right)^{1/4} \frac{1}{\sqrt{2^{n+1} (n+1)!}} \left( 2xH_n - 2nH_{n-1} \right) \\
             & = x\sqrt{\frac{2}{n+1}}\psi_n - \sqrt{\frac{n}{n+1}}\psi_{n-1} 
\end{align}
%
Or 
%
\begin{align}
  \psi_{n} &= x\sqrt{\frac{2}{n}}\psi_{n-1} - \sqrt{\frac{n-1}{n}}\psi_{n-2} 
\end{align}
%
Where
%
\begin{align}
  \psi_0 & = \left(\frac{1}{\pi}\right)^{1/4} e^{-\frac{ x^2}{2}} \\
  \psi_1 & = \left(\frac{1}{\pi}\right)^{1/4} \frac{1}{\sqrt{2}} 2x e^{-\frac{ x^2}{2}}
\end{align}
%
The algorithm for generating the harmonic oscillator wavefunctions and evalute at $x$, $\beta = 1$
%
\begin{lstlisting}
double wf(int n, double x) {
    double *W = new double[n + 2];
    double a = (1.0 / pow(pi, 0.25));
    double b = (1.0 / sqrt(2));

    W[0] = 1.0 * a; //psi_0(x)

    if (n == 0) {
        return W[0] * exp(-0.5 * x * x);
    } else {
        W[1] = sqrt(2) * x * a; //psi_1(x)
        for (int i = 1; i <= n - 1; i++) {

            W[i + 1] = (sqrt(2.0 / (i + 1)) * x * W[i] - sqrt(double(i) / (i + 1)) * W[i - 1]);
        }
        return W[n] * exp(-0.5 * x * x);
    }

    delete [] W;
\end{lstlisting}
%
The coefficients are then calculated by the following algorithm
%
\begin{lstlisting}
    for (int nxprime = 0; nxprime < nxprime_max; nxprime++) {
        for (int n = 0; n < nmax; n++) {
            innerprod = 0;
            for (int i = 0; i < N; i++) {
                innerprod += xstep * Evecs[i][nxprime] * H[i][n];
            }
            Coeff[n][nxprime] = innerprod;
        }
    }
\end{lstlisting}

\subsubsection{Validation}
%
One way to validate the coefficients is to check if it preserve the probability.
\begin{equation}
  P(n_x') = \sum_{i}^N C_{i,n_x'}^*C_{i,n_x'} \approx 1
\end{equation}
%
Another check is to plot the diagonalized curve against the linear expansion and see if the $L_2$-norm of the function is a small number
\begin{equation}
  ||\psi_{n_x'} - \psi_{n_x}||_{2} = \sqrt{ \sum_i^N |\psi_{n_x'}(x_i) - \psi_{n_x}(x_i)|^2} < \epsilon
\end{equation}

For the $L_x = 5, \omega = 1$ we need $n_x > 30$ in order to get reasonable amount of non-zero coefficients $C_{n_x',n_x}$. But our calculation is with a max $n_x=20$ due to the heavy time consumption of generating the transformed interaction elements. Besides we are studing small model spaces and only need $n_x' < 10$, in that regard $C_{n_x',n_x}$ falls off quickly and it is sufficient to have maximal $n_x = 20$ 

\begin{figure}[h!]
\centering
$\begin{array}{cc}
\includegraphics[width=2.5in]{wf_fit_nx=4} & \includegraphics[width=2.5in]{wf_fit_nx=10}
\end{array}$
\caption{Plot of the fitted curve for $L_x=5$, $\omega=1$, and the max value of $n_x=20$. The $L_2$-norm for $n_x'=10$  was 5.8E-4, while the $L_2$-norm for $n_x'=10$ was 0.0118. Which is a reasonable fit}
\end{figure}


\subsection{Transformation from polar to cartesian representation}
\label{sec:cartesian basis transform}

Harmonic oscillator wavefunctions in two dimensions for som selected values
%
\begin{align}
\psi^{\text{\tiny HO}}_{n_x,n_y}(x,y) & = \psi^{\text{\tiny HO}}_{n_x}(x)\psi^{\text{\tiny HO}}_{n_y}(y)   
  \label{Harmonic Oscillator Wavefunctions Cartesian}
\end{align}
%
\begin{align}
\psi^{\text{\tiny HO}}_{n_x,n_y}(x,y) = H_{n_x}(x)H_{n_y}(y) e^{-x^2} e^{-y^2}  
  \label{Harmonic Oscillator Wavefunctions Cartesian Explicit}
\end{align}
In cartesian representation (not normalized): 
\begin{align*}
\psi^{\text{\tiny HO}}_{0,0}(x,y) & =  e^{-x^2}e^{-y^2}   & (\epsilon = 1) \\
\psi^{\text{\tiny HO}}_{1,0}(x,y) & =  xe^{-x^2}e^{-y^2} & (\epsilon = 2) \\
\psi^{\text{\tiny HO}}_{0,1}(x,y) & =  ye^{-x^2}e^{-y^2} & (\epsilon = 2) \\
\psi^{\text{\tiny HO}}_{1,1}(x,y) & =  yxe^{-x^2/2}e^{-y^2} & (\epsilon = 3) \\ 
\end{align*}
%
Where $\epsilon = n_x + n_y + 1$.
%
In polar coordinate representation (not normalize):
\begin{align}
\psi^{\text{\tiny HO}}_{n,m}(r,\theta) & = R(r)\phi(\theta)  
  \label{Harmonic Oscillator Wavefunctions Polar}
\end{align}
%
\begin{align}
\psi^{\text{\tiny HO}}_{n,m}(r,\theta) & = r^{|m|}e^{-r^2}L_n^{|m|}e^{im\theta}  
  \label{Harmonic Oscillator Wavefunctions Polar Explicit}
\end{align}
\begin{align*}
\psi^{\text{\tiny HO}}_{0,0}(r,\theta)   & =  e^{-r^2}   & (\epsilon = 1) \\
\psi^{\text{\tiny HO}}_{0,1}(r,\theta)   & =  r e^{-r^2}e^{i\theta} & (\epsilon = 2) \\
\psi^{\text{\tiny HO}}_{0,-1}(r,\theta)  & =  r e^{-r^2}e^{-i\theta} & (\epsilon = 2) \\
\psi^{\text{\tiny HO}}_{1,0}(r,\theta)   & =  (-r^2+1)e^{-r^2} & (\epsilon = 3) \\ 
\end{align*}
%
Where $\epsilon = 2n + |m| + 1$. We transformation between the representation by the following
\begin{align*}
x & = r \cos(\theta) \\
y & = r \sin(\theta) \\
r^2 & = x^2 + y^2
  \label{def:CartesianPolarMapping}
\end{align*}
%
We have that
%
\begin{equation}
\cos(\theta) = \frac{e^{i\theta} + e^{i\theta}}{2} \qquad \cos(\theta) = \frac{e^{i\theta} - e^{i\theta}}{2i}
  \label{cosinus relation}
\end{equation}
%
We see that (normalized):
\begin{align}
\psi^{\text{\tiny HO}}_{0,0}(x,y) & = \psi^{\text{\tiny HO}}_{0,0}(r,\theta)\\
\psi^{\text{\tiny HO}}_{1,0}(x,y) & = \frac{1}{\sqrt{2}}\psi^{\text{\tiny HO}}_{1,1}(r,\theta) +   \frac{1}{\sqrt{2}}\psi^{\text{\tiny HO}}_{1,-1}(r,\theta)\\
\psi^{\text{\tiny HO}}_{0,1}(x,y) & =  \frac{i}{\sqrt{2}}\psi^{\text{\tiny HO}}_{1,-1}(r,\theta) -   \frac{i}{\sqrt{2}}\psi^{\text{\tiny HO}}_{1,1}(r,\theta)\\
  \label{we see mapping} 
\end{align}
%
The Transformation matrix $T^{N}$, $N = n_x + n_y$ :

\begin{align}
T^{0}  & = 1 \\
T^{1}  & = \begin{bmatrix*}[r] \frac{1}{\sqrt{2}} & \frac{1}{\sqrt{2}} \\ \frac{1}{\sqrt{2}} & -\frac{1}{\sqrt{2}}
 \end{bmatrix*} \\
T^{2}  & = \begin{bmatrix*}[r] \frac{1}{2} & \frac{1}{\sqrt{2}} & \frac{1}{2} \\ \frac{1}{\sqrt{2}} & 0 & -\frac{1}{\sqrt{2}} \\ \frac{1}{2} & -\frac{1}{\sqrt{2}} & \frac{1}{2} \end{bmatrix*} \\
T^{N} & = ...
\end{align}
%
These transformation matrices are obtained from Simen Kvaals PhD-thesis \cite{kvaalphd}
Our transformation from a polar basis to a cartesian one are defined,
\begin{equation}
\ket{a} \equiv \ket{n_x'(a), n_y'(a)}^{CAR} = \sum_{n_x(a)} C_{n_x(a), n_x'(a)} \sum_{n_+(a)}^{N_a} i^{-n_y'(a)} T_{n_+(a),n_x(a)}^{N_a}
  \label{DoubleDot}
\end{equation}
%
Where $N_a = n_x(a) + n_y'(a)$. 
%
Then the interaction elements in the cartesian space represented by 
%
\begin{align*}
\bra{ab}v\ket{cd}^{CAR} &= \sum_{n_x(a)}\sum_{n_x(b)}\sum_{n_x(c)}\sum_{n_y(d)} C_{n_x(a), n_x'(a)} C_{n_x(b), n_x'(b)} C_{n_x(c), n_x'(c)} C_{n_x(d), n_x'(d)} \\
&\times \sum_{n_+(a)}^{N_a}\sum_{n_+(b)}^{N_b}\sum_{n_+(c)}^{N_c}\sum_{n_+(d)}^{N_d} T_{n_+(a),n_x(a)}^{N_a} T_{n_+(b),n_x(b)}^{N_b} T_{n_+(c),n_x(c)}^{N_c} T_{n_+(d),n_x(d)}^{N_d} \\
&\times \bra{\underbrace{n_+(a)n_-(a)}_{a}|\underbrace{n_+(b)n_-(b)}_b}v\ket{\underbrace{n_+(c)n_-(c)}_c|\underbrace{n_+(d)n_-(d)}_d}^{POL} \\
&\times i^{n_y'(a)} i^{n_y'(b)} i^{-n_y'(c)} i^{-n_y'(d)}
\label{cartesian interaction element representation}
\end{align*}
%
Where 
\begin{align}
  m = n_+ - n_- \qquad n = \frac{n_+ + n_- - |m|}{2} 
\end{align}
%
Think of it as a ladder, where number of $n_+$, defines how high to the right you move on the ladder, and $n_-$ defines how high to the left.
%
Therefore each of the subshells will be defined by a set of $\{n_+,n_-\}$. Notice that we have a imaginare part. Our interaction elements in the polar basis are real, and the transformation will not change this, therefore the transformations will have the condition. 
%
\begin{equation}
\left[ n_y(a) + n_y(b) - n_y(c) - n_y(d) \right] = 2n+1 \iff  \bra{ab}v\ket{cd}^{CAR} = 0, \qquad n\in \mathcal{N}
\end{equation}


%
\begin{figure}
\centering
\input{simenmap}
\caption{Shell structure of the quantum dot, and the relations to the quantum number $n_+$, $n_-$. There is a 1-1 map between the radial quantum number $n$ and angular momentum quantum number $m_s$ to $\{n_+,n_-\}$ }
\label{fig:simenmap}
\end{figure}


We can do a simple example using the parabolic harmonic oscillator potential, i.e. $L_x=0$ and $V_0 = 0$. Our system is a harmonic oscillator with 2 shell level, and the basis states are:
\begin{align*}
\ket{0} \rightarrow &n_x'=0, n_y'=0, m_s' = -0.5, &E = 0.99988 \\  
\ket{1} \rightarrow &n_x'=0, n_y'=0, m_s' = 0.5,  &E = 0.99988 \\  
\ket{2} \rightarrow &n_x'=1, n_y'=0, m_s' = -0.5, &E = 1.99943 \\  
\ket{3} \rightarrow &n_x'=1, n_y'=0, m_s' =  0.5, &E = 1.99943 \\  
\ket{4} \rightarrow &n_x'=0, n_y'=1, m_s' = -0.5, &E = 1.99988 \\  
\ket{5} \rightarrow &n_x'=0, n_y'=1, m_s' =  0.5, &E = 1.99988  
\end{align*}
Polar basis:
\begin{align*}
  \ket{0} & = \{n=0, m=0, m_s = -0.5\} \\ 
  \ket{1} & = \{n=0, m=0, m_s = 0.5\} \\
  \ket{2} & = \{n=0, m=-1, m_s = -0.5\} \\
  \ket{3} & = \{n=0, m=-1, m_s = 0.5\} \\
  \ket{4} & = \{n=0, m=+1, m_s = -0.5\} \\
  \ket{5} & = \{n=0, m=+1, m_s = 0.5\}   
\end{align*}
Talmi matrix:
\begin{align}
  T^{(1)} &= 1 \\
  T^{(2)} &= \begin{bmatrix}\frac{1}{\sqrt{2}} & \frac{1}{\sqrt{2}} \\ \frac{1}{\sqrt{2}} &- \frac{1}{\sqrt{2}} \end{bmatrix}                             
\end{align}
%
Want to find the interaction element $\bra{45}v\ket{45}^{CAR}$:
%
\begin{equation}
  \ket{n_x' n_y'}^{CAR} = \sum_{n_x} C_{n_x',n_x} \sum_{n_+}^{n_x+n_y'} T_{n_+,n_x}^{n_x+n_y'} \ket{n_+,n_-}^{POL}
\end{equation}

\begin{align}
  \ket{4}^{CAR} & = T_{0,0}^{(1)} \ket{(n_+=0,n_-=1)}^{POL} + T_{1,0}^{(1)} \ket{(n_+=1,n_-=0)}^{POL} \\
  \ket{5}^{CAR} & = T_{0,0}^{(1)} \ket{(n_+=0,n_-=1)}^{POL} + T_{1,0}^{(1)} \ket{(n_+=1,n_-=0)}^{POL}  
\end{align}
%
If we translate these states into the Polar basis quantum numbers we get, and remember that $\ket{4}$ is a spin down states, so $\{n_+,n_-\}$ must map to a state which is spin down:
\begin{align}
  \ket{4}^{CAR} & = T_{0,0}^{(1)} \ket{4}^{POL} + T_{1,0}^{(1)} \ket{2}^{POL} \\
  \ket{5}^{CAR} & = T_{0,0}^{(1)} \ket{5}^{POL} + T_{1,0}^{(1)} \ket{3}^{POL}  
\end{align}
%
Then 
\begin{align}
  \ket{45}^{CAR} & = (T_{0,0}^{(1)})^2 \ket{4 5}^{POL} + T_{0,0}^{(1)}T_{1,0}^{(1)} \ket{43}^{POL} + T_{1,0}^{(1)}T_{0,0}^{(1)} \ket{25}^{POL} + (T_{1,0}^{(1)})^2 \ket{23}^{POL}  
\end{align}
%
Then 
\begin{align}
  \bra{45}v\ket{45}^{CAR} & = \frac{1}{2}\frac{1}{2} \bra{45}v\ket{45}^{POL} + \frac{1}{2}\frac{1}{2} \bra{43}v\ket{25}^{POL}  + \frac{1}{2}\frac{1}{2} \bra{43}v\ket{43}^{POL} \\
  &+ \frac{1}{2}\frac{1}{2} \bra{25}v\ket{25}^{POL} + \frac{1}{2}\frac{1}{2} \bra{25}v\ket{43}^{POL}+ \frac{1}{2}\frac{1}{2} \bra{23}v\ket{23}^{POL}  
\end{align}
%
List of interaction elements from \texttt{tabulate.x} with $\omega = 1.0$ and standard v.v.:
%
\begin{align}
\bra{45}v\ket{45}^{POL} & =  0.86165\\
\bra{43}v\ket{25}^{POL} & =  0.23499\\
\bra{43}v\ket{43}^{POL} & =  0.86165\\
\bra{25}v\ket{25}^{POL} & =  0.86165\\
\bra{25}v\ket{43}^{POL} & =  0.23499\\
\bra{23}v\ket{23}^{POL} & =  0.86165
\end{align}
%
The our interaction element becomes
\begin{align}
  \bra{45}v\ket{45}^{CAR} & = 0.86165 + \frac{1}{2}0.23499 = 0.97915
\end{align}
%
\begin{align}
  \bra{45}v\ket{45}^{POL} & = 0.86165
\end{align}

\subsection{Tabulating new quantum numbers}
In our cartesian basis, the two-particle interaction still does not change spin. Therefore the total spin is conserved.
 
\begin{equation}
  \bra{M_s}v\ket{M_s'} = \delta_{M_s,M_s'} \qquad M_s = m_{1s} + m_{2s} 
\end{equation}

\begin{table}[H]
\centering
\begin{minipage}[b]{0.3\linewidth}\centering 
\begin{tabular}{ccc}
\toprule
$M_s$ & $\alpha$ & State\\
\midrule
 -1 &  0  &$\ket{0,2}$ \\
 -1 &  0  &$\ket{0,4}$ \\
 -1 &  0  &$\ket{2,4}$ \\
\bottomrule
\end{tabular}
\caption{Tabulated two-particle basis for $M_s=-1$}
\end{minipage}
\hspace{0.2cm}
\begin{minipage}[b]{0.3\linewidth}
\begin{tabular}{ccc}
\toprule
$M_s$ & $\alpha$ & State\\
\midrule
  0 &  1  &$\ket{0,1}$ \\
  0 &  1  &$\ket{0,3}$ \\
  0 &  1  &$\ket{0,5}$ \\
  0 &  1  &$\ket{1,2}$ \\
  0 &  1  &$\ket{1,4}$ \\
  0 &  1  &$\ket{2,3}$ \\
  0 &  1  &$\ket{2,5}$ \\
  0 &  1  &$\ket{3,4}$ \\
  0 &  1  &$\ket{4,5}$ \\
  0 &  1  &$\ket{1,3}$ \\
\bottomrule
\end{tabular}
\caption{Tabulated two-particle basis for $M_s=0$}
\end{minipage}
\hspace{0.2cm}
\begin{minipage}[b]{0.3\linewidth}
\begin{tabular}{ccc}
\toprule
$M_s$ & $\alpha$ & State\\
\midrule
  1 &  2  &$\ket{3,4}$ \\
  1 &  2  &$\ket{1,5}$ \\
  1 &  2  &$\ket{3,5}$ \\
\bottomrule
\end{tabular}
\caption{Tabulated two-particle basis for $M_s=1$}
\end{minipage}
\end{table}
%
Every odd number is now a spin up state, and even number are spin down. The energy is thereby degenerate in spin.
%
\begin{equation}
\ket{a} = \left\{ \begin{array}{lcl} a = 2n & \mbox{if} & m_s=-0.5\\ 
 a = 2n+1 & \mbox{if} & m_s = 0.5
\end{array}\right.
\end{equation}
%
In section \ref{parabolic quantum dot} we discussed the shell structure of the parabolic quantum dot. We have the same situasjon now but with a different kind of subshells. We define a subshell with the set of quantum numbers $\{n_x',n_y'\}$, where each subshell have the energy

\begin{equation}
E_{sub} = E_{n_x'} + E_{n_y'}
  \label{subshell energy}
\end{equation}
%
And $E_{n_x'}$ are the eigenvalues found by diagonalization Eq.~\ref{eigenvalues found by diag} and Fig.~\ref{fig:Vgauss2D}, Fig. ~\ref{fig:Vgauss3D}. While $E_{n_y'} = (0.5 + n_y')$ are the one-dimensional solution to the harmonic oscillator potential. Example for the potential \ref{Vgauss}, $V=20,\sigma=1$, with $N=500$.

\begin{table}[H]
\centering
\begin{minipage}[b]{0.4\linewidth}\centering 
\begin{tabular}{cc}
\toprule
${n_x'}$ & $E_{n_x'}$\\
\midrule
  0 &  0  5.68161 \\
  1 &  0  5.68162 \\
  2 &  0  7.42767 \\
\bottomrule
\end{tabular}
\caption{Energyeigenvalues for $n_x'$}
\end{minipage}
\hspace{1cm}
\begin{minipage}[b]{0.4\linewidth}
\begin{tabular}{cc}
\toprule
$n_y'$ & $E_{n_y'}$\\
\midrule
  0 & 0.5   \\
  1 &  1   \\
  2 &  2   \\
\bottomrule
\end{tabular}
\caption{Energyeigenvalues for $n_y'$}
\end{minipage}
\end{table}
The first 6 spinorbitals are tabulated as

\begin{align*}
  \ket{0} & = \{n_x'=0, n_y'=0, m_s = \downarrow\} & E =  5.68161 \\ 
  \ket{1} & = \{n_x'=0, n_y'=0, m_s = \uparrow \}  & E =  5.68161 \\
  \ket{2} & = \{n_x'=1, n_y'=0, m_s = \downarrow\} & E = 5.68161 \\
  \ket{3} & = \{n_x'=1, n_y'=0, m_s = \uparrow \}  & E =  5.68161 \\
  \ket{4} & = \{n_x'=1, n_y'=1, m_s = \downarrow\} & E = 6.68162\\
  \ket{5} & = \{n_x'=1, n_y'=1, m_s = \uparrow \}  & E = 6.68162 
\end{align*}



\subsection{Validation of the Code}
One of the test we can do is to drop the barrier in the $x$-direction, i.e. $L_x = 0$ in Eq.~(\ref{doublepotential abs}). Then we get a harmonic oscillator potential. And we want to do a diagonalization for the 2-shell case. We get 6 single-particle states with eigenvalues in the cartesian basis:

\begin{align*}
\ket{0} \rightarrow &\{n_x'=0, n_y'=0, m_s' = \downarrow\} &E = 0.99988 \\  
\ket{1} \rightarrow &\{n_x'=0, n_y'=0, m_s' = \uparrow \} &E =  0.99988 \\  
\ket{2} \rightarrow &\{n_x'=1, n_y'=0, m_s' = \downarrow\} &E = 1.99943  \\  
\ket{3} \rightarrow &\{n_x'=1, n_y'=0, m_s' =  \uparrow\} &E =  1.99943  \\  
\ket{4} \rightarrow &\{n_x'=0, n_y'=1, m_s' = \downarrow\} &E = 1.99988 \\  
\ket{5} \rightarrow &\{n_x'=0, n_y'=1, m_s' =  \uparrow\} &E =  1.99988  
\end{align*}

If we start with the case $M_s = -1$
\begin{equation*}
\mathbf{H} = \begin{bmatrix}
\bra{\Phi_{02}}H_0\ket{\Phi_{02}} + \bra{\Phi_{02}}V\ket{\Phi_{02}} & 
\bra{\Phi_{02}}H_0\ket{\Phi_{04}} + \bra{\Phi_{02}}V\ket{\Phi_{04}} &  
\bra{\Phi_{02}}H_0\ket{\Phi_{24}} + \bra{\Phi_{02}}V\ket{\Phi_{24}} \\
\bra{\Phi_{04}}H_0\ket{\Phi_{02}} + \bra{\Phi_{04}}V\ket{\Phi_{02}} & 
\bra{\Phi_{04}}H_0\ket{\Phi_{04}} + \bra{\Phi_{04}}V\ket{\Phi_{04}} &  
\bra{\Phi_{04}}H_0\ket{\Phi_{24}} + \bra{\Phi_{04}}V\ket{\Phi_{24}} \\
\bra{\Phi_{24}}H_0\ket{\Phi_{02}} + \bra{\Phi_{24}}V\ket{\Phi_{02}} & 
\bra{\Phi_{24}}H_0\ket{\Phi_{04}} + \bra{\Phi_{24}}V\ket{\Phi_{04}} &  
\bra{\Phi_{24}}H_0\ket{\Phi_{24}} + \bra{\Phi_{34}}V\ket{\Phi_{24}} \\
\end{bmatrix}
  \label{validationDiagonalizationoftheHarmonicOscillator}
\end{equation*}
%
Where the Slater determinant is defined as
\begin{equation}
\ket{\Phi_{ab}} \equiv \frac{1}{\sqrt{2}}\left(\ket{ab}-\ket{ba}\right)
  \label{slater}
\end{equation}
%
Then we get
\begin{align*}
\mathbf{H} & \approx \begin{bmatrix}
2 + \bra{02}v\ket{02}_{\text{AS}} & 
    \bra{02}v\ket{04}_{\text{AS}} &  
    \bra{02}v\ket{24}_{\text{AS}} \\
    \bra{04}v\ket{02}_{\text{AS}} & 
4 + \bra{04}v\ket{04}_{\text{AS}} &  
    \bra{04}v\ket{24}_{\text{AS}} \\
    \bra{24}v\ket{02}_{\text{AS}} & 
    \bra{24}v\ket{04}_{\text{AS}} &  
4 + \bra{24}v\ket{24}_{\text{AS}} \\
\end{bmatrix}\\
&= \begin{bmatrix}
    3.6260   &      0    &     0 \\
         0  &  3.6264     &    0 \\
         0   &      0   & 4.6260 \\
\end{bmatrix}\\
  \label{validationDiagonalizationoftheHarmonicOscillator}
\end{align*}
%
Which is what we get when we diagonalize for $M_s=0$ in polar basis. The case $M_s=1$ is excatly the same, and for 
the $Ms = 0$ we get the eigvalues 

\begin{align*}
E_0 &= 3.1523 \qquad E_1 = 3.6260 \qquad E_2 = 3.6264 \\
E_3 &= 4.2533 \qquad E_4 = 4.2533 \qquad E_5 = 4.6267 \\
E_6 &= 4.8617 \qquad E_7 = 4.8617 \qquad E_8 = 5.1976 
\end{align*}

Which is the same as for the polar basis. 

\chapter{Results}
\begin{quote}
{\it "I've very often made mistakes in my physics by thinking the theory isn't as good as it really is, thinking that there are lots of complications that are going to spoil it - an attitude that anything can happen, in spite of what you're pretty sure should happen." (From ``Surely you're joking Mr. Feynman'') }
\\ 
\mbox{R. P. Feynman}
\end{quote}

\section*{Optimazation results}
We have done the following runs for the optimized code and compared to the code from \cite{lohne}

\begin{table}[H]
\centering 
\begin{tabular}{cccc}
\toprule
Run case & Regular code time [s] & Optimized code time [s] & Improvement [\%]\\
\midrule 
$nh=2$,\, $R = 10$ & 1897 & 4 & 4,75E+4 \\
$nh=6$,\, $R = 7$  & 1595 & 3 & 5,32E+4 \\
$nh=6$,\, $R = 8$  & 4331 & 6 & 7,22E+4\\
$nh=6$,\, $R = 9$  & 9697 & 11 & 8,82E+4 \\
$nh=6$,\, $R = 10$  & 21356  & 20 & 10,68E+4  \\
$nh=12$,\, $R = 7$  & 11775 & 23 & 5,12E+4 \\
$nh=12$,\, $R = 8$  & 30834 & 52 & 5,93E+4 \\
\bottomrule
\end{tabular}
\caption{Here the $R$ is defined as the major oscillator shells, see section \ref{parabolic quantum dot}. We have turned on the OpenMP in the optimized code and compared with the regular code from \cite{lohne}. But most of the improvement was due to the fact that we changed from Blitz++ to pointer matrices. We have on average improved the code by a 500 times speed up! And it improves with number of particles.}
\end{table}

\clearpage
\subsection*{Benchmark results for $L_x = 0$ and $\omega = 1$}
We want to test the code for the parabolic harmonic oscillator potential case and see if we get the same results as in \cite{lohne}
\begin{table}[H]
\centering 
\begin{tabular}{cccc|cccc}
\toprule
\multicolumn{4}{c|}{$nh=2$} & \multicolumn{4}{c}{$nh=6$} \\
$R_m$ & $R$ & $HF$ & $CCSD$ & $R_m$ & $R$ & $HF$ & $CCSD$\\
\midrule 
2 & 3	  & 3.253314	  &  3.152329 & 2 & 3  &  22.219813 &  22.219813  \\
4 & 10	  & 3.162691	  &  3.025231 & 4 & 10 &  20.766919 &  20.421321  \\
6 & 21	  & 3.161921	  &  3.013626 & 6 & 21 &  20.720257 &  20.260893 \\
8 & 36	  & 3.161909	  &  3.009236 & 8 & 36 &  20.719248 &  20.221750  \\
10 &55	  & 3.161908      &  3.006938 & 10& 55 &  20.719217 &  20.204344  \\
\bottomrule
\end{tabular}
\caption{Ground-state energies for holes states $nh=2$ and $nh=6$ in a circular quantum dot using a standard Coulomb interaction for $L_x=0$. The interaction elements are obtained from a cartesian basis transformation, section \ref{sec:cartesian basis transform}. HF is the Hartree- Fock energy and CCSD is the coupled-cluster energy Eq.~(\ref{Vabs}). $R_m$ is the total number of subshells while $R_m$ stands for the number of major oscillator shells Fig. \ref{fig:shellnumber}}
\label{table:benchmark}
\end{table}

\begin{table}[H]
\centering 
\begin{tabular}{cccc|cccc}
\toprule
\multicolumn{4}{c|}{$nh=2$} & \multicolumn{4}{c}{$nh=6$} \\
$R_m$ & $R$ & $HF$ & $CCSD$ & $R_m$ & $R$ & $HF$ & $CCSD$\\
\midrule 
2 & 3	  & 3.253314	  &  3.152329 & 2 & 3  &  22.212981 &  22.219813  \\
4 & 10	  & 3.162691	  &  3.025232 & 4 & 10 &  20.766419 &  20.421325  \\
6 & 21	  & 3.161921	  &  3.013627 & 6 & 21 &  20.720157 &  20.260893  \\
8 & 36	  & 3.161909	  &  3.009237 & 8 & 36 &  20.719248 &  20.221750  \\
10 &55	  & 3.161909      &  3.006938 & 10& 55 &  20.719217 &  20.204345  \\
\bottomrule
\end{tabular}
\caption{Ground-state energies for holes states $nh=2$ and $nh=6$ in a double quantum dot. These results are taken from \cite{lohne} for us to compare with. As we can see they are almost identical with Table~\ref{table:benchmark}}
\label{table:benchmark2}
\end{table}
%
We shall define a new shell structure, each subshell consist of quantum numbers $n_x'$ and $n_y'$ and they can be occupied with at most 2 electrons. In the special case where we dont have any perturbation in the middle $L_x=0$ or $V_0=0$, some subshells are part of the major oscillator shells which we have discussed in section \ref{parabolic quantum dot}. The relations are
%
\begin{equation}
R = \frac{R_m(R_m+1)}{2}.
\label{RmvsR}
\end{equation}
%
We define the total number of single-particle states in our basis as
\begin{equation}
  nh+np = 2R,
  \label{spenergyR}
\end{equation}
%
$nh$ are number of hole states and $np$ are number of particle states in our model.
\begin{table}[H]
\small
\centering 
\begin{tabular}{cccrc|cccrc}
\hline
$p$ & $n_x'$ & $n_y'$ & $m_s'$ & $E$ & $p$ & $n_x'$ & $n_y'$ & $m_s'$ & $E$\\
\hline
0  & 0 & 0 & -1 & 1 & 55 & 0 & 6 &  1 & 7   \\
1  & 0 & 0 &  1 & 1 & 56 & 7 & 0 & -1 & 8   \\
2  & 1 & 0 & -1 & 2 & 57 & 7 & 0 &  1 & 8   \\
3  & 1 & 0 &  1 & 2 & 58 & 6 & 1 & -1 & 8   \\
4  & 0 & 1 & -1 & 2 & 59 & 6 & 1 &  1 & 8   \\
5  & 0 & 1 &  1 & 2 & 60 & 5 & 2 & -1 & 8   \\ 
6  & 2 & 0 & -1 & 3 & 61 & 5 & 2 &  1 & 8   \\
7  & 2 & 0 &  1 & 3 & 62 & 4 & 3 & -1 & 8   \\
8  & 1 & 1 & -1 & 3 & 63 & 4 & 3 &  1 & 8   \\
9  & 1 & 1 &  1 & 3 & 64 & 3 & 4 & -1 & 8   \\
10 & 0 & 2 & -1 & 3 & 65 & 3 & 4 &  1 & 8   \\
11 & 0 & 2 &  1 & 3 & 66 & 2 & 5 & -1 & 8   \\
12 & 3 & 0 & -1 & 4 & 67 & 2 & 5 &  1 & 8   \\
13 & 3 & 0 &  1 & 4 & 68 & 1 & 6 & -1 & 8   \\ 
14 & 2 & 1 & -1 & 4 & 69 & 1 & 6 &  1 & 8   \\
15 & 2 & 1 &  1 & 4 & 70 & 0 & 7 & -1 & 8   \\
16 & 1 & 2 & -1 & 4 & 71 & 0 & 7 &  1 & 8   \\
17 & 1 & 2 &  1 & 4 & 72 & 8 & 0 & -1 & 9   \\
18 & 0 & 3 & -1 & 4 & 73 & 8 & 0 &  1 & 9   \\
19 & 0 & 3 &  1 & 4 & 74 & 7 & 1 & -1 & 9   \\
20 & 4 & 0 & -1 & 5 & 75 & 7 & 1 &  1 & 9   \\
21 & 4 & 0 &  1 & 5 & 76 & 6 & 2 & -1 & 9   \\
22 & 3 & 1 & -1 & 5 & 77 & 6 & 2 &  1 & 9   \\
23 & 3 & 1 &  1 & 5 & 78 & 5 & 3 & -1 & 9   \\
24 & 2 & 2 & -1 & 5 & 79 & 5 & 3 &  1 & 9   \\
25 & 2 & 2 &  1 & 5 & 80 & 4 & 4 & -1 & 9   \\ 
26 & 1 & 3 & -1 & 5 & 81 & 4 & 4 &  1 & 9   \\
27 & 1 & 3 &  1 & 5 & 82 & 3 & 5 & -1 &	9   \\
28 & 0 & 4 & -1 & 5 & 83 & 3 & 5 &  1 & 9   \\
29 & 0 & 4 &  1 & 5 & 84 & 2 & 6 & -1 & 9   \\
30 & 5 & 0 & -1 & 6 & 85 & 2 & 6 &  1 & 9   \\
31 & 5 & 0 &  1 & 6 & 86 & 1 & 7 & -1 & 9   \\
32 & 4 & 1 & -1 & 6 & 87 & 1 & 7 &  1 & 9   \\
33 & 4 & 1 &  1 & 6 & 88 & 0 & 8 & -1 & 9   \\ 
34 & 3 & 2 & -1 & 6 & 89 & 0 & 8 &  1 & 9   \\
35 & 3 & 2 &  1 & 6 & 90 & 9 & 0 & -1 & 10  \\
36 & 2 & 3 & -1 & 6 & 91 & 9 & 0 &  1 & 10  \\
37 & 2 & 3 &  1 & 6 & 92 & 8 & 1 & -1 & 10  \\
38 & 1 & 4 & -1 & 6 & 93 & 8 & 1 &  1 & 10  \\
39 & 1 & 4 &  1 & 6 & 94 & 7 & 2 & -1 & 10  \\
40 & 0 & 5 & -1 & 6 & 95 & 7 & 2 &  1 & 10  \\
41 & 0 & 5 &  1 & 6 & 96 & 6 & 3 & -1 & 10  \\
42 & 6 & 0 & -1 & 7 & 97 & 6 & 3 &  1 & 10  \\
43 & 6 & 0 &  1 & 7 & 98 & 5 & 4 & -1 & 10  \\
44 & 5 & 1 & -1 & 7 & 99 & 5 & 4 &  1 & 10  \\
45 & 5 & 1 &  1 & 7 &100 & 4 & 5 & -1 & 10  \\
46 & 4 & 2 & -1 & 7 &101 & 4 & 5 &  1 & 10  \\
47 & 4 & 2 &  1 & 7 &102 & 3 & 6 & -1 & 10  \\
48 & 3 & 3 & -1 & 7 &103 & 3 & 6 &  1 & 10  \\
49 & 3 & 3 &  1 & 7 &104 & 2 & 7 & -1 & 10  \\
50 & 2 & 4 & -1 & 7 &105 & 2 & 7 &  1 & 10  \\
51 & 2 & 4 &  1 & 7 &106 & 1 & 8 & -1 & 10  \\
52 & 1 & 5 & -1 & 7 &107 & 1 & 8 &  1 & 10  \\
53 & 1 & 5 &  1 & 7 &108 & 0 & 9 & -1 & 10  \\
54 & 0 & 6 & -1 & 7 &109 & 0 & 9 &  1 & 10  \\
\hline
\end{tabular}
\caption{$L_x=0$,\, $\omega=1.0$, here we have defined $m_s' = 2m_s$, and $m_s=\pm1/2$. $E$ are the energy eigenvalues for the singel-particles.}
\label{table:L0}
\end{table}

\begin{figure}[h!]
\centering
$\begin{array}{cc}
\includegraphics[width=3in]{nh=2_w=1_Lx=0} & \includegraphics[width=3in]{nh=6_w=1_Lx=0}\end{array}$
\caption{CCSD and HF energy as function of major oscillator shells $R_m$ for the 2- and 6-electron system with $L_x=0$,  $\omega=1$. The energy is measured in effective Hartrees ($E_H$)}
\label{fig:L=0 results}
\end{figure}

\begin{figure}[h!]
\centering
$\begin{array}{cc}
\includegraphics[width=3in]{1D_Lx=0_0} & \includegraphics[width=3in]{1D_Lx=0_1} \\ \includegraphics[width=3in]{1D_Lx=0_2} & \includegraphics[width=3in]{1D_Lx=0_3}
\end{array}$
\caption{The probability density of the four lowest eigenvalues for $L_x = 0$, $\omega = 1$ and $V(x,0)$.}
\label{fig:wf_L=0_1D}
\end{figure}

\begin{figure}[h!]
\centering
$\begin{array}{cc}
\includegraphics[width=3in]{2D_Lx=0_0} & \includegraphics[width=3in]{2D_Lx=0_1} \\ \includegraphics[width=3in]{2D_Lx=0_2} & \includegraphics[width=3in]{2D_Lx=0_3}
\end{array}$
\caption{The probability density of the four lowest eigenvalues for $L_x = 0$, $\omega = 1$ and $V(x,y)$. The energy $E$ is given in units of $[\hbar\omega]$}
\label{wf_L=0_2D}
\end{figure}

\clearpage
\section*{Standard Interaction for $L_x=2.5 \,,\, \omega=0.8$}

\begin{table}[h!]
\small
\centering 
\begin{tabular}{cccrc}
\hline
$p$ & $n_x'$ & $n_y'$ & $m_s'$ & $E$\\
\hline
0  & 0 & 0 & -1 & 0.9977\\
1  & 0 & 0 &  1 & 0.9977\\
2  & 1 & 0 & -1 & 1.0025\\
3  & 1 & 0 &  1 & 1.0025\\
4  & 2 & 0 & -1 & 1.9662\\
5  & 2 & 0 &  1 & 1.9662\\ 
6  & 0 & 1 & -1 & 1.9970\\
7  & 0 & 1 &  1 & 1.9970\\
8  & 1 & 1 & -1 & 2.0025\\
9  & 1 & 1 &  1 & 2.0025\\
10 & 3 & 0 & -1 & 2.0245\\
11 & 3 & 0 &  1 & 2.0245\\
12 & 4 & 0 & -1 & 2.8722\\
13 & 4 & 0 &  1 & 2.8722\\ 
14 & 2 & 1 & -1 & 2.9662\\
15 & 2 & 1 &  1 & 2.9662\\
16 & 0 & 2 & -1 & 2.9970\\
17 & 0 & 2 &  1 & 2.9970\\
18 & 1 & 2 & -1 & 3.0025\\
19 & 1 & 2 &  1 & 3.0025\\
20 & 3 & 1 & -1 & 3.0245\\
21 & 3 & 1 &  1 & 3.0245\\
22 & 5 & 0 & -1 & 3.8722\\
23 & 5 & 0 &  1 & 3.8722\\
\hline
\end{tabular}
\caption{single-particle states for $L_x=2.5$, $E$ is the energy eigenvalues and $m_s' = 2m_s$, where $m_s = \pm 1/2$ }
\label{table:L2.5}
\end{table}

\begin{figure}[H]
\centering
\input{w=0_8_Lx=2_5_R=12_map}
\caption{Shellstructure for $L_x=2.5$}
\label{fig:L2.5map}
\end{figure}

% \begin{figure}[H]
% \centering
% \scalebox{0.7}{\includegraphics{nh=2_w=0_8_Lx=2_5_R=1-10.eps}}
% \caption{Hartree Fock and Coupled Cluster energy for $w=0.8, \, L_x=2.5$}
% \end{figure}

\begin{table}[h!]
\centering 
\begin{tabular}{cccl|cccl}
\toprule
\multicolumn{4}{c|}{$nh=2$} & \multicolumn{4}{c}{$nh=4$} \\
$R$ & $nh+np$ & $HF$ & $CCSD$ & $R$ & $nh+np$ & $HF$ & $CCSD$ \\
\midrule 
1 & 2	  & 2.8476	  &  2.8466    &1 & 2	  & -		  &  -   	\\
2 & 4	  & 2.8466	  &  2.7185    &2 & 4	  & 9.5711	  &  9.5719     \\
3 & 6	  & 2.8337	  &  2.6566    &3 & 6	  & 9.5125	  &  9.6027  	\\
4 & 8	  & 2.8337	  &  2.6608    &4 & 8	  & 9.5125	  &  9.9226   	\\
5 & 10	  & 2.8337	  &  2.6407    &5 & 10	  & 9.5125	  &  not conv.  \\
6 & 12	  & 2.8337	  &  2.6376    &6 & 12	  & 9.3498	  &  not conv.  \\  
7 & 14	  & 2.8337	  &  2.6373    &7 & 14	  & 9.3498	  &  not conv.  \\
8 & 16	  & 2.7477	  &  2.5870    &8 & 16	  & 9.0801	  &  not conv.  \\
9 & 18	  & 2.6978        &  2.5199    &9 & 18	  & 8.8066        &  8.3057  	\\
10 &20	  & 2.6978        &  not conv. &10 &20	  & 8.6306        &  8.1228  	\\
11 &22	  & 2.6978        &  not conv. &11 &22	  & 8.6306        &  8.1119  	\\
12 &24	  & 2.6978        &  not conv. &12 &24	  & 8.6306        &  8.1095  	\\
\bottomrule
\end{tabular}
\caption{Ground-state energies for hole states $nh=2$ in a double quantum dot using a standard Coulomb interaction for $L_x=2.5$, $\omega=0.8$, where $nh+np$ are the number of holes plus particles in our system.}
\label{table:w=0.8_L=2.5}
\end{table}


\begin{figure}[h!]
\centering
$\begin{array}{cc}
\includegraphics[width=3in]{nh=2_w=0_8_Lx=2_5} & \includegraphics[width=3in]{nh=4_w=0_8_Lx=2_5}\end{array}$
\caption{CCSD and HF energy as function of subshells $R$ for the 2- and 4-electron system with $L_x=2.5$,  $\omega=0.8$. The energy is measured in effective Hartrees ($E_H$)}
\label{fig:L=2_5 w=0_8 results}
\end{figure}




\clearpage
\section*{Standard Interaction for $L_x=2.5 \,,\, \omega=1$}


\begin{table}[h!]
\centering 
\begin{tabular}{cccl|cccl}
\toprule
\multicolumn{4}{c|}{$nh=2$} & \multicolumn{4}{c}{$nh=4$} \\
$R$ & $nh+np$ & $HF$ & $CCSD$ & $R$ & $nh+np$ & $HF$ & $CCSD$ \\
\midrule 
1 & 2	  & 3.3932    &  3.3931   & 1 & 2	  &	 -	  &  -  	\\
2 & 4	  & 3.3931    &  3.2508   & 2 & 4	  & 11.1229	  &  11.1229  	\\ 
3 & 6	  & 3.3808    &  3.1865*  & 3 & 6	  & 11.0769	  &  not conv.  \\
4 & 8	  & 3.3808    &  3.1893*  & 4 & 8	  & 11.0769	  &  not conv.  \\
5 & 10	  & 3.3808    &  3.1634*  & 5 & 10	  & 11.0769	  &  not conv.  \\
6 & 12	  & 3.3808    &  3.1633*  & 6 & 12	  & 10.9337	  &  not conv.  \\  
7 & 14	  & 3.3808    &  3.1637*  & 7 & 14	  & 10.9337	  &  not conv.  \\
8 & 16	  & 3.2918    &  3.1119*  & 8 & 16	  & 10.6236	  &  not conv.  \\
9 & 18	  & 3.2400    &  3.0716*  & 9 & 18	  & 10.3394       &  not conv.  \\
10 &20	  & 3.2400    &  3.0715*  & 10 &20	  & 10.1495       &  not conv.  \\
11 &22	  & 3.2400    &  not conv.& 11 &22	  & 10.1495       &  not conv.  \\
12 &24	  & 3.2400    &  not conv.& 12 &24	  & 10.1495       &  not conv.  \\
\bottomrule
\end{tabular}
\caption{Ground-state energies for $nh=2$ in a double quantum dot using a standard Coulomb interaction for $L_x=2.5$, $\omega=1$. I have marked the CCSD energies with (*) where I have extrapolated due to poor convergence.}
\label{table:w=1_L=2.5}
\end{table}

\begin{figure}[h!]
\centering
$\begin{array}{cc}
\includegraphics[width=3.2in]{nh=2_w=1_Lx=2_5} & \includegraphics[width=3.2in]{nh=4_w=1_Lx=2_5}\end{array}$
\caption{CCSD and HF energy as function of subshells $R$ for the 2- and 4-electron system with $L_x=2.5$,  $\omega=1$. The energy is measured in effective Hartrees ($E_H$)}
\label{fig:L=2_5 w=1 results}
\end{figure}


\begin{figure}[h!]
\centering
$\begin{array}{cc}
\includegraphics[width=2.7in]{1D_Lx=2_5_0} & \includegraphics[width=2.7in]{1D_Lx=2_5_1} \\ \includegraphics[width=2.7in]{1D_Lx=2_5_2} & \includegraphics[width=2.7in]{1D_Lx=2_5_3}
\end{array}$
\caption{The probability density of the six lowest eigenstates for $L_x = 2.5$, $\omega = 1$ and $V(x,0)$.}
\label{fig:wf_Lx=2_5_1D}
\end{figure}


\begin{figure}[h!]
\centering
$\begin{array}{cc}
\includegraphics[width=2.5in]{2D_Lx=2_5_0} & \includegraphics[width=2.5in]{2D_Lx=2_5_1} \\ \includegraphics[width=2.5in]{2D_Lx=2_5_2} & \includegraphics[width=2.5in]{2D_Lx=2_5_3}
\end{array}$
\caption{The probability density of the six lowest eigenstates for $L_x = 2.5$, $\omega = 1$ and $V(x,0)$.}
\label{fig:wf_Lx=2_5_2D}
\end{figure}

\newpage
\section*{Standard Interaction for $V_0=8 \,,\, \omega=1$}

\begin{table}[h!]
\small
\centering 
\begin{tabular}{cccrc|cccrc}
\hline
$p$ & $n_x'$ & $n_y'$ & $m_s'$ & $E$ & $p$ & $n_x'$ & $n_y'$ & $m_s'$ & $E$\\
\hline
0  & 0 & 0 & -1 & 4.5510 & 20 & 0 & 3 & -1 & 7.5506\\ 
1  & 0 & 0 &  1 & 4.5510 & 21 & 0 & 3 &  1 & 7.5506\\
2  & 1 & 0 & -1 & 4.5536 & 22 & 1 & 3 & -1 & 7.5532\\
3  & 1 & 0 &  1 & 4.5536 & 23 & 1 & 3 &  1 & 7.5532\\
4  & 0 & 1 & -1 & 5.5510 & 24 & 4 & 0 & -1 & 7.8936\\
5  & 0 & 1 &  1 & 5.5510 & 25 & 4 & 0 &  1 & 7.8936\\
6  & 1 & 1 & -1 & 5.5536 & 26 & 5 & 0 & -1 & 8.1525\\
7  & 1 & 1 &  1 & 5.5536 & 27 & 5 & 0 &  1 & 8.1525\\
8  & 2 & 0 & -1 & 6.3731 & 28 & 2 & 2 & -1 & 8.3715\\
9  & 2 & 0 &  1 & 6.3731 & 29 & 2 & 2 &  1 & 8.3715\\
10 & 3 & 0 & -1 & 6.4135 & 30 & 3 & 2 & -1 & 8.3715\\
11 & 3 & 0 &  1 & 6.4135 & 31 & 3 & 2 &  1 & 8.3715\\
12 & 0 & 2 & -1 & 6.5510 & 32 & 0 & 4 & -1 & 8.4118\\
13 & 0 & 2 &  1 & 6.5510 & 33 & 0 & 4 &  1 & 8.4118\\
14 & 1 & 2 & -1 & 6.5536 & 34 & 1 & 4 & -1 & 8.5506\\
15 & 1 & 2 &  1 & 6.5536 & 35 & 1 & 4 &  1 & 8.5506\\
16 & 2 & 1 & -1 & 7.3731 & 36 & 4 & 1 & -1 & 8.8936\\
17 & 2 & 1 &  1 & 7.3731 & 37 & 4 & 1 &  1 & 8.8936\\
18 & 3 & 1 & -1 & 7.4135 & 38 & 5 & 1 & -1 & 9.1525\\
19 & 3 & 1 &  1 & 7.4135 & 39 & 5 & 1 &  1 & 9.1525\\
\hline
\end{tabular}
\caption{single-particle states for $V_0=8$, $\omega=1$. $E$ is the eigenvalues, and $m_s'$ is defined as $m_s' = 2m_s$, where $m_s=\pm1/2$.}
\label{table:V8}
\end{table}

\begin{figure}[h!]
\centering
\input{V=8_map}
\caption{Shellstructure for $V_0=8$, the levels are drawn with respect to Table~\ref{table:V8}}
\label{fig:V8map}
\end{figure}

\begin{table}[h!]
\centering 
\begin{tabular}{cccl|cccl}
\toprule
\multicolumn{4}{c|}{$nh=4$} & \multicolumn{4}{c}{$nh=8$} \\
$R$ & $nh+np$ & $HF$ & $CCSD$ & $R$ & $nh+np$ & $HF$ & $CCSD$ \\
\midrule 
2 & 4	  & 24.2844	  &  24.2844   & 2 & 4	  & -	  &  -  	      \\
3 & 6	  & 24.2844	  &  24.2039   & 3 & 6	  & -	  &  -  	      \\
4 & 8	  & 24.2844	  &  24.0921   & 4 & 8	  & 63.1758	  &  63.1758  \\
5 & 10	  & 24.2662	  &  24.0452  & 5 & 10	  & 63.1757	  &  63.0789  \\
6 & 12	  & 24.2661	  &  23.9688  & 6 & 12	  & 63.0706	  &  62.8668  \\  
7 & 14	  & 23.9971	  &  23.6772  & 7 & 14	  & 62.6058	  &  62.4029  \\
8 & 16	  & 23.8460	  &  23.4980  & 8 & 16	  & 62.3627	  &  62.1220  \\
9 & 18	  & 23.8460       &  23.4950  & 9 & 18	  & 62.3066       &  61.9723  \\
11 &22	  & 23.8460       &  23.4888  & 11 &22	  & 61.7437       &  61.8037  \\
12 &24	  & 23.8460       &  23.4857  & 12 &24	  & 61.3687       &  61.2238  \\
13 &26	  & 23.6929       &  23.3370  & 13 &26	  & 60.9041       &  60.8306  \\
14 &28	  & 23.6074       &  23.2520  & 14 &28	  & 60.6206       &  59.9635  \\
15 &30	  & 23.6074       &  23.2510  & 15 &30	  & 60.6206       &  59.9283  \\
16 &32	  & 23.6074       &  23.2497  & 16 &32	  & 60.6200       &  59.8997  \\
17 &34	  & 23.6074       &  23.2481  & 17 &34	  & 60.6012       &  59.8606  \\
18 &36	  & 23.6067       &  23.2463  & 18 &36	  & 60.5834       &  59.8198  \\
\bottomrule
\end{tabular}
\caption{Ground-state energies for 4 and 8 electrons in a double quantum dot using a standard Coulomb interaction for $V_0=8$, $\omega=1$.}
\label{table:w=1_V=8}
\end{table}

\begin{figure}[h!]
\centering
$\begin{array}{cc}
\includegraphics[width=3in]{nh=4_w=1_V=8} & \includegraphics[width=3in]{nh=8_w=1_V=8}\end{array}$
\caption{CCSD and HF energy as function of subshells $R$ for the 4- and 8-electron system with $V_0=8$,  $\omega=1$. The energy is measured in effective Hartrees ($E_H$)}
\label{fig:V=8 w=1 results}
\end{figure}

\begin{figure}[h!]
\centering
$\begin{array}{cc}
\includegraphics[width=2.5in]{1D_V=8_0} & \includegraphics[width=2.5in]{1D_V=8_1} \\ \includegraphics[width=2.5in]{1D_V=8_2} & \includegraphics[width=2.5in]{1D_V=8_3}
\end{array}$
\caption{The probability density of the six lowest eigenstates for $V_0 = 8$, $\omega = 1$ and $V_G(\bar{x },0)$.}
\label{fig:wf_V=8}
\end{figure}

\begin{figure}[h!]
\centering
$\begin{array}{cc}
\includegraphics[width=2.5in]{2D_V=8_0} & \includegraphics[width=2.5in]{2D_V=8_1} \\ \includegraphics[width=2.5in]{2D_V=8_2} & \includegraphics[width=2.5in]{2D_V=8_3}
\end{array}$
\caption{The probability density of the six lowest eigenstates for $V_0 = 8$, $\omega = 1$ and $V_G(\bar{x},\bar{y})$.}
\end{figure}


\clearpage
\section*{Standard Interaction for $V_0=10.66 \,,\, \omega=1$}

\begin{table}[h!]
\small
\centering 
\begin{tabular}{cccrc|cccrc}
\hline
$p$ & $n_x'$ & $n_y'$ & $m_s'$ & $E$ & $p$ & $n_x'$ & $n_y'$ & $m_s'$ & $E$\\
\hline
0  & 0 & 0 & -1 & 4.9123 & 20 & 0 & 3 &  1 & 7.9123   \\
1  & 0 & 0 &  1 & 4.9123 & 21 & 0 & 3 & -1 & 7.9123   \\
2  & 1 & 0 & -1 & 4.9127 & 22 & 1 & 3 &  1 & 7.9126   \\
3  & 1 & 0 &  1 & 4.9127 & 23 & 1 & 3 & -1 & 7.9126   \\
4  & 0 & 1 & -1 & 5.9123 & 24 & 4 & 0 &  1 & 8.7184   \\
5  & 0 & 1 &  1 & 5.9123 & 25 & 4 & 0 & -1 & 8.7184   \\ 
6  & 1 & 1 & -1 & 5.9127 & 26 & 5 & 0 &  1 & 8.7628   \\
7  & 1 & 1 &  1 & 5.9127 & 27 & 5 & 0 & -1 & 8.7628   \\
8  & 2 & 0 & -1 & 6.9018 & 28 & 2 & 2 &  1 & 8.9018   \\
9  & 2 & 0 &  1 & 6.9018 & 29 & 2 & 2 & -1 & 8.9018   \\
10 & 3 & 0 & -1 & 6.9071 & 30 & 3 & 2 &  1 & 8.9071   \\
11 & 3 & 0 &  1 & 6.9071 & 31 & 3 & 2 & -1 & 8.9071   \\
12 & 0 & 2 & -1 & 6.9123 & 32 & 0 & 4 &  1 & 8.9123   \\
13 & 0 & 2 &  1 & 6.9123 & 33 & 0 & 4 & -1 & 8.9123   \\ 
14 & 1 & 2 & -1 & 6.9127 & 34 & 1 & 4 &  1 & 8.9127   \\
15 & 1 & 2 &  1 & 6.9127 & 35 & 1 & 4 & -1 & 8.9127   \\
16 & 2 & 1 & -1 & 7.9018 & 36 & 4 & 1 &  1 & 9.7184   \\
17 & 2 & 1 &  1 & 7.9018 & 37 & 4 & 1 & -1 & 9.7184   \\
18 & 3 & 1 & -1 & 7.9071 & 38 & 5 & 1 &  1 & 9.7628   \\
19 & 3 & 1 &  1 & 7.9071 & 39 & 5 & 1 & -1 & 9.7628   \\
\hline
\end{tabular}
\caption{Tabulated single-particle basis for $V_0 = 10.66,\,\omega=1.0$, here we have defined $m_s' = 2m_s$, and $m_s=\pm1/2$.}
\label{table:V10.66}
\end{table}

\begin{figure}[h!]
\centering
\input{V=10_66_map}
\caption{Shellstructure for $\omega=1,\,V_0=10.66$ with respect to the single particle energies in Table \ref{table:V10.66}}
\label{fig:V10.66map}
\end{figure}


\begin{table}[h!]
\centering 
\begin{tabular}{cccl|cccl}
\toprule
\multicolumn{4}{c}{$nh=8$} & \multicolumn{4}{c}{$nh=4$}  \\
$R$ & $nh+np$ & $HF$ & $CCSD$ & $R$ & $nh+np$ & $HF$ & $CCSD$ \\
\midrule 
2 & 4	  & 65.6857 	  &  65.6857   &  2 & 4	  & 25.6188 	  &  25.6188  \\ 
4 & 8	  & 65.5764	  &  65.3979   &  4 & 8	  & 25.6188	  &  25.4430  \\
6 & 12	  & 64.8985	  &  64.6806   &  6 & 12  & 25.5957	  &  25.3295  \\
8 & 16	  & 63.9310	  &  63.4484   &  8 & 16  & 25.1959	  &  24.8764  \\
12 & 24	  & 63.9310	  &  62.6858   &  12 & 24 & 25.1959	  &  24.8653  \\
14 & 28	  & 63.2738	  &  62.6309   &  14 & 28 & 24.9871	  &  24.6589  \\  
16 & 32	  & 63.2729       &  62.5532   &  16 & 32 & 24.9872       &  24.6569  \\
18 & 36	  & 63.2372	  &  62.5532   &  18 & 36 & 24.9858	  &  24.6535  \\
20 & 40	  & 62.8530       &  62.1496   &  20 & 40 & 24.9858       &  24.6511  \\
\bottomrule
\end{tabular}
\caption{Ground-state energies for $nh=4$ and $nh=8$ in a double quantum dot using a standard Coulomb interaction for $V_0=10.66$, $\omega=1$.}
\label{table:w=1_V=10.66}
\end{table}

\begin{figure}[h!]
\centering
$\begin{array}{cc}
\includegraphics[width=3in]{nh=4_w=1_V=10_66} & \includegraphics[width=3in]{nh=8_w=1_V=10_66}\end{array}$
\caption{CCSD and HF energy as function of subshells $R$ for the 4- and 8-electron system with $V_0=10.66$,  $\omega=1$. The energy is measured in effective Hartrees ($E_H$)}
\label{fig:V=10.66 w=1 results}
\end{figure}


\clearpage
\section*{Standard Interaction for $V_0=20 \,,\, \omega=1$}

\begin{table}[h!]
\small
\centering 
\begin{tabular}{cccrc|cccrc}
\hline
$p$ & $n_x'$ & $n_y'$ & $m_s'$ & $E$ & $p$ & $n_x'$ & $n_y'$ & $m_s'$ & $E$\\
\hline
0  & 0 & 0 & -1 & 5.6822 & 16 & 0 & 3 & -1 & 8.6822\\
1  & 0 & 0 &  1 & 5.6822 & 17 & 0 & 3 &  1 & 8.6822\\
2  & 1 & 0 & -1 & 5.6822 & 18 & 1 & 3 & -1 & 8.6822\\
3  & 1 & 0 &  1 & 5.6822 & 19 & 1 & 3 &  1 & 8.6822\\
4  & 0 & 1 & -1 & 6.6822 & 20 & 2 & 1 & -1 & 8.9303\\ 
5  & 0 & 1 &  1 & 6.6822 & 21 & 2 & 1 &  1 & 8.9303\\
6  & 1 & 1 & -1 & 6.6822 & 22 & 3 & 1 & -1 & 8.9303\\
7  & 1 & 1 &  1 & 6.6822 & 23 & 3 & 1 &  1 & 8.9303\\
8  & 0 & 2 & -1 & 7.6822 & 24 & 0 & 4 & -1 & 9.6822\\
9  & 0 & 2 &  1 & 7.6822 & 25 & 0 & 4 &  1 & 9.6822\\
10 & 1 & 2 & -1 & 7.6822 & 26 & 1 & 4 & -1 & 9.6822\\
11 & 1 & 2 &  1 & 7.6816 & 27 & 1 & 4 &  1 & 9.6816\\
12 & 2 & 0 & -1 & 7.9303 & 28 & 2 & 2 & -1 & 9.9303\\
13 & 2 & 0 &  1 & 7.9303 & 29 & 2 & 2 &  1 & 9.9303\\
14 & 3 & 0 & -1 & 7.9303 & 30 & 3 & 2 & -1 & 9.9303\\
15 & 3 & 0 &  1 & 7.9303 & 31 & 3 & 2 &  1 & 9.9303\\
\hline
\end{tabular}
\caption{single-particle states for $V_0=8$, and $E$ is the eigenvalues, $m_s'=2m_s$ where $m_s$ is defined as $m_s=\pm1/2$.}
\label{table:V20}
\end{table}

\begin{figure}[h!]
\centering
\input{V=20_map}
\caption{Shellstructure for $V_0=20$, the levels are drawn with respect to Table~\ref{table:V20}}
\label{fig:V20map}
\end{figure}

\begin{table}[h!]
\centering 
\begin{tabular}{cccl|cccl}
\toprule
\multicolumn{4}{c|}{$nh=4$} & \multicolumn{4}{c}{$nh=8$} \\
$R$ & $nh+np$ & $HF$ & $CCSD$ & $R$ & $nh+np$ & $HF$ & $CCSD$ \\
\midrule 
2  & 4	  & 28.5229	  &  28.5229   & 2  & 4	  & -	  	  &  -  \\
4  & 8	  & 28.5229	  &  28.3691   & 4  & 8	  & 71.2361	  &  71.2361  \\
6  & 12	  & 28.1446	  &  27.9108   & 6  & 12  & 70.5689	  &  70.4711  \\  
8  & 16	  & 28.1171	  &  27.8315   & 8  & 16  & 70.4506	  &  70.2593  \\
10 & 20   & 28.1171	  &  27.8315   & 10 & 20  & 69.5845       &  69.2638  \\
12 & 24	  & 28.1171       &  27.8256   & 12 & 24  & 69.5121       &  69.0910  \\
14 & 28	  & 28.1153       &  27.8217   & 14 & 28  & 69.4645       &  68.9951  \\
16 & 32	  & 28.1153       &  27.8204   & 16 & 32  & 69.4616       &  68.9481  \\
\bottomrule
\end{tabular}
\caption{Ground-state energies for 4 and 8 electrons in a double quantum dot using a standard Coulomb interaction for $V_0=20$, $\omega=1$.}
\label{table:w=1_V=20}
\end{table}

\begin{figure}[h!]
\centering
$\begin{array}{cc}
\includegraphics[width=3.1in]{nh=4_w=1_V=20} & \includegraphics[width=3.1in]{nh=8_w=1_V=20}\end{array}$
\caption{CCSD and HF energy as function of subshells $R$ for the 4-and 8-electron system with $V_0=20$,  $\omega=1$. The energy is measured in effective Hartrees ($E_H$)}
\label{fig:V=20 w=1 results}
\end{figure}

\clearpage
\section*{Potentials}

\begin{figure}[h!]
\centering
$\begin{array}{cc}
\includegraphics[width=5in]{VCpotentials} \\ \includegraphics[width=5in]{VGpotentials}\end{array}$
\caption{The top figure is a plot of the potential Eq.~(\ref{Vabs}) for different $\omega$ values and $L_x=2.5$, while the bottom figure show the potential Eq.~(\ref{Vgauss}) for different $V_0$ values and $\omega=1$.}
\label{fig:VC and VG potentials}
\end{figure}


\clearpage
\section*{Standard Interaction for $V_0=10$, $\omega=1$, $a=1.449$, $\sigma=1$}
We want to look at a special potential where the may have 6 electrons in the ground state, a triple quantum dot. 

\begin{equation}
  V_S = \frac{1}{2}\bar{\omega}\left[\bar{x}^2 + \bar{y}^2 \right] + V_0 \exp\left(-\frac{(\bar{x}-a)^2}{2\sigma}\right)+V_0 \exp\left(-\frac{(\bar{x}+a)^2}{2\sigma}\right)
  \label{Vspecial}
\end{equation}
%
Where we have set: $V_0=10$, $\omega=1$, $a=1.449$, $\sigma=1$.

\begin{figure}[h!]
\centering
\scalebox{0.6}{\includegraphics{tripledot}}
\caption{Plot of the potential Eq.~(\ref{Vspecial}) with $V_G(\bar{x},0)$ and $V_0=10$, $\omega=1$, $a=1.449$, $\sigma=1$.}
\end{figure}

\begin{table}[h!]
\small
\centering 
\begin{tabular}{cccrc|cccrc}
\hline
$p$ & $n_x'$ & $n_y'$ & $m_s'$ & $E$ & $p$ & $n_x'$ & $n_y'$ & $m_s'$ & $E$\\
\hline
0  & 0 & 0 & -1 & 8.8186 & 20 & 0 & 3 &  1 & 10.8567   \\
1  & 0 & 0 &  1 & 8.8186 & 21 & 0 & 3 & -1 & 10.8567   \\
2  & 1 & 0 & -1 & 8.8397 & 22 & 1 & 3 &  1 & 11.1972   \\
3  & 1 & 0 &  1 & 8.8397 & 23 & 1 & 3 & -1 & 11.1972   \\
4  & 0 & 1 & -1 & 8.8504 & 24 & 4 & 0 &  1 & 11.7849   \\
5  & 0 & 1 &  1 & 8.8504 & 25 & 4 & 0 & -1 & 11.7849   \\ 
6  & 1 & 1 & -1 & 9.8186 & 26 & 5 & 0 &  1 & 11.8186   \\
7  & 1 & 1 &  1 & 9.8186 & 27 & 5 & 0 & -1 & 11.8186   \\
8  & 2 & 0 & -1 & 9.8397 & 28 & 2 & 2 &  1 & 11.8397   \\
9  & 2 & 0 &  1 & 9.8397 & 29 & 2 & 2 & -1 & 11.8397   \\
10 & 3 & 0 & -1 & 9.8504 & 30 & 3 & 2 &  1 & 11.8504   \\
11 & 3 & 0 &  1 & 9.8504 & 31 & 3 & 2 & -1 & 11.8504   \\
12 & 0 & 2 & -1 & 10.7849 & 32 & 0 & 4 &  1 & 11.8567   \\
13 & 0 & 2 &  1 & 10.7849 & 33 & 0 & 4 & -1 & 11.8567   \\ 
14 & 1 & 2 & -1 & 10.8186 & 34 & 1 & 4 &  1 & 12.1972   \\
15 & 1 & 2 &  1 & 10.8186 & 35 & 1 & 4 & -1 & 12.1972   \\
16 & 2 & 1 & -1 & 10.8397 & 36 & 4 & 1 &  1 & 12.2962   \\
17 & 2 & 1 &  1 & 10.8397 & 37 & 4 & 1 & -1 & 12.2962   \\
18 & 3 & 1 & -1 & 10.8504 & 38 & 5 & 1 &  1 & 12.7849   \\
19 & 3 & 1 &  1 & 10.8504 & 39 & 5 & 1 & -1 & 12.7849   \\
\hline
\end{tabular}
\caption{Tabulated single-particle basis for $V_0 = 10,\,\omega=1.0,\,a=1.449$, $E$ is the eigenvalues, and $m_s'$ is defined as $m_s=\pm1/2$.}
\label{table:tripledot}
\end{table}

\begin{figure}[h!]
\centering
$\begin{array}{cc}
\includegraphics[width=2.5in]{1D_a=1_4_0} & \includegraphics[width=2.5in]{1D_a=1_4_1} \\ \includegraphics[width=2.5in]{1D_a=1_4_2} & \includegraphics[width=2.5in]{1D_a=1_4_3}
\end{array}$
\caption{The probability density of the six lowest eigenstates for $V_0 = 10$, $\omega = 1$ and $a=1.449$ for the potential $V_S(\bar{x},0)$.}
\label{fig:wf_V=8_1D}
\end{figure}

\begin{figure}[h!]
\centering
$\begin{array}{cc}
\includegraphics[width=2.5in]{2D_a=1_4_0} & \includegraphics[width=2.5in]{2D_a=1_4_1} \\ \includegraphics[width=2.5in]{2D_a=1_4_2} & \includegraphics[width=2.5in]{2D_a=1_4_3}
\end{array}$
\caption{The probability density of the six lowest eigenstates for $V_0 = 10$, $\omega = 1$ and  $a=1.449$ for the potential $V_S(\bar{x},\bar{y})$.}
\label{wf_V=8_2D}
\end{figure}

\begin{table}[h!]
\centering 
\begin{tabular}{ccl|ccl|ccl}
\toprule
\multicolumn{3}{c|}{$nh=2$} & \multicolumn{3}{c|}{$nh=4$} & \multicolumn{3}{c}{$nh=6$}\\
$R$ & $nh+np$ & $CCSD$ & $R$ & $nh+np$ & $CCSD$ & $R$ & $nh+np$ & $CCSD$\\
\midrule 
2  & 4	  & 18.6711  & 2  & 4  & 41.5660  &  2 &  4 & -         \\
4  & 8	  & 18.7212  & 4  & 8  & no conv. &  4 &  8 & 67.2539   \\  
6  & 12	  & 18.6911  & 6  & 12 & no conv. &  6 & 12 & no conv.  \\
8  & 16	  & 18.3337  & 8  & 16 & no conv. &  8 & 16 & no conv.  \\
10 & 20	  & 22.8665  & 10 & 20 & no conv. & 10 & 20 & no conv.  \\
14 & 28	  & 24.5351  & 14 & 28 & no conv. & 10 & 20 & no conv.  \\
16 & 32	  & 18.4136  & 16 & 32 & no conv. & 16 & 32 & no conv.  \\
20 & 40	  & 26.4021  & 20 & 40 & no conv. & 20 & 40 & no conv.  \\  
\bottomrule
\end{tabular}
\caption{Ground-state energies for 2,4,6-electrons in a triple quantum dot using a standard Coulomb interaction for $V_0=10$, $\omega=1$, $a=1.449$. Here $nh+np$ are the total number of particles for a given $R$, Eq.~(\ref{spenergyR}). We have no convergence in most of the cases.}
\label{table:special}
\end{table}


% \begin{figure}[h!]
% \centering
% \scalebox{0.7}{\includegraphics{nh=2_w=1_V=20_R=1-16_std.eps}}
% \caption{$w=1.0, \, V_0=20$}
% \end{figure}

% \begin{figure}[h!]
% \centering
% \scalebox{0.7}{\includegraphics{V=8_vs_V=10_66_vs_V=20.eps}}
% \caption{The potential for $w=1.0, \, V_0=8$ and $V_0 = 10.66$}
% \end{figure}
% 
% \begin{figure}[h!]
% \centering
% \scalebox{0.7}{\includegraphics{nh=2_w=1_V=0-0_5_R=10_std.eps}}
% \caption{$w=1.0$,$V_0 = 0$ to $V_0 = 0.5$}
% \end{figure}
% 
% \begin{figure}[h!]
% \centering
% \scalebox{0.7}{\includegraphics{nh=2_w=1_V=8-20_R=12_std.eps}}
% \caption{$w=1.0$,$V_0 = 8$ to $V_0 = 20$}
% \end{figure}
% 
% \begin{figure}[h!]
% \centering
% \scalebox{0.7}{\includegraphics{nh=2_w=0_25-1_V=15_R=12_std.eps}}
% \caption{$w=1.0$,$V_0 = 0$ to $V_0 = 0.5$}
% \end{figure}
% 
% \begin{figure}[h!]
% \centering
% \scalebox{0.7}{\includegraphics{nh=4_vs_8_w=0_25-1_V=15_R=12_std.eps}}
% \caption{$w=1.0$,$V_0 = 0$ to $V_0 = 0.5$}
% \end{figure}
% 
% \begin{figure}[h!]
% \centering
% \scalebox{0.7}{\includegraphics{nh=4_w=0_25-1_V=15_R=12_std.eps}}
% \caption{$w=1.0$,$V_0 = 0$ to $V_0 = 0.5$}
% \end{figure}

\clearpage
\section{General Analysis and Discussion}
We start by considering the HF and CCSD results for 2-electron and 6-electron quantum dot in Tables \ref{table:benchmark} and \ref{table:benchmark2}, and observe that the differences are of the order 1E-6. This could be caused by the tolerence in the CCSD and HF calculation. I have used a tolerence of 1E-8 in my HF calculation and a tolerence of 1E-7 in the CCSD calculation. While I'm uncertain of what the tolerence of the CCSD and HF calculation are in \cite{lohne}. But overall the code for  the transformations of the interactions elements from a polar to a cartesian basis seem to work. 

Next we consider the double quantum dot potential Eq.~(\ref{Vabs}), with $L_x=2.5$ and $\omega=1$, see Table \ref{table:w=1_L=2.5}. For the 2-electron system we see that the CCSD calculation have a poor convergence and does not converge at all for $R=11$ and $R=12$. If we look at the possibilites for placing two electron in the lowest single-particle states, Table \ref{table:L2.5} and Fig. \ref{fig:L2.5map}. In addition we have assumed that the subshells $R_1$ and $R_2$ are degenerate, but in reality they differ in the fourth decimal place.

\begin{figure}[h!]
\centering
\input{electrons}
\caption{Different SD for the single-particle basis in Table \ref{table:L2.5}. The blue spheres represent the electrons and the spin configurations of those are denoted beneath.}
\label{fig:electrons}
\end{figure}
%
The we get a total of six slater determinants for this system. A proper way of dealing with this is using a multireference coupled-cluster theory \cite{multiCCSD}. Which means that the ground state slater determinant are a linear combination of those six possibilites with a probability amplitude for each. Our calculation in Table \ref{table:w=1_L=2.5} for the 2-electron system are with the SD no. 3 in Fig. \ref{fig:electrons}. 
The same problem applies to the HF energy as well. In the calculation in Table \ref{table:w=1_L=2.5} we are assuming that the ground-state slater determinant is no. 3 in Fig. \ref{fig:electrons}. But we have to take into account the other slater determinants as well which then gives us new coefficients for us to minimize. This is the multi-configurational Hartree-Fock method \cite{multiHF}. 

For the 4-electron system we have only one possible slater determinant which have total spin 0, see Table \ref{table:w=1_L=2.5}. We have convergence for the HF energy after $R=10$. But there is no convergence for the CCSD energy at all. This is beacuse that we dont have a big enough harmonic oscillator basis due to the cusp of the absolute value, remember that our maximal $n_x$ is only 20. Together with the fact that we are using Newton's method when we search for the energy minimum in CCSD makes the calculation unstable. However, when we apply a lower $\omega$-value we see that the ground-state energy goes a down for HF-energy and the 4-electron system. This makes sense since an increased frequency $\omega$, means a steeper potential Fig. \ref{fig:VC and VG potentials}, which means that the electrons are pushed closer together, giving them more kinetic energy. We see this effect in the ciruclar quantum dot ($L_x=0$), see \cite{lohne} and \cite{lars}.

The CCSD energies have good convergence for the potential Eq.~(\ref{Vgauss}) and we see that the ground-state energy increase with the hight of the gaussian $V_0$, see Tables \ref{table:w=1_V=8}, \ref{table:w=1_V=8} and \ref{table:w=1_V=20}. Comparing the HF energy with CCSD energy, we see that the difference does not change much after the energies have converged. This is what we expect since the HF-calculations take into account 1p-1h excitation, while the CCSD have the 2p-2h contribution in addition, this will lower the CCSD-energy compared to the HF-energy with a constant value. 
%
We have done an analysis of the ground state energies for a model space with $R=12$. We look at the 4- and 8-electron system with increasing $V_0$ for the gaussian potential Eq.~(\ref{Vgauss})

\begin{figure}[h!]
\centering
$\begin{array}{cc}
\includegraphics[width=2.5in]{nh=4_Vgauss} & \includegraphics[width=2.5in]{nh=8_Vgauss} \\ \includegraphics[width=2.5in]{nh=4_Vgauss_delta} & \includegraphics[width=2.5in]{nh=8_Vgauss_delta}
\end{array}$
\caption{}
\label{HFvsCCSD}
\end{figure}
%
For the 8-electron system, the CCSD energy does not converge for $V_0=2.5$ to $V_0=5.5$. The explanation for this could be that the slater determinant in that range is wrong. The same problem we saw earlier that we may need a multireference calculation. If we look at the single-particle energies in Fig. \ref{fig:singleparticleenergies}, we see that the single particle states in the range  $V_0=2.5$ to $V_0=5.5$ is overlapping many other excited single-particle states. This is not the case with the 4-electron system. 

\begin{figure}[h!]
\centering
$\begin{array}{c}
\includegraphics[width=5in]{single-particle_energies} 
\end{array}$
\caption{}
\label{fig:singleparticleenergies}
\end{figure}

\begin{figure}[h!]
\centering
$\begin{array}{cc}
\includegraphics[width=3in]{single-particle_hole_energies} & \includegraphics[width=3in]{correlation_energies} \\ 
\end{array}$
\caption{}
\end{figure}
%
We define the correlation energy by 
\begin{equation}
   E_{corr} = E_{CCSD} - \sum_{ij} \bra{i} \hat{h} \ket{j}
\end{equation}
%
The contribution to the energy that comes from the correlation energy falls off with increasing $V_0$. This is because when $V_0$ increases the tunneling barrier becomes bigger and the electrons will interact less with each other.
%
Finally we have looked at the triple quantum dot system. This consists of  essentially two gaussian functions in the harmonic oscillator well. We tried to adjust the parameters so that we had 6 possible single-orbital states in the ground-state, see Table \ref{table:tripledot}. Here we have the states $p=0$ and $p=1$ that are exactly degenerate, because of spin. But $p=1$ and $p=2$ are not excatly degenerate, they differ in the second decimal place. However $p=2$ and $p=3$ are degenerate because of spin. This makes it difficult to define a possible slater determinant, therefore we have tried to run the CCSD code with 2,4 and 6-electrons. And as we can see they do not converge at all, except for the 2-electron system that convergences poorly. The conclusion must be that the ground state slater we start off with is not the appropriate one. The coupled-cluster calculation seem to be very dependent on how close the energy levels are, as we have seen in the case of 8-electrons in Fig. \label{HFvsCCSD}, when $V_0$ is between 2.5 and 5. It would have been interesting to see if we could tune these parameters so that the 6 states in the triple well become more degenerate. 

\chapter{Conclusions}
\begin{quote}
\it
"My interest in science is to simply find out about the world, and the more I find out the better it is, like, to find out.`` (From ''The pleasure of finding things out``)
\end{quote}
\begin{quote}
R. P. Feynman
\end{quote}

We have developed a method for studying several interacting electrons confined in two-dimensional potential wells. In these type of quantum wells the electrons can tunnel to other excited states and in the literature such systems are dubbed two-dimensional double quantum dot systems. 
They have never been studied with an advanced many-body quantum theory like coupled-cluster theory. 
The model that we have considered is a two-dimensional double quantum well, with a circular quantum dot with a \emph{bump} in one of the dimensions and a standard one-dimensional harmonic in the other dimension. 
One of the first goals was to find the single-particle eigenfunctions of these potentials and expand them in a harmonic oscillator basis. The reason we do this is because of the known two-particle interaction elements Simen Kvaal \ref{kvaalOpenFCI} developed in his PhD thesis for the harmonic oscillator eigenfuntions. The next goal was to transformation his two-particle interaction elements from a polar basis to a cartesian basis representation. His code gives us effective interaction elements as well. We have implemented this but did not use it to produce our results. We have focused on the bare interaction only. 
But the effective interaction elements would most likely give us a better estimate for the energy. The main aim was not to give a precise description of our model but prove that our method of transformation and coupled-cluster theory together with these new interaction elements actually produce meaningful resuls. We have also developed a Restriced Hartree-Fock (RHF) program to check that the coupled-cluster  results are not totally wrong, since RHF gives a variational energy. 

We have tried to limit ourselves to closed-shell systems. We have considered ground-state slater determinants based on the lowest single-particle (in energy) states, but sometimes it is difficult to define properly what is a good closed shell systems, like in the triple dot. But for a systems where the ground-state has a single well-defined slater determinant the coupled-cluster calculations perform well in sense of convergence properties. 
Other methods for solving the double well have been cumbersome and the dimensionality quickly raises when we try to diagonalize the system \ref{doubledot}. A VMC calculation for this system is possible, but then again, finding the Jastrow-factor is an open issue. We could instead use our CCSD method use the found ground-state energies as an input in a DMC calculation.    

\section{Continuation of this thesis}
The first improvement we could do is to do a HF-calculation and use this as input to our coupled-cluster code. This  gives us new single-particle states and two-particle interaction elements. This will most likely give a better convergence for our coupled-cluster implementation. Another improvement is to include the option of chosing the analytical coulomb interaction elements \ref{rontani}, when using the standard interaction. This would speed up the part for generating the interaction elements. Furthermore, the next step is to use an effective Coulomb interaction as this will most likely result in a better estimate of the energies.
The final aim would be to extend our stationary coupled cluster approach to include time-dependent external probes. This would bring in close contact with experiments on quantum dot systems.

%We have used a simple difference method when finding the singleparticle eigenvectors and eigenvalues, but we could used a the better Richardson's deferred extrapolation method \ref{computerjensen}.



\clearpage
\begin{figure}
\begin{quote}
\it
"There's beauty not just at the dimension of one centimeter; there's also a beauty at smaller dimensions."   
\end{quote}
\begin{quote}
R. P. Feynman
\end{quote}
\end{figure}


\appendix
\chapter{Derivation of the Baker-Campbell-Hausdorff formula }
\label{dix:derivation of baker-campbell-hausdorff}

We will now derive the following relation:

\begin{align}
e^{-\lambda \hh{A}}\hh{B}e^{-\lambda\hh{A}} &= \hh{B} + \lambda \ql[\hh{B}, \hh{A}  \qr] + \frac{\lambda^2}{2!} \ql[\ql[\hh{B}, \hh{A}  \qr],\hh{A} \qr] + \frac{\lambda^3}{3!} \ql[ \ql[\ql[\hh{B}, \hh{A}  \qr],\hh{A} \qr],\hh{A} \qr] +...
\nonumber \\
&= \sum_{n=0}^{\infty} \frac{1}{n!} (\text{ad}(\hh{B}))^n \hh{A}, \qquad (\text{ad}(\hh{B}))^n \hh{A} = \ql[ \ql[\ql[\hh{B}..., \hh{A}  \qr],\hh{A} \qr],\hh{A} \qr]\,\,\,(\text{n times})
\label{app:show this relation} 
\end{align}
%
We define 

\begin{equation}
\hh{G}(\lambda) = e^{-\lambda \hh{A}}\hh{B}e^{\lambda \hh{A}}
 \label{app:G(lambda)}
\end{equation}
%
We want to do a Taylor series expansion with this function around $\lambda = 0$,

\begin{equation}
\hh{G}(\lambda) = \hh{G}(0) + \ql .\frac{d\hh{G}}{d\lambda} \qr|_{\lambda = 0} (\lambda-0) + \frac{1}{2!} \ql .\frac{d^2\hh{G}}{d^2\lambda} \qr|_{\lambda = 0} (\lambda-0)^2 + \frac{1}{3!} \ql .\frac{d^3\hh{G}}{d^3\lambda} \qr|_{\lambda = 0} (\lambda-0)^3 + ...
 \label{app:Taylor expansion}
\end{equation}
%
{\bf{0th order}}:
\begin{equation}
\hh{G}(0) = \hh{B} 
 \label{app:0th order}
\end{equation}
%
{\bf{1th order}}:
\begin{equation}
\ql.\frac{d\hh{G}}{d\lambda} \qr|_{\lambda = 0} = \ql[ -\hh{A}e^{-\lambda \hh{A}}\hh{B}e^{\lambda \hh{A}} + e^{-\lambda \hh{A}}\hh{B}e^{\lambda \hh{A}}\hh{A} \qr]_{\lambda = 0} = \ql[\hh{G}(0),\hh{A}\qr] = \ql[\hh{B},\hh{A}\qr] 
 \label{app:1th order}
\end{equation}
%
{\bf{2nd order}}:
\begin{equation}
\frac{d^2\hh{G}}{d^2\lambda} = \frac{d\hh{G}}{d\lambda} \ql[\hh{G},\hh{A}\qr] = \ql[ \frac{d\hh{G}}{d\lambda},\hh{A}\qr] =  
\ql[ \ql[ \hh{G},\hh{A}\qr] ,\hh{A}\qr] \label{app:2th order} 
\end{equation}
%
\begin{equation}
\Rightarrow \ql.\frac{d^2\hh{G}}{d^2\lambda}\qr|_{\lambda=0} = \ql[ \ql[ \hh{B},\hh{A}\qr] ,\hh{A}\qr] \label{defapp:2th order 2} 
\end{equation}
%
{\bf{3rd order}}:
\begin{equation}
\frac{d^3\hh{G}}{d^3\lambda} = \frac{d\hh{G}}{d\lambda} \ql[\ql[\hh{G},\hh{A}\qr],\hh{A}\qr] = \ql[\ql[\frac{d\hh{G}}{d\lambda} ,\hh{A}\qr],\hh{A}\qr] =  
\ql[ \ql[ \ql[ \hh{G},\hh{A}\qr],\hh{A}\qr] ,\hh{A}\qr] \label{app:3th order} 
\end{equation}
%
\begin{equation}
\Rightarrow \ql.\frac{d^3\hh{G}}{d^3\lambda}\qr|_{\lambda=0} = \ql[ \ql[\ql[ \hh{B},\hh{A}\qr],\hh{A}\qr] ,\hh{A}\qr] \label{defapp:3th order 2} 
\end{equation}
%
{\bf{n-th order}}:
\begin{equation}
\frac{d^n\hh{G}}{d^n\lambda} = \frac{d\hh{G}}{d\lambda} \underbrace{\ql[ \ql[\ql[\hh{G}..., \hh{A}  \qr],\hh{A} \qr],\hh{A} \qr]}_{\text{$(n-1)$ nested commutators}} =  \underbrace{\ql[ \ql[\ql[\hh{G}..., \hh{A}  \qr],\hh{A} \qr],\hh{A} \qr]}_{\text{$n$ nested commutators}} \nonumber 
\end{equation}







\bibliographystyle{plain}% sorting of plainnat A = Author, C = Where in the text its sorted
\bibliography{minbib}


\end{document}

